% Framework v3.069: verified five-core cubic/56-state/orientation finite comparison; physical decay open.
% Framework v3.032: correct finite two-state trial-transfer versus full physical pole identification using the already-derived source-visible Schur resolvent.
% Framework v3.026: five-core sourced spectral-weight residual and explicit physical embedding guard.
% Framework v3.025: action-sensitive reference Gram is not a physical hadron pole.
% Current integrated paper source: Paper7B-QCDReadouts-v1_159.tex (v3_017).
% Prior source: Paper7B-QCDReadouts-v1_158.tex. New cross-DAG mathematics is scoped, not physical-occurrence certified.
% Paper7B-QCDReadouts-v1_158: polyhedral full-2O working carrier and exact 2I meson/baryon geometry; physical occurrence remains conditional.
% Paper-series v2.381: complete algebraic exceptional torsion atlas and all-cocycle twist-stability; original claims retained.
% Paper-series v2.380: unified bibliography, exact exceptional double/length/gauge mathematics in I-B; previous theorem labels retained.
% Paper-series v2.379: full equivariant Dirac/Hodge/eta/torsion atlas in Paper I-B; paper-root references synchronized.
% Paper-series v2.378: 2026 exceptional literature consolidation and reference synchronization. The mathematical additions are located in Paper I-B.
% Paper-series v2.377: reference-only rebind to current Paper I-B v1.1; original mathematics retained except new Paper I-B comparison chapter.
% Framework v2.370: derived exceptional-root/Hurwitz comparison in I-A/I-C; other papers reference-only synchronized.
% Framework v2.368: preserve prior II-B convention/unital fixes and verified II-D corrections; current references synchronized.
% Framework v2.366: exact chamber--transporter/Hurwitz integration; physical action and hadron decay gates remain conditional.
% Framework v2.366: scoped Hurwitz/600-cell hadron-source integration; physical APBI/decay frontier remains open.
% Framework v2.358: current 30-root source reference synchronization; selectively updated III-E, III-F, VI-B, VI-C.
% Framework v2.357: normal-jet/Clifford/hadron-shell observable proof additions at scoped mathematical, conditional, or QCD-source levels.
% Framework v2.356: new source-owned conditional results and source-registry synchronization.
% Framework v2.350: repaired Paper VII mathematics; exact finite identities separated from analytic/physical certificates.
% Framework v2.349: integrated exact conditional mathematical closures; G0 and GP1--GP3 retained as axioms.
% Framework v2.345: F40 native three-readout and source-visibility integration.
% Framework v2.344: Twisted Hurwitz C3/Petersen bridge, native baryon source and same-action curvature firewall.
% Framework v2.338: exact Petersen 120-carrier, oriented two-source and source-ownership results. No hadron mass prediction.
% Framework v2.327: exact two-core graded transport, conditional meson spin source, native mass-spectrum firewall.
% Framework v2.323: Petersen finite strong-source coherence, exact multi-core mirror spectra, QCD action-incidence and continuum firewalls. Existing source theorems preserved.
% Framework v2.275: exact positive dense-twin word cubature, finite repeatability and eighth-moment classification, nested sampling frames, joint action and POVM-jet exports; preserves v2.274.
% Framework v2.274: two-scalar positive convolution certificates, first-leaked exceptional coset detection, branch-symmetric twin mixing, and finite endpoint moments; all v2.273 branches preserved.
% Framework v2.273: restored missing Problem 3 theorems from v2.268; exact exceptional frames, twin geometry, convolution repeatability and projector/occurrence firewalls; merged on v2.272 without overwriting independent branches.
% Framework v2.272: merged v2.270 covariant native-readout mathematics with v2.271 Problem 4 finite-action Hopf hypercharge/holonomy corrections; no universal bulk-gauge uniqueness claim.
% Independent branches: v2.270 readout/guard exports and v2.270 K16 action-consumed charge closure; all prior scoped assumptions retained.
% Framework v2.267: exact Paper II-A converse/observable certificates, III-F action tests, V-B right orbit and V-C native gauge scope; prior v2.266 preserved.
% Framework v2.266: exact Paper II-C outcome-count/readout converse results and scoped downstream imports; 30-root references synchronized.
% Framework v2.265: II-C coordinate-stage constitution and exact first-symbol host inheritance; prior mathematical domains retained.
% Framework v2.264: II-B reverse and action-bridge integration; preserves prior mathematical domains; active 30-root references synchronized.
% Framework v2.263: Paper I-D adds bridge-controlled root-carry parity and source-relative old-cap cancellation; other root mathematical bodies retained; active 30-root references synchronized.
% Framework v2.262: merges Paper II-F three-gap closure from v2.260 into v2.261, preserving newer II-B/II-D repairs; 30-root references synchronized.
% Framework v2.261: merged Paper II-B GP2-level/geometric repairs; preserves v2.260 II-D repairs; active references synchronized.
% Framework v2.260: Paper II-D signed RS, finite-symbol typing, and conditional curvature/Jordan repairs; references synchronized.
% Framework v2.259: conditional native SU(2) atlas and explicit later discharge; v2.258 mathematical bodies retained elsewhere.
% Framework v2.258: preserves parallel II-A/dynamics update and adds complete readout/frequency closures.
% Framework v2.257: exact minimal complete position readout and arbitrary-preparation frequency law.
% Framework v2.256: exact native inertia and sharp bounded ordered-source closure; preserves concurrent v2.255 work.
% Framework v2.255: preserves parallel detector integration and adds exact finite readout closures.
% Framework v2.254: shared bibliography/reference sync; substantive Pasterski detector-Ward additions only IV-B/IV-C/VII-C.
% Framework v2.253: typed raw coordinate defect in I-A and synchronized VI-A import; active references synchronized.
% Framework v2.252: optional shell-frame source-return calculation in III-G; active references synchronized.
% Framework v2.251: forward common-action execution, joint source jets and scoped coupled control.
% Framework v2.251: exact full-heat variations, fermion source jets and explicit physical completion frontier.
% Framework v2.251: certified interacting source jets, pole stability and scoped completion proofs.
% Framework v2.251: interacting finite-source evaluation and coupled local vacuum classification.
% Framework v2.251: exact finite source execution and native causal common-parent bridge.
% Framework v2.251: constructive parent and scoped source-incidence audit; see release report.
% Framework v2.251: common-action two-jet, corrected source/response closure and synchronized references.
% Framework v2.244: scoped rooted source-order Schur diagnostics and synchronized paper references.
% Framework v2.242: signed rooted phase lift, complete transverse return and exact family cubic.
% Framework v2.241: proved rooted conormal geometry, conditional cap realization and scoped action tests.
% Framework v2.239: coherent 30-paper reference rebinding; selective new proofs on established owners.
% Framework v2.236: source registry synchronized; existing mathematical body retained except declared insertions.
\documentclass[12pt,a4paper]{article}
\PassOptionsToPackage{hidelinks}{hyperref}

% Paper-specific support packages not supplied by the shared series macro file.
\usepackage{microtype}
\usepackage{etoolbox}

% Macros-v2_203: v3.077 source-current bindings for unconditionality source integration. No mathematical macro changes.
%% Macros-v2_201: downstream full-carrier geometry update, five current root bindings, no mathematical macro change.
%% Macros-v2_200: Paper III-F v1.280 current-source binding for framework v3.074; mathematical macro definitions unchanged.
%% Macros-v2_198: current project v3.072 III-D and bibliography binding; mathematical macros unchanged.
%% Current project v3.070 source-reference bindings; previous macro bodies unchanged.
% Current II-D reference root for PapersMinimal v3.063; mathematical macros otherwise unchanged.
% Macros-v2_188: re-based current paper identities and A--G family; no change to geometry.
% Macros-v2_184: active Paper III-E v1.200/III-F v1.267 binding for framework v3.031; mathematical macros unchanged.
% Macros-v2_183: updated shared cross-reference bindings for Paper1B/III-E/III-F v3.030; mathematical macros unchanged.
% Macros-v2_182: active Paper II-F/V-A cross-references for framework v3.029; definitions unchanged.
% Paper-series v2.381: new I-B torsion/modular results, bibliography and current source roots.
% Paper-series v2.380: new exceptional bibliography and synchronized paper references.
% Paper-series v2.379 current 31-paper macro/aux registry.
% Paper-series v2.378: current 31-paper AUX/bibliography registry updated for Paper I-B v1.2 literature additions.
% Framework v2.371: current-source reference registry refreshed; old theorem labels retained.
% Framework v2.370: current 30-paper source/AUX root rebinding; native mathematical macros unchanged.
% Framework v2.368: synchronized merged 30-root registry; mathematical notation unchanged.
% Framework v2.358: active 30-paper source registry synchronized; existing mathematical macros unchanged.
% Framework v2.357: selective new source-owned certificates; active reference roots updated for five revised papers.
% Framework v2.356: active 30-root registry synchronized to the 30 papers in this release.
% Framework v2.275: exact positive dense-twin word cubature, finite repeatability and eighth-moment classification, nested sampling frames, joint action and POVM-jet exports; preserves v2.274.
% Framework v2.274: filename and root-source synchronization.
% Framework v2.273: restored missing Problem 3 theorems from v2.268; exact exceptional frames, twin geometry, convolution repeatability and projector/occurrence firewalls; merged on v2.272 without overwriting independent branches.
% Framework v2.272: merged v2.270 covariant native-readout mathematics with v2.271 Problem 4 finite-action Hopf hypercharge/holonomy corrections; no universal bulk-gauge uniqueness claim.
% Independent branches: v2.270 readout/guard exports and v2.270 K16 action-consumed charge closure; all prior scoped assumptions retained.
% Framework v2.269: synchronized 30 active AUX roots; mathematics and bibliography preserved from Macros-v2_161.
% Framework v2.266: Paper II-C converse-derived outcome resources; II-E, III-C, IV-A/B, V-A and VI-B downstream mathematical closures; 30-root registry synchronized.
% Framework v2.265: II-C coordinate-stage constitution and exact first-symbol host inheritance; prior mathematical domains retained.
% Framework v2.264: paper II-B reverse/audit and IV-A observable-minimal exchange; synchronized external references to 30 active roots.
% Framework v2.263: Paper I-C adds bridge-controlled root-carry parity and source-relative old-cap cancellation; other root mathematical bodies retained; active 30-root references synchronized.
% Framework v2.262: integrates Paper II-F radius, Casimir and readout results; active 30-root references synchronized.
% Framework v2.261: merged Paper II-B GP2-level/geometric repairs; preserves v2.260 II-D repairs; active references synchronized.
% Framework v2.260: Paper II-D signed RS, finite-symbol typing, and conditional curvature/Jordan repairs; references synchronized.
% Framework v2.259: conditional native SU(2) atlas and explicit later discharge; v2.258 mathematical bodies retained elsewhere.
% Framework v2.258: preserves parallel II-A/dynamics update and adds complete readout/frequency closures.
% Framework v2.257: exact minimal complete position readout and arbitrary-preparation frequency law.
% Framework v2.256: exact native inertia and sharp bounded ordered-source closure; preserves concurrent v2.255 work.
% Framework v2.255: preserves parallel detector integration and adds exact finite readout closures.
% Framework v2.254: Pasterski bibliography and synchronized 30-paper external-reference registry.
% Framework v2.253: typed raw coordinate defect in I-A and synchronized VI-A import; active references synchronized.
% Framework v2.252: synchronized source registry; optional III-G shell-frame calculation.
% Framework v2.251: forward common-action execution, joint source jets and scoped coupled control.
% Framework v2.251: exact full-heat variations, fermion source jets and explicit physical completion frontier.
% Framework v2.251: certified interacting source jets, pole stability and scoped completion proofs.
% Framework v2.251: synchronized active source registry; mathematical notation unchanged.
% Framework v2.251: synchronized active source registry; mathematical notation unchanged.
% Framework v2.251: synchronized active source registry; mathematical notation unchanged.
% Framework v2.251: synchronized source registry; no mathematical notation changes.
% Framework v2.244: synchronized 30-paper AUX registry and rooted phase-source action-order tests.
% Framework v2.242: signed rooted phase lift, complete transverse return and exact family cubic.
% Framework v2.241: proved rooted conormal geometry, conditional cap realization and scoped action tests.
% Framework v2.239: all 30 active source roots synchronized; no mathematical macro changes.
% Framework v2.238: Paper IV A/B/C split and 30-root reference registry.
% Framework v2.237: Paper VII-D technical split and 28-root reference registry.
% Framework v2.236: Paper VI-C technical-volume reference added; all current root bindings advanced.
% Framework v2.235: 26-root AUX and bibliography bindings synchronized; mathematical macros unchanged.
% BEGIN INLINE SOURCE: Paper00-TeXBuildCompatibility-v1_0.tex
% Rendering-only repairs for inherited long appendices and bitmap text fonts.
% Integer counters through 26 retain the original labels; later values use AA,
% AB, ... . This file changes no action coefficient, state or theorem premise.
\RequirePackage{alphalph}
\makeatletter
\let\BHPrimitiveAlph\@Alph
\newalphalph{\BHExtendedAlph}{\BHPrimitiveAlph}{26}
\let\@Alph\BHExtendedAlph
\AtBeginDocument{%
  \@ifpackageloaded{microtype}{\microtypesetup{expansion=false}}{}%
}
\makeatother

% END INLINE SOURCE: Paper00-TeXBuildCompatibility-v1_0.tex

% Framework v2.164: current reference/bibliography profile; mathematical definitions retained.
% Framework v2.162: finite metric, functional Hilbert and scoped Ward integration.
% Framework v2.161: inherits Macros-v2_57.tex exactly except the current external-reference binding.
% Macros v2.57 / framework v2.156: current reference registry only; mathematical definitions unchanged.
% Framework v2.143 / Macros v2.55: current external-reference bindings; mathematical notation unchanged.
% Framework v2.110 / Macros v2.54: current external-reference registry and bibliography; mathematical notation inherited.
% Framework v2.67 / Macros-v2_53: adds Born-action-cell and joint-EH readout symbols; absolute unit realization is metrological; MATCH_{J/scal}^{EH} remains the sole G24 normalization bridge.
% Framework v2.35 / Macros-v2_52: grouped paper numbering, active AUX registry, and Born-v6_20 bibliography binding synchronized; notation unchanged.
% Framework v2.26: active-source AUX and bibliography bindings synchronized; notation unchanged.
% Framework v2.09: reviewed early-reciprocity integration; scoped physical imports.
% Framework v2.08: Born--Euler mathematics and conditional complete-action potential criterion.
% Framework v2.07: active AUX registry includes Paper I-B v1.2; shared bibliography Born-v6_18.
% Macros-v2_46 / framework v2.04: mathematical atlas and downstream realization correction.
% Macros-v2_46 / framework v2.03: scoped structure integration and current dependency bindings.
% Macros-v2_46 / framework v2.01: shared source bindings; prior revision comments are historical.
% v2.42 / framework v1.99: Synchronized RG/AS owner and bibliography bindings.
% v2.41 / framework v1.99: active manuscript AUX bindings and geometric-closure bibliography synchronized.
% Release-v1.80 centralized source synchronization for scalar-circle/comparison-band integration.
% Release v1.79: Hurwitz hierarchy and kinetic-normalized action response integration; physical mass matching remains conditional.
% Release v1.78: approved M8 closure-batch integration; physical and analytic domains retained.
% Current v1.78 registry: all active TeX roots are synchronized below.
% Earlier release comments are historical scope records, not current status.
% Release-v1.77 synchronization: native same-owner Noether closure, generator--flow/history typing, analytic-certification firewalls, scalaron/K16, LF/QCD, quantum-Ward, and RAR boundary-owner transport are registered; shared mathematical notation is unchanged.
% Release-v1.76 synchronization: registry points to the v1.76 active paper roots.
% Release-v1.76 synchronization: registers the local A+Phi+g Noether/constraint closure and dynamics-frontier propagation; shared mathematical notation is unchanged.
% Release-v1.75 synchronization: registry points to the v1.75 active paper roots.
% Release-v1.75 synchronization: integrates the joint-Dirac Einstein bulk real form and the local odd-scalar heat-coefficient action on explicit analytic first-jet/interior domains; v1.74 Yang--Mills closure is retained. Fermionic reality, pure-metric curvature-squared/exact-nonlocal terms, noncanonical boundaries, physical vacuum selection, and actual solution-history tokens remain open.
% Release-v1.74 synchronization: registry points to the v1.74 active paper roots and retains the v1.73 theorem labels.
% Release-v1.74 synchronization: canonical Macros-v2_34.tex; generated local nonlinear Yang--Mills action, constrained connection-field phase space, source-free Gauss momentum map, and Legendre closure are integrated on their proved domains. The remaining non-gauge ACT_cl sectors and actual solution-history selection remain open.
% Release-v1.71 synchronization: closes the former JOINT_x/JOINT_5
% connection-occurrence gap by canonical structure-group extension of the
% nonflat Born connection, faithful associated representation, and the local
% cotangent--Clifford operator constructor.  Actual fluctuation members,
% actions, vacua, solutions, calibrations, and analytic QFT domains remain
% separately typed.  The active-AUX registry is advanced once.
% Release-v1.71 canonical macro and active-AUX registry.
% ======================================================================
% Macros-v2_39.tex -- canonical publication macro file
% Release v1.71 retains the v1.70 null-ray/QCD boundaries and additionally
% removes the duplicate structural connection actuality premise.  It does not
% claim global bundle topology, select a gauge/Higgs member, or promote the
% finite construction to an interacting continuum theory.
% Release v1.68 constructs the finite light-front total-fraction operators and
% parton number/momentum operator measures, proves their simplex and Hausdorff-limit
% theorems, and separates positive cut measures from scheme-renormalized
% light-ray distributions.  It also separates finite real Schur poles, stable
% norm-resolvent/Riesz pole-residue limits, and continued resonance
% pole/residue limits.  The Paper-VII-C scattering package is claim-indexed:
% local QFT and spectral claims import no S-matrix or infrared premise, while
% asymptotic amplitudes, LSZ, S-matrix poles, completeness, and cross sections
% import only their named massive or massless/long-range clauses.
% Release v1.67 factorizes BND48_H into a same-action Hessian/Schur
% descendant and still-open typed boundary-trace, critical-lift, and
% refinement-limit conditions.  It proves the finite-collar Dirichlet-to-
% Neumann law sqrt(y)tanh(L sqrt(y)), its nonuniform critical limit, and the
% finite-return ultra-infrared turnover.  The Paper-I-A Hardy constant mode is
% an exact geometric near-hit, not yet the physical metric boundary mode.
% The current local f(R) Hessian remains massive and therefore still refutes
% a RAR prediction on the minimal action branch.
% Release v1.66 derives the transverse depth-six normalizer bridge from the
% already generated C6 Reynolds image and Paper-I-A binary-dihedral normalizer.
% The bridge is the unique least 2D3-stable scalar-containing realization of
% the faithful quotient 2D3/C6=C2, has trace defect six and hence K=48, but is
% transverse to the nested active carrier.  Exact K=44 and K=52 contraction,
% semigroup, and energy-midpoint identities are registered as structural
% readouts rather than projector events.  The active E48 carrier gap remains
% unchanged.  A same-action boundary-symbol/Calderon intertwiner is isolated
% as the necessary and still unproved owner bridge to the static DtN sort;
% representation equivalence alone neither derives RAR nor creates a K=48
% action junction, pole, occupation, or cosmological transition.
% Release v1.65 derives the native E_48-centered carrier mass gap from the
% GP1--GP3-generated active spectrum, distinguishes its carrier-free interval
% from a particle located at E_48, and withdraws the unsupported promotion to
% a positive continuum Yang--Mills mass gap.  The latter remains a valid
% conditional transfer theorem on an enlarged continuum branch, but it is not
% an active framework prediction or an unfinished native-gap proof.
% Release v1.64 synchronizes Derivation Batches 14--22.  It registers the
% GP1--GP3 heat-profile orbit, same-Hessian electroweak response and pole
% ownership, K16/K48 action-response corrections, perturbation state and
% non-Gaussianity typing, finite-vacuum/gap/QCD limit repairs, factorized
% position-area ownership, the actual static metric response, and the ordered
% Wilson/string-tension limits.  Strict continuum convergence, positive
% physical gap, surviving higher cumulants, and positive tension remain open.
% Release v1.61 closes TIMEWIDTH16 as an independent direct-production datum.
% JUNC_{2|7} continues to own the finite intrinsic microscopic profile/width,
% while the actual same-action background/production history determines its
% physical-time pullback, support, and duration.  The current Paper-I-A--VII-C AUX
% registry is advanced once; shared mathematical notation is unchanged.
% Release v1.60 derives the final-octet primitive-remnant depletion, canonical
% Hopf-current representative, exact Hopf descent to dOmega/(4pi), and
% closed-FLRW a^{-3} neutral-current shape.  It retypes K=48 as the onset
% of the final one-unit geometric transfer and records cDM/RAR as distinct
% physical projections of that common neutral parent while preserving
% occurrence, Hilbert-owner, stress-matching, and static-kernel premises.
% Release v1.59 transports the v1.58 exceptional-shell/K16 closure through
% Papers III-F and VI-B--VII-C.  It replaces Paper VI-B's obsolete mixed-Hessian
% output typing by the third-variation pair vertex and Bogoliubov occupation
% support, derives the post-seed Higgs-pair feedback vertex and its exact
% alpha--delta kernel factors, propagates the unique K16/K48 action-junction
% split, identifies the QCD delta count with the upstream 2O shell atom only
% as a common-scalar numerical compatibility, and narrows THERM16 to dynamical
% approach/beta selection because the finite Gibbs/KMS target is already
% unique on a supplied simple observable algebra once its Hamiltonian and
% inverse temperature are supplied; direct-sum central weights remain typed.
% Release v1.58 registers the exceptional nearest-shell closure, the derived
% transverse K=16 action event, the unique interior action junction, the
% type-correct ordered pair vertex, the thin-layer squeeze law, and the
% radiation-like/matter-like kinematic split.  REG_0, transfer magnitude,
% physical layer width, and thermalization remain explicitly conditional.
% Release v1.57 registers the active-depth-six obstruction, the native
% K=48 carrier desert, the typed 7|9 owner reduction, and the dynamical
% reheating handoff across Papers III-B--VI-B.
% Release v1.56 registers the least-support n_2 inflation dressing of the
% persistent n_7 Laurent response, the analytic global-convexity and positive
% endpoint-measure proofs, the derived inflation core and K=16 coefficient
% freeze, and the no-particle dark-sector split into static RAR plus the
% conditional homogeneous effective-curvature response at reflected K=48.
% There is no n_6 action, scalaron-DM mass, or second Einstein--Hilbert term.
% Release v1.55 registers the protected-octet resolved/residual reflection
% q_sp+Z_W=8, its j<->8-j and K<->64-K scale covariance, and the derived
% rank-one-exit quadratic grade C_1^2/U=C_2 in Papers III-C, III-E, and VIII.
% The reflection is a typed trace/scale theorem, not a time reversal or a
% physical event-occurrence theorem; reheating transfer remains conditional.
% Release v1.54 registers the exact positive-part Laurent indices, the
% active-domain/terminal-endpoint action split, the no-second-EH firewall, and
% the q_DM=6/K_DM=48 scalaron readout typing across Papers III-D/E/F, VI,
% VIII, and IX.
% Release v1.53 registers the common nonnegative Laurent obstacle projection
% for n_1^+ and n_7^+, discharges the three former exit/freeze assumptions,
% and advances the Paper-III-D/E/F, VIII, and IX consumers.  Shared notation
% and bibliography ownership are unchanged.
% Release v1.51 registers the PQTI derivation batch and its updated
% Paper-III-A/C/D/E/F, VIII, and IX consumers.  Shared notation and
% bibliography ownership are unchanged.
% Release v1.50 advances the Paper-III-B AUX target to v1.40 after the
% rigorous audit repair; shared notation and bibliography ownership are unchanged.
% Release v1.49 advances the Paper-II-E/F AUX target to v3.14 after the
% minimal audit repair; shared notation and bibliography ownership are unchanged.
% Release v1.46 synchronizes Paper II-A's minimal frozen constructor-origin
% theorem and its five dependent AUX targets.  Shared notation is unchanged.
% Release v1.45 updates the Paper-II-B AUX target to v2.23 after the
% GP3-rigidity and appendix-placement repair; shared notation is unchanged.
% Release v1.44 advances the centralized Paper-I-A--VII-C AUX registry after the
% Paper-I-A v5.16 protected-completion scope repair.  Every current
% paper imports this one file; no paper-level opt-in or \externaldocument list
% is permitted.  The bibliography remains centralized here through Born-v6_13.
% Shared bibliography, theorem infrastructure, colors, notation alignment,
% and current Paper-I-A--VII-C external-reference registration.
% This file centralizes common imports, bibliography sorting, colors, theorem
% infrastructure, and the notation shared by the complete series.
% v2.9 (release v1.43) made the Paper-I-A--VII-C AUX registry automatic and
% centralized the bibliography and all external-reference ownership.
% v2.8 release-registry refresh (v1.33) registers the rooted generated-parent
% cascade and the POL/REG/free-law/11+4/Born-group closure revisions.  No
% shared mathematical notation or Paper-II-A fundamental principle is changed.
% v2.8 release-registry refresh (v1.32) registers the universal generated-
% continuation revisions in Papers II-A, III-A, III-C, IV, and VIII.  No
% shared mathematical notation or Paper-II-A fundamental principle is changed.
% v2.8 release-registry refresh (v1.30) registers the CRP4/projective-state
% source revisions and their direct registry-only consumers.  No shared
% mathematical notation or Paper-II-A axiom is changed.
% v2.7 release-registry refresh (v1.27) updates only the current external
% Paper-I-A--VII-C basenames after the placement/assumption synchronization; no
% shared mathematical notation or Paper-II-A axiom is changed.
% v2.7 synchronizes the geometric defect-repair revision: Paper-II-A v3.42,
% Paper-III-C v1.26, Paper-III-D v1.27, Paper-III-E v1.32, and Paper-III-F
% v1.34.  It registers the common trace-defect/spectral-filtration theorem,
% its owner-typed response pullback, and the reduced GRED and TERM interfaces.
% No shared mathematical notation is changed.
% v2.6 synchronizes Paper-III-E v1.31's cone-type repair, minimal-response
% classification, and explicit classical action-role boundary, together with
% Paper-III-F v1.32's downstream response-atlas synchronization.  No shared
% mathematical notation is changed.
% v2.5 synchronizes the Paper-III-D v1.26 faithful-trace-duality proof,
% exact ordered-owner joint locus, and typed alternative-presentation
% bookkeeping repair, together with its Paper-III-E and closure-frontier
% references.  No shared mathematical notation is changed.
% v2.4 synchronizes the Paper-III-C v1.25 orbit/process/owner repair, its
% downstream Papers III-B--F, V, VI, VIII, and IX revisions, and Born-v6_13.
% v2.3 synchronizes the Paper-II-A v3.41 witness/support/principal inheritance
% layer, the Paper-III-A v1.24 joint-realizability/readout firewall, all dependent
% host-rate revisions, and the current Paper-I-A--VII-B source registry.
% v2.2 synchronizes the generated active-realization, determinant-depth,
% dimension-weighted hierarchy, and typed cosmological-event interfaces.
% v2.1 synchronizes the exceptional-support interface revisions: Paper I-A
% v5.9, Paper II-A v3.30, Paper II-E/F v3.7, and Paper IV v4.16.
% v2.0 introduces the stage-safe G0 occurrence notation, synchronizes the
% active Paper-I-A--IV source registry (current rigorous Paper-III closure:
% III-A v1.17, III-B v1.22, III-C--E v1.17, III-F v1.14, IV v4.14), and makes
% external-reference self-skip depend on logical paper role rather than a
% versioned job name.
% ======================================================================

% ======================================================================
% macros.tex — Shared macros for the Born Geometry series
% ======================================================================

% ======================================================================
% 1. PACKAGES, BIBLIOGRAPHY, COLORS, AND UNIT DEFINITIONS
% ======================================================================
% Core packages shared across the complete Paper-I-A--VII-C series.  Loading them here makes
% the macro file a genuine series preamble while remaining compatible with existing
% paper preambles that load the same packages before \input.
\usepackage{amsmath}
\usepackage{amssymb}
\usepackage{amsthm}
\usepackage{mathrsfs}
\usepackage{array}
\usepackage{booktabs}
\usepackage{xcolor}
\makeatletter
\@ifpackageloaded{geometry}{}{\usepackage{geometry}\geometry{margin=1in}}
\makeatother
\usepackage{mdframed}

% Shared support packages required by the macros below and by most papers.
\usepackage{ragged2e}
\usepackage{longtable}
\usepackage{tabularx}
\usepackage{enumitem}
\usepackage[most]{tcolorbox}

% Series-wide numeric bibliography convention using the portable BibTeX path.
% All Papers I-A--VII-C use the ordinary \cite interface.
\usepackage[numbers,sort&compress]{natbib}
\providecommand{\printbibliography}{%
  \bibliographystyle{plainnat}%
  \bibliography{Born-v6_60}%
}

% Shared colors.  \providecolor avoids clobbering a paper-specific override.
\providecolor{postulatecolor}{RGB}{230,242,255}
\providecolor{compare}{RGB}{230,245,230}
\providecolor{definitioncolor}{RGB}{230,242,255}

\providecolor{openriskcolor}{RGB}{255,235,235}
\providecolor{bridgecolor}{RGB}{240,240,255}
\providecolor{gold}{RGB}{255,215,0}

% ======================================================================
% 2. CUSTOM ENVIRONMENTS
% ======================================================================
% Common table column types used in all papers.  Guarded so that a paper may
% define its own versions before loading this file.
\makeatletter
\@ifundefined{NC@find@L}{%
  \newcolumntype{L}[1]{>{\RaggedRight\arraybackslash}p{#1}}%
}{}
\@ifundefined{NC@find@C}{%
  \newcolumntype{C}[1]{>{\Centering\arraybackslash}p{#1}}%
}{}
\@ifundefined{NC@find@Y}{%
  \newcolumntype{Y}{>{\RaggedRight\arraybackslash}X}%
}{}
\makeatother

% ======================================================================
% 3. THEOREM STYLES AND ENVIRONMENTS
% ======================================================================
% Custom underlined theorem style (useful for remarks or notes)
\newtheoremstyle{underlined}
  {3pt}   % Space above
  {3pt}   % Space below
  {}      % Body font (empty means normal/roman text, typical for remarks)
  {}      % Indent amount
  {}      % Theorem head font (can be \itshape, \bfseries, etc.)
  {.}     % Punctuation after theorem head
  { }     % Space after theorem head
  {\underline{\thmname{#1}\thmnumber{ #2}}\thmnote{ (#3)}} % Theorem head formatting

% Standard plain theorem style environments.  Each environment is guarded so
% the shared macro file can be loaded after a paper-specific theorem setup.
\theoremstyle{plain}
\makeatletter
\@ifundefined{theorem}{\newtheorem{theorem}{Theorem}[section]}{}
\@ifundefined{lemma}{\newtheorem{lemma}{Lemma}[section]}{}
\@ifundefined{proposition}{\newtheorem{proposition}{Proposition}[section]}{}
\@ifundefined{corollary}{\newtheorem{corollary}{Corollary}[section]}{}
\makeatother

% Definition style environments (typically non-italicized body)
\theoremstyle{definition}
\makeatletter
\@ifundefined{definition}{\newtheorem{definition}{Definition}[section]}{}
\@ifundefined{postulate}{\newtheorem{postulate}{Postulate}[section]}{}
\@ifundefined{closure}{\newtheorem{closure}{Closure}[section]}{}
\makeatother

% Applying the custom underlined style
\theoremstyle{underlined}
\makeatletter
\@ifundefined{remark}{\newtheorem{remark}{Remark}[section]}{}
\@ifundefined{convention}{\newtheorem{convention}{Convention}[section]}{}
\@ifundefined{principle}{\newtheorem{principle}{Principle}[section]}{}
\makeatother

% Framed proof environment.
% Every \begin{proof} ... \end{proof} is automatically framed/shaded.
% This avoids mdframed/trivlist nesting, which can drop the top border after
% theorem-like environments.
\makeatletter
\renewenvironment{proof}[1][\proofname]{%
  \par\addvspace{6pt}%
  \noindent
  \begin{tcolorbox}[%
    enhanced,
    breakable,
    colback=black!4,
    colframe=black!35,
    boxrule=0.35pt,
    arc=1pt,
    left=6pt,
    right=6pt,
    top=4pt,
    bottom=4pt,
    boxsep=0pt,
    before skip=0pt,
    after skip=6pt,
    parbox=false
  ]%
  \pushQED{\qed}%
  \normalfont\itshape #1\@addpunct{.}\enspace
  \normalfont\ignorespaces
}{%
  \unskip\nobreak\hfill\popQED
  \end{tcolorbox}%
  \par\addvspace{2pt}%
}
\makeatother

\mdfdefinestyle{postulatebox}{
  backgroundcolor=yellow!10,
  linecolor=black,
  linewidth=0.4pt,
  skipabove=\baselineskip,
  skipbelow=\baselineskip,
  nobreak=true
}

% ======================================================================
% PROOF-READABILITY HELPERS
% ======================================================================

\newcommand{\ProofPlan}[1]{%
  \par\smallskip
  \noindent\textbf{Proof plan. }#1%
  \par\smallskip
}

\newlist{proofsteps}{enumerate}{2}
\setlist[proofsteps,1]{
  label=\textbf{(\alph*)},
  ref=\alph*,
  leftmargin=2.2em,
  itemsep=0.45\baselineskip,
  topsep=0.35\baselineskip
}
\setlist[proofsteps,2]{
  label=\textbf{(\roman*)},
  ref=\roman*,
  leftmargin=2.0em,
  itemsep=0.25\baselineskip,
  topsep=0.2\baselineskip
}


% ======================================================================
% 4. NUMBER SYSTEMS AND ALPHABETS
% ======================================================================
\providecommand{\RR}{\mathbb{R}} % Real numbers
\providecommand{\PP}{\mathbb{P}} % Projective space / Probability
\providecommand{\CC}{\mathbb{C}} % Complex numbers
\providecommand{\ZZ}{\mathbb{Z}} % Integers
\providecommand{\NN}{\mathbb{N}} % Natural numbers
\providecommand{\HH}{\mathbb{H}} % Quaternions / Hyperbolic space
\providecommand{\LFRayGeom}{\mathsf{LF\!\mbox{-}RAY\!\mbox{-}GEOM}}
\providecommand{\LFRep}{\mathsf{LF\!\mbox{-}REP}}
\providecommand{\LFState}{\mathsf{LF\!\mbox{-}STATE}}
\providecommand{\LFLim}{\mathsf{LF\!\mbox{-}LIM}}
\providecommand{\LFScat}{\mathsf{LF\!\mbox{-}SCAT}}
% Number systems / projective spaces
\providecommand{\CP}{\mathbb{CP}}
\providecommand{\RP}{\mathbb{RP}}

% ======================================================================
% 5. GROUPS AND LIE THEORY
% ======================================================================
\providecommand{\SO}{\ensuremath{\mathrm{SO}}}     % Special Orthogonal group
\providecommand{\SU}{\ensuremath{\mathrm{SU}}}     % Special Unitary group
\providecommand{\Spin}{\ensuremath{\mathrm{Spin}}} % Spin group
\providecommand{\CS}{\ensuremath{\mathrm{CS}}}     % Chern-Simons (capitalized)
\providecommand{\cs}{\ensuremath{\mathrm{cs}}}     % Chern-Simons (lowercase)

% ======================================================================
% 6. COMMON MANIFOLDS AND TOPOLOGY
% ======================================================================
\providecommand{\Sone}{S^1}                 % Standard circle

% ======================================================================
% 7. MATH OPERATORS (Declared via \DeclareMathOperator)
% ======================================================================
\DeclareMathOperator{\End}{End}   % Endomorphisms
\DeclareMathOperator{\Aut}{Aut}   % Automorphisms
\DeclareMathOperator{\spec}{spec} % Lowercase spectrum
\DeclareMathOperator{\tr}{tr}     % Lowercase trace (if you already have \Tr)
\DeclareMathOperator{\supp}{supp} % Support
\DeclareMathOperator{\TrA}{Tr_{A}} % Partial trace over subsystem A
\DeclareMathOperator{\TrB}{Tr_{B}} % Partial trace over subsystem B
% Lie/algebra operators
\DeclareMathOperator{\Ad}{Ad}
\DeclareMathOperator{\Sym}{Sym}
\DeclareMathOperator{\su}{su}
% ======================================================================
% 8. STANDARD OPERATOR-LIKE SYMBOLS (Safe definitions via \providecommand)
% ======================================================================
\providecommand{\vol}{\operatorname{vol}}   % Volume form
\providecommand{\Vol}{\operatorname{Vol}}   % Total volume
\providecommand{\Ric}{\operatorname{Ric}}   % Ricci tensor
\providecommand{\Rm}{\operatorname{Rm}}     % Riemann tensor
\providecommand{\Tr}{\operatorname{Tr}}     % Uppercase trace
\providecommand{\Spec}{\operatorname{Spec}} % Uppercase spectrum
\providecommand{\diag}{\operatorname{diag}} % Diagonal matrix
\providecommand{\rank}{\operatorname{rank}} % Rank
\providecommand{\scal}{\mathcal R} % Scalar curvature
\providecommand{\Cl}{\operatorname{Cl}}     % Clifford algebra
\providecommand{\diam}{\operatorname{diam}} % Diameter

% ======================================================================
% 9. FUNCTIONS, TRIGONOMETRY, AND CALCULUS
% ======================================================================

% Codifferential naming conventions
% (Avoid \ddag: commonly defined by physics.sty and symbol packages)
\providecommand{\codiff}{\mathrm{d}^{\dagger}} % Codifferential operator

% Quantum mechanics bra-ket notation (needed for basis tables)
\providecommand{\ket}[1]{\lvert #1 \rangle}
\providecommand{\bra}[1]{\langle #1 \rvert}
\providecommand{\braket}[2]{\langle #1 \mid #2 \rangle}

% ======================================================================
% 10. PHYSICAL SCALES AND CONSTANTS (Paper 1c / Born Geometry)
% ======================================================================
\providecommand{\Rp}{\ensuremath{R_{\mathrm P}}}     % matched common EH curvature scale (not defined from observed G)
\providecommand{\hBorn}{\ensuremath{h_{\mathrm B}}} % physical representative 2\pi\hbar of one primitive Born scalar action cell
\providecommand{\UEHjoint}{\ensuremath{U_{J}}} % joint-Dirac EH curvature-coordinate scale 8 kappa_J
\providecommand{\Gread}{\ensuremath{G_{\rm scal}}} % scalar-action coordinate c_op^3/(h_B U); physical matching explicit
\providecommand{\GreadJoint}{\ensuremath{G_{J,\mathrm{raw}}}} % raw ordinary-trace coordinate = G_scal/96
\providecommand{\Rds}{R_{\mathrm{dS}}}    % de Sitter Curvature
\providecommand{\Mpl}{M_{\mathrm{Pl}}}    % Planck mass
\providecommand{\lp}{\ell_{\mathrm{Pl}}}  % Planck length
\providecommand{\Pmax}{P_{\mathrm{max}}}  % Maximum momentum
\providecommand{\Mplbar}{\bar{M}_{\mathrm{Pl}}} % Reduced Planck mass
\providecommand{\eff}{\mathrm{eff}}       % Subscript for effective quantities
\providecommand{\Cosmo}{\ensuremath{\Lambda}}                % Cosmological constant (use \Cosmo in source)
\providecommand{\LCDM}{\ensuremath{\Cosmo\mathrm{CDM}}}  % \LambdaCDM (typesets with \Cosmo)
\providecommand{\Lcut}{\ensuremath{\Lambda_{\mathrm{cut}}}}  % Cutoff scale (brace when used as script)
\providecommand{\Rcs}{\ensuremath{R_{\mathrm{cs}}}}      % Configuration-space radius
\providecommand{\CosmoEff}{\ensuremath{\Cosmo_{\eff}}}   % Effective cosmological constant

% Disambiguation: \Rmom is the momentum-space (S^3) radius. 
% Reserve plain R for the 4D curvature scalar in f(R).
\providecommand{\Rmom}{\ensuremath{R_{\mathrm{mom}}}} 
\providecommand{\Punit}{P_{0}} % Momentum unit set by \Rmom (= \hbar/\Rmom)
\providecommand{\Rvec}{\ensuremath{R_{\mathrm{vec}}}} % invariant vector-space radius in the reduced round sector
% Specific numeric values and expressions
\providecommand{\Vlogcap}{L_{\mathrm{Hopf}}^{(1)}}
\providecommand{\Hopf}{\operatorname{Hopf}}
% Scalar phase (S^1) sector + universal cover notation (Paper I-A)
\providecommand{\tper}{T_\star}      % Time period/scale
\providecommand{\Eper}{E_\star}      % Energy period/scale
\providecommand{\ttilde}{\tilde t}   % Universal cover time
\providecommand{\Etilde}{\tilde E}   % Universal cover energy

% ======================================================================
% 10A. SERIES-WIDE LOGICAL / NOTATIONAL DISAMBIGUATIONS
% ======================================================================
% Use these symbols to avoid collisions among: (i) the Paper--III
% no-insertion/geometric-closure class, (ii) Paper--I-A trimer diagonals,
% (iii) Paper--V-A/V-C family-space diagonals, (iv) scalar-index running,
% and (v) protected shell-clock labels.
\providecommand{\Rprim}{\ensuremath{\mathsf R_{\mathrm{prim}}^{\Sigma}}} % primitive finite structural readouts
\providecommand{\Rloop}{\ensuremath{\mathsf R_{\circ}^{\Sigma}}} % closed-loop readout role
\providecommand{\Rbdry}{\ensuremath{\mathsf R_{\partial}^{\Sigma}}} % boundary-conditioned readout role
\providecommand{\Rgen}{\ensuremath{\mathsf R_{\mathrm{fin}}^{\Sigma}}} % finite generated closure of primitive roles
\providecommand{\PhysOcc}{\operatorname{PhysOcc}}
\providecommand{\Actual}{\operatorname{Actual}}
\providecommand{\GeomLic}{\operatorname{GeomLic}}
\providecommand{\InvDist}{\operatorname{InvDist}}
\providecommand{\Dcl}{\ensuremath{\mathsf{D}_1}}      % unmarked positive Dirac-even closed spectral-readout form class
\providecommand{\Dheat}{\ensuremath{\mathsf D_{\mathrm{heat}}}} % one-parameter heat-semigroup subclass
\providecommand{\DtriI}{\ensuremath{D^{\mathrm{tri}}_1}} % Paper--I-A cyclic-trimer diagonal
\providecommand{\DtriII}{\ensuremath{D^{\mathrm{tri}}_2}}
\providecommand{\DfamI}{\ensuremath{D^{\mathrm{fam}}_1}} % Family-space diagonal selector
\providecommand{\DfamII}{\ensuremath{D^{\mathrm{fam}}_2}}
\providecommand{\alpharun}{\ensuremath{\alpha_{\mathrm{run}}}} % scalar spectral-index running
\providecommand{\Ksh}{\ensuremath{K_{\mathrm{sh}}}} % protected shell-clock coordinate
\providecommand{\ddef}{\ensuremath{\delta_{\mathrm{def}}}} % universal defect scalar after Paper--III realization

% ======================================================================
% 10B. SERIES-WIDE LOGICAL STATUS TAGS
% ======================================================================
% Upright small-caps wrapper for audit/status labels.  It prevents LaTeX from
% trying to use unavailable small-caps italic font shapes inside theorem bodies
% such as TU/lmr/m/scit.
\providecommand{\AuditLabel}[1]{}
\renewcommand{\AuditLabel}[1]{{\normalfont\scshape #1}}

% These tags are bookkeeping labels for consistency audits. They do not
% create new mathematical hypotheses; they record how a datum is being used
% in the proof chain.
\providecommand{\StatusPrimitive}{}
\renewcommand{\StatusPrimitive}{\AuditLabel{Primitive}}
\providecommand{\StatusDerived}{}
\renewcommand{\StatusDerived}{\AuditLabel{Derived}}
\providecommand{\StatusImported}{}
\renewcommand{\StatusImported}{\AuditLabel{Imported}}
\providecommand{\StatusConditional}{}
\renewcommand{\StatusConditional}{\AuditLabel{Conditional}}
\providecommand{\StatusBridge}{}
\renewcommand{\StatusBridge}{\AuditLabel{Bridge-derived}}
\providecommand{\StatusComparison}{}
\renewcommand{\StatusComparison}{\AuditLabel{Comparison-only}}

% ======================================================================
% 10C. SERIES-WIDE CLAIM-WORDING AUDIT LABELS
% ======================================================================
% These are textual audit labels used to keep modal language consistent
% across the series.  They introduce no new mathematical assumptions.
\providecommand{\ClaimForced}{}
\renewcommand{\ClaimForced}{\AuditLabel{forced}}
\providecommand{\ClaimUnique}{}
\renewcommand{\ClaimUnique}{\AuditLabel{unique}}
\providecommand{\ClaimCanonical}{}
\renewcommand{\ClaimCanonical}{\AuditLabel{canonical}}
\providecommand{\ClaimMinimal}{}
\renewcommand{\ClaimMinimal}{\AuditLabel{minimal}}
\providecommand{\ClaimClosed}{}
\renewcommand{\ClaimClosed}{\AuditLabel{closed}}
\providecommand{\ClaimPrediction}{}
\renewcommand{\ClaimPrediction}{\AuditLabel{prediction}}

% ======================================================================
% 10D. SERIES-WIDE THEOREM-CHAIN AUDIT LABELS
% ======================================================================
% These labels mark the direction of load-bearing theorem chains. They are
% audit annotations only; they introduce no new mathematical assumptions.
\providecommand{\ChainInput}{}
\renewcommand{\ChainInput}{\AuditLabel{Chain input}}
\providecommand{\ChainOutput}{}
\renewcommand{\ChainOutput}{\AuditLabel{Chain output}}
\providecommand{\ChainGuard}{}
\renewcommand{\ChainGuard}{\AuditLabel{Chain guard}}
\providecommand{\ChainNoBackflow}{}
\renewcommand{\ChainNoBackflow}{\AuditLabel{No backflow}}

% ======================================================================
% 10E. SERIES-WIDE LOCAL THEOREM-AUDIT LABELS
% ======================================================================
% These labels mark a statement-level dependency audit around individual
% theorem/definition nodes. They are bookkeeping labels only.
\providecommand{\LocalAudit}{}
\renewcommand{\LocalAudit}{\AuditLabel{Local audit}}
\providecommand{\UsesOnly}{}
\renewcommand{\UsesOnly}{\AuditLabel{Uses only}}
\providecommand{\Exports}{}
\renewcommand{\Exports}{\AuditLabel{Exports}}
\providecommand{\DoesNotUse}{}
\renewcommand{\DoesNotUse}{\AuditLabel{Does not use}}

% ======================================================================
% 10F. SERIES-WIDE NUMERICAL-ROW AUDIT LABELS
% ======================================================================
% These labels classify individual rows in numerical, benchmark, status,
% and comparison tables.  They are audit annotations only and do not change
% any theorem, value, formula, or comparison number.
\providecommand{\RowAudit}{}
\renewcommand{\RowAudit}{\AuditLabel{Numerical-row audit}}
\providecommand{\RowTheorem}{}
\renewcommand{\RowTheorem}{\AuditLabel{Theorem-fixed row}}
\providecommand{\RowBridge}{}
\renewcommand{\RowBridge}{\AuditLabel{Bridge-fixed row}}
\providecommand{\RowRepresentative}{}
\renewcommand{\RowRepresentative}{\AuditLabel{Representative row}}
\providecommand{\RowComparison}{}
\renewcommand{\RowComparison}{\AuditLabel{Comparison row}}
\providecommand{\RowCalibrated}{}
\renewcommand{\RowCalibrated}{\AuditLabel{Calibrated row}}
\providecommand{\RowEstimate}{}
\renewcommand{\RowEstimate}{\AuditLabel{Estimate row}}
\providecommand{\RowOpen}{}
\renewcommand{\RowOpen}{\AuditLabel{Open row}}
\providecommand{\RowNoSelector}{}
\renewcommand{\RowNoSelector}{\AuditLabel{No selector}}

% ======================================================================
% 10G. TABLE HELPERS
% ======================================================================
% Table helpers used by Paper V-C.  The unused readout-label preset has been
% removed; only active table helpers remain.
\providecommand{\gpsection}[1]{%
  \addlinespace[0.55em]
  \multicolumn{3}{@{}l}{\textsc{#1}}\\[-0.25em]
  \midrule
}
\providecommand{\mathcell}[1]{%
  \(\begin{gathered}#1\end{gathered}\)%
}

\newcommand{\DerBy}[1]{%
  \begin{tabular}[t]{@{}l@{}}#1\end{tabular}%
}

\newcommand{\TableNotes}[1]{%
  \par\medskip
  \begin{minipage}{0.94\linewidth}
  \footnotesize\emph{Notes.} #1
  \end{minipage}\par
}
% ======================================================================
% 10H. PAPER VI-A / QCD NOTATION
% ======================================================================
% Shared symbols used by Paper VI-A and downstream QCD/Yang--Mills handoffs.
\providecommand{\Pqcd}{\mathfrak P_{\rm VIII}^{\rm QCD}}
\providecommand{\Iqcd}{\mathfrak I_{\rm QCD}^{(7)}}
\providecommand{\muQCD}{\mu_{\rm QCD}^{\rm int}}
\providecommand{\LambdaQCD}{\Lambda_{\rm QCD}}
\providecommand{\LambdaBH}{\Lambda_{\rm BH}^{\rm QCD}}
\providecommand{\HQCD}{\mathcal H_{\rm QCD}^{\rm BH}}
\providecommand{\HQCDfin}{H_{\rm QCD}^{\rm fin}}
\providecommand{\Hraw}{\mathcal H_{\rm raw}}
\providecommand{\HphysQCD}{\mathcal H_{\rm phys}^{\rm QCD}}
\providecommand{\Vcol}{V_C}
\providecommand{\Ps}{P_{\rm sing}}
\providecommand{\PiOct}{\Pi_8}
\providecommand{\MSbar}{\overline{\rm MS}}
\providecommand{\YM}{\mathrm{YM}}
\providecommand{\Ktri}{K_{\triangle}}
\providecommand{\KYM}{K_{\YM}^{\rm BH}}
\providecommand{\aBH}{a_{\rm BH}}
\providecommand{\epsBH}{\varepsilon_{\rm BH}}
\newcommand{\RowConstructive}{\textnormal{\textsc{Constructive}}}
\newcommand{\RowObligation}{\textnormal{\textsc{Bridge obligation}}}
\newcommand{\RowTarget}{\textnormal{\textsc{Target row}}}
\newcommand{\RowValidation}{\textnormal{\textsc{Validation row}}}
\providecommand{\slashed}[1]{\not{#1}}

% ======================================================================
% 10I. PAPER VI-B / CONTINUOUS-QFT NOTATION
% ======================================================================
% Shared symbols used by Paper VI-B and any downstream continuous-QFT or
% Yang--Mills handoff paper.  All definitions are guarded so individual
% papers may override them locally before loading this file if needed.
\providecommand{\A}{\mathfrak A}
\providecommand{\Hphys}{\mathcal H_{\rm phys}}
\providecommand{\QCD}{\mathrm{QCD}}
\providecommand{\BH}{\mathrm{BH}}
\providecommand{\loc}{\mathrm{loc}}

\providecommand{\RBH}{\mathcal R_{\rm BH}}
\providecommand{\DG}{\mathfrak D_G}
\providecommand{\PW}{\mathcal{P\!W}} % Peter--Weyl spectral-symbol subspace
\providecommand{\Toe}{\operatorname{Toe}} % Toeplitz compression/representative
\providecommand{\Hol}{\operatorname{Hol}} % Holonomy

% ======================================================================
% 11. LATE-LOADED PACKAGES (Must be at the end of the preamble)
% ======================================================================
% Load late packages only if a paper has not already loaded them, avoiding
% option clashes in paper-specific preambles.
\makeatletter
\@ifpackageloaded{hyperref}{}{\usepackage{hyperref}}
\makeatother
% --- Series-wide disambiguations (Pass 3) ---
\providecommand{\diracdelta}{\delta_{\rm D}} % Dirac delta distribution (delta reserved for hierarchy)
\providecommand{\Hg}{H_{\mathrm g}} % Hodge harmonic projector (keep H for Hubble)
\providecommand{\Gr}{G_{\mathrm r}} % Green operator (keep G for Newton)
\providecommand{\FD}{\mathrm{D}} % functional/variational derivative symbol (delta reserved)
\providecommand{\Id}{\mathbf{1}} % identity/Kronecker tensor (avoid \delta_{ij})


% BEGIN INLINE SOURCE: Paper00-ActiveExternalReferences-v2_130.tex
% Framework v2.191: current reference binding; existing definitions retained.
% Framework v2.190: current22-root registry with native clock-evolution proofs.
% Framework v2.189: exact current22-root external registry, including retained v2.188 hardening owners.
% Framework v2.187: current registry synchronization only; source release v2.186 preserved.
% Framework v2.186: authoritative current 22-paper registry.
% Framework v2.182: all current 22 logical paper AUX bindings.
% Framework v2.170: current I-B v1.43 external AUX binding.
% Framework v2.167: exact current 22-root registry after CP--BL2 integration.
% Framework v2.166: exact current 22-root compact-release AUX registry.
% Framework v2.165: authoritative current root bindings.
% Framework v2.164: exact current 22 manuscript AUX bindings.
% Framework v2.162: finite metric, functional Hilbert and scoped Ward integration.
% Framework v2.161: current 22 manuscript AUX bindings; source only.
% Framework v2.156 / registry v2.113: synchronized current 22 manuscript roots; source-only registry.
% Framework v2.143 / registry v2.111: active manuscript AUX basenames; no generated AUX or PDF included.
% ======================================================================
% Central external-reference imports, numbers only, no PDF jumping
% ======================================================================

% xr-hyper is loaded centrally so that papers may retain their existing
% \usepackage{xr-hyper} lines without redefining \externaldocument after this
% registry is installed.  If plain xr was loaded earlier, retain it.
\makeatletter
\@ifpackageloaded{xr-hyper}{}{%
  \@ifpackageloaded{xr}{}{\usepackage{xr-hyper}}%
}
\makeatother

% Make external-document registration idempotent.  The registry identifies a
% logical paper by its Paper<number><optional letter> prefix.  Prefix-qualified
% imports remain independent.  Current paper roots contain no local
% \externaldocument calls; this guard also makes accidental duplicates harmless.
\ExplSyntaxOn
\prop_new:N \g_bh_external_refs_prop
\cs_new_eq:NN \BHOriginalExternalDocument \externaldocument

\cs_new_protected:Npn \bh_external_document:nnnn #1#2#3#4
 {
  \tl_set:Nn \l_tmpa_tl {#3}
  \regex_replace_once:nnN
   { \A (?: xref/ )? (Paper[0-9]+[A-G]?) .* \Z }
   { \1 }
   \l_tmpa_tl
  \tl_set:Nx \l_tmpb_tl { #1 | \tl_use:N \l_tmpa_tl }
  \prop_if_in:NVTF \g_bh_external_refs_prop \l_tmpb_tl
   {
    \PackageInfo{BornHopf}
     {Skipping~duplicate~external~refs~for~#3}
   }
   {
    \prop_gput:NVn \g_bh_external_refs_prop \l_tmpb_tl {#3}
    \IfNoValueTF{#4}
     { \BHOriginalExternalDocument[#1][#2]{#3} }
     { \BHOriginalExternalDocument[#1][#2]{#3}[#4] }
   }
 }

\RenewDocumentCommand{\externaldocument}{O{} O{} m o}
 { \bh_external_document:nnnn {#1}{#2}{#3}{#4} }

\NewDocumentCommand{\BHExternalRefs}{m}
 {
  % Do not import the aux file of the logical paper currently being compiled,
  % even when the job and registry use different version suffixes.
  \tl_set:Nx \l_tmpa_tl { \c_sys_jobname_str }
  \regex_replace_once:nnN
   { \A (Paper[0-9]+[A-G]?) .* \Z }
   { \1 }
   \l_tmpa_tl
  \tl_set:Nn \l_tmpb_tl { #1 }
  \regex_replace_once:nnN
   { \A (Paper[0-9]+[A-G]?) .* \Z }
   { \1 }
   \l_tmpb_tl
  \str_if_eq:eeTF { \tl_use:N \l_tmpa_tl } { \tl_use:N \l_tmpb_tl }
   {
    \PackageInfo{BornHopf}
     {Skipping~external~refs~for~#1~because~this~is~the~current~document}
   }
   {
    \IfFileExists{xref/#1.aux}
     { \externaldocument[][nocite]{xref/#1}[#1.pdf] }
     {
      \IfFileExists{#1.aux}
       { \externaldocument[][nocite]{#1} }
       {
        \PackageInfo{BornHopf}
         {External~aux~file~#1.aux~not~found;~no~labels~imported}
       }
     }
   }
 }
\ExplSyntaxOff

% External references are automatic.  This file is the sole owner of both the
% bibliography and the current Paper-I-A--VII-D AUX registry (30 roots including I-C).  Missing auxiliaries
% are informational on the first pass and are populated by the source-only
% two-pass series build.  Basenames deliberately omit temporary download/copy
% suffixes such as ``(1)'' or ``-1''.
  % Paper I-A: round S3 compendium.
  \BHExternalRefs{Paper1A-Geometry-v5_181}

  % Paper I-B: complete exceptional geometry and finite spectral structures.
  \BHExternalRefs{Paper1B-ExceptionalGeometry-v1_27}
  % Paper I-D: technical companion of I-A through I-C.
  \BHExternalRefs{Paper1D-TechnicalMaterial-v1_69}
  % Paper I-C: reciprocal M8 compendium.
  \BHExternalRefs{Paper1C-M8Geometry-v1_146}

\ifdefined\BHMathOnly
  % Papers I-A and I-B import only the mathematical atlas labels.
\else
  % Paper II: five physical-mathematical main papers plus one technical proof volume.
  \BHExternalRefs{Paper2A-GeometricConstitution-v3_227}
  \BHExternalRefs{Paper2B-MomentumCarrier-v2_147}
  \BHExternalRefs{Paper2C-ReciprocalLocalization-v2_177}
  \BHExternalRefs{Paper2D-SymmetrySpectra-v2_145}
  % The former II-F export labels remain in II-E; current II-F is the technical volume.
  \BHExternalRefs{Paper2E-KinematicExports-v3_140}
  % Paper II-F: complete technical proof and calculation owner for II-A--II-E.
  \BHExternalRefs{Paper2F-TechnicalMaterial-v1_71}

  % Paper III: Dynamics and Intrinsic Geometric Constitution (six components).
  \BHExternalRefs{Paper3A-GeneratedDynamics-v1_181}
  \BHExternalRefs{Paper3B-SpectralAction-v1_192}
  \BHExternalRefs{Paper3C-ProtectedHierarchy-v1_196}
  \BHExternalRefs{Paper3D-RunningResponse-v1_168}
  \BHExternalRefs{Paper3E-GravityAction-v1_204}
  \BHExternalRefs{Paper3F-CompleteDynamics-v1_281}
  % Paper III-G: complete preserved technical proof owner.
  \BHExternalRefs{Paper3G-TechnicalMaterial-v1_58}

  % Paper IV-A: generated finite quantum mechanics.
  \BHExternalRefs{Paper4A-FiniteQuantumMechanics-v4_217}
  % Paper IV-B: complete finite quantum processes and physical response.
  \BHExternalRefs{Paper4B-QuantumResponse-v1_55}
  % Paper IV-C: extended proofs and finite technical calculations.
  \BHExternalRefs{Paper4C-TechnicalMaterial-v1_50}

  % Paper V-A: The Standard Model.
  \BHExternalRefs{Paper5A-FiniteMatterGeometry-v3_219}

  % Paper V-B: Finite Spectrum and Flavor Closure.
  \BHExternalRefs{Paper5B-TerminalSpectrum-v5_202}

  % Paper V-C: Finite Interaction Boundaries of the Standard Model.
  \BHExternalRefs{Paper5C-InteractionResponse-v2_185}
  % Paper V-D: preserved technical proof and calculation owner.
  \BHExternalRefs{Paper5D-TechnicalMaterial-v1_58}

  % Paper VI-A: Gravity and Closed Cosmology.
  \BHExternalRefs{Paper6A-GravityCosmology-v2_268}

  % Paper VI-B: Geometric Dark Sector and Cosmological Readouts.
  \BHExternalRefs{Paper6B-CosmologyReadouts-v1_215}

  % Paper VI-C: complete technical proof and comparison owner.
  \BHExternalRefs{Paper6C-TechnicalMaterial-v1_53}

  % Paper VII-A: Finite QCD Kernel and Nonperturbative Closure.
\BHExternalRefs{Paper7A-QCDKernel-v1_185}

  % Paper VII-B: Hadronic, Thermal, and Nuclear QCD Readouts.
  \BHExternalRefs{Paper7B-QCDReadouts-v1_169}

  % Paper VII-C: Continuous Quantum Field Theory.
  \BHExternalRefs{Paper7C-ContinuousQFT-v2_208}

  % Paper VII-D: complete technical material for VII-A through VII-C.
  \BHExternalRefs{Paper7D-TechnicalMaterial-v1_60}

\fi

% END INLINE SOURCE: Paper00-ActiveExternalReferences-v2_130.tex



\providecommand{\LambdacomparisonBH}{\Lambda_{\rm BH}^{\rm comp}}
\providecommand{\deltageom}{\delta_{\rm geom}}
\providecommand{\alphasseed}{\alpha_s^{\rm seed}}


% Paper-specific layout refinements.
\setlength{\tabcolsep}{2pt}
\AtBeginEnvironment{longtable}{\footnotesize}
\ifcsname c@biburlnumpenalty\endcsname\setcounter{biburlnumpenalty}{9000}\fi
\ifcsname c@biburllcpenalty\endcsname\setcounter{biburllcpenalty}{9000}\fi
\ifcsname c@biburlucpenalty\endcsname\setcounter{biburlucpenalty}{9000}\fi


\title{\textbf{Born--Hopf Geometry in a Closed Universe VII-B: Hadronic and Strong-Interaction Readouts}\\[0.4em]
\large Physical spectral, current, light-ray, thermal and nuclear observables}
\author{Peter Kugelmann\thanks{\texttt{pkugelmann@proton.me}}}
\date{Version 1.169 --- 20 November 2026 \\ Framework v3.077}
\setcounter{tocdepth}{2}

% Current release citation-only xref imports (no material DAG dependency).
\BHExternalRefs{Paper1A-Geometry-v5_175}
\BHExternalRefs{Paper7A-QCDKernel-v1_185}
% Cross-DAG current-series import for the new source-owned mathematical bridge.
\BHExternalRefs{Paper1B-ExceptionalGeometry-v1_27}
\BHExternalRefs{Paper2C-ReciprocalLocalization-v2_170}
\BHExternalRefs{Paper3C-ProtectedHierarchy-v1_190}

\begin{document}
% Publication typesetting: controlled prose line-breaking, without scaling mathematics.
\emergencystretch=2em
\maketitle
\begin{abstract}
Paper VII-B owns finite hadron spectra and pole/residue calculations,
 current matrix elements, native physical null and parton distributions,
 thermal and nuclear response, and source/detector contracts. Color carrier
 arithmetic or Hurwitz support geometry does not itself determine a hadron
 mass, binding energy, form factor or cross section. The derived split grammar
 remains in the main series; the earlier $E_{36}$/$E_{40}$ phenomenological
 numeric comparisons are preserved in Paper VII-D as open operator-derivation
 targets, not discarded or treated as completed predictions. The continuum
 and scattering conditions are referred to Paper VII-C.
\end{abstract}
\tableofcontents
\listoftables
\clearpage

\section{Inherited QCD operators, mass inputs, and readout contracts}
\label{p11:sec:v237-b1}
Hadronic observables consume the supplied strong Hamiltonian, already generated quark masses, physical states and source/detector maps; Paper~VII-B does not select those inputs from measured spectra.

% Source Paper VII-B, original line 101; inherited proof bodies unchanged.
% Former logical part: Paper-VII-A Interface and Paper-VII-B Readout Scope; integrated into the new thematic argument.
\label{p11:part:paper-x-interface}


% Source Paper VII-B, original line 104; inherited proof bodies unchanged.
\subsection{Imported finite QCD core}
\label{p11:sec:imported-finite-qcd-core}

\begin{definition}[Imported repaired Paper-VII-A core datum]
\label{p11:def:imported-paper-x-core-datum}
Paper~VII-B imports exactly the repaired Paper~VII-A datum
\[
\mathfrak D_{7A}^{\rm QCD,core}
=\left(\begin{gathered}
C,SU(3)_c,\Ps,\mathcal K_H,\Delta_{2,H},\\
\{H_L,\widehat K_L(t),\mathfrak A_{L,{\rm str}},
\mathfrak A'_{L,{\rm str}},\mathcal O_L,W_L[J]\}_L,\\
\mathfrak T_{\rm chiral},\mathfrak T_{\rm top}
\end{gathered}\right)
\]
of Definition~\ref{p10:def:finite-qcd-core-datum}, with the acyclic handoff of Theorem~\ref{p10:thm:paper-x-xi-boundary}.  The source-native common-heat gauge normalization is inherited at its native matching datum.  A named low-scale \(\alphasseed\) remains a scale/RG comparison coordinate until matching is supplied.  The charged finite mass/source member is already inherited and \(M_q^{\rm phys}(\mu)\) consumes its applicable G16 transport.  The topological \(\theta\) is an auxiliary Fourier source; Paper~V-C fixes the native finite sector character and rooted determinant orientation, hence its finite \(\bar\theta=0\) invariant. Regulated physical transport is retained only by claims that consume it.  A finite nonlinear action is regulator/presentation data rather than a physical continuum selector.  Paper~VII-C supplies the bounded phase-cluster state and the ordinary local-QFT
export conditional on its actual reconstruction certificate, without a full
action-universality theorem; such a theorem is retained only for representative-independent or canonical comparison claims.  No fixed \(Q_{\rm wall}^{\rm BH}\), fixed nonperturbative scale \(\Lambda_{\rm BH}\), count-derived shell amplitude, positive continuum gap, physical string tension, \(F_\pi\), \(B_0\), or topological-susceptibility decimal is part of the unconditional import.
\end{definition}

\begin{corollary}[Finite strong readout maps are generated descendants]
\label{p11:cor:v147-generated-readout-descendants}
For the generated finite strong owner imported from Paper~VII-A, every hadronic,
current, distribution, thermal, topological, and multi-baryon readout map that
Paper~VII-B constructs by its licensed finite spectral/source operations is a
generated structural/readout descendant.  No second premise is required to
assert physicality of the map itself.  States or produced histories remain inputs where the readout consumes them.
The charged finite mass/source packet is inherited as an already generated
output; numerical execution consumes its applicable G16 transport.
Hadronization/process certificates, a selected unscreened area-law target,
exact light-front representation only for its optional fraction realization,
and stronger continuum comparison/universality hypotheses remain separately typed; bounded phase-cluster existence and the conditional ordinary local-QFT
reconstruction certificate are inherited from Paper~VII-C.  In particular, a
positive finite carrier gap is not thereby promoted to the continuous
Yang--Mills mass-gap theorem.
\end{corollary}
\begin{proof}
The imported owner is the finite structural datum of
Definition~\ref{p11:def:imported-paper-x-core-datum}.  Applying the Paper-II-A
structural-descendant rule to Paper~VII-B's licensed spectral and source
constructors co-realizes their output maps.  The excluded items are member,
state/history, calibration, representation-certificate, or analytic-limit
claims, and therefore lie outside that finite structural induction.
\end{proof}


% BEGIN INLINE SOURCE: Paper7B-SpectralPopulationScope-Main-v1_0.tex
% Framework v2.187: positive finite spectral extraction and its state scope.
\begin{corollary}[Finite spectra precede population readouts]
\label{p11:cor:v2187-spectrum-before-population}
Retain an actual finite self-adjoint singlet-sector Hamiltonian
$H_{\chi,L}$ from the complete QCD source-transfer packet.  Its spectral
values and projectors,
\[
 H_{\chi,L}=\sum_{\lambda}\lambda P_{\lambda,L},
\]
are licensed finite readouts without a produced population or preparation
history.  A state is consumed by a probability
$\operatorname{Tr}(\rho_LP_{\lambda,L})$, a state-weighted source response,
or a named production/abundance observable.  A named physical particle or
resonance additionally retains the sector dictionary, isolation and only the
pole/refinement/channel certificates its interpretation uses.

These finite spectral maps do not consume positive asymptotic string tension,
an ordered massless condensate, a positive continuum sector gap, or exact
finite species/Fock light-front correspondence.  The exact Goldstone/GMOR,
unscreened Wilson-area, distribution and asymptotic process claims retain
their separately printed hypotheses.
\end{corollary}
\begin{proof}
Finite functional calculus depends on $H_{\chi,L}$ alone.  Changing a density
matrix changes the Born weights but not this operator's eigenvalues or
projectors.  The inherited source packet fixes its normalized source maps;
physical identification, state/history and limiting constructions enter only
the later readout that consumes them.
\end{proof}

% END INLINE SOURCE: Paper7B-SpectralPopulationScope-Main-v1_0.tex


\begin{remark}[Assumption-minimized strong-interface scope]\label{p11:rem:conditional-strong-interface-scope}
Paper~VII-A proves the exact ancestry certificate
\[
 \deltageom=\omega_{2O}=\frac1{6\sqrt2},
 \qquad
 \deltageom^{-2}=\operatorname{rank}(E_6)h(E_6)=6\cdot12=8\cdot9=72.
\]
This common scalar is not a QCD transition probability and is not the
source-native gauge-normalization owner.  The common-heat branch fixes the
absolute gauge matching coefficient at its own scale; a low-scale
\(\alphasseed\) is obtained only after scale/RG transport.  The charged finite quark-mass/source operator is inherited; its named-scale transport remains typed.  The auxiliary topological Fourier source is distinct from the fixed native trivial sector character and rooted finite strong-CP invariant; no new physical \(\bar\theta_{\rm QCD}\) member is selected here.  The exact finite action formula is regulator data.  Paper~VII-C v2.83 requires no full universality test for the base phase-selected bounded-local export; action universality is tested only when representative-independence or canonical same-phase comparison is claimed.  Paper~VII-B therefore reports dimensionless geometric spectra
unconditionally and reports dimensionful QCD observables only as eigenvalues,
correlator residues, response derivatives, or continuum limits of the supplied
finite Hamiltonian family.  Comparison data never select these upstream
members.
\end{remark}

\begin{center}
\fbox{\parbox{0.92\textwidth}{
\textbf{Paper-VII-B theorem scope.}  Every physical number must be obtained from a specified finite operator/source member; a state, preparation history or limiting certificate is retained only where that readout consumes it.  Finite counts may classify representations and exact Hodge spectra, but may not be promoted to amplitudes, current residues, topology weights, or physical length/energy scales without an operator theorem.  When the required constructor is absent, Paper~VII-B records a typed interface obligation instead of inserting a selector.
}}
\end{center}

\begin{definition}[Assumption-minimized finite readout grammar]
Let \(P_{\chi,L}\ne0\) be a spectral or conserved-sector orthogonal projector commuting with the supplied self-adjoint \(H_L\), and let \(E_{0,L}\) be the full finite vacuum energy. All compressed operators below are understood \emph{on} \(\operatorname{Ran}P_{\chi,L}\), never on the ambient carrier with its artificial zero complement. Define
\begin{align}
\Delta E_{\chi,L}&:=\inf\operatorname{spec}\bigl((P_{\chi,L}H_LP_{\chi,L})|_{\operatorname{Ran}P_{\chi,L}}\bigr)-E_{0,L},\\
\widehat K_{\chi,L}(t)&:=P_{\chi,L}e^{-t(H_L-E_{0,L})}P_{\chi,L},\\
\epsilon_{\chi,L}(t)&:=1-e^{-t\Delta E_{\chi,L}}.
\end{align}
If \(\mathcal O\) is a supplied source operator, its pole residue is
\[
Z_{\mathcal O,n,L}:=|\langle\Omega_L|\mathcal O|n,L\rangle|^2.
\]
If \(J\) is a current whose conservation has been established on the stated finite action/state branch, the corresponding finite Ward identity fixes its conserved charge at zero momentum; slopes, magnetic moments, axial/tensor charges, and other transverse form factors remain the corresponding matrix elements of \(J\).  Topology is read from \(Z_L(\theta)\) and the computed charge distribution/covariance, while chiral constants are read from the double-limit order parameter and current correlators exported as \(\mathfrak T_{\rm chiral}\).  These are deterministic readouts once the owning source/operator member is supplied; they are not new phenomenological coefficients introduced by Paper~VII-B.
\end{definition}


% Source Paper VII-B, original line 5300; inherited proof bodies unchanged.
\subsection{inherited chiral spurion and selector retirement}
\label{p11:sec:v154-inherited-chiral-spurion}

\begin{corollary}[Kaplan--Manohar freedom is not a second physical mass selector]
\label{p11:cor:v154-km-selector-retired}
On the native branch in which an actual nonzero G0 quark-source member occurs
and the G16 RG/threshold transport to \(\mu\) is certified, the microscopic
spurion \(M_q^{\rm phys}(\mu)\) is inherited from the injective upstream
source-owner map.  The Kaplan--Manohar transformation therefore acts only as a
reparametrization of the truncated chiral description and supplies no further
physical member choice.
\end{corollary}

The remaining chiral/readout inputs are claim-specific: controlled-regime
occurrence, low-energy constants and correlator limits, electromagnetic or
anomaly inputs when consumed, and the accuracy/order of the effective
expansion.  None is a replacement quark-mass selector.
% END TRANSITIVE-CHIRAL-SPURION v1.54

% BEGIN SAFEGUARD-QCD-MASS-SPURION v1.55

% Source Paper VII-B, original line 5320; inherited proof bodies unchanged.
\subsection{synchronized QCD mass-spurion boundary}
\label{p11:sec:v155-safeguard-qcd-mass-spurion}
All QCD/chiral mass-dependent constructions in this paper consume the inherited
running microscopic spurion \(M_q^{\rm phys}(\mu)\).  The useful chiral and
Kaplan--Manohar algebra is retained, but it cannot reselect the physical
microscopic member.  Hence Paper VII-B has no separate G5Q mass selector and no
new QCD mass-constructor debt.
% END SAFEGUARD-QCD-MASS-SPURION v1.55

% BEGIN CHARGED-SOURCE-READOUT-SYNC v1.57

% Source Paper VII-B, original line 5330; inherited proof bodies unchanged.
\subsection{charged source/readout synchronization}
\label{p11:sec:v157-charged-source-sync}
The charged noncentral finite source member is now generated upstream before
G16 transport.  All quark-mass-dependent readouts in this paper therefore
consume that member after the claimed running/threshold handoff; none retains
an independent same-action family-assembly occurrence premise.

% END CHARGED-SOURCE-READOUT-SYNC v1.57


% BEGIN CHARGED-COMPARISON-READOUT-SYNC v1.58

% Source Paper VII-B, original line 5341; inherited proof bodies unchanged.
\subsection{charged comparison/readout synchronization}
\label{p11:sec:v158-charged-comparison-sync}
The inherited microscopic quark source now carries the derived finite charged
comparison coefficient of Paper~V-B v5.91.  Chiral and QCD readouts consume
that coefficient after the existing G16/RG/threshold handoff.  They do not
supply a second $24/25$ factor, an exponent selector, or a replacement quark
mass source.  Named-scheme values and pole quantities retain their prior
claim-scoped evaluation status.
% END CHARGED-COMPARISON-READOUT-SYNC v1.58




% BEGIN GEOMETRY-FIRST CONSEQUENCES v2.201
\setlength{\emergencystretch}{2em}


\section{Color-singlet hadron supports and channel classification}
\label{p11:sec:v237-b2}
Exact Hurwitz and color-singlet support decompositions specify allowed finite channels, not numerical transition strengths or dimensionful hadron masses.

% Source Paper VII-B, original line 221; inherited proof bodies unchanged.
% Former logical part: Hadron, Distribution, Thermal, and Nuclear Readout Packages; integrated into the new thematic argument.
\label{p11:part:hadron-distribution-thermal-nuclear-readouts}


% Source Paper VII-B, original line 224; inherited proof bodies unchanged.
\subsection{Primitive singlet-channel geometry and supports}
\label{p11:sec:hadron-spectrum}


% Source Paper VII-B, original line 227; inherited proof bodies unchanged.
\subsubsection{Hadron support classification and primitive channel hierarchy}
\label{p11:subsec:hadron-support-primitive-channel-hierarchy}


\begin{table}[htbp]
\centering
\footnotesize
\begin{tabular}{>{\raggedright\arraybackslash}p{0.34\textwidth} >{\raggedright\arraybackslash}p{0.56\textwidth}}
\toprule
Claim & Discharged by / status \\
\midrule
Finite singlet hadron states & constructed inside the physical finite QCD singlet space \\
Hadron masses as finite spectral gaps & read from the native Hamiltonian/transfer spectrum \\
Baryon and magnetic-moment residues & computed as finite current and residue readouts \\
PDFs, DAs, and fragmentation functions & represented by finite matrix elements and handoff operators \\
Validation-only spectrum comparison & restricted to comparison after native computation \\
\bottomrule
\end{tabular}
\caption[Hadron-spectrum and distribution readouts]{Hadron-spectrum and distribution readouts.}
\label{p11:tab:hadron-spectrum-ledger}
\TableNotes{This ledger separates the native finite hadron/readout problem from downstream spectroscopy and global-fit comparison.  Masses, residues, and distributions enter as finite spectral or matrix-element outputs, not fitted inputs.}
\end{table}

The allowed hadron support classes are derivable from \(C\cong\CC^3\) and the singlet projector.  Their masses are not new parameters: they are the singlet eigenvalue differences of the native nonperturbative Hamiltonian.

\begin{theorem}[Singlet support classification]
\label{p11:thm:singlet-support-classification}
Given \(C\cong\CC^3\) and \(\Ps\), the primitive color-singlet hadronic supports are exactly the singlet subspaces generated by the following tensor classes:
\begin{align}
\label{p11:eq:meson-support-class}
\text{mesons:}\qquad & C\otimes C^\ast\cong\mathbf1\oplus\mathbf8,\\
\label{p11:eq:baryon-support-class}
\text{baryons:}\qquad & \wedge^3C\cong\mathbf1,\\
\label{p11:eq:antibaryon-support-class}
\text{antibaryons:}\qquad & \wedge^3C^\ast\cong\mathbf1,\\
\label{p11:eq:glueball-support-class}
\text{glueballs:}\qquad & \Ps(\mathbf8^{\otimes n}),\quad n\ge2,\\
\label{p11:eq:hybrid-support-class}
\text{hybrids:}\qquad & \Ps((C\otimes C^\ast)_{\mathbf8}\otimes\mathbf8),
\end{align}
together with higher singlet tensor powers such as tetraquark and pentaquark supports.  No new color carrier is introduced for exotic hadrons.
\end{theorem}
\begin{proof}
A quark transforms in \(C\), an antiquark in \(C^\ast\), and a gluon in the adjoint \(\mathbf8\).  The trace contraction gives the meson singlet in \(C\otimes C^\ast\).  Since \(\dim C=3\), the invariant determinant tensor gives \(\wedge^3C\cong\mathbf1\), forcing minimal quark-only singlets to contain three quarks.  A single adjoint is not a singlet, while adjoint tensor powers contain singlet contractions.  Hybrids require the octet part of \(C\otimes C^\ast\) to combine with the gluon octet to a singlet.  Higher exotics are admissible exactly when their color tensor product contains \(\mathbf1\).\qedhere
\end{proof}


\begin{theorem}[Hadronic fermion parity]
\label{p11:thm:hadronic-fermion-parity}
Every physical color-singlet QCD sector obeys
\begin{equation}
\label{p11:eq:fermion-parity-baryon-parity}
(-1)^F=(-1)^B.
\end{equation}
Thus sectors with even baryon number are bosonic.  Sectors with odd baryon
number are fermionic.
\end{theorem}
\begin{proof}
Quarks and antiquarks are fermions, while gluons are bosons.  Hence a raw tensor sector has fermion parity \((-1)^{n_q+n_{\bar q}}\).  Modulo two,
\begin{equation}
 n_q+n_{\bar q}\equiv n_q-n_{\bar q}\pmod2,
\end{equation}
because their difference is \(2n_{\bar q}\).  On physical sectors, Corollary~\ref{p10:cor:integer-baryon-number} gives \(n_q-n_{\bar q}=3B\).  Therefore \((-1)^F=(-1)^{3B}=(-1)^B\).\qedhere
\end{proof}

\begin{corollary}[Minimal baryon exchange symmetry]
\label{p11:cor:minimal-baryon-exchange-symmetry}
For the minimal baryon support \(\wedge^3C\cong\mathbf1\), the color wavefunction is antisymmetric.  Since quarks are fermions, the remaining spin/flavor/orbital wavefunction is symmetric under exchange of identical quarks, modulo the usual restriction to actually identical quark labels.
\end{corollary}
\begin{proof}
The unique color-singlet contraction of three fundamental colors is
\(\epsilon_{abc}q^aq^bq^c\), i.e. the exterior product \(\wedge^3C\).  This
factor changes sign under exchange of two identical quark color labels.  The
full fermionic state must also change sign under such an exchange, so the
non-color factor must be symmetric for identical quarks.  Nonidentical flavor
labels are not subject to an exchange-symmetry constraint between distinct
particles.
\end{proof}


\begin{theorem}[QCD use of imported color-Casimir attraction channels]
\label{p11:thm:color-casimir-attraction-channels}
Paper~V-C, the proposition titled \emph{Primitive color attraction channels}, fixes the
primitive \(SU(3)_c\) quark color factors.  In Paper~VII-B notation,
\begin{equation}
\label{p11:eq:casimir-attractive-channels}
(q\bar q)_{\mathbf1}:-\frac43,
\qquad
(q\bar q)_{\mathbf8}:+\frac16,
\qquad
(qq)_{\bar{\mathbf3}}:-\frac23,
\qquad
(qq)_{\mathbf6}:+\frac13,
\qquad
(gg)_{\mathbf1}:-3 .
\end{equation}
Thus meson singlets, antitriplet diquark subchannels, and pure-glue singlets
are the primitive attractive color channels used by the finite-QCD readout closure.
\end{theorem}

\begin{proof}
The quark entries are imported from Paper~V-C.  The pure-glue singlet entry is
the same Casimir identity
\[
T_1\cdot T_2=\frac12(C_R-C_{R_1}-C_{R_2})
\]
with \(R=\mathbf1\), \(R_1=R_2=\mathbf8\), and \(C_A=3\), giving
\(T_1\cdot T_2=-3\).  Negative values are attractive for one-gluon exchange.
\end{proof}


\begin{theorem}[Computed channel response replaces Casimir loss]
\label{p11:thm:primitive-channel-loss-hierarchy}
\label{p11:thm:primitive-strong-amplitude-loss-law}

The color factors of Theorem~\ref{p11:thm:color-casimir-attraction-channels} classify representation contractions but do not determine a transfer amplitude.  For every invariant channel projector \(P_{\mathcal C,L}\), the rigorous finite response is
\[
\Delta E_{\mathcal C,L}=\inf\operatorname{spec}(P_{\mathcal C,L}H_LP_{\mathcal C,L})-E_{0,L},
\qquad
\epsilon_{\mathcal C,L}(t)=1-e^{-t\Delta E_{\mathcal C,L}}.
\]
No universal relation \(\epsilon_{\mathcal C}=C_{\mathcal C}^2\deltageom\) and no carrier-count cutoff follow from the current corpus.
\phantomsection\label{p11:eq:primitive-channel-loss-rule}
\phantomsection\label{p11:eq:primitive-channel-weights}
\phantomsection\label{p11:eq:primitive-channel-amplitude-ratio}
\phantomsection\label{p11:eq:primitive-channel-loss-ratio}
\phantomsection\label{p11:eq:primitive-n-carrier-loss}
\end{theorem}
\begin{proof}

For any fixed color representation, multiplying an allowed channel interaction by a positive scalar leaves the representation and Casimir unchanged while changing the spectral gap and transfer eigenvalue.  Hence the Casimir cannot determine the response coefficient.  The displayed quantities are instead exact spectral functionals of the supplied self-adjoint Hamiltonian.
\end{proof}

\begin{corollary}[No universal count-derived hadron floor]
\label{p11:cor:numerical-primitive-channel-floors}

The finite spectral theorem gives only \(\Delta E_{\mathcal C,L}\ge0\).  A strictly positive channel floor is obtained exactly when the corresponding restricted spectral gap is positive; its value must be computed from \(H_L\).  The predecessor decimals based on a fixed wall scale and \(C^2\delta\) are withdrawn.

\end{corollary}
\begin{proof}

Positivity follows from subtracting the ground energy.  A counterfamily \(H_L(c)=E_{0,L}I+cP_{\mathcal C,L}\), \(c>0\), has identical finite carrier and projector data but arbitrary channel gap \(c\).  Therefore no nonzero universal numerical floor is implied by the geometry alone.
\end{proof}



% BEGIN FRAMEWORK v2.323: PETERSEN HADRON SUPPORT AND SPECTRAL READOUT BOUNDARY
\section{Petersen multi-core visibility and hadronic spectral tests}
\label{p11:sec:v323-petersen-hadron-test}
\noindent\textit{Status: exact finite geometric spectra; the map to actual
color-singlet constituent supports and the physical poles remains open.}
For an occurring embedded $2I$, use the Petersen edge space
$W=\bigoplus_{m=1}^5V_m$ and the self-adjoint mirror contrast
$D=\Pi_3-\Pi_{3'}$ of Paper~I-A and Paper~VII-A.  The selected-core
projector $P_S$ has rank $3|S|$ and $D_S=P_SDP_S$.

\begin{theorem}[Exact multi-core mirror-compression spectrum]
\label{p11:thm:v323-petersen-k-core-spectrum}
The spectrum depends only on $k=|S|$ and is
\[
\begin{array}{c|c|l}
 k&\dim P_SW&\operatorname{Spec}(D_S)\\\hline
 0&0&\varnothing\\
 1&3&0^{(3)}\\
 2&6&\bigl(\pm\sqrt{15}/10\bigr)^{(2)},\ 0^{(2)}\\
 3&9&\pm\sqrt5/5,\ \bigl(\pm\sqrt{35}/10\bigr)^{(2)},\ 0^{(3)}\\
 4&12&\bigl(\pm\sqrt{15}/5\bigr)^{(3)},\ 0^{(6)}\\
 5&15&(+1)^{(3)},\ (-1)^{(3)},\ 0^{(9)}.
\end{array}
\]
For $\pm a$ entries, the displayed multiplicity belongs to \emph{each}
sign.  All $\binom5k$ supports of a fixed size are unitarily
conjugate under the native signed $A_5$ action.
\end{theorem}
\begin{proof}
Use the exact orientation-independent rank-three projectors
$\Pi_3=K^{\mathsf T}K/5$ and
$\Pi_{3'}=(K^\sigma)^{\mathsf T}K^\sigma/5$ on the fifteen oriented
Petersen edges.  Substitution into $D_S$ for one representative of
each size gives the characteristic polynomials written in
Paper~VII-D.  The icosahedral $A_5$ action is transitive on the
$k$-subsets of five labels and preserves $D$, hence preserves all
other spectra.  The exact certificate enumerates all 60 even
permutations and checks the required signed conjugations.
\end{proof}

\begin{corollary}[Two-core and complementary three-core supports are
not spectrally identical for this observable]
\label{p11:cor:v323-2-3-contrast}
For two selected cores, the mirror-contrast operator has extreme
magnitudes $\sqrt{15}/10$; its complementary three-core support has
extreme magnitude $\sqrt{35}/10$.  These are distinct geometric
source-visibility bounds, even though the separately constructed
ordinary-to-conjugate maps obey
$J_{S^c}=-J_S$ and have identical quadratic polar metrics.
No such distinction exists in one isolated local core.  None of the
numbers is an energy or a meson/baryon mass ratio.
\end{corollary}
\begin{proof}
Take the operator norms from the preceding spectra.  The
complement identity concerns the different transfer map $J_S$,
not the compression $P_SDP_S$, so the conclusions are compatible.
\end{proof}

\begin{proposition}[Conditional physical hadron readout criterion]
\label{p11:prop:v323-actual-hadron-spectrum-criterion}
The geometric contrast yields a hadronic observable only after a
source-owned identification of an actual color-singlet QCD subspace
with the selected core channels, its actual Hamiltonian, and an
interpolator with nonzero matrix elements on the relevant poles.
The finite hadron masses remain the eigenvalue gaps or same-channel
transfer poles of the supplied $H_L$ from Paper~VII-A and this paper.
If a proposed Petersen subspace is nonreducing under $H_L$, its
compression spectrum cannot be substituted for physical pole data.
In particular the labels $k=2$ and $k=3$ are not quark number or
baryon number; neither the meson assignment nor the baryon assignment
follows from this theorem.
\end{proposition}
\begin{proof}
The proof is the finite spectral/resolvent definition already used by
the same-channel transfer construction: the pole positions and
residues are functions of the \emph{actual} Hamiltonian, vacuum and
source.  The geometric projector $P_S$ supplies none of those maps.
\end{proof}
% END FRAMEWORK v2.323: PETERSEN HADRON SUPPORT AND SPECTRAL READOUT BOUNDARY



% BEGIN FRAMEWORK v2.327: MESON SQUARED-GAP OBSERVABLE AND GRAM COHERENCE
\section{Meson spin-pair diagnostics without postselected mass triplets}
\label{p11:sec:v327-meson-mirror-readouts}
\noindent\textit{Status: exact conditional algebra and measured targets,
not a predicted mass spectrum.  Known heavy-quark spin symmetry explains
an approximate part of the observed pattern; it is not claimed here
as a new empirical discovery.}

\begin{proposition}[Gram rule for orthogonal versus coherent supports]
\label{p11:prop:v327-meson-gram-support}
For any genuine source amplitudes $u_i$ with
$\|u_i\|^2=n_i$ and real normalized overlaps
$\gamma_{ij}=\operatorname{Re}\langle u_i,u_j\rangle/
\sqrt{n_in_j}$, one has
\[
 \left\|\sum_i u_i\right\|^2
  =\sum_i n_i+2\sum_{i<j}\sqrt{n_in_j}\,\gamma_{ij}.
\]
The overlap matrix must be positive semidefinite; its entries
cannot be selected independently. For three equal supports of size
$40$ with common real overlap $\gamma$, positivity requires
$-\tfrac12\le\gamma\le1$, and
\[
 \operatorname{Spec}(G)=
 \{40(1+2\gamma),40(1-\gamma)^{(2)}\},\qquad
 \operatorname{Tr}G=120,
 \qquad \left\|u_1+u_2+u_3\right\|^2=120+240\gamma.
\]
The limits $\gamma=0$ and $\gamma=1$ reproduce the orthogonal
$\sqrt{120}$ and aligned $3\sqrt{40}$ amplitudes already
used for finite terminal-support reasoning. The squared norms are
\emph{not} physical hadron masses without the native same-action
source and spectral contraction. No Koide equal-norm constraint is
introduced.
\end{proposition}
\begin{proof}
Expand the squared norm. The equicorrelation Gram matrix is
$40[(1-\gamma)I_3+\gamma\mathbf1\mathbf1^{\mathsf T}]$.
Its eigenvalues are displayed; their nonnegativity gives the range.
\end{proof}

\begin{definition}[Source-owned pseudoscalar/vector squared-gap test]
\label{p11:def:v327-meson-squared-gap-observable}
For matched ground-state flavor sources $f$ define the \emph{measured}
spin-squared difference and average
\[
 \Delta_f^{(2)}=M_{V,f}^2-M_{P,f}^2,\qquad
 \overline{M_f^2}=\frac{M_{P,f}^2+3M_{V,f}^2}{4}.
\]
These are to be evaluated from actual, source-visible pseudoscalar and
vector poles of the same finite $H_L$ (including the controlled
refinement and resonance pole convention where necessary).
If and only if the conditional model of the mirror--spin spectrum theorem in Paper~VII-A
embeds with corresponding \emph{active states of the same mirror
parity}, it yields $\Delta_f^{(2)}=4\kappa_f$ and
$\overline{M_f^2}=a_f\pm b_f$.
It yields no universal $\kappa_f$ and no $M_{P,f}$ or $M_{V,f}$
number without the actual QCD coefficients.
\end{definition}

\begin{remark}[Benchmark spectrum and established context]
\label{p11:rem:v327-meson-matched-spectrum-benchmark}
Approximate experimental squared-mass differences in $\mathrm{GeV}^2$
for ground-state spin-partner channels are\footnote{Particle Data Group,
\emph{Review of Particle Physics} (2026), meson listings,
\url{https://pdg.lbl.gov/2026/listings/contents_listings.html}.}:
\[
\begin{array}{c|cccccccc}
 \text{pair} & \pi/\rho & K/K^* & D/D^* & D_s/D_s^*
             & B/B^* & B_s/B_s^* & \eta_c/J\!\!/\!\psi & \eta_b/\Upsilon\\\hline
 \Delta^{(2)} &0.58&0.55&0.55&0.59&0.48&0.52&0.69&1.16
\end{array}
\]
The broad $\rho$ requires a specified pole/Breit--Wigner
mass convention, so its number is a coarse diagnostic.
For heavy--light $D$ and $B$ families, an approximately shared
$M_V^2-M_P^2$ is already a consequence of heavy-quark spin symmetry
at leading order, with chromomagnetic and inverse-heavy-mass
corrections.\footnote{J.~L. Rosner and M.~B. Wise,
\emph{Phys. Rev. D} \textbf{47} (1993) 343,
\url{https://doi.org/10.1103/PhysRevD.47.343};
P. Ball and V. M. Braun, \emph{Phys. Rev. D} \textbf{49} (1994) 2472,
\url{https://doi.org/10.1103/PhysRevD.49.2472}.}
The $\eta_b/\Upsilon$ pair is an explicit counterexample to any
universal light/heavy squared-gap identity. No adjustable signed
square-root-mass triplet is used as a model input.
\end{remark}

\begin{proposition}[No mass prediction from the two-core projector alone]
\label{p11:prop:v327-no-automatic-meson-spectrum}
The existence and ranks of $Q_{mn}$, the grading algebra of
$U_{mn}$, and the conditional $\mathbf1\oplus\mathbf3$ spin
multiplicities do not imply $\Delta_f^{(2)}=\mathrm{const}$,
any identified pseudoscalar/vector pair, a universal $E_{32}$-based
squared gap, or an exact Pythagorean physical meson-mass identity.
Such statements require a specified $\iota_L$, verified source
support, the leakage/resolvent test and a computed coefficient
$\kappa_f$ in the native QCD action. In the absence of these,
$\kappa_f$ can vary without changing the geometric projector.
\end{proposition}
\begin{proof}
In the conditional model
$\Delta_f^{(2)}=4\kappa_f$. Distinct choices of $\kappa_f$
with the same $Q_{mn}$ and $U_{mn}$ preserve every finite graph
identity but change this difference. The physical embedding is also
an independent necessary hypothesis by the compression--leakage identity in Paper~VII-A.
\end{proof}
\paragraph{Required next executable test.}
The next nonnumerological task is to construct a color-singlet
intertwiner between the native Petersen two-core sector and the
physical spin-resolved source; evaluate $B$ and source residues in
the \emph{same} finite $H_L$; then compute
$\Delta_f^{(2)}$ and $\overline{M_f^2}$ without using measured
meson masses to select $b_f$ or $\kappa_f$.  The graph projector is
an eligibility test, not an independent force term.
% END FRAMEWORK v2.327: MESON SQUARED-GAP OBSERVABLE AND GRAM COHERENCE


\section{Hadron spectral masses, poles, and native finite blocks}
\label{p11:sec:v237-b3}
Finite eigenvalues, GEVP, Schur residues and stable finite poles are kept distinct from resonance widths, scattering sheets and physically occupied source states.

% Source Paper VII-B, original line 591; inherited proof bodies unchanged.
\subsubsection{Native hadron eigenvalue problem and mass closure}
\label{p11:subsec:native-hadron-eigenvalue-mass-closure}

\begin{definition}[Finite hadron eigenvalue problem]
\label{p11:def:hadronic-eigenvalue-problem}
For a supplied finite physical member \(H_L\), a hadronic sector \(\chi\) is a color-singlet invariant block.  Its finite excitation energies satisfy
\begin{equation}
\label{p11:eq:hadron-eigenvalue-problem}
P_{\chi,L}(H_L-E_{0,L})P_{\chi,L}\Psi_{h,L}=M_{h,L}\Psi_{h,L},
\qquad M_{h,L}\ge0.
\end{equation}
A physical hadron mass is a controlled refinement/continuum limit of such isolated spectral values when that limit exists.  No preferred one-step transfer time or separate wall scale is part of the definition.
\end{definition}


\begin{theorem}[Spectral-correlator and transfer-energy bridge]
\label{p11:thm:native-spectral-correlator-mass-bridge}

For a supplied self-adjoint \(H_L\) and source \(\mathcal O_h\),
\[
C_{h,L}(t)=\sum_n |\langle\Omega_L|\mathcal O_h|n,L\rangle|^2e^{-t(E_{n,L}-E_{0,L})}.
\]
If the lightest nonzero overlap in the channel has gap \(M_{h,L}\), then
\[
M_{h,L}=-\lim_{t\to\infty}\frac{d}{dt}\log C_{h,L}(t).
\]
For the Paper-VII-A transfer \(\widehat K_L(t_0)\), an exact transfer eigenvalue obeys \(M_{h,L}=-t_0^{-1}\log\lambda_h(t_0)\).
\phantomsection\label{p11:eq:spectral-correlator-decomposition}
\phantomsection\label{p11:eq:mass-from-correlator}
\phantomsection\label{p11:eq:mass-from-transfer-eigenvalue}
\end{theorem}
\begin{proof}

Insert the spectral resolution of \(H_L-E_{0,L}\).  The smallest exponent with nonzero source overlap dominates the large-time logarithmic derivative.  Functional calculus gives \(\lambda_h(t_0)=e^{-t_0M_{h,L}}\).  No independent wall scale appears.
\end{proof}

% BEGIN FRAMEWORK v3.026: source-owned five-core spectral budget, conditional.
\begin{theorem}[Conditional five-core spectral-weight and operator-residual budget]
\label{p11:thm:v3026-five-core-spectral-budget}
Fix an \emph{already realized finite physical action member} with
actual Hilbert sector pair and an off-diagonal Hamiltonian coupling
$K:\mathcal H_{\rm loc}\to\mathcal H_{\rm rec}$.
Assume the same action and source/test definitions justify all sector
projectors, and that the symmetry-compatible \emph{isometric}
transport $U_{\rm loc},U_{\rm rec}$ required in
Proposition~\ref{p10:prop:v3026-five-core-physical-matching-gate}
has been separately supplied. Put
$\widetilde C_m=U_{\rm rec}C_mU_{\rm loc}^\dagger$,
$b_m=\operatorname{Tr}(\widetilde C_m^\dagger K)$, and
$K_{\rm geom}=(40/69)\sum_m b_m\widetilde C_m$.
For any finite Hermitian receiver Hamiltonian $H_{\rm rec}$ and
$\eta>0$, define the \emph{finite model spectral weight}
\[
S_\eta(E)={\eta\over\pi}\operatorname{Tr}
 \left[K^\dagger\big((E-H_{\rm rec})^2+\eta^2 I\big)^{-1}K\right].
\]
Then, without fitted edge weights,
\[
\boxed{\int_{-\infty}^{\infty}S_\eta(E)\,dE
 -{40\over69}\sum_{m=1}^5|b_m|^2
 =\|K-K_{\rm geom}\|_{\rm HS}^2\ge0.}
\]
The defect vanishes exactly when the \emph{actual supplied} block
lies in the transported four-dimensional five-core span.
If the geometry is not physically embedded, this is only an
operator-space comparison statement and is \emph{not} a measured
hadronic line shape, pole, cross section, partial width or lifetime.
The resonance/Schur continuation and source-visible limits of this
paper remain independent requirements.
\end{theorem}
\begin{proof}
The isometric transport preserves the five-core Gram
$G=(69/40)(I_5-J_5/5)$, whose Moore--Penrose inverse is
$G^+=(40/69)(I_5-J_5/5)$. The sum $\sum_m b_m=0$ follows
from $\sum_m\widetilde C_m=0$; therefore $K_{\rm geom}$ is the
Hilbert--Schmidt orthogonal projection of $K$ onto the core span and
\[
\|K\|_{\rm HS}^2={40\over69}\sum_m|b_m|^2
 +\|K-K_{\rm geom}\|_{\rm HS}^2.
\]
The finite spectral resolution $H_{\rm rec}=\sum_jE_jP_j$ and
$\int_{\mathbb R}(\eta/\pi)[(E-E_j)^2+\eta^2]^{-1}
\,dE=1$ give
$\int_{\mathbb R}S_\eta(E)\,dE=\sum_j
\operatorname{Tr}(K^\dagger P_jK)=\|K\|_{\rm HS}^2$.
Combine the identities. No action occurrence is inferred from the
existence of a finite algebraic projection.
\end{proof}
% END FRAMEWORK v3.026: conditional geometric spectral budget.

% Framework v2.217 helper integration.
\begin{remark}[Finite cyclic source and the physical readout clock]
\label{p11:rem:v2217-temporal-source-clock}
Paper~VII-A, Theorem~\ref{p10:thm:v2217-temporal-gw-source}, proves an
actual finite nonzero off-clock pseudoscalar source, its $72a_R^2$
coincident contact and its summed axial-divergence identity. Those cyclic
Euclidean clock labels are not yet the positive-time Hamiltonian used by
the preceding spectral bridge. Already its uniform-sector source has
$C(0)=3a_R^2|1+\sigma\zeta|^2/|1-\sigma\zeta|^2$, whereas
$C(1)=12a_R^2/|1-\sigma\zeta|^2>C(0)$ for either spin sign.
It therefore cannot itself be read as a one-sided positive vacuum Laplace
correlator with nonnegative excitation energies. This comparison concerns
the explicitly projected uniform sector, not an asserted formula for the
whole 864-mode source. The actual transfer/current/clock matching remains
necessary before extracting a pion pole or axial residue.
\end{remark}


\begin{proposition}[Dimensionless representation relative to an arbitrary reference unit]
\label{p11:prop:dimensional-hadron-mass-form}

If a supplied finite Hamiltonian carries physical energy units and \(\mu_{\rm ref}>0\) is any reference unit, then
\[
M_{h,L}=\mu_{\rm ref}\,\mathcal M_{h,L}(\text{dimensionless ratios of the supplied member}).
\]
Changing \(\mu_{\rm ref}\) changes the coordinates but not \(M_{h,L}\).  The current corpus does not single out a universal \(\Lambda_{\rm BH}\) for this purpose.
\phantomsection\label{p11:eq:hadron-mass-dimensional-form}
\end{proposition}
\begin{proof}

This is dimensional bookkeeping: divide the finite matrix by \(\mu_{\rm ref}\), diagonalize the dimensionless matrix, and multiply the eigenvalue back by \(\mu_{\rm ref}\).  Since the reference is arbitrary, it cannot be promoted to a prediction.
\end{proof}

\begin{theorem}[Nucleon mass as a baryon-sector spectral gap]
\label{p11:thm:native-nucleon-mass-eigenvalue}

Let \(P_{N,L}\ne0\) project onto a baryon-number-one, spin/isospin nucleon sector of a supplied full-QCD member, \emph{provided} the actual Hamiltonian preserves that sector: \([P_{N,L},H_L]=0\). In particular, a strict isospin label requires an isospin-preserving mass/source branch; it does not follow merely from baryon conservation on unequal-mass branches. Then the finite nucleon gap is
\[
M_{N,L}=\inf\operatorname{spec}\bigl((P_{N,L}H_LP_{N,L})|_{\operatorname{Ran}P_{N,L}}\bigr)-E_{0,L}.
\]
Its physical value requires the typed quark-mass/source member and a controlled refinement limit; no count-derived additive mass formula is implied.
\phantomsection\label{p11:eq:nucleon-sector-labels}
\phantomsection\label{p11:eq:nucleon-mass-transfer-eigenvalue}
\end{theorem}
\begin{proof}

The stated \([P_{N,L},H_L]=0\) makes \(\operatorname{Ran}P_{N,L}\) invariant, so its restricted self-adjoint Hamiltonian has an exact finite spectrum. Subtracting the vacuum gives the gap. Without this commutator certificate, compression gives a trial/Ritz eigenvalue, not necessarily a pole of the actual full Hamiltonian; a source-visible spectral projector is an alternative only after its own identification.
\end{proof}


\begin{definition}[Finite baryon flux-junction support]
\label{p11:def:finite-baryon-flux-junction-support}
The minimal \(B=1\) color-singlet support is
\begin{equation}
\epsilon_{abc}q^aq^bq^c .
\end{equation}
Equivalently, in Wilson-transfer language it is represented by three
fundamental center-flux legs meeting at a \(Z_3\)-neutral baryon junction.  The
total state is a color singlet, but each leg carries fundamental \(Z_3\)-N-ality
before meeting at the junction.
\end{definition}

\begin{lemma}[Commutant criterion for baryon-junction superselection]
\label{p11:lem:no-incoherent-baryon-junction-branch}

A coherent baryon-junction branch is unique only if the commutant of the compiled baryon-sector algebra is scalar.  More generally, nontrivial projections in \(\mathfrak A'_{L,{\rm str}}\) are dynamically generated superselection projectors and must be retained.

\end{lemma}
\begin{proof}

By the finite Wedderburn decomposition of Paper~VII-A, the commutant is the full algebra of multiplicity-space operators.  A projection commutes with all compiled observables exactly when it lies in this commutant.  Thus absence of an upstream label does not exclude a dynamically generated branch; scalar commutant is the necessary and sufficient finite criterion for irreducibility.
\end{proof}

\begin{theorem}[Exact determinant Casimir contraction and dynamical response boundary]
\label{p11:thm:baryon-flux-junction-contraction}

For a color-singlet three-quark determinant,
\[
\sum_{i<j}T_i\!\cdot T_j=-2.
\]
This is the exact color contraction.  The associated baryon interaction energy is the matrix element of the supplied pair/junction operator; it is not fixed to \(2\sqrt{\deltageom}\) or to a universal transfer loss.
\phantomsection\label{p11:eq:coherent-baryon-junction-loss}
\end{theorem}
\begin{proof}

Expand \(0=(T_1+T_2+T_3)^2=3C_F+2\sum_{i<j}T_i\cdot T_j\) with \(C_F=4/3\).  The contraction follows.  Multiplying the radial/junction operator by any positive response coefficient preserves this color identity while changing the energy, proving that the latter is independent dynamical information.
\end{proof}

\begin{theorem}[Exact criterion for a positive nucleon gap]
\label{p11:thm:positive-native-nucleon-mass-floor}

The nucleon-sector gap satisfies \(M_{N,L}\ge0\).  A positive level-independent nucleon floor exists precisely when
\[
\inf_L M_{N,L}>0
\]
along the specified admissible refinement.  No positive universal lower bound follows from the determinant Casimir alone.
\phantomsection\label{p11:eq:positive-native-nucleon-mass-floor}
\end{theorem}
\begin{proof}

Finite positivity follows from ground-energy subtraction.  The uniform statement is equivalent to positivity of the displayed infimum by definition.  The family \(H_L=E_{0,L}I+L^{-1}P_{N,L}\) gives positive finite gaps but zero uniform infimum, so finite positivity alone is insufficient.
\end{proof}

\begin{corollary}[Scale-free baryon gap ratio is a computed observable]
\label{p11:cor:dimensionless-baryon-anchor-floor}

For any nonzero computed reference gap \(\Delta E_{\rm ref,L}\), the ratio \(M_{N,L}/\Delta E_{\rm ref,L}\) is dimensionless and determined by the supplied spectrum.  Without such a computed reference, no universal baryon anchor ratio is defined.
\phantomsection\label{p11:eq:dimensionless-coherent-baryon-floor}
\end{corollary}
\begin{proof}

Both numerator and denominator are spectral outputs.  Their quotient is invariant under a common change of energy units, but it varies under independent changes of admissible Hamiltonian matrix elements.  Hence geometry alone does not fix it.
\end{proof}

\begin{theorem}[Finite hadron-spectrum closure]
\label{p11:thm:finite-hadron-spectrum-closure}
For every conserved charge block \((B,Q,\vec N,J^P,\ldots)\), the finite Born--Hopf QCD Hamiltonian restricted to \(\Ps\Hraw\) has a finite singlet spectrum.  The hadron masses and mixings in that block are the corresponding eigenvalue differences and eigenvectors.  Therefore the hadron spectrum is a native spectral output of \(\HQCDfin\), not an imported mass table.
\end{theorem}
\begin{proof}
The finite quantum bridge makes \(\HQCDfin\) self-adjoint on a finite-dimensional state space.  The exact selection laws proved above decompose this space into finite conserved charge blocks.  The spectral theorem gives a complete eigenbasis in each block, and physical masses are differences from the vacuum energy in the chosen comparison convention.  No hadron mass is inserted into the Hamiltonian.\qedhere
\end{proof}

% BEGIN M8 v1.78 XI_IMPORT
\begin{remark}[Hadron and light-front readout imports]
\label{p11:rem:m8-batch-qcd-readout-inputs}
The finite field/dual identities in IV refine the
algebraic input to X; they supply no new hadron
Hamiltonian, physical species assignment or state.
The already derived number and momentum sum rules
remain algebraic readouts on their stated finite
representation.  Their PDF moment limits retain the
state, refinement and convergence hypotheses.
Native action symmetries are tested on the actual
claim-complete owner before any reduced spectral or
light-front interpretation is applied.
\end{remark}
% END M8 v1.78 XI_IMPORT




% Source Paper VII-B, original line 767; inherited proof bodies unchanged.
\subsubsection{Same-channel mixing and pole extraction}
\label{p11:subsec:same-channel-mixing-pole-extraction}

\begin{lemma}[Scheme-invariant isolated pole on a common analytic chart]
\label{p11:lem:v153-scheme-invariant-isolated-pole}
Let two renormalized inverse-propagator presentations of the same physical
two-point owner satisfy
\[
 \Gamma'(p)=A(p)\Gamma(p)B(p)
\]
on one analytic sheet, with \(A,B\) analytic and pointwise invertible near an
isolated zero \(p_*\) of \(\det\Gamma\).  Then \(p_*\) is an isolated zero
of \(\det\Gamma'\) with the same multiplicity.  For a \emph{simple} inverse-propagator pole, residue tensors transform
covariantly by the value-at-pole invertible channel maps. For a higher-order
pole, the full Laurent principal part transforms, and individual simple-pole
coefficients also receive derivatives of the channel maps. A physical
normalized residue is fixed only after the channel normalization is declared.
\end{lemma}

\begin{proof}
Since
\(\det\Gamma'=\det A\,\det B\,\det\Gamma\) and the first two factors
are analytic and nonzero on the chart, they do not change the divisor of the
isolated zero.  For a simple pole of \(\Gamma^{-1}\), the identity
\(\Gamma'^{-1}=B^{-1}\Gamma^{-1}A^{-1}\) gives
\(\operatorname{Res}_{p_*}\Gamma'^{-1}=B(p_*)^{-1}(\operatorname{Res}_{p_*}\Gamma^{-1})A(p_*)^{-1}\).
For a higher-order pole the same product must be expanded to full Laurent order; nonconstant analytic factors mix the coefficients.
\end{proof}

\begin{definition}[Same-channel finite transfer block]
\label{p11:def:same-channel-finite-transfer-block}
Fix level \(L\), transfer time \(t>0\), and a channel \(\chi\).  For interpolators \(\mathcal S_\chi=\{\mathcal O_{\chi,a}\}\), set \(|a\rangle=\mathcal O_{\chi,a}^\dagger|\Omega_L\rangle\) and delete null Gram directions.  Define
\begin{equation}
\label{p11:eq:same-channel-gram-transfer}
G^{(\chi)}_{ab}=\langle a|b\rangle,
\quad T^{(\chi)}_{ab}(t)=\langle a|e^{-t(H_L-E_{0,L})}|b\rangle,
\quad K_{\rm eff}^{(\chi)}=G^{-1/2}T(t)G^{-1/2}.
\end{equation}
For every positive eigenvalue \(\lambda_n^{(\chi)}(t)\), define the transfer-energy approximant
\begin{equation}
\label{p11:eq:same-channel-pole-masses}
M_n^{(\chi)}(L,\mathcal S_\chi;t)=-t^{-1}\log\lambda_n^{(\chi)}(t).
\end{equation}
The parameter \(t\) is the transfer time already present in Paper~VII-A; no additional dimensional normalizer is introduced.
\end{definition}

\begin{theorem}[Computed same-channel mixing and Rayleigh--Ritz control]
\label{p11:thm:no-parameter-same-channel-mixing}

At fixed \(L\), \(t>0\), and interpolator set, every entry of \(G^{(\chi)}\) and \(T^{(\chi)}(t)\) is a matrix element of the supplied finite state/operator package.  The Gram reduction therefore introduces no fit parameter.  On nested same-channel trial spaces, the min--max principle applied to the positive transfer operator gives monotone Ritz approximants to its largest eigenvalues, hence nonincreasing trial energies at fixed transfer time.  These are Ritz values of the \emph{transfer operator}, not in general eigenvalues of a compressed Hamiltonian or poles of the full resolvent. Identification with finite-volume physical pole energies requires a reducing spectral sector or a separate source-visible full-resolvent certificate. Continuum pole claims additionally require their stated refinement/continuum hypotheses.

\end{theorem}
\begin{proof}

Finite matrix multiplication fixes all entries before comparison.  Put $K=e^{-t(H_L-E_{0,L})}>0$. For any nested same-channel subspaces $V\subset W$ with Gram-orthogonal projectors, the Courant--Fischer formula for the compressed \emph{same} Hermitian operator $K$ gives $\lambda_j(K|_V)\leq\lambda_j(K|_W)$ for every ordinal $j\leq\dim V$, with the eigenvalues sorted decreasingly. Since $-\log(\cdot)/t$ reverses this order, the corresponding trial energies are nonincreasing under enlargement. In general $P e^{-tH_L}P\ne e^{-t P H_L P}$; the statement is not an identification of either compressed spectrum with the full pole set. If $P$ reduces $H_L$ the equality holds on $\operatorname{Ran}P$ and the full-resolvent spectral projection certifies the visible poles. Without reduction the full-resolvent test of Proposition~\ref{p11:prop:two-support-full-resolvent-pole-certificate} is necessary. Norm/strong-resolvent convergence and isolation of the relevant eigenvalue are separate analytic hypotheses for a continuum limit.
\end{proof}

\begin{theorem}[Readout chronology eliminates fitted normalizers without inventing a shell scale]
\label{p11:thm:no-additional-shell-normalizer-qcd}

Once a finite physical member \(H_L\) and transfer time \(t\) are supplied, all spectral energies and residues are fixed by matrix operations.  Paper~VII-B introduces no additional multiplicative normalizer fitted to hadron data.  This statement does not assert that Paper~VII-A has already derived an absolute QCD scale or quark-mass operator.
\phantomsection\label{p11:eq:qcd-extra-shell-carrier-test}
\phantomsection\label{p11:eq:qcd-extra-shell-source-zero}
\end{theorem}
\begin{proof}

Diagonalization and matrix elements are deterministic functions of the supplied finite matrices.  A posteriori rescaling to match comparison data would alter the upstream member and violate the acyclic handoff.  Conversely, if an upstream physical member is missing, the numerical execution is undefined rather than normalized by Paper~VII-B.
\end{proof}

\begin{corollary}[First two-support transfer extraction and conditional pole identification]
\label{p11:cor:first-two-support-pole-extraction}
Let a channel \(\chi\) have two non-null Gram-orthonormal native support states
\(|A\rangle,|B\rangle\).  In that basis the first nontrivial transfer block is
\begin{equation}
\label{p11:eq:first-two-support-transfer-block}
K^{(\chi)}_{AB}(t)=\begin{pmatrix}
\lambda_A&\kappa_{AB}\\
\overline{\kappa}_{AB}&\lambda_B
\end{pmatrix},
\qquad
\kappa_{AB}(t)=\langle A|\widehat K_L(t)|B\rangle .
\end{equation}
Its two exact \emph{trial-transfer} eigenvalues are
\begin{equation}
\label{p11:eq:first-two-support-eigenvalues}
\lambda_{\pm}^{(\chi)}
=\frac{\lambda_A+\lambda_B}{2}
\pm
\sqrt{\left(\frac{\lambda_A-\lambda_B}{2}\right)^2+
|\kappa_{AB}|^2},
\end{equation}
and the corresponding transfer-energy approximants (retaining the legacy
cross-reference label) are
\begin{equation}
\label{p11:eq:first-two-support-pole-masses}
M_{0,1}^{(\chi)}(t)=-t^{-1}\log\lambda_{\pm}^{(\chi)}(t) .
\end{equation}
Thus same-channel mixing lowers the lightest \emph{trial} energy and raises the
orthogonal trial partner whenever \(\kappa_{AB}\ne0\), while the shift is
computed from the native overlap, not from a tunable mixing angle.  The two
values are exact finite \emph{physical} pole energies if this non-null
support space reduces the physical \(H_L\) in its selected conserved
channel.  Without that certificate, their actual physical-pole content is
determined instead by the source-visible full resolvent of
Proposition~\ref{p11:prop:two-support-full-resolvent-pole-certificate};
a trial-transfer eigenvalue alone is not a pole assertion.
\end{corollary}
\begin{proof}
The displayed spectrum is the exact spectrum of the complex Hermitian
two-by-two compression.  A relative rephasing makes \(\kappa_{AB}\)
nonnegative real but is not needed for the formula. Positivity of the
finite transfer makes both eigenvalues positive on the non-null span.
The larger eigenvalue increases and the smaller decreases with
\(|\kappa_{AB}|\), so applying \(-t^{-1}\log\) yields the claimed trial
energy shifts. If the support projector \(P_{AB}\) commutes with \(H_L\),
then \(P_{AB}e^{-t(H_L-E_{0,L})}P_{AB}\) is the transfer of the actual
restricted Hamiltonian. Its eigenvalues are \(e^{-t(E_n-E_{0,L})}\)
for genuine full-Hamiltonian eigenvectors in that subspace, and their
nonzero spectral residues are visible to the spanning interpolators.
Without this commutation, no spectral-correlator pole bridge applies to
the two compressed eigenvalues in general.
\end{proof}

\begin{proposition}[Exact full-resolvent certification of two-support physical poles]
\label{p11:prop:two-support-full-resolvent-pole-certificate}
Let \(P=P_{AB}\) be the orthogonal projector onto the two-source support,
\(Q=I-P\), and \(\widetilde H_L=H_L-E_{0,L}I\). The exact compressed
source-visible resolvent is
\begin{equation}
\label{p11:eq:two-support-schur-pole-certificate}
P(z-\widetilde H_L)^{-1}P
=\Bigl[zP-P\widetilde H_LP-\Sigma_P(z)\Bigr]_{\operatorname{Ran}P}^{-1},
\quad
\Sigma_P(z)=P\widetilde H_LQ
\bigl[zQ-Q\widetilde H_LQ\bigr]_{\operatorname{Ran}Q}^{-1}
Q\widetilde H_LP,
\end{equation}
for \(z\) in the joint resolvent domain where the indicated inverses
exist. A real finite spectral value \(\epsilon\) is a pole in the
\(P\)-visible correlator if and only if
\(P\Pi_{\epsilon}(\widetilde H_L)P\ne0\); then its residue is precisely
\(P\Pi_{\epsilon}(\widetilde H_L)P\). In particular, the bare two-state
transfer eigenvalues do not determine full physical poles whenever the
source leaks into \(Q\). For a selected physical finite \(H_L\), the
full-resolvent computation is an exact non-fitted pole readout; it does
not independently establish that the required physical owner has been
realized or that a continuum pole limit exists.
\end{proposition}
\begin{proof}
Block inversion of \(z-\widetilde H_L\) relative to
\(\operatorname{Ran}P\oplus\operatorname{Ran}Q\) proves the Schur identity.
The finite spectral theorem gives
\(P(z-\widetilde H_L)^{-1}P=
\sum_{\epsilon}P\Pi_{\epsilon}(\widetilde H_L)P/(z-\epsilon)\).
For distinct real \(\epsilon\) there is no cancellation between
these simple pole terms, so the residue is nonzero exactly under the
stated source-visibility test. If \([P,\widetilde H_L]=0\), the leakage
term vanishes, and the transfer eigenvalues match the visible full
Hamiltonian poles. Otherwise the self-energy must be retained. This is
the two-support specialization of the previously established full
source-resolvent interface also used in
Proposition~\ref{p11:prop:v3023-mobile-2i-pole-interface}.
\end{proof}


% Source Paper VII-B, original line 1956; inherited proof bodies unchanged.
% Former logical part: Constructed QCD Blocks and Downstream Readout Packages; integrated into the new thematic argument.
\label{p11:part:constructed-qcd-blocks}


% Source Paper VII-B, original line 1959; inherited proof bodies unchanged.
\subsection[Constructed QCD blocks and downstream closure packages]
{Constructed QCD blocks and closure packages}
\label{p11:sec:validation}
\label{p11:sec:constructed-qcd-blocks}


\begin{theorem}[Inclusive color-counting export]
\label{p11:thm:inclusive-color-counting-export}
The derived color carrier \(C\simeq\mathbb C^3\) fixes \(N_c=3\).  For high-energy inclusive annihilation away from threshold and resonance regions,
the electromagnetic current-correlator/OPE hard coefficient is
\(3\sum_fQ_f^2[1+\alpha_s(s)/\pi+O(\alpha_s^2)]\).  On a physical inclusive
annihilation domain satisfying \textnormal{(SM1)--(SM3)} and
\textnormal{(SM7)} of Definition~\ref{p12:def:claim-indexed-smat-package},
\textnormal{(SM4)} if an LSZ representation is used, and a selected QED
infrared route from Definition~\ref{p12:def:claim-indexed-infrared-package},
the optical-theorem/cross-section identification gives
\begin{equation}
\label{p11:eq:eplus-eminus-r-ratio}
R(s)=\frac{\sigma(e^+e^-\to{\rm hadrons})}{\sigma(e^+e^-\to\mu^+\mu^-)}
=3\sum_{f:2m_f<\sqrt s}Q_f^2\left[1+\frac{\alpha_s(s)}{\pi}+O(\alpha_s^2)\right].
\end{equation}
Thus
\begin{equation}
\label{p11:eq:r-ratio-threshold-values}
R_{uds}=2,
\qquad
R_{udsc}=\frac{10}{3},
\qquad
R_{udscb}=\frac{11}{3}
\end{equation}
at leading hard order.  The coefficient is a direct export of the derived
color multiplicity, electric charges, and strong running.  Its interpretation
as the displayed physical inclusive observable is claim-scoped; detailed
fragmentation cancels in the inclusive sum to leading power on that domain.
\end{theorem}
\begin{proof}
The electromagnetic current OPE has flavor weight \(Q_f^2\), while the color
trace gives \(N_c=3\); real and virtual gluon corrections give the displayed
first perturbative coefficient.  These are local/current-correlator statements.
The optical theorem and incident/final-state normalization convert the
spectral discontinuity into the cross-section ratio only under the named
scattering and infrared clauses.  Inclusivity then removes dependence on
detailed exclusive hadronization functions to the stated perturbative order.
\end{proof}

\begin{longtable}{>{\raggedright\arraybackslash}p{0.23\textwidth} >{\raggedright\arraybackslash}p{0.38\textwidth} >{\raggedright\arraybackslash}p{0.28\textwidth}}
\caption[QCD validation and failure-test ledger]{QCD validation and failure-test ledger.}\label{p11:tab:validation-ledger}\\
\toprule
Target & Native quantity to derive & Comparison layer \\
\midrule
\endfirsthead
\toprule
Target & Native quantity to derive & Comparison layer \\
\midrule
\endhead
Strong running & conditional perturbative source/threshold branch and scheme conversion & lattice/PDG running-coupling summaries \\
Inclusive color counting & \(R(s)=3\sum_fQ_f^2[1+\alpha_s(s)/\pi+\cdots]\) & high-energy \(e^+e^-\to\) hadrons \\
Mass gap/string tension & computed finite gap sequence; conditional area-law definition after pairing & lattice pure-gauge/full-QCD comparison \\
Chiral sector & correlator-derived \(\Sigma,F_\pi,B_0\) and pseudoscalar spectrum & lattice/chiral phenomenology \\
Nucleon & spectral mass differences, current residues and sigma terms & experiment/\allowbreak lattice/\allowbreak isospin+EM decompositions \\
Hadron spectrum & low meson/baryon/glueball towers & PDG/lattice spectroscopy \\
PDF moments & typed \(\langle x^n\rangle_{i/h}(\mu)\) and evolution & DIS/Drell--Yan/global PDF fits \\
DAs and FFs & pion-DA parity, DA moments, fragmentation/hadronization map & exclusive processes and fragmentation data \\
Matrix elements & weak, scalar, electromagnetic, tensor, gravitational, and two-baryon contact matrix elements & lattice/flavor/global-fit/nuclear comparisons \\
Topology & \(\chi_t=V^{-1}\operatorname{Var}(\nu)\), higher cumulants \(b_2,b_4\), \(U(1)_A\), and the \(\eta'\) anomaly contribution from the computed sector measure & lattice/topological and spectroscopy comparisons \\
Thermal QCD & Goldstone pion-gas endpoint, Stefan--Boltzmann endpoint, finite-trace crossover susceptibilities and, only after the thermodynamic/refinement limit, a possible chiral transition temperature & lattice thermodynamics/heavy-ion constraints \\
Cosmology export & hadronic/nuclear reaction kernels, OPE boundary, thermal trace-anomaly operator readout & BBN/CMB/structure propagation in Paper~VI-B \\
\bottomrule
\end{longtable}
\TableNotes{These targets test the completed finite Paper~VII-B readout map.  They are validation and failure-test channels, not selection data for the finite strong-kernel construction.}

\begin{theorem}[Acyclic validation hierarchy]
\label{p11:thm:qcd-validation-hierarchy}

The rigorous hierarchy is logical rather than numerical: (i) verify exact finite geometry and algebra; (ii) verify the supplied finite source/action member; (iii) compute spectra, residues, traces and cumulants; (iv) prove required refinement/continuum limits; only then (v) compare with lattice, experiment, or phenomenology.  String tension, chiral constants, and topology measure remain validation targets until computed.  A uniform continuous gap is tested only on an enlarged comparison branch that independently claims one; it is not an active native prediction.
\phantomsection\label{p11:eq:minimal-validation-target-triplet}
\phantomsection\label{p11:eq:strong-validation-targets}
\end{theorem}
\begin{proof}

Each stage is a function of the preceding data.  Allowing comparison data to choose an earlier coefficient reverses the dependency graph and constitutes calibration.  This proves the stated ordering while making no claim that the open observables are already solved.
\end{proof}


\begin{longtable}{@{}>{\raggedright\arraybackslash}p{0.22\textwidth} >{\raggedright\arraybackslash}p{0.72\textwidth}@{}}
\caption[Constructed QCD block status]{Constructed QCD block status.}\label{p11:tab:constructed-block-status}\\
\toprule
Block & Constructed source, output, and remaining obligation \\
\midrule
\endfirsthead
\toprule
Block & Constructed source, output, and remaining obligation \\
\midrule
\endhead
Primitive support transfer & \textbf{Construction:} compression of \(e^{-t(H_L-E_0)}\) to computed invariant sectors\par\textbf{Output:} finite spectral gaps/overlaps\par\textbf{Remaining:} execute supplied physical member  \\
Topology & \textbf{Construction:} \(Z_L(\theta)\), charge distribution and covariance\par\textbf{Output:} finite cumulants/susceptibility\par\textbf{Remaining:} admissible index/continuum measure  \\
Distribution evolution & \textbf{Construction:} derived \(\LFRayGeom\) plus source-owned light-ray/twist-two operators; exact fraction POVMs are an optional stronger finite branch\par\textbf{Output:} renormalized operator moments, DGLAP/ERBL and current/stress sum rules\par\textbf{Remaining:} prove \(\LFState_{h,L}\) and only the consumed \(\LFLim\)/factorization clauses; prove \(\LFRep_L\) only when the optional finite POVM branch is claimed  \\
Pseudoscalar/chiral & \textbf{Construction:} mass-source Hessian plus chiral correlators\par\textbf{Output:} symbolic/finite mass and current blocks\par\textbf{Remaining:} physical \(M_q^{\rm phys},F_\pi,B_0,\Sigma\), SSB limit  \\
Resonance/\(p^4\) & \textbf{Construction:} current/source correlator poles, continued matrix residues, and low-momentum expansion\par\textbf{Output:} \(L_i\), pole residues, pole-width parameters\par\textbf{Remaining:} execute correlators/continuation; physical rates are claim-scoped  \\
Baryon currents & \textbf{Construction:} determinant color support plus computed current operators\par\textbf{Output:} charges, axial/tensor/magnetic/scalar readouts\par\textbf{Remaining:} execute baryon eigenvectors/current matrix elements  \\
Pure glue & \textbf{Construction:} exact Hodge/incidence geometry plus Yang--Mills quadratic form\par\textbf{Output:} dimensionless geometry, finite spectral targets\par\textbf{Remaining:} conditional continuum-gap certificate only if claimed; area-law/pairing proof  \\
Nuclear & \textbf{Construction:} OPE with computed chiral inputs plus QCD-matched contact operators\par\textbf{Output:} two-/many-baryon kernels\par\textbf{Remaining:} contact/multi-baryon execution  \\
Thermal & \textbf{Construction:} finite analytic \(Z_L(T,\mu)\)\par\textbf{Output:} responses/crossover diagnostics\par\textbf{Remaining:} thermodynamic/refinement limit  \\
Physical resonances & \textbf{Construction:} Schur/Feshbach continuation of specified blocks\par\textbf{Output:} complex poles/residues\par\textbf{Remaining:} open-channel continuum construction  \\
\bottomrule
\end{longtable}
\TableNotes{No row upgrades an unresolved source or limit to a numerical prediction.  Each physical output is computed from the supplied operator member or remains an explicit obligation.}

\paragraph{Controlled bounded Wilson source readouts.}
\label{p11:par:v202-bounded-source-handoff}
For the actual positive heat blocks, real and imaginary parts of normalized
Wilson traces are bounded by one. Theorem~\ref{p12:thm:v202-source-errors}
therefore controls every fixed connected source derivative, with the
degree-17 finite-block margin in
Corollary~\ref{p12:cor:v202-native-cumulants}. Once bounded QFT4 state
comparison is established, Theorem~\ref{p12:thm:v202-bounded-source-inheritance}
also supplies their source limits. This does not assign a topological
sector weight, physical susceptibility, unbounded current normalization,
thermodynamic limit or physical pole.


% Source Paper VII-B, original line 2079; inherited proof bodies unchanged.
\subsubsection{Derived finite-readout upgrades from the closure ledger}
\label{p11:subsec:derived-finite-readout-upgrades}

The following package separates exact finite/geometric consequences from conditional standard-QCD readouts.  No transfer-loss parameter, chiral normalization, or determinant-baryon residue is presumed fixed by counting; such quantities enter only through their owning matrix/correlator computations.

\begin{corollary}[Perturbative comparison readouts from the imported evolution interface]
\label{p11:cor:derived-perturbative-comparison-package}
On the explicitly selected Paper~V-C comparison member \(t=\tfrac12\), for which
\(\alpha_s((M_Z^{\rm BH})^2)=\delta\), Paper~VII-B may combine the threshold chart with the imported DGLAP/ERBL/timelike evolution operator.  It may therefore compute branch-conditional scheme scales, inclusive electromagnetic color-counting ratios, evolved moments, and hard-process limits without owning or rederiving the perturbative kernels.  In particular,
\begin{equation}
\label{p11:eq:lambda-msbar-five-two-loop-bh}
\Lambda_{\MSbar,5}^{(2),\rm BH}=206.9499404\,\mathrm{MeV}
\end{equation}
is conditional on the explicit \emph{external} comparison reference
\(M_Z^{\rm cmp}=91.1876\,\mathrm{GeV}\), the chart
\(\alpha_s((M_Z^{\rm cmp})^2)=\delta=1/\sqrt{72}\), and the conventional asymptotic two-loop \(n_f=5\) \(\overline{\rm MS}\) definition
\[
\Lambda_5^{(2)}=M_Z^{\rm cmp}\exp\!\left(-\frac{2\pi}{\beta_0\alpha_s}\right)
\left(\frac{\beta_0\alpha_s}{4\pi}\right)^{-\beta_1/(2\beta_0^2)},
\quad \beta_0=\frac{23}{3},\quad\beta_1=\frac{116}{3}.
\]
The decimal is a \(t=\tfrac12\) comparison-scheme output, \emph{not} a native Born--Hopf prediction of the dimensionful input or a nonperturbative scale, and
\begin{equation}
\label{p11:eq:derived-uv-momentum-angular-fixed-point}
Q_\infty=\frac9{25},\qquad G_\infty=\frac{16}{25},\qquad
J_{q,\infty}=\frac9{50},\qquad J_{g,\infty}=\frac8{25}
\end{equation}
are imported-evolution fixed-point consequences for a Paper-VII-B boundary state.
\end{corollary}
\begin{proof}
Apply the Paper-V-C perturbative clock and evolution semigroup to the native
Paper-VII-B boundary data.  The decimal follows by substituting the declared external
\(91.1876\,\mathrm{GeV}\) reference and \(\delta=1/\sqrt{72}\) into the displayed two-loop convention (yielding \(206.94994037\ldots\,\mathrm{MeV}\)); no absolute scale arises from the dimensionless coupling alone. The other numerical outputs are deterministic conditional comparison readouts; the evolution machinery remains owned by Paper~V-C.
\end{proof}

\begin{corollary}[Exact support identities versus dynamical baryon matrix elements]
\label{p11:cor:primitive-support-and-baryon-matrix-elements}

The color support/Casimir identities are exact.  Axial, tensor, magnetic, scalar and sigma-term readouts are instead matrix elements or source derivatives of the supplied baryon eigenstates.  The symmetric spin-flavor values such as \(g_A=5/3\) may be retained only as explicitly labeled representation benchmarks, not as finite-QCD predictions.
\phantomsection\label{p11:eq:primitive-support-floors-closed}
\phantomsection\label{p11:eq:primitive-axial-tensor-charges}
\phantomsection\label{p11:eq:primitive-baryon-magnetic-moments}
\phantomsection\label{p11:eq:primitive-sigma-term-floor}
\phantomsection\label{p11:eq:primitive-ope-coefficient}
\phantomsection\label{p11:eq:primitive-pion-nucleon-coupling-floor}
\end{corollary}
\begin{proof}

Group theory fixes transformation laws and Clebsch--Gordan coefficients but leaves reduced matrix elements.  Wigner--Eckart therefore separates exact representation factors from dynamical reduced matrix elements.  Feynman--Hellmann similarly expresses scalar responses as derivatives of the computed eigenvalue rather than a count-derived coefficient.
\end{proof}

\begin{theorem}[Topological variance is the computed charge covariance]
\label{p11:thm:orthogonal-hurwitz-adjoint-winding-variance}
\label{p11:cor:full-primitive-topological-winding-block}

Let the regulated topological charge decompose as \(\nu_L=w\cdot\boldsymbol\nu_L\).  Then
\[
\operatorname{Var}(\nu_L)=w^TC_Lw,
\qquad (C_L)_{ij}=\operatorname{Cov}(\nu_{i,L},\nu_{j,L}).
\]
Neither the number of channels nor a finite 24/25 arithmetic factor fixes \(C_L\).
\phantomsection\label{p11:eq:derived-topological-enhancement-32}
\phantomsection\label{p11:eq:full-primitive-Xt-closed}
\end{theorem}
\begin{proof}

Variance of a linear combination is bilinear in the covariance matrix.  Positive-semidefinite matrices \(C_L\), \(2C_L\), and correlated deformations can share the same channel labels/count while giving different variances.  Hence the covariance must be computed from the topological transfer measure.
\end{proof}

\begin{corollary}[Hard-exclusive asymptotic kernels with dynamical normalization]
\label{p11:cor:hard-exclusive-da-moments-and-kernels}

The inverse moments implied by the asymptotic pion and baryon DAs are exact
integrals of those asymptotic shapes.  On the process domain of
Theorem~\ref{p11:thm:native-qcd-factorization-bridge}, hard-exclusive amplitudes
factorize into perturbative kernels times the relevant decay constants/DAs.
Their absolute normalization remains conditional on the computed \(F_\pi\),
baryon normalization, and source matrix elements.
\phantomsection\label{p11:eq:pion-da-inverse-moments-closed}
\phantomsection\label{p11:eq:pion-transition-form-factor-asymptotic-bh}
\phantomsection\label{p11:eq:pion-em-form-factor-hard-asymptotic-bh}
\phantomsection\label{p11:eq:baryon-da-inverse-moments-closed}
\end{corollary}
\begin{proof}

The inverse moments are elementary beta-function integrals.  The separately
declared factorization/scattering package supplies the convolution formula, but
the current correlator residue \(F_\pi\) is not fixed by the 24-cell geometry.
Therefore the shape-dependent integrals survive while predecessor absolute
numerical normalizations are withdrawn.
\end{proof}

\begin{corollary}[Thermal endpoint susceptibilities]
\label{p11:cor:thermal-bqs-susceptibility-endpoint}
In the massless three-flavor Stefan--Boltzmann endpoint, with
\begin{equation}
\mu_u=\frac13\mu_B+\frac23\mu_Q,
\quad
\mu_d=\frac13\mu_B-\frac13\mu_Q,
\quad
\mu_s=\frac13\mu_B-\frac13\mu_Q-\mu_S,
\end{equation}
the quadratic \((B,Q,S)\) susceptibility matrix is
\begin{equation}
\label{p11:eq:thermal-bqs-chi2-matrix}
\chi_2^{BQS}=
\begin{pmatrix}
1/3 & 0 & -1/3\\
0 & 2/3 & 1/3\\
-1/3 & 1/3 & 1
\end{pmatrix}.
\end{equation}
The fourth-to-second diagonal ratios are
\begin{equation}
\label{p11:eq:thermal-bqs-kurtosis-ratios}
\frac{\chi_4^B}{\chi_2^B}=\frac{2}{3\pi^2},
\qquad
\frac{\chi_4^Q}{\chi_2^Q}=\frac{2}{\pi^2},
\qquad
\frac{\chi_4^S}{\chi_2^S}=\frac{6}{\pi^2}.
\end{equation}
Together with Theorems~\ref{p11:thm:native-qcd-stefan-boltzmann-limit}
and~\ref{p11:thm:native-low-temperature-pion-gas}, this fixes the exact
high-temperature normalization of the charge susceptibilities.  The interacting
crossover is not an independent thermal input.  It is the finite cumulant
readout of the same spectrum and charge operators.
\end{corollary}
\begin{proof}
For a massless Dirac flavor, \(p_f/T^4=N_c[7\pi^2/180+(\mu_f/T)^2/6+(\mu_f/T)^4/(12\pi^2)]\).  Differentiating this expression with the displayed charge assignments gives \eqref{p11:eq:thermal-bqs-chi2-matrix} and \eqref{p11:eq:thermal-bqs-kurtosis-ratios}.
\end{proof}


\begin{corollary}[Pseudoscalar ideal-gas fingerprint as a symbolic endpoint]
\label{p11:cor:native-pseudoscalar-gas-fingerprint}

For a specified pseudoscalar mass set \(\{m_a\}\), the noninteracting gas pressure and susceptibilities are the standard Bose sums and are exact within the ideal-gas approximation.  They are functions of \(m_a/T\); no fixed Paper-VII-B crossover temperature or pseudoscalar mass decimal is implied.
\phantomsection\label{p11:eq:ps-gas-pressure-fingerprint}
\phantomsection\label{p11:eq:ps-gas-energy-fingerprint}
\phantomsection\label{p11:eq:ps-gas-susceptibility-fingerprint}
\phantomsection\label{p11:eq:ps-gas-eos-numerical-fingerprint}
\phantomsection\label{p11:eq:ps-gas-bqs-numerical-fingerprint}
\phantomsection\label{p11:eq:ps-gas-fourth-cumulant-numerical-fingerprint}
\end{corollary}
\begin{proof}

The finite-volume Bose partition function factorizes over one-particle modes, and differentiating its logarithm gives pressure and charge cumulants.  Because the masses and any crossover temperature are independent dynamical inputs/readouts, substituting predecessor decimals is not a theorem of the current corpus.
\end{proof}



% Source Paper VII-B, original line 2588; inherited proof bodies unchanged.
\subsubsection{Native resonance and \texorpdfstring{$p^4$}{p4} chiral closure blocks}
\label{p11:subsec:native-resonance-p7-closures}

Given a physical quark-mass operator and controlled chiral limit, the light
pseudoscalar layer is governed by Ward identities and correlator-derived
low-energy constants.  The next layer concerns non-Goldstone poles and the
analytic \(p^4\) boundary.  Pole bases, low-energy constants, slopes, and
complex pole-width parameters are defined by the corresponding continued QCD
correlators; physical partial rates additionally use the claim-scoped
scattering package.  Resonance saturation, when used, is an approximation to
those correlators rather than a uniqueness theorem.

\begin{theorem}[Vector/scalar resonance supports are spectral sectors, not count energies]
\label{p11:thm:leading-transverse-contact-support-selection}

The finite carrier and quantum numbers define vector and scalar invariant source sectors.  Their lowest energies are the corresponding spectral gaps of \(H_L\); symmetry alone does not set \(M_V\) to a sum of count-derived floors or \(M_S=4\pi F_\pi\).
\phantomsection\label{p11:eq:leading-vector-support-selection}
\phantomsection\label{p11:eq:leading-scalar-contact-selection}
\end{theorem}
\begin{proof}

Project \(H_L\) onto the specified \(J^{PC}\) and flavor sectors and apply the min--max principle.  Multiplying an allowed interaction in either sector by a positive scalar changes the pole while preserving the support labels and symmetries, disproving a unique count-based energy.
\end{proof}

\begin{theorem}[Vector nonet as a finite correlator eigenproblem]
\label{p11:thm:native-leading-vector-nonet}

In the approximate \(SU(3)_f\) regime, vector sources decompose into octet plus singlet.  The physical vector poles are eigenvalues of the corresponding vector-current correlator/transfer block.  First-order flavor breaking is described by the allowed mass-source matrix elements; its reduced coefficients must be computed and are not fixed by reusing \(B_0\) from the pseudoscalar sector.
\phantomsection\label{p11:eq:native-vector-base-mass}
\phantomsection\label{p11:eq:native-vector-mass-square-operator}
\phantomsection\label{p11:eq:native-vector-nonet-numeric}
\phantomsection\label{p11:eq:native-vector-gmo}
\end{theorem}
\begin{proof}

Representation theory fixes the octet/singlet decomposition and the number of reduced matrix elements.  Wigner--Eckart does not equate reduced matrix elements belonging to different operators (vector pole Hamiltonian versus pseudoscalar chiral Hessian).  Hence the pole matrix is a computed source block.
\end{proof}

\begin{corollary}[Vector residues and radii as current-correlator readouts]
\label{p11:cor:native-vector-current-constants}

Vector decay constants, electromagnetic form-factor slopes and radii are residues/derivatives of the corresponding conserved-current matrix elements.  Vector-meson dominance may be used as an explicitly labeled approximation after the vector spectral density is known, but it is not a uniqueness theorem.
\phantomsection\label{p11:eq:native-L9-vector}
\phantomsection\label{p11:eq:native-pion-charge-radius-vector}
\phantomsection\label{p11:eq:native-kl3-vector-slopes}
\end{corollary}
\begin{proof}

Ward identities determine \(F(0)\) for conserved charges.  The Taylor coefficient \(F'(0)\) and pole residues are independent transverse matrix elements.  Analytic families with the same \(F(0)\) and arbitrary slope provide a direct non-uniqueness proof.
\end{proof}

\begin{theorem}[Axial/vector spectral constraints and non-unique saturation]
\label{p11:thm:native-axial-pole-L10}

The vector--axial correlator determines \(L_{10}\) through its low-momentum expansion (with the stated renormalization convention).  If one additionally imposes narrow-resonance saturation and the first two Weinberg sum rules, the solutions form at least the one-parameter family
\[
F_A^2=c,\quad F_V^2=F_\pi^2+c,\quad
M_A^2=\frac{F_V^2}{F_A^2}M_V^2,
\qquad c>0.
\]
Therefore the sum rules do not uniquely fix \(M_A\), \(F_V\), or \(L_{10}\).
\phantomsection\label{p11:eq:native-axial-pole-mass}
\phantomsection\label{p11:eq:native-L10}
\phantomsection\label{p11:eq:native-axial-strange-operator}
\end{theorem}
\begin{proof}

Substitution verifies both Weinberg sum rules for every \(c>0\).  Hence any unique choice such as \(F_V^2=2F_\pi^2\) requires another dynamical input.  The exact object is the current-correlator coefficient itself.
\end{proof}

\begin{theorem}[Scalar low-energy constants are source-response coefficients]
\label{p11:thm:native-scalar-L5-L8}

The chiral constants \(L_5\) and \(L_8\) are coefficients in the low-energy expansion of the QCD generating functional and are computable from scalar/pseudoscalar source correlators.  Relations obtained by integrating out a narrow scalar resonance are conditional saturation formulas with independently computed \(M_S,c_d,c_m\); symmetry does not force \(M_S=4\pi F_\pi\) or \(c_d=c_m=F_\pi/2\).
\phantomsection\label{p11:eq:native-scalar-contact-scale}
\phantomsection\label{p11:eq:native-L5-L8}
\phantomsection\label{p11:eq:native-2L8-minus-L5-zero}
\phantomsection\label{p11:eq:native-scalar-kl3-slopes}
\end{theorem}
\begin{proof}

The chiral operator basis fixes which coefficients occur but not their numerical values.  The family \(M_S=c_SF_\pi, c_d=aF_\pi,c_m=bF_\pi\) respects the same dimensions/symmetries for arbitrary positive dimensionless constants, proving non-uniqueness without correlator input.
\end{proof}

\begin{theorem}[Four-Goldstone low-energy constants from the generating functional]
\label{p11:thm:native-L1-L2-L3}

The constants \(L_1,L_2,L_3\) are the corresponding \(O(p^4)\) coefficients of the renormalized chiral generating functional.  Resonance-exchange formulas are optional approximations conditional on computed vector/scalar poles and residues; they are not fixed by finite support counts.
\phantomsection\label{p11:eq:native-L123-numeric}
\end{theorem}
\begin{proof}

Matching the low-momentum four-current/four-Goldstone amplitudes to the chiral operator basis uniquely defines the renormalized coefficients in a chosen scheme.  Different admissible resonance residues produce different saturation values while leaving the finite carrier unchanged, so the carrier alone does not determine them.
\end{proof}

\begin{theorem}[Anomalous singlet contribution to L7 is conditional on the computed singlet pole]
\label{p11:thm:native-L7-anomalous}

When a sufficiently heavy pseudoscalar singlet is integrated out in the chiral domain, its exchange gives the standard negative contribution to \(L_7\) expressed in terms of the singlet mass and decay normalization.  Those quantities must be computed from the topological/chiral correlator block; no fixed anomaly mass follows from channel counts.
\phantomsection\label{p11:eq:native-L7}
\end{theorem}
\begin{proof}

Tree-level elimination of a quadratic singlet field produces the usual inverse-mass-squared local operator.  The sign and structure follow algebraically, while its magnitude depends on the singlet propagator and coupling.  Current Paper~VII-A exports those as topological/chiral targets rather than fixed decimals.
\end{proof}

\begin{corollary}[Low-energy-constant status summary]
\label{p11:cor:native-leading-p7-constants-summary}

The renormalized \(L_i^r(\mu)\) are QCD correlator/matching outputs.  Symmetry may force some operator structures to vanish in special limits, but absence of a finite count label does not imply \(L_4=L_6=0\), and no resonance-saturation numerical summary is theorem-level without computed residues.
\phantomsection\label{p11:eq:native-leading-p7-constants-summary}
\end{corollary}
\begin{proof}

The chiral basis contains all symmetry-allowed operators.  Their coefficients are independent reduced response data unless an additional dynamical relation is proved.  A zero coefficient is therefore a substantive dynamical statement, not a consequence of bookkeeping.
\end{proof}


% BEGIN INLINE SOURCE: Paper7B-ChiralCertificateHandoff-Main-v1_0.tex
% Framework v2.164: chiral execution consumes the complete same-action certificate.

% Source Paper VII-B, original line 3106; inherited proof bodies unchanged.
\subsubsection{Proof-audited downstream closure blocks B--F}
\label{p11:sec:block-b-through-f-closures}

The preceding constructions remove the native QCD inputs.  The remaining downstream objects are physical pole dressing, the pseudoscalar topology/glue interface, distribution functions, the thermal equation of state, and nuclear physics.  This subsection records the proof-audited closure status of these blocks.  It also separates theorem-level readouts from finite matrix executions still required for numerical precision.

\begin{theorem}[Block B: exact finite Schur spectrum]
\label{p11:thm:block-b-finite-pole-dressing}
Let \(\mathcal H_\chi=P_\chi\mathcal H_\chi\oplus Q_\chi\mathcal H_\chi\)
be a fixed finite quantum-number channel split into compact and complementary
supports.  With
\begin{equation}
H_\chi=\begin{pmatrix}H_{PP}&H_{PQ}\\ H_{QP}&H_{QQ}\end{pmatrix},
\end{equation}
define, for \(z\notin\sigma(H_{QQ})\),
\begin{equation}
\label{p11:eq:block-b-feshbach-effective-hamiltonian}
H_{{\rm eff},L}^{(\chi)}(z)
=H_{PP}+H_{PQ}(z-H_{QQ})^{-1}H_{QP}.
\end{equation}
The exact projected resolvent identity is
\begin{equation}
\label{p11:eq:block-b-projected-resolvent}
P_\chi(z-H_\chi)^{-1}P_\chi
=\bigl[z-H_{{\rm eff},L}^{(\chi)}(z)\bigr]^{-1},
\end{equation}
so finite eigenvalues away from \(\sigma(H_{QQ})\) are the real zeros of
\begin{equation}
\label{p11:eq:block-b-pole-equation}
\det\!\bigl[z\mathbf 1-H_{{\rm eff},L}^{(\chi)}(z)\bigr]=0.
\end{equation}
At this finite self-adjoint level every pole is real and every resolvent residue
is the compression of an orthogonal spectral projector.  There is no
\(+i0\) boundary value, branch cut, second sheet, or nonzero width.  Those
objects arise only after a separately controlled open-channel/refinement limit
and analytic continuation.
\end{theorem}
\begin{proof}
Block Gaussian elimination gives
Equation~\eqref{p11:eq:block-b-projected-resolvent}.  The poles of the full
finite resolvent are eigenvalues of the self-adjoint matrix \(H_\chi\), hence
are real.  The finite spectral theorem gives
\((z-H_\chi)^{-1}=\sum_\lambda(z-\lambda)^{-1}P_{\lambda,L}\), so the residue
at \(\lambda\) is \(P_\chi P_{\lambda,L}P_\chi\).  An \(i0\) prescription
cannot be generated before a continuous spectrum and its boundary values
exist.
\end{proof}


\begin{corollary}[Pole lowering and retained residue share one eliminated channel]
\label{p11:cor:schur-positive-gap-weight}
In Theorem~\ref{p11:thm:block-b-finite-pole-dressing}, suppose the
canonical complementary Hamiltonian at the pole obeys
$0<\delta I\preceq H_{QQ}-EI\preceq DI$.
For a normalized retained null vector $u$ put
\[
 s=u^\dagger H_{PQ}(H_{QQ}-EI)^{-1}H_{QP}u
   =u^\dagger H_{PP}u-E.
\]
The normalized full eigenvector has retained population
\[
 w_{\rm ret}
 =\frac{1}{1+u^\dagger H_{PQ}(H_{QQ}-EI)^{-2}H_{QP}u},
 \qquad
 \boxed{\frac{\delta}{\delta+s}\le w_{\rm ret}
                         \le\frac{D}{D+s}.}
\]
For a simple visible pole its retained matrix residue is
$w_{\rm ret}uu^\dagger$.  For a degenerate pole let $U$ span the full
retained nullspace isometrically, and put
\[
 S=U^\dagger H_{PQ}(H_{QQ}-EI)^{-1}H_{QP}U,\qquad
 T=U^\dagger H_{PQ}(H_{QQ}-EI)^{-2}H_{QP}U.
\]
The complete retained residue is $U(I+T)^{-1}U^\dagger$, with
$(I+S/\delta)^{-1}\preceq(I+T)^{-1}\preceq(I+S/D)^{-1}$.
\end{corollary}
\begin{proof}
Apply Theorem~\ref{p1b:thm:response-positive-gap-weight} and the
inherited Hamiltonian norm of
Corollary~\ref{p10:cor:schur-positive-gap-weight} to the fixed channel.
The full eigenvector lift gives the norm; positive-gap spectral calculus
gives the scalar and matrix bounds.
\end{proof}
The lowering and residue normalization are correlated outputs of this
same channel, while a current source contributes its actual overlap.
Different eliminated supports cannot exchange these certificates, and
neither an indefinite pivot nor an open-channel resonance automatically
satisfies the hypotheses.

\begin{theorem}[A complete response determines the least source-visible sector]
\label{p11:thm:response-visible-inverse}
For a fixed retained QCD source frame and finite canonical Hermitian
split, the entire matrix response
\[
 M(z)=zI-H_{PP}
       +H_{PQ}(H_{QQ}-zI)^{-1}H_{QP}
     =zI-H_{PP}+\sum_\lambda\frac{R_\lambda}{\lambda-z},
 \quad R_\lambda\succeq0,
\]
determines its minimal visible Hermitian environment up to unitary
intertwining of the eliminated block and couplings.  Its dimension is
\[
 \boxed{\dim\mathcal E_{\min}
          =\sum_\lambda\operatorname{rank}R_\lambda.}
\]
An environment reducing subspace annihilated by $H_{PQ}$ is invisible to
this response and is not reconstructed.  A degenerate pole contributes
the rank of its full positive coefficient matrix.
\end{theorem}
\begin{proof}
Use the positive residue-factor construction and cyclic visible subspace
of Theorem~\ref{p1b:thm:response-minimal-visible-realization}.
Its hypotheses hold for this finite same-owner Schur split.
\end{proof}
This is an inverse certificate for complete exact matrix response data.
A few masses, diagonal correlators, selected source contractions, or
finite noisy samples do not constitute that data and leave an inverse
family.  The theorem does not reconstruct source-invisible QCD states,
determine a thermal ensemble, or replace the separate continuum and
resonance limits.  The minimal sector is relative to the declared tests;
a mode consumed by another observable remains in that observable's
owner.

\begin{corollary}[Complex resonance poles are Schur-denominator zeros]
\label{p11:cor:block-b-leading-vector-complex-poles}

After an open-channel limit, boundary values, and a specified analytic
continuation are supplied, resonance poles are the zeros of the continued
Schur/Feshbach denominator.  Their widths are determined by the limiting
imaginary self-energy.  No KSRF-derived fixed pole is asserted.
\phantomsection\label{p11:eq:block-b-vector-complex-poles}
\end{corollary}
\begin{proof}

Theorem~\ref{p11:thm:block-b-finite-pole-dressing} gives the finite Schur
identity.  A proved limiting boundary value and continuation across its cut
define the resonance sheet; pole positions are determinant zeros there.
Couplings entering the self-energy remain QCD matrix elements.
\end{proof}

\begin{theorem}[Topology/anomaly/pseudoscalar separation from the generating function]
\label{p11:thm:block-c-topology-pseudoscalar-separation}

The topological sector is specified by \(Z_L(\theta)\) and its cumulants; the pseudoscalar anomaly block is then built from the computed \(\chi_{\YM}\), chiral decay normalization, and neutral flavor mass matrix.  Channel counts do not determine \(\chi_{\YM}\).  Glue/pseudoscalar poles follow from the resulting coupled correlator matrix.
\phantomsection\label{p11:eq:block-c-topological-variance-factor}
\phantomsection\label{p11:eq:block-c-anomaly-mass-bridge}
\phantomsection\label{p11:eq:block-c-pseudoscalar-glue-pole-equation}
\end{theorem}
\begin{proof}

Differentiation of \(\log Z_L(\theta)\) gives topological cumulants.  The anomalous Ward identity couples the singlet pseudoscalar current to this sector.  Diagonalizing the combined two-point function yields poles.  The covariance counterfamily of Theorem~\ref{p11:thm:orthogonal-hurwitz-adjoint-winding-variance} excludes a universal count multiplier.
\end{proof}

\begin{theorem}[Block D: physical light-ray owner and optional finite fraction realization]
\label{p11:thm:block-d-finite-distribution-spectral-measures}
The physical distribution owner is the Wilson-completed light-ray/twist-two
family of Theorem~\ref{p11:thm:lf-wilson-line-identification}; its current and
stress-tensor moments obey Theorem~\ref{p11:thm:native-pdf-sum-rules} without
exact R3.  If the stronger exact \(\LFRep_L\) certificate is additionally
supplied, Appendix~\ref{p11:app:v264-exact-lfrep-positive-measures} constructs
the positive finite number/momentum measures and their exact simplex sum rule.
The latter are a finite realization and are not substituted for the
renormalized physical PDF.
\phantomsection\label{p11:eq:block-d-pdf-spectral-measure}
\phantomsection\label{p11:eq:block-d-moment-readout}
\end{theorem}
\begin{proof}
The first statement is Theorems~\ref{p11:thm:lf-wilson-line-identification} and
\ref{p11:thm:native-pdf-sum-rules}.  The optional finite statement is
Theorem~\ref{p11:thm:light-front-fraction-povms} together with the exact R3
hypothesis stated in Appendix~\ref{p11:app:v264-exact-lfrep-positive-measures}.
The support/positivity firewall follows because a renormalized light-ray
operator-valued distribution need not be the spectral measure of the finite
commuting fraction operators.
\end{proof}

\begin{theorem}[Interacting thermal block is finite and analytic]
\label{p11:thm:block-e-thermal-finite-trace}

For each finite \(L\), the interacting QCD thermal block is the exact finite trace of the supplied Hamiltonian/source member.  It is analytic for \(T>0\); susceptibilities are derivatives of \(\log Z_L\).  A crossover may be extracted from these derivatives, while a true transition requires the controlled infinite-level/volume limit.  No fixed shell temperature is part of the block.
\phantomsection\label{p11:eq:block-e-finite-thermal-trace}
\end{theorem}
\begin{proof}

Diagonalizing the finite Hamiltonian expresses the partition function as a finite sum of exponentials.  Positivity and analyticity follow immediately, as do derivative formulas.  Nonanalyticity can only arise from an accumulation of levels in the specified limit.
\end{proof}

\begin{theorem}[Multi-baryon transfer cumulants from the supplied Hamiltonian]
\label{p11:thm:block-f-nuclear-transfer-cumulants}

Given the finite multi-baryon Hamiltonians in sectors of baryon number \(A\), connected nuclear interactions are defined by the usual inclusion--exclusion/Moebius cumulants of their energies or logarithmic transfer amplitudes.  The algebraic cumulant identities are exact; numerical contact/OPE entries are not assumed and must be computed from the underlying QCD matrix elements.

\end{theorem}
\begin{proof}

Moebius inversion on the partition lattice separates connected from disconnected contributions uniquely.  Applying it to computed sector energies or log transfers yields the connected kernels.  Since the operation is algebraic, it introduces no new coefficient, but it cannot supply missing two-body matrix elements.
\end{proof}


% Source Paper VII-B, original line 5059; inherited proof bodies unchanged.
\subsection{finite spectrum and source-response tests}
\label{p11:sec:v2210-lit-readout-tests}
This section uses the complete finite QCD member and source dictionary
already exported by Paper~VII-A. It supplies numerical-error and reduction
tests for the actual spectrum and source jets. A certificate does not fill
in an unexecuted hadron, heavy-source or multi-baryon matrix.


% Source Paper VII-B, original line 5066; inherited proof bodies unchanged.
\subsubsection{What a finite correlation GEVP proves}
\begin{theorem}[Finite GEVP compression and certified source-frame gap]
\label{p11:thm:v2210-lit-gevp-frame-gap}
On the actual finite physical excited quotient let
$A=H-E_0\ge0$ have ordered eigenvalues
$\Delta_1\le\cdots\le\Delta_d$. Let $B:\mathbb C^N\to\mathcal H_{\rm exc}$
have the actual vacuum-orthogonal source columns, and put
\[
 G(t)=B^\dagger e^{-tA}B,
 \quad G(t_0)\succ0,\quad
 V=e^{-t_0A/2}BG(t_0)^{-1/2},\quad \tau=t-t_0>0.
\]
Then $V^\dagger V=I_N$ and the generalized eigenvalues of
$(G(t),G(t_0))$ are the eigenvalues of
$M(\tau)=V^\dagger e^{-\tau A}V$. Ordered decreasingly, they obey
\[
 \boxed{0<\lambda_j(M)\le e^{-\tau\Delta_j},\qquad
 -\tau^{-1}\log\lambda_j(M)\ge\Delta_j,\quad j\le N.}
\]
Thus these effective energies are upper bounds on the corresponding ordered
physical energies, not lower bounds on all unseen states.

For a full physical gap lower bound, require the actual complete source frame
$W=BB^\dagger\ge w_-I$ with certified $w_->0$. If
\[
 C_{\rm tot}(T)=\operatorname{Tr}(We^{-TA}),\qquad
 |\widehat C_{\rm tot}(T)-C_{\rm tot}(T)|\le\varepsilon_C,
 \quad C_+=\widehat C_{\rm tot}(T)+\varepsilon_C>0,
\]
then
\[
 \boxed{\Delta_1\ge
 \max\{0,T^{-1}\log(w_-/C_+)\}.}
\]
If a Hermitian computed compression has
$\|\widehat M-M\|\le\varepsilon_M$, each actual generalized eigenvalue is
within $\varepsilon_M$ of its ordered computed eigenvalue. Whenever
$0<\widehat\lambda_j-\varepsilon_M$, the corresponding effective energy
lies in
\[
 \left[-\tau^{-1}\log(\widehat\lambda_j+\varepsilon_M),
       -\tau^{-1}\log(\widehat\lambda_j-\varepsilon_M)\right].
\]
\end{theorem}
\begin{proof}
The Gram identity gives $V^\dagger V=I$, and ordinary congruence transforms
the GEVP into $M$. The min--max principle for compression of the positive
matrix $e^{-\tau A}$ gives the first bounds. A normalized lowest eigenvector
$e_1$ has $\langle e_1,We_1\rangle\ge w_-$, hence
$C_{\rm tot}(T)\ge w_-e^{-T\Delta_1}$; taking logarithms and using the
certified upper bound gives the gap bound. Hermitian eigenvalue perturbation
and monotonicity of $-\log$ give the final interval.
\end{proof}
This is a finite counterpart of correlation-matrix extraction
technology~\cite{LitQCDBlossier2009GEVP} and sharpens the error-carrying
execution of Section~\ref{p11:sec:v2201-derived-d23}. An incomplete frame
has no strictly positive lower frame bound on the entire excited quotient.
A large-volume gap conclusion needs uniform physical frame/correlation
estimates and the actual transfer clock; finite source ranks supply neither.
The framework's native carrier-support gap and optional continuous
Yang--Mills gap remain distinct.

\begin{proposition}[Nonreducing source compression has a measurable return]
\label{p11:prop:v2210-lit-semigroup-leakage}
For the isometry $V$ above, let $P=VV^\dagger$, $Q=I-P$ and
$\delta=\|PAQ+QAP\|$. Then
\[
 \boxed{\|V^\dagger e^{-\tau A}V
       -e^{-\tau V^\dagger AV}\|\le\tfrac12\tau^2\delta^2.}
\]
This compares the actual source compression with the exponential of its
compressed Hamiltonian; it does not replace the former by the latter.
\end{proposition}
\begin{proof}
Put $A_0=PAP+QAQ\ge0$ and $E=A-A_0=PAQ+QAP$.
Both semigroups are contractions. Duhamel's identity gives
\[
 e^{-\tau A}-e^{-\tau A_0}
 =-\int_0^\tau e^{-(\tau-s)A_0}Ee^{-sA}\,ds.
\]
Expand $e^{-sA}$ once more by the same identity. The single-insertion term
vanishes between $P$ on both sides because $E$ is off diagonal. The
remaining double integral has norm at most
$\int_0^\tau\int_0^s\delta^2\,dr\,ds=\tau^2\delta^2/2$.
On $\operatorname{ran}P$, $A_0$ is unitarily represented by
$V^\dagger AV$, proving the assertion.
\end{proof}
Together with Proposition~\ref{p11:prop:v199-compression-defect}, this is an
explicit test of a proposed source-range reduction. Small geometric rank
alone says nothing about $\delta$ or the physical axial/source defect.


% Source Paper VII-B, original line 5356; inherited proof bodies unchanged.
\subsection{Geometry-first consequences and physical scope}
\label{p11:sec:v2201-geometry-first-consequences}
\paragraph{Geometry ancestry.}
Finite color/matter operator and physical pairing $\to$ specified transfer/source or thermal density $\to$ spectral and thermodynamic certificates.
\paragraph{Status.}
These are consequences on the domains stated below.
A mathematical certificate or an executed restriction retains every
physical source, state, clock and readout hypothesis of its inputs.


% Source Paper VII-B, original line 5365; inherited proof bodies unchanged.
\subsubsection{D23: A source-visible gap has a finite positive certificate}
\label{p11:sec:v2201-derived-d23}
\noindent\textit{Status: proved result on stated domain.}

For a positive physical transfer Hamiltonian $H-E_0$, a source correlation
has spectral form $C(T)=\sum_rw_re^{-\Delta_rT}$, $w_r\ge0$.
Its effective gap obeys
\[
 m_{\rm eff}(T)=-\partial_T\log C(T),\qquad
 m_{\rm eff}'(T)=-\operatorname{Var}_T(\Delta_r)\le0.
\]
Thus an early effective gap is generally an upper bound on the lowest
visible energy, not a lower gap proof. If the lowest visible weight is
independently at least $w_0>0$, then
\[
 \boxed{\Delta_{\rm vis}\ge
 \frac1T\log\frac{w_0}{C(T)}.}
\]
\begin{proof}
The effective gap is the tilted positive spectral mean.
For the second inequality use $C(T)\ge w_0e^{-\Delta_{\rm vis}T}$.
\end{proof}
For a full gap claim, a complete source frame
$W=\sum_j|j\rangle\langle j|\ge w_0I$ on the excited physical quotient
gives the same lower bound with
$C_{\rm tot}(T)=\operatorname{Tr}(We^{-T(H-E_0)})$.
A uniform large-volume conclusion consumes a uniform frame bound and
correlation estimate; a finite count does not supply either.

On the conditional single-loop Casimir transfer, the smallest nontrivial
singlet class-function sources $\chi_3,\chi_{\bar3}$ have
$C(n\tau)=e^{-2nt/3}$, hence exact visible gap $2t/(3\tau)$.
This executes a minimal gauge-invariant transfer block when its
cotangent/transfer owner is licensed. It is not a hadron mass, an area law
or the full QCD gap, and $t$ and $\tau$ remain distinct parameters.
\par\noindent\textit{Inherited inputs:} \ref{p10:thm:finite-qcd-matrix-compiler}, \ref{p10:rem:three-distinct-heat-parameters}, \ref{p10:rem:v2159-electric-heat-scope}, \ref{p10:thm:canonical-haar-heat-representative}.



\section{Baryon spin, currents, chiral residues, and form factors}
\label{p11:sec:v237-b4}
Source visibility and tensor geometry do not determine an unsupplied current coefficient. Charged and baryonic physical pole readouts require the actual Hamiltonian and source matrix elements.

% Source Paper VII-B, original line 371; inherited proof bodies unchanged.
\subsubsection{Nucleon scalar, axial, and nuclear bridges}
\label{p11:subsec:nucleon-scalar-axial-nuclear-bridges}

\begin{theorem}[Nucleon mass trace-anomaly bridge]
\label{p11:thm:nucleon-trace-anomaly-bridge}
On the matched relativistic one-nucleon comparison branch, with the
declared normalization and with the full finite scale residual of
Theorem~\ref{p11:thm:thermal-trace-anomaly-operator-closure} proved absent
or absorbed into the matched operators, the mass is fixed by the trace of
the QCD energy-momentum tensor:
\begin{equation}
\label{p11:eq:nucleon-trace-anomaly-bridge}
M_N=\frac1{2M_N}\left[
\left\langle N\left|\frac{\beta(g_s)}{2g_s}F_{\mu\nu}^aF^{a\mu\nu}\right|N\right\rangle
+\sum_q m_q(1+\gamma_m)\langle N|\bar q q|N\rangle
\right].
\end{equation}
Without that residual certificate, the displayed bracket additionally
contains $\langle N|\mathcal R_{\rm scale}|N\rangle$ in the same
normalization. The native finite mass remains the independently defined
baryon-sector spectral gap. On the matched branch the bridge reduces it to
computed trace and scalar matrix elements~\cite{CollinsDuncanJoglekar1977,ShifmanVainshteinZakharov1978}.
\end{theorem}
\begin{proof}
The QCD scale Ward identity gives \(T^\mu{}_{\mu}=\frac{\beta(g_s)}{2g_s}F^2+\sum_qm_q(1+\gamma_m)\bar q q\).  For a relativistically normalized one-nucleon state, \(\langle N(p)|T^{\mu\nu}|N(p)\rangle=2p^\mu p^\nu\).  Taking the trace gives \(\langle N|T^\mu{}_{\mu}|N\rangle=2M_N^2\), hence \eqref{p11:eq:nucleon-trace-anomaly-bridge}.\qedhere
\end{proof}

\begin{theorem}[Native scalar-response bridge]
\label{p11:thm:native-scalar-response-bridge}
Let the finite hadron gap be \(M_h(m_q)=E_h(m_q)-E_0(m_q)>0\) on differentiable isolated state/vacuum branches. The exact finite scalar response is the \emph{vacuum-subtracted} Feynman--Hellmann identity
\[
\frac{\partial M_h}{\partial m_q}
=\langle h|\partial_{m_q}H|h\rangle_{\rm norm}
-\langle\Omega|\partial_{m_q}H|\Omega\rangle_{\rm norm}.
\]
When the scalar insertion is renormalized and normalized covariantly to the same vacuum-subtracted mass observable, the relativistically normalized form is:
\begin{equation}
\label{p11:eq:feynman-hellmann-scalar-bridge}
\frac{\partial M_h}{\partial m_q}=\frac1{2M_h}\langle h|\bar q q|h\rangle_{\rm conn,vac\,sub},
\qquad
\sigma_{qh}=m_q\frac{\partial M_h}{\partial m_q}.
\end{equation}
Therefore sigma terms cannot be inserted as independent phenomenological coefficients; they are native response coefficients of the hadron eigenvalue problem~\cite{Feynman1939,Hellmann1937}.
\end{theorem}
\begin{proof}
For a finite quark-mass source, \(\partial_{m_q}H=\int d^3x\,\bar q q\) on its declared renormalized branch. Apply Feynman--Hellmann separately to \(H\Psi_h=E_h\Psi_h\) and \(H\Omega=E_0\Omega\), then subtract: \(\partial_{m_q}(E_h-E_0)=\langle h|H_q'|h\rangle-\langle\Omega|H_q'|\Omega\rangle\) in unit normalization. The displayed covariant matrix element is the connected, vacuum-subtracted version multiplied by the compatible external-state normalization \(2M_h\); raw \(\langle h|\bar q q|h\rangle\) without vacuum subtraction would be incorrect if \(E_0\) depends on \(m_q\). Multiply by \(m_q\) for the sigma term.\qedhere
\end{proof}


\begin{theorem}[RG-invariant scalar bridge]
\label{p11:thm:rg-invariant-scalar-bridge}
The renormalized quark mass and scalar density have opposite anomalous
dimensions,
\begin{equation}
\label{p11:eq:mass-scalar-opposite-running}
\mu\frac{d m_q}{d\mu}=-\gamma_m m_q,
\qquad
\mu\frac{d}{d\mu}(\bar q q)_\mu=+\gamma_m(\bar q q)_\mu.
\end{equation}
Hence the RG-invariant scalar products are
\[
  m_q(\mu)(\bar q q)_\mu,\qquad
  m_q(\mu)\Sigma_q(\mu),\qquad
  \sigma_{qh}=m_q\frac{\partial M_h}{\partial m_q}.
\]
The individual factors remain convention dependent.
\end{theorem}
\begin{proof}
Equation~\eqref{p11:eq:mass-scalar-opposite-running} gives \(\mu d[m_q(\bar q q)_\mu]/d\mu=0\). GMOR and the scalar part of the trace-anomaly bridge are therefore scheme-safe only when written in invariant products. The full trace \(T^\mu{}_{\mu}\) is physical, while its quark/gluon split is a renormalization-convention decomposition.
\end{proof}

\begin{theorem}[Goldberger--Treiman consequence of computed current residues]
\label{p11:thm:native-goldberger-treiman-bridge}
Assume the controlled chiral regime, an isolated nucleon pole, PCAC, and a pion pole in the axial channel.  Then the computed correlator readouts obey
\begin{equation}
\label{p11:eq:goldberger-treiman-bridge}
g_A M_N=f_\pi g_{\pi NN}+O(m_q).
\end{equation}
Hence, once \(M_N\), \(f_\pi\), and \(g_A\) have been obtained from the supplied QCD Hamiltonian/current correlators, the leading pion--nucleon coupling is a dependent readout rather than an additional coefficient~\cite{GoldbergerTreiman1958}.  This identity does not determine any of those three owning correlator quantities from finite counts.
\end{theorem}
\begin{proof}
The nucleon axial matrix element defines \(g_A=G_A(0)\). The divergence of the axial current has a pion pole with residue proportional to \(f_\pi g_{\pi NN}\). Matching this pole to the zero-momentum axial response gives \eqref{p11:eq:goldberger-treiman-bridge}, up to explicit chiral-breaking corrections.
\end{proof}


\begin{theorem}[Symmetric spin--flavor benchmark for the nucleon axial charge]
\label{p11:thm:leading-native-nucleon-axial-charge}

If the lowest nucleon state is additionally assumed to lie in the completely symmetric two-flavor \(SU(6)\)-type S-wave spin--flavor representation, then
\[
\Delta u=\frac43,\qquad \Delta d=-\frac13,\qquad g_A=\frac53.
\]
This is a representation benchmark, not a theorem that the compiled QCD ground state realizes that representation.  The physical \(g_A\) is the axial-current matrix element of the computed nucleon state.
\phantomsection\label{p11:eq:leading-native-ga}
\phantomsection\label{p11:eq:leading-native-ga-five-thirds}
\end{theorem}
\begin{proof}

The symmetric three-quark spin--flavor wavefunction gives the standard matrix element of \(\sum_k\sigma_3^{(k)}\tau_3^{(k)}\), yielding the displayed values.  QCD dynamics can mix other orbital/gluonic components, so realization of this benchmark is an extra state hypothesis and must be checked spectrally.
\end{proof}

\begin{theorem}[One-pion-exchange boundary from computed chiral inputs]
\label{p11:thm:native-one-pion-exchange-nuclear-boundary}
In the two-nucleon sector, the long-range isovector force is the pion pole of
the native axial current.  At leading chiral order its momentum-space kernel is
\begin{equation}
\label{p11:eq:native-one-pion-exchange-potential}
V_{1\pi}(\vec q)
=
-\frac{g_A^2}{4F_\pi^2}
\frac{(\vec\sigma_1\!\cdot\!\vec q)(\vec\sigma_2\!\cdot\!\vec q)}
{\vec q^{\,2}+m_\pi^2}
(\vec\tau_1\!\cdot\!\vec\tau_2),
\end{equation}
with \(F_\pi\), \(m_\pi\), and \(g_A\) obtained from the owning chiral/current readouts.  The shorter-range two-baryon constants are the analytic threshold
coefficients of the native \(BB\to BB\) transfer amplitude,
\begin{equation}
\label{p11:eq:native-nuclear-contact-readouts}
C_\alpha^{NN}
=
\mathcal P_\alpha
\left.
\mathcal T_{NN\to NN}(\vec p\,',\vec p)
\right|_{\rm analytic\ at\ \vec q=0},
\end{equation}
and are not independent nuclear phenomenology inputs.
\end{theorem}
\begin{proof}
PCAC and the Goldberger--Treiman bridge identify the pion--nucleon vertex with
\(g_A/(2F_\pi)\) times the derivative pion field.  Exchanging
the light pion between two native nucleon eigenstates gives the pole factor in
\eqref{p11:eq:native-one-pion-exchange-potential}, with spin and isospin matrices
fixed by the two-nucleon representation.  After the pion pole and other resolved
light exchanges are subtracted, the remaining low-momentum amplitude is analytic
and decomposes into the allowed contact operator basis; the projection
\(\mathcal P_\alpha\) extracts the coefficients.  Since the amplitude is a
matrix element of the supplied QCD Hamiltonian between computed baryon states, those contacts are derived readouts once the regulator/action and matching prescription are fixed; they are not fixed by geometry alone.
\end{proof}



% Source Paper VII-B, original line 505; inherited proof bodies unchanged.
\subsubsection{Baryon topology and multiplet quantization}
\label{p11:subsec:baryon-topology-multiplet-quantization}

\begin{theorem}[Microscopic--topological baryon matching under a current-matched chiral map]
\label{p11:thm:microscopic-topological-baryon-matching}
The microscopic and chiral descriptions of baryon number agree \emph{provided} a same-action low-energy matching map has been supplied that sends the normalized microscopic baryon current to the chiral winding current, with matching charge orientation and no unmatched baryon-current contact.  The microscopic color-singlet baryon support is
\begin{equation}
\label{p11:eq:microscopic-baryon-color-singlet}
\wedge^3 C\cong\mathbf1,
\end{equation}
while the aligned chiral sector has a Goldstone field
\begin{equation}
\label{p11:eq:chiral-winding-map}
U:S^3\to SU(3).
\end{equation}
Since \(\pi_3(SU(3))\cong\ZZ\), the chiral theory has integer topological charge
\begin{equation}
\label{p11:eq:topological-baryon-number}
B_{\rm top}
=
\frac{1}{24\pi^2}
\int d^3x\,\epsilon^{ijk}
\operatorname{Tr}\!\bigl[
(U^{-1}\partial_iU)(U^{-1}\partial_jU)(U^{-1}\partial_kU)
\bigr].
\end{equation}
On the stated current-matched branch, the low-energy chiral image of the microscopic baryon current is the normalized winding current, so its charge is \(B_{\rm top}\). Thus \(B=(n_q-n_{\bar q})/3=B_{\rm top}\) holds \emph{on that branch}; the determinant color invariant and \(\pi_3(SU(3))\cong\ZZ\) separately do not construct the matching map.
\end{theorem}
\begin{proof}
The microscopic statement is the determinant invariant of the three-dimensional color carrier.  The macroscopic statement follows from finite-energy compactification of the three-dimensional configuration domain and from \(\pi_3(SU(3))\cong\ZZ\), as in the Skyrme/WZW construction~\cite{Skyrme1961,Witten1983Baryon,AdkinsNappiWitten1983}.  The microscopic conserved current is \(J_B^\mu=\frac13\sum_f\bar q_f\gamma^\mu q_f\); the effective winding current is identically conserved. These independent conservation statements and integer classifications do not identify their charges. The theorem's conditional conclusion follows only after the declared same-action/current matching carries the microscopic charge into the winding-current class with the stated normalization.
\end{proof}

\begin{theorem}[WZW baryon multiplet quantization]
\label{p11:thm:wzw-baryon-multiplet-quantization}
For a chiral soliton of winding number \(B\), collective \(SU(3)\) quantization in the presence of the WZW term with coefficient \(N_c=3\) imposes
\begin{equation}
\label{p11:eq:wzw-right-hypercharge-constraint}
Y_R=\frac{N_cB}{3}.
\end{equation}
For \(B=1\), this gives \(Y_R=1\).  The lowest admissible flavor representations are
\begin{equation}
\label{p11:eq:baryon-octet-decuplet}
\mathbf8_{J=1/2}
\qquad\text{and}\qquad
\mathbf{10}_{J=3/2}.
\end{equation}
Thus the baryon octet/decuplet representation structure is fixed by the same derived color multiplicity that fixes the WZW anomaly coefficient~\cite{Witten1983WZW,Witten1983Baryon}.
\end{theorem}
\begin{proof}
Write the slowly rotating soliton as \(U(x,t)=A(t)U_0(x)A(t)^{-1}\), with \(A(t)\in SU(3)\).  The WZW term contributes a first-order collective-coordinate term whose coefficient is \(N_cB\).  Canonical quantization gives the right-hypercharge constraint.  For \(N_c=3\) and \(B=1\), the lowest \(SU(3)\) representations containing the required state are \((p,q)=(1,1)\), the octet, and \((p,q)=(3,0)\), the decuplet.  The collective right-isospin assignments give \(J=1/2\) and \(J=3/2\), respectively.
\end{proof}

\begin{corollary}[Leading baryon mass relations]
\label{p11:cor:leading-baryon-mass-relations}
Treating the light-quark mass matrix as an \(SU(3)\)-breaking spurion, first-order flavor breaking in the WZW-quantized baryon multiplets gives the octet relation
\begin{equation}
\label{p11:eq:baryon-gell-mann-okubo}
2(M_N+M_\Xi)=3M_\Lambda+M_\Sigma,
\end{equation}
and the decuplet equal-spacing rule
\begin{equation}
\label{p11:eq:decuplet-equal-spacing}
M_{\Sigma^\ast}-M_\Delta
=
M_{\Xi^\ast}-M_{\Sigma^\ast}
=
M_\Omega-M_{\Xi^\ast}.
\end{equation}
These are representation-theoretic mass relations; the absolute masses and higher-order splittings remain native spectral targets~\cite{GellMann1961Eightfold,Okubo1962}.
\end{corollary}
\begin{proof}
In the isospin limit, the leading octet mass operator has the form \(M_8=M_0+aY+b[I(I+1)-Y^2/4]\).  Substitution of \((Y,I)=(1,1/2),(0,0),(0,1),(-1,1/2)\) for \(N,\Lambda,\Sigma,\Xi\) gives the octet relation.  For the decuplet, the leading breaking is linear in hypercharge, and the hypercharges of \(\Delta,\Sigma^\ast,\Xi^\ast,\Omega\) differ by one, giving equal spacing.
\end{proof}

\begin{theorem}[Large-color scaling is asymptotic, not an exact Nc=3 identity]
\label{p11:thm:nc3-large-color-scaling}

The usual large-\(N_c\) power counting is an asymptotic organization of a family of gauge theories as \(N_c\to\infty\).  At \(N_c=3\) it supplies an ordering estimate only; it does not fix an absolute hadronic mass, residue, or a Paper-VII-B normalization.
\phantomsection\label{p11:eq:nc-scaling-specialized}
\end{theorem}
\begin{proof}

An asymptotic statement \(X(N_c)=O(N_c^p)\) leaves an arbitrary dimensionless coefficient and subleading terms.  Evaluating the order symbol at \(N_c=3\) therefore cannot yield an equality.  Current Paper~VII-A also exports no universal nonperturbative scale with which such an equality could be normalized.
\end{proof}



% Source Paper VII-B, original line 2708; inherited proof bodies unchanged.
\subsubsection{Chiral execution and the same-action nonvanishing certificate}
\label{p11:sec:chiral-certificate-handoff}
Paper~VII-A now supplies the weighted mode-window lower bound
of Theorem~\ref{p10:thm:weighted-chiral-mode-window} and the direct
energy/source criterion of
Theorem~\ref{p10:thm:chiral-energy-source-certificate}.
On the compatible pressure-limit domain,
Theorem~\ref{p10:thm:chiral-pressure-energy-equivalence} makes
positive scalar expectation at vanishing excess energy density an
equivalent test. The source-moment theorem is a sufficient route
to that test, not another constitutive requirement.

Execution consumes the actual compiled singlet Hamiltonian, scalar
source, state, normalization and ordered reciprocal readout sequence.
The finite Schur response retains
Theorem~\ref{p10:thm:complete-schur-scalar-source}'s four source
terms and inherited norm. Source/determinant transport retains its
existing contacts and sea response. No reference to a flavor-orbit
trial selects a physical QCD state.
The actual positive limiting response remains to be established.
The axial residues $F_\pi$, the ratio $B_0=\Sigma/F_\pi^2$ and the
matched chiral constants remain downstream correlator outputs.
The renormalization convention in
Theorem~\ref{p11:thm:native-chiral-running} supplies no missing
condensate or pole normalization.

% END INLINE SOURCE: Paper7B-ChiralCertificateHandoff-Main-v1_0.tex


\begin{theorem}[Renormalization-group running with arbitrary matching scale]
\label{p11:thm:native-chiral-running}

For each renormalized chiral constant,
\[
L_i^r(\mu_2)=L_i^r(\mu_1)+\frac{\Gamma_i}{16\pi^2}\log\frac{\mu_1}{\mu_2}
\]
with the standard convention-dependent coefficients \(\Gamma_i\).  The matching scale \(\mu_1\) is a renormalization convention inside the EFT window; it is not uniquely fixed to the first vector pole.
\phantomsection\label{p11:eq:native-chiral-scale}
\phantomsection\label{p11:eq:native-Li-running}
\phantomsection\label{p11:eq:native-chiral-logs-numeric}
\end{theorem}
\begin{proof}

Renormalization-scale independence of physical amplitudes gives the RG equation and its integrated form.  Choosing a convenient matching scale reorganizes loops and local terms but cannot be selected by symmetry alone.
\end{proof}

\begin{corollary}[One-loop decay-constant ratio as a conditional chiral formula]
\label{p11:cor:native-one-loop-FK-Fpi}

The one-loop expression for \(F_K/F_\pi\) is retained in its standard form as a function of the renormalized \(L_5^r(\mu)\), pseudoscalar masses, \(F_0\), and chiral logarithms.  No native decimal follows until those correlator-derived inputs and the renormalization convention are supplied.
\phantomsection\label{p11:eq:native-FK-Fpi-one-loop}
\phantomsection\label{p11:eq:native-FK-Fpi-one-loop-numeric}
\end{corollary}
\begin{proof}

This is direct evaluation of the one-loop chiral generating functional.  The scale dependence of the logarithms cancels that of \(L_5^r\); therefore a numerical value is meaningful only for a complete matched input set.
\end{proof}

\begin{corollary}[One-loop pseudoscalar masses retain analytic low-energy constants]
\label{p11:cor:native-one-loop-mass-logs}

The one-loop pseudoscalar mass formulas contain both chiral logarithms and symmetry-allowed analytic \(L_i^r\) combinations.  These analytic terms cannot be set to zero from finite support geometry; they remain computed/matched QCD response coefficients.
\phantomsection\label{p11:eq:native-one-loop-mass-logs-numeric}
\end{corollary}
\begin{proof}

Renormalization of the loop divergences requires the local \(O(p^4)\) counterterms.  Setting their finite parts to special values is a renormalization/matching choice unless independently derived.  Therefore a log-only numerical mass is not a theorem.
\end{proof}

\begin{corollary}[Vector strong widths from computed couplings]
\label{p11:cor:native-vector-strong-widths}

For a vector resonance with a specified on-shell coupling to two pseudoscalars,
the standard phase-space formula determines its partial width on the
\textnormal{(SM1)--(SM3)} plus \textnormal{(SM5)} open-channel domain, together
with the relevant phase-space and channel-normalization parts of
\textnormal{(SM7)}, in
Definition~\ref{p12:def:claim-indexed-smat-package}; \textnormal{(SM4)} is also
used when the on-shell final-state coupling is identified by LSZ.  The vector
is the continued pole, not a stable asymptotic input; \textnormal{(SM2)} applies
to the declared stable external channel particles.  KSRF and ideal-flavor
relations may be tested as dynamical/saturation approximations but are not
consequences of the finite carrier.
\phantomsection\label{p11:eq:native-g-rho-pi-pi}
\phantomsection\label{p11:eq:native-vector-widths-numeric}
\end{corollary}
\begin{proof}

On the named scattering domain, the rate follows from the squared on-shell
matrix element and two-body phase space.  The coupling is obtained from a
residue/three-point function of the supplied QCD source algebra, but the LSZ and
channel normalization are additional claim-local hypotheses.  Current
conservation and representation labels do not fix its magnitude.
\end{proof}

\begin{remark}[Open vector three-point execution]
\label{p11:rem:vector-vertex-reclosure-targets}
The phase-space width relations in
Corollary~\ref{p11:cor:native-vector-strong-widths} are the leading conditional
channel formulas.  Flavor-dependent corrections to \(g_{K^*K\pi}\) and
\(g_{\phi KK}\) require the corresponding finite three-point current/transfer
matrix elements.  No additional strange-leg factor is assigned before that
calculation.
\end{remark}


% Source Paper VII-B, original line 2813; inherited proof bodies unchanged.
\subsubsection[Native baryon spin--gluon and dressed-residue closure blocks]
{Baryon spin--gluon and dressed residues}
\label{p11:subsec:native-baryon-dressed-closures}


\begin{lemma}[Geometric 8/3 second moment is not an energy weight]
\label{p11:lem:minimal-transverse-spin-glue-sheet}

The finite Hurwitz/color geometry may give an exact isotropic second-moment coefficient \(8/3\).  This coefficient belongs to the relevant tensor/source contraction.  It does not by itself equal an additive Hamiltonian energy or a probability weight.
\phantomsection\label{p11:eq:minimal-transverse-spin-glue-weight}
\end{lemma}
\begin{proof}

Multiplying the physical operator associated with that tensor structure by any scalar leaves the normalized geometric second moment unchanged while rescaling its energy contribution.  Therefore an operator/Hamiltonian coefficient must be computed separately.
\end{proof}

\begin{lemma}[Wigner--Eckart form of baryon current projections]
\label{p11:lem:dressed-determinant-current-projections}

Within a fixed flavor/spin multiplet, vector, axial, tensor and magnetic current matrix elements factor into known Clebsch--Gordan coefficients and independent reduced matrix elements.  Charge conservation fixes only the conserved vector charge normalization.  Fractions such as \(3/4\), \((3/4)^2\), or \(1-\delta\) are not fixed by support counting.
\phantomsection\label{p11:eq:dressed-determinant-current-projections}
\end{lemma}
\begin{proof}

Apply Wigner--Eckart to each current operator.  Different irreducible operators have independent reduced matrix elements.  Ward identities constrain the conserved vector generator at zero momentum but not axial/tensor overlaps or magnetic slopes, proving the stated separation.
\end{proof}

\begin{theorem}[Baryon spin/flavor spectrum as a finite block eigenproblem]
\label{p11:thm:native-baryon-spin-gluon-ladder}

The determinant color singlet fixes the baryon color support.  For each spin/flavor sector \(\chi\), the physical finite poles are the eigenvalues of \(P_{\chi,L}(H_L-E_{0,L})P_{\chi,L}\).  Spin--glue and hyperfine splittings are off-diagonal/diagonal matrix elements of the supplied Hamiltonian and are not fixed by count-derived rungs.
\phantomsection\label{p11:eq:baryon-determinant-spin-glue-components}
\phantomsection\label{p11:eq:native-baryon-anchor}
\phantomsection\label{p11:eq:native-delta-anchor}
\phantomsection\label{p11:eq:octet-strange-recoupling-operator}
\phantomsection\label{p11:eq:octet-strange-recoupling-projectors}
\phantomsection\label{p11:eq:octet-strange-recoupling-eigenvalues}
\phantomsection\label{p11:eq:native-octet-mass-ladder}
\phantomsection\label{p11:eq:decuplet-strange-recoupling-operator}
\phantomsection\label{p11:eq:native-decuplet-mass-ladder}
\phantomsection\label{p11:eq:native-baryon-ladder-numeric}
\phantomsection\label{p11:eq:native-baryon-isospin-splittings}
\end{theorem}
\begin{proof}

Conserved quantum numbers block-diagonalize the self-adjoint Hamiltonian.  The spectral theorem gives the poles in each block.  Independent admissible changes of the hyperfine operator preserve the determinant color support while changing splittings, so no universal ladder follows from support alone.
\end{proof}

\begin{remark}[Finite baryon recoupling boundary]
\label{p11:rem:baryon-recoupling-status}
The rational coefficients in \eqref{p11:eq:native-octet-mass-ladder} and \eqref{p11:eq:native-decuplet-mass-ladder} are leading eigenvalues of the finite spin--flavor recoupling block.  Physical pole masses require the full finite baryon Hamiltonian diagonalization of Theorem~\ref{p11:thm:finite-hadron-spectrum-closure}.
\end{remark}

\begin{theorem}[Baryon residues and scalar responses from source derivatives]
\label{p11:thm:native-dressed-baryon-residues}

For a normalized baryon pole \(|B,L\rangle\), every current residue is a matrix element of its supplied current operator, and every scalar response is a Feynman--Hellmann derivative of \(M_{B,L}\) with respect to the physical source.  No universal determinant dressing factor multiplies all currents.
\phantomsection\label{p11:eq:native-dressed-axial-tensor}
\phantomsection\label{p11:eq:native-g-piNN-dressed}
\phantomsection\label{p11:eq:native-magnetic-zs-zv}
\phantomsection\label{p11:eq:native-dressed-magnetic-isoscalar-isovector}
\phantomsection\label{p11:eq:native-dressed-magnetic-moments}
\phantomsection\label{p11:eq:native-dressed-anomalous-magnetic-moments}
\phantomsection\label{p11:eq:native-dressed-sigma-terms}
\phantomsection\label{p11:eq:native-dressed-scalar-fractions}
\end{theorem}
\begin{proof}

This is the definition of the source-generated correlator and the Feynman--Hellmann theorem.  Since distinct currents are distinct operators, a common multiplicative factor would require an additional operator identity; none follows from color antisymmetry.
\end{proof}

\begin{theorem}[Ward identity fixes charge but not radius]
\label{p11:thm:finite-em-current-radius-reclosure}

For a hadron electromagnetic form factor \(F_h(q^2)\), the Ward identity fixes \(F_h(0)=Q_h\).  The charge radius is the independent slope (with the usual sign convention), e.g.
\[
\langle r_h^2\rangle=6\,F_h'(0)/Q_h
\]
for nonzero charge after the chosen \(q^2\) convention.  It must be computed from the current matrix element.
\phantomsection\label{p11:eq:finite-em-current-radius-reclosure-form-factor}
\phantomsection\label{p11:eq:finite-em-current-radius-xi-values}
\phantomsection\label{p11:eq:em-radius-reclosure-derived}
\phantomsection\label{p11:eq:em-radius-reclosure-numerics}
\end{theorem}
\begin{proof}

Analytic families \(F(q^2)=Q+cq^2+O(q^4)\) satisfy the same Ward identity for arbitrary \(c\).  Therefore charge normalization does not determine the radius.  Neutral states analogously permit arbitrary leading slope.
\end{proof}


\begin{corollary}[Octet axial currents have two reduced SU(3) couplings]
\label{p11:cor:native-octet-axial-current-block}

In exact flavor \(SU(3)\), octet-baryon axial matrix elements are determined by Clebsch--Gordan coefficients multiplying two reduced couplings conventionally denoted \(F\) and \(D\).  These reduced couplings are QCD current matrix elements and must be computed; no determinant-support factor fixes their ratio.
\phantomsection\label{p11:eq:native-FD-block}
\phantomsection\label{p11:eq:native-hyperon-axial-ratios}
\phantomsection\label{p11:eq:native-BBP-axial-coefficients}
\phantomsection\label{p11:eq:native-BBP-GT-extension}
\end{corollary}
\begin{proof}

The octet appears twice in \(8\otimes8\), giving two independent reduced octet matrix elements.  Wigner--Eckart therefore requires two constants.  Support counting cannot reduce this multiplicity without an additional dynamical identity.
\end{proof}


% Source Paper VII-B, original line 5434; inherited proof bodies unchanged.
\subsection{baryon spin--flux operator versus physical poles}
\label{p11:sec:v2220-baryon-flux-readouts}
\noindent\textit{Status: precise source/transfer readout, with the physical
spectral and continuum calculations explicitly still outstanding.}

Paper~VII-A constructs the bounded, gauge-invariant Hurwitz face
observable $\mathcal F_H$ in
\eqref{p10:eq:v2220-global-observable}.  The exact conditional
fixed-link comparison in
\eqref{p10:eq:v2220-first-spinflux} is a coefficient of a
\emph{squared one-body-based three-quark kinetic transition amplitude},
not an energy eigenvalue.  Its sign can vary with the $SU(3)$ links.
The pure-Haar values $1/9$ and $11/15552$, the sharp local bounds
$[-3/4,9/2]$, and the comparison coefficient $x^8$ are therefore
not baryon masses or physical hyperfine coefficients.

\begin{proposition}[Physical singlet spectral readout]
\label{p11:prop:v2220-physical-baryon-readout}
Let $H_L$ be the actually instantiated, self-adjoint color-singlet
finite QCD Hamiltonian of Paper~VII-A, including the inherited physical
mass sources and scale/coupling transport, and let
$K_L(T)=e^{-T(H_L-E_{0,L})}$ be its positive Euclidean transfer.
For a gauge-invariant baryon interpolator $\mathcal B_{J,f}$, the
finite-volume correlator is
\begin{equation}
\label{p11:eq:v2220-physical-correlator}
 C_{J,f}(T)=\langle0|\mathcal B_{J,f}K_L(T)
 \mathcal B_{J,f}^{\dagger}|0\rangle
 =\sum_n |Z_{n,J,f}|^2e^{-T(E_{n,J,f}-E_{0,L})}.
\end{equation}
If its lowest visible eigenlevel is isolated and has nonzero source
overlap, its energy is selected by the large-$T$ exponential.
When two such stable levels in the same flavor sector can be isolated,
\begin{equation}
\label{p11:eq:v2220-stable-hyperfine}
 E_{3/2,f}-E_{1/2,f}
 =-\lim_{T\to\infty}\frac1T
 \log\!\frac{C_{3/2,f}(T)}{C_{1/2,f}(T)}.
\end{equation}
For an unstable resonance, this real finite-volume limit need not
select the resonance mass: the continued scattering-channel pole and
residue prescription of this paper remains compulsory.
\end{proposition}
\begin{proof}
Use the spectral theorem for the finite self-adjoint $H_L$ and expand
each source-created state into the physical eigenbasis.  The positive
weights in \eqref{p11:eq:v2220-physical-correlator} select the lowest
nonzero-overlap level.  Taking the logarithmic ratio cancels the vacuum
energy.  An unstable resonance contributes through finite-volume
multi-particle eigenstates; its pole is a different analytic readout.
\end{proof}

\begin{proposition}[Bounded geometric source and Feynman--Hellmann test]
\label{p11:prop:v2220-flux-source-test}
On a finite gauge-link representation for which $\mathcal F_H$ is the
bounded self-adjoint multiplication operator from Paper~VII-A, introduce
the \emph{bookkeeping} source
\[
 H_L(\lambda)=H_L+\lambda\mathcal F_H.
\]
For any isolated, differentiable baryon eigenbranch in the $J,f$
sector (or its normalized isolated spin multiplet),
\begin{equation}
\label{p11:eq:v2220-flux-fh}
 \left.\partial_\lambda E_{J,f}(\lambda)\right|_{0}
 =\langle \mathcal F_H\rangle_{J,f}.
\end{equation}
In particular, whenever both spin branches are differentiable,
\begin{equation}
\label{p11:eq:v2220-flux-hyperfine}
 \left.\partial_\lambda(E_{3/2,f}-E_{1/2,f})\right|_0
 =\langle\mathcal F_H\rangle_{3/2,f}
  -\langle\mathcal F_H\rangle_{1/2,f}.
\end{equation}
These matrix elements are not fixed by the pure-gauge variance or by
the conditional comparison transfer.
\end{proposition}
\begin{proof}
Finite-dimensional analytic spectral perturbation and the
Feynman--Hellmann identity give the derivatives for a simple isolated
level.  For an isolated spin multiplet invariant under the scalar
source, use its normalized spectral-projector trace.  Taking the
difference proves \eqref{p11:eq:v2220-flux-hyperfine}.
\end{proof}



% BEGIN FRAMEWORK v2.235: SIGNED HURWITZ FACE-SOURCE CONSUMPTION
\begin{theorem}[Exact signed-mode compression of the declared face source]
\label{p11:thm:v235-signed-face-source-compression}
Work on the untruncated product-Haar link Hilbert space of
Proposition~\ref{p10:prop:v235-signed-haar-norms}, with the
orthogonal \(96\)-mode projections \(P_\chi,P_\psi\).
For the \emph{already declared bookkeeping source}
\(\mathcal F_H=96^{-1}\sum_fF(U_f)\) of
Proposition~\ref{p11:prop:v2220-flux-source-test}, the complete
\(192\)-mode compression is
\[
 \boxed{P_\chi\mathcal F_HP_\chi=\tfrac19P_\chi,
 \quad P_\psi\mathcal F_HP_\psi=\tfrac{85}{768}P_\psi,
 \quad P_\chi\mathcal F_HP_\psi=0.}
\]
The exact difference between the two projected coefficients is
\(1/9-85/768=1/2304\).  These are finite reference-source
compressions, not physical hadron-pole slopes of an undeformed action.
\end{theorem}
\begin{proof}
On any edge fan, the six detour links form \(K_{2,3}\) and contain
no complete triangular face.  The signed \(\chi_e\) therefore has
no nontrivial single-face compression.  Adding the central edge
produces precisely the three fan faces for \(\psi_e\).
For each fan face \(V_i\), Haar/Schur orthogonality gives
\[
 \begin{gathered}
 \mathbb E_0[\psi_e^2\mid V_i]
    =\frac{8-\chi_8(V_i)}{81},\qquad
 \mathbb E_0\psi_e^2=\frac8{81},\\
 \operatorname{Cov}_0(\psi_e^2,F(V_i))=-\frac1{729}.
 \end{gathered}
\]
The other \(93\) faces have covariance zero.  Dividing by the
norm and averaging over all \(96\) faces yields
\(\tfrac19+\tfrac{3(-1/729)}{96(8/81)}=\tfrac{85}{768}\).
Distinct signed edge modes have at least four exclusive links on
each side; a potential supported on one triangle cannot bridge
them, so the off-diagonal matrix elements vanish.
Color-conjugation parity eliminates the mixed blocks.
\end{proof}

\begin{corollary}[Adjoint-character selector for finite face potentials]
\label{p11:cor:v235-signed-adjoint-selection}
For any bounded real inversion-even, color-conjugation-even
\(SU(3)\) class function \(v\), define
\[
 \mathcal V_v=\frac1{96}\sum_fv(U_f),\qquad
 \bar v=\int_{SU(3)}v(U)\,dU,\qquad
 a_8(v)=\int_{SU(3)}v(U)\chi_8(U)\,dU.
\]
The signed Haar-mode projections obey
\[
 \boxed{P_\chi\mathcal V_vP_\chi=\bar v P_\chi,
 \quad P_\psi\mathcal V_vP_\psi
 =\bigl(\bar v-\tfrac1{256}a_8(v)\bigr)P_\psi,
 \quad P_\chi\mathcal V_vP_\psi=0.}
\]
For \(v=F\), \(\bar F=a_8(F)=1/9\), recovering
Theorem~\ref{p11:thm:v235-signed-face-source-compression}.
For the existing compact heat representative
\(v_t=-\log(K_t(U)/K_t(I))\) at \(t>0\), this identity determines
the projected reference response through one adjoint-character
integral.  It does not identify face heat time with physical
Hamiltonian evolution time.
\end{corollary}
\begin{proof}
Only the three fan faces have nontrivial \(\psi_e\) covariance,
and \(\mathbb E_0[\psi_e^2\mid V_i]
=(8-\chi_8(V_i))/81\).
After division by \(\mathbb E_0\psi_e^2=8/81\), each such face
contributes \(-a_8(v)/8\) relative to the Haar mean; averaging
three faces among \(96\) gives \(-a_8(v)/256\).
The same support and color-conjugation arguments as in the
preceding theorem establish the other blocks.
\end{proof}

\begin{corollary}[Vertex-gradient sources are redundant on the finite link owner]
\label{p11:cor:v235-hurwitz-edge-source-ward}
For the represented bounded edge observables
\(\widehat D_e=\sum_{f\supset e}F(U_f)\), and the incidence
maps of Theorem~\ref{p1:thm:v235-hurwitz-chain-complex},
\[
 \partial_1\widehat{\mathbf D}=0,\qquad
 H_L(\boldsymbol\lambda+\partial_1^{\mathsf T}\phi)
 =H_L(\boldsymbol\lambda),\quad
 H_L(\boldsymbol\lambda)=H_L+\sum_e\lambda_e\widehat D_e.
\]
Thus exactly \(23\) independent formal vertex-gradient source
directions are redundant wherever this \emph{same-owner} bounded
link-source deformation is admitted.  Its finite response and
retarded susceptibilities inherit the same null directions.
\end{corollary}
\begin{proof}
The multiplication-operator identity follows from
\(\widehat{\mathbf D}=B\widehat{\mathbf F}\) and
\(\partial_1B=0\).  Its inner product with the gradient
\(\partial_1^{\mathsf T}\phi\) vanishes identically, including
under arbitrary time evolution of the admitted Hamiltonian.
The image of \(\partial_1^{\mathsf T}\) has rank \(23\).
\end{proof}

\begin{remark}[Compressed source response versus actual QCD consumption]
\label{p11:rem:v235-source-vs-poles}
On the additional same-owner \emph{pure-electric} branch with common
\(\alpha>0\), the signed packet compression of
\(H_E+\lambda\mathcal F_H\) is
\[
 \bigl(9\hbar^2\alpha+\lambda/9\bigr)P_\chi+
 \bigl(11\hbar^2\alpha/2+85\lambda/768\bigr)P_\psi.
\]
This is not an invariant spectral subspace of the face multiplication
operator, and its diagonal difference
\(7\hbar^2\alpha/2+\lambda/2304\) is not an interacting baryon gap.
The remaining physical-consumption calculation is the actual
same-owner fermionic source jet.  On an admitted affine invertible
branch \(D_F(u)=D_F(0)+u\,\iota_e(\mathcal B_e)\), it has the
diagnostic identity
\[
 \left.\frac{d^3}{du^3}[-\log\det D_F(u)]\right|_{u=0}
 =-2\operatorname{Tr}\!\left[
   \bigl(D_F(0)^{-1}\iota_e(\mathcal B_e)\bigr)^3\right].
\]
Neither the insertion \(\iota_e\) nor the noncancellation of its
spin/color, clock and boundary contractions follows from the
comparison cubic of Theorem~\ref{p10:thm:v235-rooted-weyl-cubic}.
The physical baryon poles, mixed action contacts and continuum
claims therefore remain guarded.
\end{remark}
% END FRAMEWORK v2.235: SIGNED HURWITZ FACE-SOURCE CONSUMPTION

% BEGIN FRAMEWORK v2.230: ANALYTIC SPIN-FLUX READOUT FIREWALL
\begin{remark}[All-link analytic coefficient does not fix a baryon level]
\label{p11:rem:v230-all-link-flux-readout}
The all-link analytic coefficient calculation in
Theorem~\ref{p10:thm:v230-all-link-spinflux} now establishes the
$O(x^8)$ comparison identity without gauge-background sampling.  It
does not identify the comparison time $x$ with the physical Euclidean
transfer time $T$ in \eqref{p11:eq:v2220-physical-correlator}.
The unequal-propagator spin--flavor channels and the interacting
baryon-state matrix elements in
\eqref{p11:eq:v2220-flux-hyperfine} remain separate action and readout
obligations.  The exact Haar average is a pure-link benchmark, not the
interacting physical expectation of this observable.
\end{remark}
% END FRAMEWORK v2.230: ANALYTIC SPIN-FLUX READOUT FIREWALL

\paragraph{Claim boundary.}
The update proves the geometric origin and invariant finite statistics
of a \emph{possible consumer observable} in the already declared QCD
action, not its nonzero matrix-element difference on physical baryon
eigenstates.  The physical strange-baryon spin-gap ratio near $8/9$,
a universal $E_{36}\sqrt{24/25}$ energy normalization, positive physical
string tension, spatial hadron radii and resonance pole locations
remain independent interacting/readout obligations.  No new spin,
color, temporal-clock or hadron-mass coefficient is selected here.
% END FRAMEWORK v2.220: BARYON SPIN--FLUX SPECTRAL READOUT BOUNDARY

% BEGIN FRAMEWORK v2.226: KINEMATIC CLOSURES


\section{Physical null, light-ray distributions and fragmentation}
\label{p11:sec:v237-b5}
The native physical null generator supplies the finite operator owner. Species/Fock positivity, finite simplex measures and asymptotic factorization consume additional, named certificates; full calculations are retained in Paper~VII-D.

% Source Paper VII-B, original line 872; inherited proof bodies unchanged.
\subsection[PDFs, distribution amplitudes, fragmentation, and matrix elements]
{PDFs, amplitudes, fragmentation, and matrix elements}
\label{p11:sec:hadronic-distributions}

Hadron masses do not by themselves determine hadronic observables.  Scattering, decay, collider, and exclusive processes require operator and distribution data.  In a native QCD completion these data must be readouts of the singlet hadron states, not fitted functions inserted later.


% Source Paper VII-B, original line 878; inherited proof bodies unchanged.
\subsubsection{Distributions, amplitudes, and perturbative evolution}
\label{p11:subsec:distributions-amplitudes-perturbative-evolution}

\begin{definition}[Null-ray geometry, physical operator owner, state, and limit package]
\label{p11:def:hadronic-distribution-package}
At level \(L\), a physical hadronic light-ray readout consumes four differently typed
packages.
\begin{enumerate}[label=\textnormal{(LF-\Alph*)},leftmargin=3.6em]
\item \textbf{Geometry.}
The derived upstream package \(\LFRayGeom\) of Paper~III-E,
Definition~\ref{p3e:def:lf-ray-geometry-package}: the null operation pair,
its coframe and local generator components, and the complete basepoint-indexed
null-path family.  It supplies neither a global null foliation nor a selected
hadron/beam alignment.
\item \textbf{Physical null/source owner.}
Paper~VII-A supplies the conserved physical null generator
\begin{equation}
\label{p11:eq:v264-physical-plus-owner}
 P_{\rm phys,L}^+
 =\frac1{\sqrt2}\left(H_L/c_{\rm op}+P_{H,L}\right),
 \qquad [H_L,P_{\rm phys,L}^+]=0,
\end{equation}
together with the QCD current/twist-two source family and the color connection
that defines Wilson transport.  Theorem~\ref{p10:thm:v264-twist-two-current-stress-owner}
fixes the \(n=1\) current and forward \(n=2\) stress-tensor normalizations from
the same action/source owner.  The exact finite certificate \(\LFRep_L\) is an
additional stronger realization only when a claim explicitly asks for the
species/Fock simplex, positive number/momentum POVMs, or an occupation-space
identification with \(P_{\rm phys,L}^+\).  In particular, clause~\textnormal{(R3)}
is not imported by the physical light-ray/twist-two definition itself.
\item \textbf{Physical state and alignment.}
Let \(\mathcal H_{L,\rm sing}^{\rm QCD}\) be the finite gauge-singlet QCD carrier
on which \(H_L\) and \(P_{\rm phys,L}^+\) act.  A certificate
\(\LFState_{h,L}\) consists of a nonzero physical fixed-plus sector
\begin{equation}
\label{p11:eq:v264-physical-fixed-plus-sector}
 E^{\rm phys}_{p_h,L}
 :=\mathbf1_{\{p_h^+\}}
 \!\left(P_{\rm phys,L}^+\big|_{\mathcal H_{L,\rm sing}^{\rm QCD}}\right),
 \qquad p_h^+>0.
\end{equation}
For the remainder of this paper write \(E_{p_h,L}:=E^{\rm phys}_{p_h,L}\).  Then
\begin{equation}
\label{p11:eq:lf-fixed-total-plus-sector}
 E_{p_h,L}P_{\rm phys,L}^+=p_h^+E_{p_h,L}.
\end{equation}
and a normalized vector
\(|h,L;p_h^+\rangle\in E^{\rm phys}_{p_h,L}\mathcal H_{L,\rm sing}^{\rm QCD}\),
and either a declared state/beam alignment with one generated conjugate of the
clock--Hopf ray or an exact covariance theorem proving the readout independent
of that choice.  The vector may be supplied as a preparation or constructed on
the stable source-overlap branch of
Corollary~\ref{p11:cor:source-selected-fixed-momentum-state}; no construction is
inferred for a resonance residue.  The representation and conservation of the
total generator do not select this vector.  For an actually realized hadron
preparation or history, the corresponding token remains G0 occurrence content.

If the optional exact \(\LFRep_L\) certificate is additionally realized, clause
\textnormal{(R3)} gives
\(P_{F,L}^+=P_{\rm phys,L}^+\), so the Fock-sector projector used in
Appendix~\ref{p11:app:v264-exact-lfrep-positive-measures} coincides with
\(E^{\rm phys}_{p_h,L}\).
\item \textbf{Claim-local limiting clauses.}
Only a claim that takes a continuum distribution limit imports the applicable
clauses of \(\LFLim\) from
Definition~\ref{p11:def:claim-local-lf-limit-package}.  Local current and
twist-two matrix elements use no exact finite Fock/simplex certificate.  The
optional finite positive fraction measures additionally consume \(\LFRep_L\),
but no continuum clause until a weak or renormalized limit is claimed.
\end{enumerate}
At a factorization scale \(\mu\), the physical readout is denoted
\begin{equation}
\label{p11:eq:hadronic-functional-package}
\mathcal F_{\rm had}^{\rm phys}(h;\mu)
:=\operatorname{Readout}_{\mu}
\left(\LFRayGeom,P_{\rm phys,L}^+,\LFState_{h,L},
      \LFLim_{\mathcal C},\mathcal O_{\rm SM}\right),
\end{equation}
where \(\mathcal C\) is the claim being evaluated.  Exact finite fraction POVMs
form an optional stronger realization, not an adjustable PDF shape or fitted
normalization.  No S-matrix or infrared-completion premise belongs to this
definition.
\end{definition}

\begin{theorem}[Physical fixed-total-plus sector preservation and state-occurrence firewall]
\label{p11:thm:v134-fixed-plus-sector-preservation}
On the gauge-singlet QCD carrier, assume the Paper~VII-A R1--R2 physical-null
generator theorem and gauge invariance of the total generator.  Let
\(H_{L,\mathrm{sing}}\) be the restriction of \(H_L\) to that carrier and let
\(E^{\rm phys}_{p_h,L}\) be
\eqref{p11:eq:v264-physical-fixed-plus-sector}.  Then
\begin{equation}
\label{p11:eq:v134-fixed-plus-sector-preservation}
 [H_{L,\mathrm{sing}},E^{\rm phys}_{p_h,L}]=0,
 \qquad
 e^{-itH_{L,\mathrm{sing}}/\hbar}E^{\rm phys}_{p_h,L}
 =E^{\rm phys}_{p_h,L}e^{-itH_{L,\mathrm{sing}}/\hbar}.
\end{equation}
Hence a supplied or actually occurring state that starts in the physical
\(p_h^+\) sector remains in that sector.  This statement uses neither the
species projectors nor clause~\textnormal{(R3)}.  If the optional exact
\(\LFRep_L\) branch is realized, the same sector is simultaneously the fixed
Fock-plus sector because \(P_{F,L}^+=P_{\rm phys,L}^+\).
\end{theorem}
\begin{proof}
Paper~VII-A gives \([H_L,P_{\rm phys,L}^+]=0\).  Gauge invariance makes
\(P_{\rm phys,L}^+\) preserve the singlet carrier, so finite spectral functional
calculus gives commutation with
\(\mathbf1_{\{p_h^+\}}(P_{\rm phys,L}^+)\), proving the first identity.
Exponentiation gives the second.  State occurrence remains distinct from
spectral-sector preservation, and species/Fock additivity is an optional
stronger representation statement.
\end{proof}

\begin{definition}[Claim-local light-front limiting clauses]
\label{p11:def:claim-local-lf-limit-package}
The symbol \(\LFLim\) denotes a menu, not a global conjunction.  A claim
imports only the clauses it uses:
\begin{enumerate}[label=\textnormal{(L\arabic*)},leftmargin=2.8em]
\item declared embeddings or refinement channels comparing the finite
carriers, states, physical null/source operators, and, only on the optional
exact-LFRep branch, the positive fraction operator measures;
\item an elapsed-cover local-window or decompactifying regime in which the
selected clock--Hopf null helix converges on every claimed compact parameter
interval to the target light ray, with curvature errors controlled;
\item convergence of the color connections and holonomies on those paths,
with gauge covariance preserved;
\item a specified ultraviolet light-ray renormalization scheme and a stated
test-function domain on which the Fourier/Mellin operation commutes with the
limit; and
\item only for TMD claims, a rapidity regulator, soft factor/subtraction, and
the corresponding rapidity-evolution and matching domain.
\end{enumerate}
Clause~\textnormal{(L1)} is sufficient for a weak limit of the finite positive
momentum measures once either
Theorem~\ref{p11:thm:hausdorff-pdf-moment-limit} or
Theorem~\ref{p11:thm:bounded-lipschitz-pdf-limit} is discharged.
Identifying that measure with a renormalized collinear PDF additionally uses
clauses~\textnormal{(L2)--(L4)}, and a TMD identification also uses
clause~\textnormal{(L5)}.  These clauses are distinct from the claim-indexed
\(\LFScat\) conditions in Paper~VII-C.
\end{definition}


\begin{remark}[Exact finite versus controlled asymptotic LF representation]
\label{p11:rem:v186-exact-asymptotic-lfrep-split}
Two certificates are now kept distinct.  The existing $\LFRep_L$ is an exact
finite representation certificate and remains the premise of the optional
finite fraction operators, exact simplex identities, and positive operator
measures collected in Appendix~\ref{p11:app:v264-exact-lfrep-positive-measures}.
It is not a premise of the main physical light-ray/twist-two owner or of the
current/stress-tensor sum-rule theorem.  A controlled local-flat certificate
$\LFRep^{\rm asymp}_{m_x}$ may instead consist of comparison maps for which the
R1--R6 defects vanish in a declared $m_x\to\infty$ or decompactifying regime,
together with the spectral-sector and zero-mode stability estimates consumed
by the claim.  Such an asymptotic certificate may feed the claim-local limit
package, but it does not replace $\LFRep_L$ in an exact finite theorem unless
that theorem is explicitly restated with quantified error terms.  Thus exact
finite LF algebra and continuum light-ray recovery are separate structural
claims rather than two names for one premise.
\end{remark}

\begin{remark}[Physical versus additive finite plus momentum]
\label{p11:rem:v252-physical-fock-plus-distinction}
Before clause~\textnormal{(R3)} is invoked, Paper~VII-A distinguishes
\[
 P_{\rm phys,L}^+=\frac1{\sqrt2}
 \left(H_L/c_{\rm op}+P_{H,L}\right)
 \qquad\text{from}\qquad
 P_{F,L}^+=\sum_iP_{i,L}^+.
\]
All exact simplex and finite fraction-operator theorems below consume R3 and
therefore work on the subdomain where
\(P_{F,L}^+=P_{\rm phys,L}^+\).  On that subdomain the shorter notation
\(P_L^+\) used below denotes this common operator.  No equality is inferred
outside that certified R3 domain.
\end{remark}

\begin{theorem}[Native light-ray/twist-two physical readout]
\label{p11:thm:lf-wilson-line-identification}
Let \(|h,L;p_h^+\rangle\) be a physical state from
Definition~\ref{p11:def:hadronic-distribution-package}.  Assume the generated null direction, the finite QCD quark/gluon fields, and
the QCD color connection define the gauge-invariant bilocals along \(z_\lambda:=\gamma_{+,z_0}(\lambda)\),
\begin{equation}
\label{p11:eq:v264-physical-light-ray-bilocal}
 \mathcal O_{q,L}(\lambda)
 =\bar q_{f,L}(z_\lambda)\slashed n\,
 W_{+,L}(\lambda,0)q_{f,L}(z_0),
\end{equation}
and the analogous gluon operator.  Their local Mellin derivatives are the
source-owned twist-two insertions \(O_{q,n}^{(L)}\) and \(O_{g,n}^{(L)}\) of
Paper~VII-A.  Under clauses~\textnormal{(L2)--(L4)} of \(\LFLim\), the
renormalized Fourier transform of these bilocals defines the physical collinear
PDFs, and their Mellin moments are the corresponding renormalized twist-two
matrix elements.  No exact finite species/Fock equality
\(P_{F,L}^+=P_{\rm phys,L}^+\), commuting fraction simplex, or finite number
POVM is used in this physical definition.

This theorem by itself does not prove support on \([0,1]\), positivity after a
chosen subtraction scheme, or small-\(x\)/zero-mode control.  Those are separate
spectral/renormalization statements.  TMDs additionally require
clause~\textnormal{(L5)}.
\end{theorem}
\begin{proof}
Gauge covariance of the QCD fields and endpoint transformation of the Wilson
transporter make \eqref{p11:eq:v264-physical-light-ray-bilocal} gauge invariant.
Covariant Taylor expansion along the same null path generates powers of
\(n\!\cdot\!D\), hence the local twist-two source family already fixed in
Paper~VII-A.  Clauses~\textnormal{(L2)--(L4)} are exactly the geometric-path,
color-holonomy, ultraviolet-renormalization, and Fourier/Mellin interchange
hypotheses needed to pass to the renormalized light-ray distribution.  None of
these steps invokes the optional one-body/Fock realization.
\end{proof}

\begin{theorem}[Current and stress-tensor PDF sum rules]
\label{p11:thm:native-pdf-sum-rules}
Assume Theorem~\ref{p11:thm:lf-wilson-line-identification} and the finite
current/spin-two owner identity of
Theorem~\ref{p10:thm:v264-twist-two-current-stress-owner}.  Let \(h\) have net
flavor charges \(N_f^h\), baryon number \(B_h\), and electric charge \(Q_h\).
Assume \emph{in addition} a separately certified physical support/crossing and zero-mode prescription: the relevant quark/antiquark/gluon light-ray distributions have the standard \([0,1]\) physical-fraction support (with the negative signed-$x$ quark branch correctly identified with antiquarks), their zero-mode and endpoint contributions are included in the moments, and no regulator/contact remainder spoils the local-current and stress-tensor Mellin identities. Whenever these conditions and the renormalized light-ray limits and finite Ward/stress identities
are preserved by the Paper~VII-C QFT4/QFT10 reconstruction,
\begin{align}
\label{p11:eq:v264-pdf-net-flavor-sum-rule}
 \int_0^1\!\bigl(q_f(x,\mu)-\bar q_f(x,\mu)\bigr)\,dx&=N_f^h,\\
\label{p11:eq:v264-pdf-baryon-sum-rule}
 \frac13\sum_f\int_0^1\!\bigl(q_f-\bar q_f\bigr)(x,\mu)\,dx&=B_h,\\
\label{p11:eq:v264-pdf-charge-sum-rule}
 \sum_f Q_f\int_0^1\!\bigl(q_f-\bar q_f\bigr)(x,\mu)\,dx&=Q_h,\\
\label{p11:eq:v264-pdf-momentum-sum-rule}
 \sum_f\int_0^1\!x\bigl(q_f+\bar q_f\bigr)(x,\mu)\,dx
 +\int_0^1\!xg(x,\mu)\,dx&=1.
\end{align}
The first three identities are current-charge statements.  The last is the
forward \(n=2\) twist-two/stress-tensor identity normalized by the physical
translation charge \(P_{\rm phys}^+\).  Exact finite R3 is not a premise. The \([0,1]\) integral boundaries are a \emph{separate certified light-ray support/crossing input}, not a consequence of the preceding physical bilocal theorem alone.
\end{theorem}
\begin{proof}
The \(n=1\) member of the source-owned quark twist-two tower is the normalized
vector current by Theorem~\ref{p10:thm:v264-twist-two-current-stress-owner}; its
forward charge gives the net-flavor identity, and baryon number and electric
charge are the displayed linear combinations.  The forward \(n=2\) quark and
gluon insertions sum to \(T^{++}_{\rm QCD}\) on physical singlet matrix elements.
The total null translation charge is therefore the sum of the corresponding
quark, antiquark, and gluon momentum moments, yielding
\eqref{p11:eq:v264-pdf-momentum-sum-rule}.  Paper~VII-C transports these finite
identities only when every contact, anomaly, regulator, and composite term
belongs to the controlled QFT4/QFT10 family.
\end{proof}

\begin{remark}[Optional exact finite positive-measure realization]
\label{p11:rem:v264-optional-finite-positive-measure-branch}
Appendix~\ref{p11:app:v264-exact-lfrep-positive-measures} retains the stronger
exact-\(\LFRep_L\) construction of commuting finite fraction operators, positive
number/momentum POVMs, the exact finite \([0,1]\)-support argument, and the
Hausdorff/bounded-Lipschitz positive-measure limits.  Those statements are
useful finite realization theorems but are not the owner of the physical
renormalized PDF.  In particular, their finite positivity and support may not
be transferred through a subtraction/continuum limit without a separate
spectral/renormalization theorem.
\end{remark}

\begin{theorem}[Application of the imported Paper-V-C QCD evolution interface]
\label{p11:thm:native-dglap-scale-evolution}
\label{p11:thm:parameter-free-perturbative-qcd-completion}
Paper~V-C, Proposition~\ref{p7v2:prop:dglap-erbl}, owns the leading-twist
spacelike DGLAP, timelike fragmentation DGLAP, and ERBL evolution generators
after a factorization scheme and the corresponding nonperturbative boundary
are supplied.  Its Proposition~\ref{p7v2:prop:qcd-running} owns perturbative
strong-coupling running between thresholds, while
Proposition~\ref{p7v2:prop:threshold-obligation} records the additional
threshold-matching data that must be supplied.  Paper~VII-B owns the
nonperturbative boundary functions and applies the imported semigroup:
\begin{equation}
\label{p11:eq:dglap-equation}
\mathcal H_{\rm QCD}^{7B}(Q)
=\mathcal U_{\rm QCD}^{5C}
\!\left(Q,\mu_{\rm QCD}^{\rm int};\alpha_s,\mathcal T_q\right)
\mathcal H_{\rm QCD}^{7B}(\mu_{\rm QCD}^{\rm int}).
\end{equation}
The scale variable $Q$ in this theorem is a renormalization/factorization
scale, not physical time.  Applying the imported scale semigroup therefore
constructs no G0 history and supplies no hadron-state selection.
The imported color factors are
\begin{equation}
\label{p11:eq:dglap-color-factors}
C_F=\frac43,\qquad C_A=3,\qquad T_F=\frac12,
\end{equation}
and the evolution preserves the Paper-VII-B flavor and momentum sum rules.  Paper~VII-B
does not rederive the beta function, splitting kernels, or ERBL fixed modes.
\end{theorem}
\begin{proof}
The perturbative semigroup and its coefficients are fixed upstream.  Once
Paper~VII-B supplies a native PDF, DA, GPD, or fragmentation boundary function,
its scale transport is the unique application of that imported linear evolution
operator.  This transport adds no new nonperturbative coefficient.
\end{proof}

\begin{theorem}[Conditional QCD factorization bridge]
\label{p11:thm:native-qcd-factorization-bridge}

For a process and scale regime in which the relevant factorization theorem
applies, the operator-level quantity separates into perturbatively calculable
hard kernels convolved with renormalized PDFs/DAs/FFs defined by the typed
light-front/out-state operators, up to the theorem's power-suppressed remainder.
The factorization scale \(\mu\) is a conventional separation scale; no Paper-VII-A
wall scale is required.  If the left-hand side is asserted to be an exclusive
asymptotic amplitude, an \(S\)-matrix element, or a scattering cross section,
the claim additionally consumes precisely the relevant clauses of the
claim-indexed package in Definition~\ref{p12:def:claim-indexed-smat-package},
including an infrared route when the channel is massless or long range.  A
light-ray operator identity or evolution equation by itself does not.
\phantomsection\label{p11:eq:gauge-invariant-pdf-definition}
\phantomsection\label{p11:eq:hard-factorization-bridge}
\end{theorem}
\begin{proof}

This is the standard factorization statement under its process-dependent
infrared-safety, scale-separation, operator-definition, and renormalization
hypotheses.  Proposition~\ref{p12:prop:scattering-nonclaim} and
Theorem~\ref{p12:thm:scattering-hypothesis-firewall} supply the additional
asymptotic claim firewall.  Paper~VII-B supplies nonperturbative matrix elements
only after the corresponding operator interfaces exist; it does not infer them
from finite color counts.
\end{proof}


\begin{theorem}[Native distribution-amplitude export]
\label{p11:thm:native-distribution-amplitude-export}
For a physical pion eigenstate, the leading-twist distribution amplitude is the gauge-invariant light-cone matrix element
\begin{equation}
\label{p11:eq:pion-da-definition}
\langle0|\bar d(z)\gamma^\mu\gamma_5W(z,0)u(0)|\pi^+(P)\rangle
=i f_\pi P^\mu\int_0^1 dx\,e^{ixP\cdot z}\phi_\pi(x,\mu).
\end{equation}
The local axial-current limit fixes
\begin{equation}
\label{p11:eq:pion-da-normalization}
\int_0^1dx\,\phi_\pi(x,\mu)=1.
\end{equation}
On a process domain satisfying the factorization and asymptotic hypotheses of
Theorem~\ref{p11:thm:native-qcd-factorization-bridge}, hard exclusive amplitudes
factorize into native distribution amplitudes and perturbative kernels.  The
operator's scale evolution itself is ERBL evolution, with kernels fixed by the
derived color data, and
\begin{equation}
\label{p11:eq:pion-da-asymptotic}
\phi_\pi(x,\mu\to\infty)=6x(1-x)
\end{equation}
for the leading pion channel~\cite{EfremovRadyushkin1980,LepageBrodsky1980}.
\end{theorem}
\begin{proof}
The Wilson line makes the bilocal light-cone operator gauge invariant.  Setting
\(z=0\) reduces the bilocal operator to the axial current used to define
\(f_\pi\), giving the normalization.  On the separately declared
scattering/factorization domain, short-distance exclusive factorization
separates hard propagators from the light-cone hadron projection.  Independently
of that asymptotic claim, collinear renormalization of the light-ray operator
gives ERBL evolution, whose asymptotic fixed form is the lowest Gegenbauer mode.
\end{proof}

\begin{theorem}[Isospin-symmetric pion DA parity]
\label{p11:thm:isospin-symmetric-pion-da-parity}
In the isospin-symmetric limit \(m_u=m_d\), the native leading-twist pion
distribution amplitude obeys the exchange symmetry
\begin{equation}
\label{p11:eq:pion-da-exchange-symmetry}
\phi_\pi^{\rm BH}(x,\mu)=\phi_\pi^{\rm BH}(1-x,\mu).
\end{equation}
Equivalently, in the Gegenbauer expansion
\begin{equation}
\label{p11:eq:pion-da-gegenbauer-native}
\phi_\pi^{\rm BH}(x,\mu)
=6x(1-x)\left[1+\sum_{n\ge1}a_n^{\pi,\rm BH}(\mu)
C_n^{(3/2)}(2x-1)\right],
\end{equation}
all odd coefficients vanish,
\begin{equation}
\label{p11:eq:pion-da-odd-gegenbauer-vanish}
a_{2k+1}^{\pi,\rm BH}(\mu)=0.
\end{equation}
Thus the first nontrivial pion-DA target is the even moment \(a_2^{\pi,\rm BH}\),
not an arbitrary asymmetric shape.
\end{theorem}
\begin{proof}
For equal up and down masses, the finite QCD transfer is invariant under the
combined flavor rotation and charge-conjugation operation that exchanges the
quark and antiquark light-cone momentum fractions in the pion.  The Wilson-line
operator in Equation~\eqref{p11:eq:pion-da-definition} is mapped to its reversed
separation with the same matrix element, giving \(x\leftrightarrow1-x\).  The
Gegenbauer polynomial \(C_n^{(3/2)}(2x-1)\) has parity \((-1)^n\) under this
exchange, so all odd coefficients vanish.  Isospin breaking may generate odd
moments for kaons or for unequal quark masses, but those are native mass-spurion
readouts, not shape inputs.
\end{proof}


\begin{corollary}[Closed pion-DA first moments]
\label{p11:cor:closed-pion-da-first-moments}
In the isospin-symmetric pion channel, the first light-cone momentum moments are
fixed before any shape fit:
\begin{equation}
\label{p11:eq:pion-da-first-moment-closure}
\int_0^1 dx\,x\,\phi_\pi^{\rm BH}(x,\mu)=\frac12,
\qquad
\int_0^1 dx\,(2x-1)\phi_\pi^{\rm BH}(x,\mu)=0 .
\end{equation}
Thus the first unresolved pion-DA shape coefficient is the even coefficient
\(a_2^{\pi,\rm BH}(\mu)\); the normalization, exchange parity, and first moment
are already closed.
\end{corollary}
\begin{proof}
Normalize by Equation~\eqref{p11:eq:pion-da-normalization}.  Equation~\eqref{p11:eq:pion-da-exchange-symmetry} makes the distribution symmetric about
\(x=1/2\), so the expectation value of \(x-1/2\) vanishes.  This gives the first
equality and the equivalent \(\xi=2x-1\) moment.
\end{proof}

\begin{corollary}[Asymptotic pion-DA moments]
\label{p11:cor:asymptotic-pion-da-moments}
At the ERBL ultraviolet fixed point, the pion distribution amplitude and all of
its moments are fixed:
\begin{equation}
\label{p11:eq:asymptotic-pion-da-fixed-shape}
\phi_{\pi,\infty}^{\rm BH}(x)=6x(1-x).
\end{equation}
For nonnegative integers \(n\),
\begin{equation}
\label{p11:eq:asymptotic-pion-x-moments}
\langle x^n\rangle_{\pi,\infty}^{\rm BH}
=
\int_0^1dx\,x^n\,6x(1-x)
=
\frac{6}{(n+2)(n+3)}.
\end{equation}
With \(\xi=2x-1\),
\begin{equation}
\label{p11:eq:asymptotic-pion-xi-moments}
\langle \xi^{2m}\rangle_{\pi,\infty}^{\rm BH}
=
\frac{3}{(2m+1)(2m+3)},
\qquad
\langle \xi^{2m+1}\rangle_{\pi,\infty}^{\rm BH}=0 .
\end{equation}
Equivalently, in the Gegenbauer expansion of
Equation~\eqref{p11:eq:pion-da-gegenbauer-native},
\begin{equation}
\label{p11:eq:asymptotic-pion-gegenbauer-vanishing}
a_n^{\pi,\rm BH}(\infty)=0\qquad(n\ge1).
\end{equation}
Thus pion-DA normalization, exchange parity, first moment, ultraviolet shape, and
all ultraviolet moments are closed; only finite matching-scale even coefficients
such as \(a_2^{\pi,\rm BH}(\mu_0)\) remain native matrix-element computations.
\end{corollary}
\begin{proof}
\ProofPlan{We evaluate the asymptotic pion-DA moments in four steps: (a) use the ultraviolet fixed shape; (b) compute ordinary \(x\)-moments; (c) compute \(\xi\)-moments and parity; (d) read off the Gegenbauer fixed point.}
\begin{proofsteps}
\item \textbf{Ultraviolet fixed shape.}
Equation~\eqref{p11:eq:pion-da-asymptotic} gives the ultraviolet fixed shape.

\item \textbf{Ordinary \(x\)-moments.}
The
ordinary moment is
\[
6\int_0^1dx\,(x^{n+1}-x^{n+2})
=6\left(\frac1{n+2}-\frac1{n+3}\right),
\]
which gives \eqref{p11:eq:asymptotic-pion-x-moments}.

\item \textbf{\(\xi\)-moments and parity.}
For the \(\xi\)-moments,
write \(x=(1+\xi)/2\), \(dx=d\xi/2\), and
\(x(1-x)=(1-\xi^2)/4\).  Then
\[
6x(1-x)\,dx=\frac34(1-\xi^2)\,d\xi .
\]
Odd moments vanish by parity and
\[
\frac34\int_{-1}^{1}d\xi\,\xi^{2m}(1-\xi^2)
=
\frac34\left(\frac{2}{2m+1}-\frac{2}{2m+3}\right)
=
\frac{3}{(2m+1)(2m+3)}.
\]

\item \textbf{Gegenbauer fixed-point readout.}
The Gegenbauer coefficients vanish because \(6x(1-x)\) is the lowest ERBL
eigenfunction alone.
\end{proofsteps}
\end{proof}

\begin{theorem}[Asymptotic determinant-baryon distribution amplitude]
\label{p11:thm:asymptotic-determinant-baryon-da}

Conditional on the standard leading-twist baryon light-cone operator and ERBL domain, the asymptotic three-quark distribution amplitude is
\[
\Phi_{\rm as}(x_1,x_2,x_3)=120x_1x_2x_3,
\qquad x_i\ge0,\quad x_1+x_2+x_3=1.
\]
Its permutation symmetry is compatible with the non-color symmetry required by the determinant color singlet.  This asymptotic shape does not fix the finite-scale boundary DA.
\phantomsection\label{p11:eq:asymptotic-baryon-da}
\phantomsection\label{p11:eq:asymptotic-baryon-da-moments}
\phantomsection\label{p11:eq:asymptotic-baryon-da-low-moments}
\end{theorem}
\begin{proof}

The lowest anomalous-dimension ERBL eigenfunction for the leading three-quark baryon operator is proportional to \(x_1x_2x_3\); normalization on the simplex gives the factor 120.  The determinant fixes color antisymmetry, while the finite-scale coefficients of higher conformal polynomials remain dynamical.
\end{proof}


\begin{theorem}[Native GPD unification bridge]
\label{p11:thm:native-gpd-unification-bridge}
Off-forward light-cone matrix elements of a physical nucleon define native generalized parton distributions \(H_q,E_q,\widetilde H_q,\widetilde E_q\).  Their projections unify PDFs, electromagnetic form factors, axial form factors, and nucleon angular momentum:
\begin{align}
H_q(x,0,0)&=\begin{cases}q(x),&0<x\le1,\\-\bar q(-x),&-1\le x<0,\end{cases}\notag\\
\widetilde H_q(x,0,0)&=\begin{cases}\Delta q(x),&0<x\le1,\\\Delta\bar q(-x),&-1\le x<0,\end{cases} \label{p11:eq:gpd-forward-limits}\\
\int_{-1}^{1}dx\,H_q(x,\xi,t)&=F_1^q(t),\notag\\
\int_{-1}^{1}dx\,E_q(x,\xi,t)&=F_2^q(t), \label{p11:eq:gpd-first-moments}\\
\int_{-1}^{1}dx\,\widetilde H_q(x,\xi,t)&=G_A^q(t),\notag\\
J_q&=\frac12\int_{-1}^{1}dx\,x\,[H_q(x,\xi,0)+E_q(x,\xi,0)]. \label{p11:eq:gpd-axial-moment-ji}
\end{align}
Together with the gluon analogue, \(\frac12=\sum_qJ_q+J_g\).  Hence PDFs, form factors, axial charge, and spin decomposition are not separate QCD sectors; they are moments and limits of the same native off-forward operator family~\cite{Ji1997,Radyushkin1997}.
\end{theorem}
\begin{proof}
The forward limit removes the momentum transfer and returns the signed-extension PDF operator, with the negative-$x$ quark branch fixed by charge conjugation as the antiquark PDF (and the helicity branch by its axial convention). A possible zero-mode at $x=0$ remains part of the separately certified distributional support prescription.  Integrating over \(x\) collapses the bilocal vector and axial operators to local currents, giving Dirac, Pauli, and axial form factors.  The first \(x\)-weighted moment gives the quark part of the energy-momentum tensor; its nucleon matrix element decomposes into total angular momentum, giving Ji's relation.  All operators are Wilson-line-completed native QCD operators.
\end{proof}

\begin{theorem}[Native Bjorken sum-rule bridge]
\label{p11:thm:native-bjorken-sum-rule-bridge}
The isovector axial charge equals the nonsinglet polarized first moment,
\begin{equation}
\label{p11:eq:ga-polarized-first-moment}
g_A=\Delta u-\Delta d.
\end{equation}
Consequently polarized DIS obeys
\begin{equation}
\label{p11:eq:bjorken-sum-rule}
\int_0^1dx\,[g_1^p(x,Q^2)-g_1^n(x,Q^2)]
=\frac{g_A}{6}C_{\rm NS}(Q^2)+O(Q^{-2}),
\end{equation}
with leading perturbative coefficient
\begin{equation}
\label{p11:eq:bjorken-coefficient-leading}
C_{\rm NS}(Q^2)=1-\frac{\alpha_s(Q^2)}{\pi}+O(\alpha_s^2).
\end{equation}
Thus the leading nonsinglet spin-structure export is fixed once the computed baryon axial response \(g_A\) and the perturbative strong running on the chosen source branch are known~\cite{Bjorken1966}.
\end{theorem}
\begin{proof}
The first moment of the nonsinglet polarized PDF is the matrix element of the nonsinglet axial current.  The operator-product expansion for two electromagnetic currents in polarized DIS gives the first moment of \(g_1^p-g_1^n\) as that axial-current matrix element multiplied by a short-distance nonsinglet coefficient.  The coefficient is perturbative and depends only on the already-derived QCD coupling and color data.
\end{proof}


\begin{theorem}[Asymptotic quark--gluon momentum partition]
\label{p11:thm:asymptotic-quark-gluon-momentum-partition}
Let
\begin{equation}
Q(\mu)=\sum_f\int_0^1dx\,x(q_f+\bar q_f),
\qquad
G(\mu)=\int_0^1dx\,xg(x,\mu).
\end{equation}
The second Mellin moment of singlet DGLAP evolution gives the ultraviolet fixed momentum fractions
\begin{equation}
\label{p11:eq:asymptotic-momentum-partition}
Q_\infty=\frac{3n_f}{16+3n_f},
\qquad
G_\infty=\frac{16}{16+3n_f}.
\end{equation}
For \(n_f=3\),
\begin{equation}
\label{p11:eq:asymptotic-momentum-partition-three-flavors}
Q_\infty=\frac9{25},
\qquad
G_\infty=\frac{16}{25}.
\end{equation}
Thus at asymptotically high scale the total quark/gluon momentum split is fixed by the derived \(SU(3)_c\) color factors.
\end{theorem}
\begin{proof}
The second moment of the singlet DGLAP equations gives
\begin{equation}
\frac{dQ}{d\log\mu^2}=\frac{\alpha_s}{2\pi}\left[-\frac43C_FQ+\frac{n_f}{3}G\right],
\qquad
Q+G=1.
\end{equation}
For \(SU(3)_c\), \(C_F=4/3\), hence the fixed point satisfies \((n_f/3)G=(16/9)Q\), or \(Q/G=3n_f/16\).  Together with momentum conservation this yields \eqref{p11:eq:asymptotic-momentum-partition}.\qedhere
\end{proof}

\begin{corollary}[No physical projector from the arithmetic 25 decomposition]
\label{p11:cor:qcd-remnant-envelope-25-support}

The LO singlet evolution fixed point \(Q_\infty:G_\infty=9:16\) is valid, but the arithmetic rewrite \(25=(1+8)+16\) does not define an operator projector.  Any remnant fraction must be the expectation value of an explicitly constructed projector or positive operator in the physical state.
\phantomsection\label{p11:eq:hidden-eta-qcd-envelope}
\end{corollary}
\begin{proof}

A number decomposition does not specify a Hilbert-space subspace.  Moreover, the current corpus contains representation-dependent finite ratios for the same 24/25 arithmetic.  Therefore no representation-independent QCD projector follows without an operator map.
\end{proof}


\begin{corollary}[Asymptotic quark--gluon angular-momentum partition]
\label{p11:cor:asymptotic-quark-gluon-angular-momentum-partition}
Let
\[
\mathcal J_q(\mu)=2\sum_fJ_f(\mu),
\qquad
\mathcal J_g(\mu)=2J_g(\mu),
\qquad
\mathcal J_q+\mathcal J_g=1 .
\]
The ultraviolet fixed point of the total quark and gluon angular momenta is
\begin{equation}
\label{p11:eq:asymptotic-angular-momentum-partition}
\mathcal J_{q,\infty}
=
\frac{3n_f}{16+3n_f},
\qquad
\mathcal J_{g,\infty}
=
\frac{16}{16+3n_f}.
\end{equation}
Equivalently,
\begin{equation}
\label{p11:eq:asymptotic-Jq-Jg}
J_{q,\infty}
=
\frac12\,\frac{3n_f}{16+3n_f},
\qquad
J_{g,\infty}
=
\frac12\,\frac{16}{16+3n_f}.
\end{equation}
For \(n_f=3\),
\begin{equation}
\label{p11:eq:asymptotic-Jq-Jg-three-flavors}
J_{q,\infty}=\frac9{50},
\qquad
J_{g,\infty}=\frac8{25}.
\end{equation}
If the ultraviolet light-flavor block is flavor symmetric, each active
quark-plus-antiquark flavor carries \(J_{f+\bar f,\infty}=3/50\).
\end{corollary}
\begin{proof}
Ji's relation in Equation~\eqref{p11:eq:gpd-axial-moment-ji} reads the total
quark angular momentum from the same twist-two spin-two energy-momentum
operator whose forward form factor gives the momentum fraction.  The quark and
gluon parts of this operator mix under the same singlet anomalous-dimension
matrix in the forward momentum and total-angular-momentum projections.  Hence
\((\mathcal J_q,\mathcal J_g)\) obeys the same ultraviolet fixed-point equation
as \((Q,G)\) in Theorem~\ref{p11:thm:asymptotic-quark-gluon-momentum-partition}.
Substituting the derived \(C_F=4/3\) gives the ratio
\(\mathcal J_q/\mathcal J_g=3n_f/16\), and normalization gives
\eqref{p11:eq:asymptotic-angular-momentum-partition}.  Dividing by two gives
\eqref{p11:eq:asymptotic-Jq-Jg}.
\end{proof}



% Source Paper VII-B, original line 1531; inherited proof bodies unchanged.
\subsubsection{Fragmentation and hadronization exports}
\label{p11:subsec:fragmentation-hadronization-exports}

\begin{proposition}[Typed hadronization-channel boundary]
\label{p11:prop:fragmentation-as-singlet-outcome}

A rigorous finite hadronization export requires a completely positive trace-preserving map \(\mathcal E_{{\rm had},L}\) from the resolved colored/pre-hadronic density algebra to the color-singlet out-state algebra, preserving all conserved charges and total four-momentum in the stated approximation.  Its Kraus/outcome weights then define inclusive hadronic probabilities.
\phantomsection\label{p11:eq:native-hadronization-map}
\end{proposition}
\begin{proof}

A CPTP map supplies positivity and total probability one.  If its adjoint fixes each conserved charge operator, the corresponding expectation values are conserved.  In particular, Paper~VII-A's internally derived conservation of the identified total $P_L^+$ fixes the compatibility condition $\mathcal E_{{\rm had},L}^*(P_{L,{\rm out}}^+)=P_{L,{\rm in}}^+$ on any hadronization channel claimed to represent the same closed approximation.  That relation constrains a supplied channel but does not construct its Kraus operators or prove that such a channel exists.  The current finite QCD carrier does not uniquely construct the map, so existence remains an explicit interface rather than a consequence of singlet projection alone.
\end{proof}


\begin{theorem}[Conditional fragmentation sum rules from the hadronization channel]
\label{p11:thm:native-fragmentation-hadronization-export}

Conditional on a light-front cut-operator constructor and a charge-preserving
hadronization channel \(\mathcal E_{\rm had}\), fragmentation functions are
matrix elements of the corresponding cut/Wilson-line operators and obey the
algebraic momentum and charge sum rules.  This operator statement consumes no
\(S\)-matrix hypothesis.  Interpreting the channel outcomes as physical
asymptotic hadrons uses \textnormal{(SM1)--(SM3)} of
Definition~\ref{p12:def:claim-indexed-smat-package}; claiming that the summed
hadronic outcomes exhaust the relevant interacting sector additionally uses
\textnormal{(SM6)}.  A massless/long-range asymptotic interpretation uses a
route from Definition~\ref{p12:def:claim-indexed-infrared-package}.  Paper~VII-B
does not claim that \(\mathcal E_{\rm had}\) is generated by finite color
projection alone.
\phantomsection\label{p11:eq:hadronization-charge-sector-decomposition}
\phantomsection\label{p11:eq:fragmentation-momentum-sum-rule}
\end{theorem}
\begin{proof}

Trace preservation and conservation by \(\mathcal E_{\rm had}\) give the
algebraic inclusive sum rules.  The claim-indexed asymptotic and, where
necessary, infrared clauses justify only their physical out-state
interpretation; \textnormal{(SM6)} is needed only for an exhaustive channel
sum.  Absent those clauses the finite singlet algebra supplies channel outcomes
but not their interpretation as asymptotic scattering probabilities.
\end{proof}


% Source Paper VII-B, original line 1574; inherited proof bodies unchanged.
\subsubsection{Distribution computation targets}
\label{p11:subsec:distribution-computation-targets}

\begin{definition}[Distribution computation targets]
\label{p11:def:distribution-computation-targets}
A minimal finite evaluation basis consists of moments and selected matrix elements,
\begin{equation}
\label{p11:eq:first-distribution-moment-targets}
\langle x^n\rangle_{q/h}(\mu_0)\ \text{beyond the closed UV endpoints},
\qquad
g_A,
\qquad
\sigma_{\pi N},
\qquad
\langle h'|\mathcal O_a|h\rangle,
\end{equation}
before attempting full \(x\)-dependent PDFs, complete fragmentation functions, or a global hadronization model.
\end{definition}


% Source Paper VII-B, original line 5193; inherited proof bodies unchanged.
\subsubsection{Spacelike source readouts before exact finite Fock identification}
\begin{theorem}[Gauge-completed spacelike jets with a finite remainder]
\label{p11:thm:v2210-lit-spacelike-source-jets}
On a licensed smooth local reciprocal-position chart, take an actual
spacelike path $\gamma(z)$ with $\gamma(0)=x_0$, the pulled-back QCD
connection $\mathcal A_z$, and finite quark field word $q_z=q(\gamma(z))$.
Use the actual transporter $W(0,z)$, in the convention
\[
 \partial_zW(0,z)=W(0,z)\mathcal A_z,\qquad W(0,0)=I,
 \qquad D_zq_z=\partial_zq_z+\mathcal A_zq_z.
\]
Assume the licensed finite words are $N+1$ differentiable along this path.
For a fixed spin/flavor source matrix $\Gamma$ commuting with color put
\[
 \mathcal O_\Gamma(z)=\bar q(x_0)\Gamma W(0,z)q_z.
\]
The color-contracted operator is gauge invariant. Its exact local expansion
is
\[
 \boxed{\mathcal O_\Gamma(z)
 =\sum_{n=0}^N\frac{z^n}{n!}
      \bar q(x_0)\Gamma D_z^nq_z|_{z=0}
    +\mathcal E_{N+1}(z),}
\]
with the operator bound
\[
 \boxed{\|\mathcal E_{N+1}(z)\|
 \le\frac{|z|^{N+1}}{(N+1)!}\|\bar q(x_0)\Gamma\|
       \sup_{s\text{ between }0,z}\|W(0,s)D_s^{N+1}q_s\|.}
\]
Every normalized actual hadron-state matrix element obeys the same bound.
No exact species projector, Fock-plus additivity or $\LFRep_L$ premise is
used in this finite source calculation.
\end{theorem}
\begin{proof}
Gauge covariance gives $q_z\mapsto g_zq_z$,
$\bar q(x_0)\mapsto\bar q(x_0)g_0^{-1}$ and
$W(0,z)\mapsto g_0W(0,z)g_z^{-1}$, so the contracted word is invariant.
The transporter equation yields
$\partial_z(Wq)=W D_zq$; induction gives
$\partial_z^n(Wq)=WD_z^nq$. Taylor's formula with its integral remainder
and the operator norm give the two displayed statements. A normalized
matrix element is bounded by the operator norm.
\end{proof}
The jets are covariant local source words before the appropriate
twist/trace projection. Curvature, path and higher-twist terms are not
discarded by calling them PDF moments. The actual stable source-selected
hadron state, physical momentum/alignment and normalization remain those of
Definition~\ref{p11:def:hadronic-distribution-package}.

\begin{corollary}[Conditional reduced Ioffe-time renormalization test]
\label{p11:cor:v2210-lit-ioffe-ratio-test}
Suppose the declared matrix elements $\mathcal M(\nu,z)$ of the preceding
same operator have a certified state-independent scalar multiplicative
renormalization with $Z(z,\mu)\ne0$
$\mathcal M^R(\nu,z)=Z(z,\mu)\mathcal M(\nu,z)$ on their comparison domain.
If the reference $\mathcal M(0,z)\ne0$, then
\[
 \frac{\mathcal M^R(\nu,z)}{\mathcal M^R(0,z)}
 =\frac{\mathcal M(\nu,z)}{\mathcal M(0,z)}.
\]
An additive contact, operator-mixing matrix or state-dependent factor does
not satisfy this cancellation premise. A physical collinear PDF conclusion
still requires the actual short-distance/large-momentum matching and its
claimed test-function, source and hadron limits.
\end{corollary}
\begin{proof}
Cancel the same nonzero scalar factor. The omitted alternatives have no
such scalar identity, which is why the premise is retained.
\end{proof}
The spacelike route follows the observable logic of
\cite{Ji2013LaMET,LitQCDJiZhangZhao2018Bilocal,LitQCDRadyushkin2017Pseudo}.
Those sources supply matching/renormalization technology on their actual
domains, not this framework's uncomputed matrix elements. It is a physical
source-observable route alongside the main light-ray/twist-two chain; it
does not reopen the optional exact finite fraction-POVM package. TMD and
fragmentation claims retain their separate soft/rapidity and out-state
clauses. A blanket separate-hadron TMD factorization transfer is invalid for
the known non-Abelian counterexample domain~\cite{LitQCDRogersMulders2010TMD}.

\paragraph{Current execution boundary after these certificates.}
QC01 still needs the actual physical quark/action member and normalization.
QC02 can use the certified frame test, but its physical transfer, uniform
frame and APBI/flux bounds remain actual executions. QC03 now has exact
blocked/Duhamel source contacts and signed sector-cumulant tests; its
weighted ordered chiral response and actual winding weights remain.
QC04 has error-carrying spectral and spacelike source tests; hadron,
heavy-source, decay and multi-baryon matrix values remain to be computed.
QC05 retains only the large-volume/operator limits consumed by the claimed
observable, with OS and a continuous gap required only on the branches that
assert them. QC06 retains every generated multiface interaction and source
jet through its actual blocking history. These reductions neither infer
numerical QCD from counts nor promote optional exact SSB, Clay or finite
$\LFRep_L$ conclusions into gates on finite massive observables.
% END VERIFIED QCD READOUT CERTIFICATES v2.210


\begingroup
\small
\RaggedRight
\setlength{\emergencystretch}{10em}
\sloppy
% Bibliography consolidated at the end of this volume
\endgroup


% BEGIN TRANSITIVE-CHIRAL-SPURION v1.54


\section{Thermal, nuclear, and finite capacity responses}
\label{p11:sec:v237-b6}
Thermal derivatives and nuclear transfer spectra are observable-directed finite operators. Native Hurwitz splits and boundary-reduced ranks are retained; the earlier $E_{36}/E_{40}$ phenomenological correlations are preserved separately in Paper~VII-D as non-derivational targets, not discarded or promoted to mass formulas.

% Source Paper VII-B, original line 1593; inherited proof bodies unchanged.
\subsection[Finite-temperature QCD and cosmological export]
{Thermal QCD and cosmological export}
\label{p11:sec:thermal-cosmology-export}

This section constructs the finite thermal QCD trace and the thermodynamic data exported to Paper~VI-B.

\begin{definition}[Native finite thermal partition function]
\label{p11:def:native-thermal-partition-function}
\leavevmode\par\noindent
The finite-temperature QCD object at regulated Born--Hopf level is
\begin{equation}
\label{p11:eq:native-thermal-partition-function}
Z_{\rm QCD}^{\rm fin}(T,\mu_B,\mu_Q,\mu_S)
=
\operatorname{Tr}_{\Ps\Hraw}
\exp\!\left[-\frac{\HQCDfin-\mu_BB-\mu_QQ-\mu_SS}{T}\right].
\end{equation}
The full thermal target replaces \(\HQCDfin\) by the nonperturbative
transfer/refinement object \(\HQCD\) only after the transfer rule has been
fixed. Equivalently, one may use flavor chemical potentials \(\mu_f\) and
conserved operators \(\widehat N_f\):
\begin{equation}
\label{p11:eq:native-thermal-partition-flavor}
Z_{\rm QCD}^{\rm fin}(T,\mu_f)
=
\operatorname{Tr}_{\Ps\Hraw}
\exp\!\left[-\frac{\HQCDfin-\sum_f\mu_f\widehat N_f}{T}\right].
\end{equation}
\end{definition}

\begin{theorem}[Finite thermal QCD export exists]
\label{p11:thm:finite-thermal-export-exists}
The partition function \(Z_{\rm QCD}^{\rm fin}\) in \eqref{p11:eq:native-thermal-partition-function} is a finite native object.  Its definition uses no lattice equation of state, heavy-ion fit, BBN table, or cosmological likelihood.
\end{theorem}
\begin{proof}
The trace is over the finite singlet strong-sector state space and the exponent is built from \(\HQCDfin\), which is finite and self-adjoint by Theorem~\ref{p10:thm:finite-strong-hamiltonian-closed}.  The conserved charges are fixed operators on the same finite state space by Corollary~\ref{p10:cor:integer-baryon-number} and Theorems~\ref{p10:thm:integer-hadronic-electric-charge} and~\ref{p10:thm:vector-flavor-charge-conservation}.  Hence the regulated partition function is defined internally before any thermal or cosmological comparison.\qedhere
\end{proof}


\begin{theorem}[Finite grand-potential and mechanical-source bridge]
\label{p11:thm:native-thermodynamic-response-bridge}
On a fixed differentiable finite source chart let $H,Q_a$ be independent of
$T$ and of the affine chemical sources, and write
$K=H-\sum_a\mu_aQ_a$, $W=\log\operatorname{Tr}e^{-K/T}$.
The complete metric/volume chart includes its capacity, measure, normalization
and transported-frame response. With $V$ its volume define
$\varpi=TW/V$, $\epsilon=\langle H\rangle/V$, $n_a=\langle Q_a\rangle/V$,
and $s=(W+\langle K\rangle/T)/V$. Then exactly
\begin{equation}
\label{p11:eq:thermal-response-basic}
 \boxed{\epsilon=Ts-\varpi+\sum_a\mu_an_a,\qquad
 p_{\rm mech}=T\partial_VW=-\langle\partial_VK\rangle.}
\end{equation}
The volume derivative is the actual full source derivative; in an isotropic
metric chart $V=V_0a^3$ it is $(a/3V)D_a$. A prescribed changing carrier is
first pulled to its common differentiable frame with all projector jets.
At a rank change the corresponding finite comparison/contact, rather than
an unjustified derivative through a changing trace space, is required.

Only on an extensive thermodynamic branch with the stated volume law does
$p_{\rm mech}=\varpi=:p$, $s=\partial_Tp$, and the usual equation-of-state
form follow. At zero chemical potentials on that branch,
\begin{equation}
\label{p11:eq:thermal-trace-anomaly-response}
 {\epsilon-3p\over T^4}=T{d\over dT}\left({p\over T^4}\right).
\end{equation}
Finite affine charge susceptibilities are derivatives of the grand-potential
density,
\begin{equation}
\label{p11:eq:bqs-susceptibilities}
 \chi_{ijk}^{BQS}=\left.
 {\partial^{i+j+k}(\varpi/T^4)\over
 \partial(\mu_B/T)^i\partial(\mu_Q/T)^j\partial(\mu_S/T)^k}
 \right|_{\mu=0}.
\end{equation}
Using $p/T^4$ here is justified after the same extensivity identification.
A sound speed uses $dp/d\epsilon$ on its declared equilibrium/adiabatic
trajectory, not an unspecified derivative at nonzero conserved density.
\end{theorem}
\begin{proof}
Duhamel differentiation gives $\partial_{\mu_a}W=\langle Q_a\rangle/T$,
$\partial_TW=\langle K\rangle/T^2$, and
$\partial_VW=-\langle\partial_VK\rangle/T$ on this fixed source chart.
The Gibbs entropy is $W+\langle K\rangle/T$, giving the energy identity.
The remaining derivatives are then ordinary differentiation. Equality of
pressure and grand-potential density additionally needs the volume law,
as stated; a finite $V$-independent two-level Hamiltonian has zero mechanical
pressure but positive $T\log Z/V$.
\end{proof}

\begin{theorem}[Complete finite thermal scale and volume insertion]
\label{p11:thm:thermal-trace-anomaly-operator-closure}
On the preceding actual source chart,
\begin{equation}
\label{p11:eq:thermal-trace-anomaly-operator-closure}
 \boxed{\epsilon-3p_{\rm mech}
 ={1\over V}\langle H+3V\partial_VK\rangle
 ={T\over V}(T\partial_T+\mu_a\partial_{\mu_a}-3V\partial_V)W.}
\end{equation}
At $\mu=0$ and $V=V_0a^3$ the complete finite trace insertion is
$\Theta_L^{\rm full}=V^{-1}(H+aD_aH)$. It retains the full metric source,
including any ceiling/capacity, central action, determinant, measure and
frame terms prescribed by that same owner. If a differentiable dimensional
presentation actually has
$H(a,\lambda)=a^{-1}H(1,a^{d_i}\lambda_i)$, differentiation gives the
finite Euler identity
\[
 H+a\partial_aH=\sum_i d_i\lambda_i\partial_{\lambda_i}H.
\]
Every retained dimensionful mark, including a cutoff mark when this
presentation contains it, belongs in that sum. No such homogeneity is inferred
from finite dimensionality.

On a matched renormalized QCD comparison branch define its ordinary insertion
$\Theta^{\rm QCD}=\beta(g_s)F^2/(2g_s)+
\sum_qm_q(1+\gamma_m)\bar q q$ and the actual residual
$\mathcal R_{\rm scale}=\Theta_L^{\rm full}-\Theta^{\rm QCD}$.
Then the same-source statement is
\[
 \epsilon-3p_{\rm mech}
 =\langle\Theta^{\rm QCD}+\mathcal R_{\rm scale}\rangle.
\]
Dropping this residual requires its actual matching/vanishing certificate.
The complete finite insertion is already a calculable native response;
an uncomputed regulator term is not set to zero by quoting the continuum
two-term anomaly formula.
\end{theorem}
\begin{proof}
The pressure formula and $\epsilon=\langle H\rangle/V$ give the first
identity. The temperature and chemical derivatives of $W$ give
$T(T\partial_T+\mu_a\partial_{\mu_a})W=\langle H\rangle$,
which proves the second. The isotropic chain rule and the displayed
homogeneity assumption give the finite Euler identity. The final equality
is the exact source decomposition with its retained residual. At a nonlinear
source pullback all first and higher contact terms are those of
Theorem~\ref{p11:thm:v2210-lit-complete-duhamel-contacts}.
\end{proof}

\begin{theorem}[Native QCD Stefan--Boltzmann limit]
\label{p11:thm:native-qcd-stefan-boltzmann-limit}
On a certified free-gas comparison window in which the active masses and
strong infrared scales are small compared with $T$, the transported coupling
is small, and $T$ is small compared with the fixed physical UV ceiling,
the matched QCD readout approaches the free adjoint-gluon/fundamental-quark
gas with its actual volume, singlet and occupation errors controlled.  The derived carrier \(N_c=3\) gives
\begin{equation}
\label{p11:eq:qcd-stefan-boltzmann-pressure}
p_{\rm SB}^{\rm BH}(T)=
\left[
\frac{8\pi^2}{45}+N_f\frac{7\pi^2}{60}
\right]T^4.
\end{equation}
For \(N_f=3\),
\begin{equation}
\label{p11:eq:qcd-stefan-boltzmann-pressure-three-flavors}
p_{\rm SB}^{(3)}=\frac{19\pi^2}{36}T^4.
\end{equation}
Moreover \(\epsilon_{\rm SB}=3p_{\rm SB}\), \(I_{\rm SB}=\epsilon-3p=0\), and \(c_s^2\to1/3\).
\end{theorem}
\begin{proof}
Massless bosons contribute \(g\pi^2T^4/90\).  Gluons have \(g_g=2(N_c^2-1)=16\).  A massless Dirac quark flavor contributes the fermion factor \(7/8\) times \(g_q=2_{\rm spin}2_{q\bar q}N_c=12\), hence \(7\pi^2T^4/60\) per active flavor.  Summing gives \eqref{p11:eq:qcd-stefan-boltzmann-pressure}; conformal free-gas thermodynamics gives the remaining statements.\qedhere
\end{proof}



\begin{proposition}[Fixed-ceiling free thermal remainder]
\label{p11:prop:v2211-fixed-ceiling-thermal-remainder}
For a free massless species with a spherical spatial ceiling $\Lambda$,
write $r=\Lambda/T>0$. Per $g$ bosonic or fermionic internal modes, the
missing pressure tails satisfy
\[
 0\le{p_{\infty,b}-p_{\Lambda,b}\over T^4}
 \le {g e^{-r}(r^2+2r+2)\over2\pi^2(1-e^{-r})},\qquad
 0\le{p_{\infty,f}-p_{\Lambda,f}\over T^4}
 \le {g e^{-r}(r^2+2r+2)\over2\pi^2}.
\]
The corresponding bosonic energy tail is bounded by
$g e^{-r}(r^3+3r^2+6r+6)/[2\pi^2(1-e^{-r})]$ in units $T^4$;
the fermionic bound omits $(1-e^{-r})^{-1}$. For positive free frequencies
$\omega_j$ and independently truncated occupations $0,\ldots,N_j$,
\[
 W_{\infty}-W_{\{N_j\}}
 =-\sum_j\log[1-e^{-(N_j+1)\omega_j/T}].
\]
These are explicit free comparison errors, not an interacting convergence
theorem. The standard Stefan--Boltzmann coefficients can therefore be used
at a declared tolerance when these tails and the actual interaction,
mass, volume and singlet-projection errors are controlled. They are not a
$T\to\infty$ theorem at fixed finite capacity: then $Z\to\dim\mathcal H$.
\end{proposition}
\begin{proof}
Integrate $gT^4x^2[-\log(1-e^{-x})]/(2\pi^2)$ or
$gT^4x^2\log(1+e^{-x})/(2\pi^2)$ from $r$ to infinity. Use
$-\log(1-e^{-x})\le e^{-x}/(1-e^{-r})$ and
$\log(1+e^{-x})\le e^{-x}$, then integrate $x^2e^{-x}$.
For energy use $x^3/(e^x\mp1)$ and the same estimate. The geometric
oscillator sum gives $Z_{N_j}=(1-e^{-(N_j+1)\omega_j/T})/
(1-e^{-\omega_j/T})$, proving the logarithmic error. A true zero mode
requires its separate licensed finite state; this oscillator formula does
not assume an unnormalizable zero-frequency Gaussian.
\end{proof}

\begin{remark}[Free local gluon count versus finite physical singlet trace]
\label{p11:rem:v245-gluon-count-singlet-thermal-firewall}
The coefficient \(g_g=16=8_{\rm adjoint}\times2_{\rm helicity}\) in the
Stefan--Boltzmann theorem is the local free massless gauge-mode count that
controls the asymptotically free continuum high-temperature endpoint.  At
finite QCD resolution the physical thermal state space remains the exact
Haar-singlet space of Paper~VII-A.  Its free-adjoint comparison trace is the
Haar-projected character trace of
Proposition~\ref{p10:prop:v245-singlet-projected-free-gluon-trace}, not the
unprojected product obtained by simply raising each oscillator factor to the
sixteenth power.  Thus the continuum \(g_g=16\) endpoint and exact finite
singlet projection are complementary statements, not competing definitions
of the physical thermal trace.
\end{remark}


\begin{theorem}[Native low-temperature Goldstone pion gas]
\label{p11:thm:native-low-temperature-pion-gas}
In the exact light-quark chiral limit and for temperatures small compared with the first non-Goldstone excitation gap, the
thermal trace is dominated by the three pions of the broken
\(SU(2)_L\times SU(2)_R\) subcoset.  Thus
\begin{equation}
\label{p11:eq:pion-gas-pressure}
\frac{p_\pi}{T^4}=\frac{\pi^2}{30},
\qquad
\frac{\epsilon_\pi}{T^4}=\frac{\pi^2}{10},
\qquad
\frac{s_\pi}{T^3}=\frac{2\pi^2}{15},
\qquad
c_s^2\to\frac13 .
\end{equation}
For physical light-quark masses this endpoint is replaced by the corresponding
massive pion gas.  At leading chiral order its mass is the correlator-owned relation
\(m_\pi^2=B_0(m_u+m_d)+O(m_q^2)\), with \(B_0=\Sigma/F_\pi^2\) determined by the controlled chiral-limit readouts rather than by a wall scale.
\end{theorem}
\begin{proof}
Conditional on occurrence of the broken chiral phase, the two-flavor subsector contains three Goldstone pions.  A massless real boson contributes \(\pi^2T^4/90\) to the pressure;
therefore three Goldstone bosons give \(p_\pi/T^4=3\pi^2/90=\pi^2/30\).  The
remaining thermodynamic identities follow from the conformal massless-gas
relations, while finite quark masses enter through the chiral mass operator and its computed low-energy constants.
\end{proof}

\begin{definition}[Thermal QCD export object]
\label{p11:def:thermal-export-object}
The thermal QCD export object is
\begin{equation}
\label{p11:eq:thermal-export-object}
\begin{aligned}
 \mathcal E_{\rm QCD}(T,\mu)=\bigl(&p,\epsilon,s,c_s^2,
 \chi_B,\chi_Q,\chi_S,\\
 &T_{c,L}(h)\ \text{or its controlled limit},\ \Gamma_{\rm had}\bigr),
\end{aligned}
\end{equation}
where
\begin{equation}
\label{p11:eq:thermal-derivatives}
\begin{aligned}
 \varpi&=\frac{T}{V}\log Z_{\rm QCD}^{\rm fin},\qquad
 p=p_{\rm mech}=T\partial_V\log Z_{\rm QCD}^{\rm fin},\\
 \epsilon&=Ts-\varpi+\mu_Bn_B+\mu_Qn_Q+\mu_Sn_S,\\
 c_s^2&=\left.{dp\over d\epsilon}\right|_{\rm declared\ path}.
\end{aligned}
\end{equation}
and conserved-charge susceptibilities are derivatives of the actual
grand-potential density $\varpi$. On the certified extensive branch
$p=\varpi$ and $s=\partial_Tp$; otherwise finite entropy and mechanical
pressure retain Theorem~\ref{p11:thm:native-thermodynamic-response-bridge}.
\end{definition}


\begin{definition}[Thermal reference coordinate; no native shell temperature]
\label{p11:def:native-scalar-shell-thermal-activation-scale}
For any chosen positive reference \(T_{\rm ref}\), one may introduce the coordinate \(u=\log(T/T_{\rm ref})\).  The choice of \(T_{\rm ref}\) is a coordinate convention and does not define a physical crossover.  The repaired Paper-VII-A core exports no theorem equating a scalar-shell count or amplitude with a temperature.
\end{definition}

\begin{definition}[Finite chiral crossover and continuum transition target]
\label{p11:def:native-thermal-crossover-readout}
At finite level \(L\), choose a \emph{dimensionful Hamiltonian} chiral source \(h\), with \(H_L(h)=H_L(0)+h\mathcal Q_L+\cdots\), \(Z_L(T,h)=\operatorname{Tr}e^{-H_L(h)/T}\), and \(V_L>0\). Define the corresponding scalar-density response
\[
\Sigma_L(T,h):=-\frac{T}{V_L}\,\partial_h\log Z_L(T,h)
=\frac1{V_L}\langle\partial_hH_L(h)\rangle_{T,h},
\]
and a chosen thermal response such as
\begin{equation}
\label{p11:eq:canonical-log-scalar-response}
\chi_{\Sigma,L}(T,h):=\left|T\,\partial_T\Sigma_L(T,h)\right|.
\end{equation}
On a compact temperature interval, a finite crossover readout is any maximizer
\begin{equation}
\label{p11:eq:canonical-native-crossover-selector}
T_{c,L}(h)\in\operatorname*{arg\,max}_{T\in I}\chi_{\Sigma,L}(T,h).
\end{equation}
\label{p11:eq:canonical-native-crossover-fixed-point}
If instead \(h\) is declared as a dimensionless source \emph{in the exponent} \(Z=\operatorname{Tr}e^{-H/T+h\mathcal Q}\), then \(V_L^{-1}\partial_h\log Z\) is the appropriate distinct convention; its susceptibility and selected maximizer need not coincide with the dimensionful-Hamiltonian-source readout above. A genuine chiral phase-transition temperature, if it exists, is a nonanalyticity arising only after the specified thermodynamic/refinement and \(h\downarrow0\) limits.  No reflection of a logarithmic coordinate selects it.
\end{definition}

\begin{theorem}[No count-derived thermal activation scale]
\label{p11:prop:thermal-activation-shell-amplitude}

The repaired Paper-VII-A core supplies no theorem equating a finite shell coefficient with a temperature.  At finite level, all thermal scales must be extracted from derivatives or spectral features of
\[
Z_L(T)=\operatorname{Tr}e^{-H_L/T}
\]
for the supplied Hamiltonian.  The predecessor relation involving \(24/15\), \(\sqrt\delta\), and a wall scale is withdrawn.
\phantomsection\label{p11:eq:thermal-activation-fpi-lock}
\end{theorem}
\begin{proof}

Multiplying a dimensionless finite count by an arbitrary external energy unit yields a continuum of temperatures with identical geometric data.  Therefore a physical temperature requires a dimensionful Hamiltonian/source member.  Thermal responses computed from that member are deterministic and need no extra shell selector.
\end{proof}

\begin{theorem}[Finite thermal analyticity and crossover no-go]
\label{p11:prop:native-transition-scaling}
\label{p11:prop:transition-scale-reconciliation}
\label{p11:thm:scalar-shell-saturation-canonical-crossover}

For finite-dimensional \(H_L\), \(Z_L(\beta)=\operatorname{Tr}e^{-\beta H_L}\) is strictly positive and real analytic for every \(\beta>0\).  Hence no genuine nonanalytic phase transition occurs at finite \(L\).  A finite crossover location may be defined by a response extremum, but a coordinate reflection \(u\mapsto-u\) does not determine that extremum.
\phantomsection\label{p11:eq:native-transition-scaling}
\phantomsection\label{p11:eq:thermal-crossover-shell-factor}
\phantomsection\label{p11:eq:thermal-selfdual-equivalence}
\phantomsection\label{p11:eq:zeta-thermal-saturation-one}
\phantomsection\label{p11:eq:canonical-crossover-numeric-symbolic}
\end{theorem}
\begin{proof}

Diagonalize \(H_L\): \(Z_L(\beta)=\sum_j e^{-\beta E_j}\), a finite sum of entire functions and positive on the positive real axis.  Thus \(\log Z_L\) and all finite derivatives are analytic.  For non-uniqueness of a coordinate-fixed crossover, the family \(\chi(u)=e^{-(u-a)^2}\) has its maximum at arbitrary \(a\) while the coordinate involution retains fixed point zero.
\end{proof}

\begin{proposition}[Hadronic rates as native thermal averages]
\label{p11:prop:hadronic-rates-native-thermal-averages}
Hadronic reaction kernels exported to cosmology must be thermal averages of native hadron masses, cross sections, and matrix elements:
\begin{equation}
\label{p11:eq:hadronic-rate-native-average}
\Gamma_{\rm had}(T,\mu)
=
\operatorname{ThermalAverage}_{T,\mu}
\left[M_h,\langle h'|\mathcal O|h\rangle,\sigma_{hh'}(s),\mathcal F_{\rm had}\right].
\end{equation}
They may not be imported from BBN or heavy-ion tables if the claim is finite QCD readout closure.
Here any \(\sigma_{hh'}\) asserted as a scattering cross section is admitted
only on the \textnormal{(SM1)--(SM3)} plus \textnormal{(SM7)} domain of
Definition~\ref{p12:def:claim-indexed-smat-package}, with
\textnormal{(SM4)} additionally if an LSZ representation is used, and after a
selected infrared route when applicable; the masses and local matrix elements
themselves do not consume those hypotheses.
\end{proposition}
\begin{proof}
A cosmological hadronic rate is obtained by summing or integrating QCD
transition probabilities over the thermal ensemble.  The masses and local
matrix elements are readouts of \(\mathfrak S_{\rm QCD}^{\rm BH}\) by
Theorem~\ref{p10:thm:master-finite-nonperturbative-qcd-closure}; the cross
sections and inclusive distributions are downstream claim-scoped limits of
those readouts.  BBN or heavy-ion tables are downstream
comparison/propagation products; using them to choose the listed QCD ingredients
would reverse the readout-chronology direction of Theorem~\ref{p10:thm:qcd-readout-chronology}.
\end{proof}

\begin{theorem}[Cosmological export boundary]
\label{p11:thm:cosmological-export-boundary}
Paper~VII-B may export \(\mathcal E_{\rm QCD}(T,\mu)\), nucleon/pion/nuclear boundary data, and hadronic rates to Paper~VI-B, but Paper~VI-B survey fits, CMB likelihoods, BBN abundance comparisons, and structure-formation data may not select the Paper~VII-B QCD kernels.
\end{theorem}
\begin{proof}
QCD determines thermodynamic and hadronic inputs.  Cosmological pipelines propagate them and compare with data.  If cosmological agreement is used to tune the QCD kernel, the QCD kernel is no longer native to the strong-sector carrier.\qedhere
\end{proof}



% Source Paper VII-B, original line 4029; inherited proof bodies unchanged.
\subsection{Nuclear closed-shell transfer outlook}
\label{p11:app:nuclear-closed-shell-transfer-outlook}

This section records a nuclear-physics outlook consequence of the selected-cell geometry of Paper~I-A and the finite Gram-reduction rule used in the Paper~VII-B transfer readout.  It is part of the finite nuclear-readout domain of this paper and records a finite closed-shell readout target generated by the same native selector machinery.  The statement is deliberately finite.  The compact Born--Hopf carrier is not moved from the selected value
\[
        m_{\mathsf B}=1,
        \qquad
        \mathcal H_{\mathsf B}=\mathcal H_1 .
\]
The later index below is therefore not a new compact truncation level.  It is a finite symmetric readout degree over the coordinate support already present in the selected cell.

\begin{definition}[Selected finite shell-readout support]
\label{p11:def:selected-finite-shell-readout-support}
Let
\[
        X:=V_{1/2}\otimes V_{1/2}
\]
be the selected coordinate sector imported from Paper~I-A.  Thus \(\dim X=4\).  For \(r\geq 0\), define the raw degree-\(r\) finite shell-readout support by
\[
        \mathcal B_r:=\operatorname{Sym}^r(X).
\]
Equivalently, \(r\) is the homogeneous symmetric transfer degree of the fixed selected coordinate sector.  It is not the truncation parameter \(m\) in the family \(\mathcal H_m\).
\end{definition}

\begin{lemma}[Raw selected-cell shell count]
\label{p11:lem:raw-selected-cell-shell-count}
For every \(r\geq0\),
\[
        \dim \mathcal B_r
        =
        \binom{r+3}{3}.
\]
In particular,
\[
        \dim\mathcal B_0,\ldots,\dim\mathcal B_6
        =
        1,4,10,20,35,56,84.
\]
\end{lemma}

\begin{proof}
This is Definition~\ref{p1:def:raw-symmetric-coordinate-degree}, namely
\(\mathcal B_r=\Sym^r(X)\) with
\(\dim\mathcal B_r=\binom{r+3}{3}\).  Paper~VII-B imports the rank sequence and
adds only the readout interpretation.
\end{proof}

\begin{remark}[Why the first three terms are already selected]
\label{p11:rem:first-three-terms-selected}
Paper~I-A identifies the mathematical level-one carrier with
\[
        \mathcal H_1
        =
        \mathcal S_0\oplus\mathcal S_{1/2}\oplus\mathcal S_1,
\]
with shell dimensions
\[
        1,
        \qquad
        4,
        \qquad
        10.
\]
Thus the first three raw readout degrees agree with the retained selected-cell shell filtration.  No boundary-rank subtraction is present before the first coordinate-closure boundary.
\end{remark}

\begin{lemma}[First forbidden coordinate-exclusion boundary]
\label{p11:lem:first-forbidden-coordinate-exclusion-boundary}
Let
\[
        \mathcal K
        :=
        (V_{1/2}\otimes V_0)
        \oplus
        (V_0\otimes V_{1/2})
\]
be the protected chiral block of the selected cell.  Coordinate tensoring by \(X=V_{1/2}\otimes V_{1/2}\) gives
\[
        X\otimes (V_{1/2}\otimes V_0)
        \cong
        (V_0\otimes V_{1/2})
        \oplus
        (V_1\otimes V_{1/2}),
\]
and
\[
        X\otimes (V_0\otimes V_{1/2})
        \cong
        (V_{1/2}\otimes V_0)
        \oplus
        (V_{1/2}\otimes V_1).
\]
The returned summands lie in \(\mathcal H_1\).  The exterior summands
\[
        V_1\otimes V_{1/2},
        \qquad
        V_{1/2}\otimes V_1
\]
have total weight \(3/2\) and therefore lie beyond the fixed selected cell.  Each exterior chiral boundary has dimension
\[
        \dim V_1\,\dim V_{1/2}=3\cdot2=6.
\]
\end{lemma}

\begin{proof}
This is Proposition~\ref{p1:prop:first-excluded-chiral-boundary}.  Paper~VII-B
imports the two exterior six-dimensional chiral boundary summands and assigns
them only their finite-transfer status.
\end{proof}

\begin{definition}[Oriented coordinate-boundary face]
\label{p11:def:oriented-coordinate-boundary-face}
Adopt the Paper~I-A finite representative and normal-line choice of Definition~\ref{p1:def:oriented-boundary-face}. Paper~I-A supplies oriented abstract three-dimensional factors \(F_\sigma\subset X\) in a \emph{chosen} splitting \(X\cong N_\sigma\oplus F_\sigma\). The Hurwitz trimer comparison spaces
\[
        E_+\oplus E_- ,\qquad\dim E_+=\dim E_-=3,
\]
are separate until an explicit orientation-compatible identification \(j_\sigma:E_\sigma\xrightarrow{\sim}F_\sigma\) is chosen. Such a comparison isomorphism is presentation data, not a canonical source-generated physical embedding.  The protected commutant line on \(\mathcal K\) is generated by
\[
        \Gamma_{\mathcal K}=P_L-P_R,
\]
so the only non-scalar native selector on the protected block is the already fixed oriented chiral split.
Fix one selected-cell orientation sign \(\sigma\in\{+,-\}\).  Let \(N_\sigma=\CC\eta_\sigma\) be the Paper-I-A selected normal-symbol line.  For \(r\geq3\), import the corresponding oriented boundary face
\[
        \mathcal L_r^\sigma
        :=
        N_\sigma\odot\operatorname{Sym}^{r-1}(j_\sigma(E_\sigma))
        =N_\sigma\odot\operatorname{Sym}^{r-1}(F_\sigma),
\]
where \(\eta_\sigma\) denotes the first excluded chiral normal symbol produced by the coordinate-closure boundary.  For \(r<3\), set \(\mathcal L_r^\sigma=0\).
The opposite sign gives the orientation-reversed selected cell; it is not an additional simultaneous deletion subspace in the same readout. The ranks are independent of a basis choice, but the \emph{embedded} subspace and its annihilation by an actual nuclear source map \(S_r\) depend on the declared splitting and \(j_\sigma\); the Gram-kernel theorem below is conditional on that \emph{same} embedded image.
\end{definition}

\begin{lemma}[Triangular rank of the coordinate boundary]
\label{p11:lem:triangular-rank-coordinate-boundary}
For \(r\geq3\),
\[
        \dim \mathcal L_r^\sigma
        =
        \binom{r+1}{2}.
\]
Thus
\[
        \dim \mathcal L_3^\sigma,
        \dim \mathcal L_4^\sigma,
        \dim \mathcal L_5^\sigma,
        \dim \mathcal L_6^\sigma
        =
        6,10,15,21.
\]
\end{lemma}

\begin{proof}
This is the boundary-rank statement in
Theorem~\ref{p1:thm:boundary-reduced-ranks}, after adopting the finite
representative and normal-line choice of
Definition~\ref{p1:def:oriented-boundary-face}.
\end{proof}

\begin{theorem}[Actual nuclear Gram-kernel certificate]
\label{p11:thm:v2211-nuclear-gram-kernel-certificate}
Let $\mathcal L_r^\sigma\subset\mathcal B_r$ be the imported geometric
face and let $S_r:\mathcal B_r\to\mathcal H_{r,\rm trial}^{\rm nuc}$
be the actual finite nuclear source/transfer map, with $G_r=S_r^\dagger S_r$.
If its inherited source map satisfies
\[
 S_r\mathcal L_r^\sigma=0,\qquad
 G_r\ge g_r P_{(\mathcal L_r^\sigma)^\perp},\quad g_r>0,
\]
then $\ker S_r=\mathcal L_r^\sigma$ exactly and its Gram reduction has
rank $\dim\mathcal B_r-\dim\mathcal L_r^\sigma$.
For an actual finite fermionic occupation realization with two independent
spin slots per retained mode, the capacity is twice that rank. Identification
of these capacities with physically ordered nuclear closed shells additionally
consumes the corresponding actual nuclear Hamiltonian ordering/gap readout.
\end{theorem}
\begin{proof}
The first condition puts the face in the kernel. For $S_rv=0$, the lower
bound gives $0=\langle v,G_rv\rangle\ge
g_r\|P_{(\mathcal L_r^\sigma)^\perp}v\|^2$, so $v$ lies in the face.
Rank-nullity gives the rank. Exterior occupation of a rank-$d$ one-body
space with two supplied spin slots has maximal occupation $2d$ by CAR.
Neither rank-nullity nor CAR orders the energies of that nuclear state.
\end{proof}

\begin{proposition}[No-insertion status of the coordinate boundary]
\label{p11:prop:no-insertion-status-coordinate-boundary}
On the actual source-map branch satisfying
Theorem~\ref{p11:thm:v2211-nuclear-gram-kernel-certificate}, passage from
$\mathcal B_r$ to the boundary-reduced support below inserts neither a new
projector nor a new truncation level. It uses the verified Gram kernel and
the following native data:
\[
        m_{\mathsf B}=1,
        \quad
        X=V_{1/2}\otimes V_{1/2},
        \quad
        \mathcal K,
        \quad
        \Gamma_{\mathcal K},
        \quad
        E_\sigma,
\]
and the finite Gram reduction of null directions in the same-channel transfer readout.  In particular, the forbidden face \(\mathcal L_r^\sigma\) is an excluded boundary of the fixed selected cell, not a retained sector of a larger \(\mathcal H_m\).
\end{proposition}

\begin{proof}
Paper~III fixes the physically realized compact momentum carrier at \(m_p=1\) and proves that the primitive branch has no momentum-truncation flow.  Paper~I-A supplies the mathematical \(\mathcal H_1\) finite-cell structure and its first excluded coordinate boundary.  Hence the exterior coordinate-exclusion summands at total weight \(3/2\) are not promoted to retained states by replacing the realized carrier with a higher mathematical truncation \(\mathcal H_m\).  No-insertion also forbids adding a new channel projector that would remove an arbitrary sector.  The only available non-scalar selector on the protected block is the commutant line generated by \(\Gamma_{\mathcal K}\), and the only available rank-three finite support is the selected oriented support \(E_\sigma\).  The face \(\mathcal L_r^\sigma\) is therefore the native boundary generated by the first excluded chiral normal direction and its tangential propagation over \(E_\sigma\).  Reversing the root orientation exchanges \(\sigma\) with \(-\sigma\) and gives an isomorphic selected-cell readout with the same rank; no-insertion does not license deleting both oriented faces at once.

In the finite transfer construction, null Gram directions are deleted before forming the effective transfer block.  The imported geometric description alone does not prove the kernel of
a later symmetric nuclear source map. On the stated branch the actual
$S_r\mathcal L_r^\sigma=0$ and positive complementary Gram bound do prove
that kernel, by Theorem~\ref{p11:thm:v2211-nuclear-gram-kernel-certificate}.
Thus deleting it is the inherited Gram reduction. If this source-map test
has not been executed, only the geometric quotient ranks below are derived;
the actual nuclear occupation identification remains conditional.
\end{proof}

\begin{definition}[Boundary-reduced finite shell support]
\label{p11:def:boundary-reduced-finite-shell-support}
The selected boundary-reduced degree-\(r\) support is
\[
        \mathcal N_r
        :=
        \begin{cases}
        \mathcal B_r, & 0\leq r\leq2,\\[3pt]
        \mathcal B_r/\mathcal L_r^\sigma, & r\geq3.
        \end{cases}
\]
The sign \(\sigma\) records the fixed selected-cell orientation.  The two signs are exchanged by the orientation-reversing interchange of the two chiral factors and have the same rank.  They are not summed or quotiented simultaneously in one selected-cell readout.
\end{definition}

\begin{theorem}[Imported boundary ranks and native occupation readout]
\label{p11:thm:native-boundary-reduced-shell-count}
The selected finite shell-readout ranks are
\[
        \dim\mathcal N_r
        =
        \begin{cases}
        \binom{r+3}{3}, & 0\leq r\leq2,\\[5pt]
        \binom{r+3}{3}-\binom{r+1}{2}, & r\geq3.
        \end{cases}
\]
Consequently,
\[
        \dim\mathcal N_0,
        \ldots,
        \dim\mathcal N_6
        =
        1,4,10,14,25,41,63.
\]
On the certified Gram-map and twofold finite fermionic occupation branch
of Theorem~\ref{p11:thm:v2211-nuclear-gram-kernel-certificate}, this gives
the finite capacities
\[
        2\dim\mathcal N_0,
        \ldots,
        2\dim\mathcal N_6
        =
        2,8,20,28,50,82,126.
\]
The first corrected count \(28\) means that the raw degree-three capacity
\(40\) is reduced to \(28\) retained occupation slots by a rank-\(12\)
doubled boundary exclusion.  It is not an amount of matter transferred to
reciprocal position and is not a gravitational activation threshold.
\end{theorem}

\begin{proof}
Paper~I-A, Theorem~\ref{p1:thm:boundary-reduced-ranks}, proves the finite quotient
ranks
\[
  \dim\mathcal N_0,\ldots,\dim\mathcal N_6=1,4,10,14,25,41,63
\]
after the stated orientation and normal-line choices.  Paper~VII-B does not
rederive those geometric ranks.  Its occupation step consumes the actual Gram-kernel certificate and the
licensed twofold finite fermionic readout; multiplying the imported ranks
then gives $2,8,20,28,50,82,126$. Without the actual map certificate this
calculation is the geometric capacity target, not an executed nuclear
closed-shell theorem.
\end{proof}

\begin{remark}[Theorem output and downstream transfer]
\label{p11:rem:boundary-reduced-shell-count-status}
The geometric rank sequence is a Paper~I-A theorem.  Paper~VII-B proves that a source map passing the actual Gram-kernel test
and a licensed twofold occupation readout use that quotient without an
inserted projector. The geometric quotient itself does not execute the test.  The imported ranks are
\[
        1,4,10,14,25,41,63.
\]
It does not by itself compute nuclear binding energies, isotope stability windows, deformation corrections, or continuum scattering data.  Those require the corresponding finite Hamiltonian or transfer block.  The point is narrower: on the verified map/occupation branch the finite
capacities require no new shell projector or momentum-carrier change.
Their arithmetic is the rank of the boundary-reduced selector support.
Their actual occurrence and energetic nuclear closed-shell ordering are
readouts of the nuclear Hamiltonian, not consequences of rank arithmetic.  By the Paper~III no-boundary-transfer theorem, the excluded face is not mapped to a position-space field; gravitational source status is owned separately by Paper~VI-A through metric variation of the generated reciprocal-position action.
\end{remark}



% Source Paper VII-B, original line 5157; inherited proof bodies unchanged.
\subsubsection{Complete thermal/source contacts}
\begin{theorem}[Duhamel response with the actual nonlinear source Hessian]
\label{p11:thm:v2210-lit-complete-duhamel-contacts}
Let $H(J)$ be the actual twice differentiable finite Hermitian member and
$Z(J)=\operatorname{Tr}e^{-\beta H(J)}$ at fixed $\beta>0$. With the
Duhamel form defined in
Theorem~\ref{p10:thm:v2210-lit-ward-defect},
\[
 \boxed{\partial_a\log Z=-\beta\langle H_a\rangle,
 \qquad \partial_a\partial_b\log Z
 =-\beta\langle H_{ab}\rangle
    +\beta^2\mathcal C_\rho(H_a,H_b).}
\]
Consequently the Hessian is positive semidefinite for an affine Hermitian
charge source $H(J)=H_0-J^aQ_a$. For a general nonlinear source or metric
pullback, positivity applies to the Duhamel term and does not discard the
contact term $-\beta\langle H_{ab}\rangle$.
\end{theorem}
\begin{proof}
Differentiate the exponential by Duhamel's integral. Dividing by $Z$ gives
the first identity; differentiating the expectation of $H_a$ gives the
printed contact and connected response. Positivity of the Duhamel quadratic
form follows from its positive logarithmic means. For the sharp countertest,
$H(h)=H_0+ch^2I$, $c>0$, has
$\partial_h^2\log Z|_0=-2\beta c<0$ although all commutators are unchanged.
\end{proof}
Section~\ref{p11:sec:v2201-derived-d24} remains the affine-source result.
The present theorem covers the complete nonlinear source/metric jets and
prevents positive finite charge susceptibility from being promoted to a
nonzero condensate, a pressure value or a nuclear contact coefficient.
Paper~VII-A's sector cumulants in
Theorem~\ref{p10:thm:v2210-lit-sector-cumulants} are ordinary charge
cumulants of their declared sector weights; their fourth cumulant need not
be positive. None of these auxiliary source operations inserts a new
physical strong-CP member.


% Source Paper VII-B, original line 5402; inherited proof bodies unchanged.
\subsubsection{D24: Complete finite thermal susceptibility is a positive Duhamel form}
\label{p11:sec:v2201-derived-d24}
\noindent\textit{Status: proved result on stated domain.}

For $Z(\beta,\mu)=\operatorname{Tr}e^{-\beta(H-\mu_aQ_a)}$,
\[
 \partial_{\mu_a}\log Z=\beta\langle Q_a\rangle,\qquad
 \partial_{\mu_a}\partial_{\mu_b}\log Z
 =\beta^2\int_0^1\operatorname{Tr}
 (\rho^s\delta Q_a\rho^{1-s}\delta Q_b)\,ds.
\]
The real symmetric susceptibility matrix is positive semidefinite.
\begin{proof}
Differentiate the exponential by its Duhamel integral.
In the density eigenbasis the quadratic form for a Hermitian $A$ is
$\sum_{ij}(p_i-p_j)/(\log p_i-\log p_j)|A_{ij}|^2$, using $p_i$ on
the diagonal. All logarithmic means are nonnegative. Subtract the first
derivative product to obtain the connected form.
\end{proof}
Ordinary charge covariance is sufficient when the charges commute with the
Hamiltonian; otherwise it is the wrong response. The independent finite
partition-function derivative check reproduces the formula.
Metric pressure still uses $-\langle\partial_aH\rangle/(3\mathcal Va^{-1})$.
A central action term changes energy and pressure despite leaving all
commutators unchanged. BBN/nuclear/thermal descendants must therefore
consume the actual action entries, temperature/history and reaction source
jets; this positive finite test cannot fill in absent nuclear matrix elements.
\par\noindent\textit{Inherited inputs:} \ref{p3f:thm:v169-finite-source-work}, \ref{p11:cor:thermal-bqs-susceptibility-endpoint}, \ref{p10:def:finite-qcd-compilation-datum}.

% END GEOMETRY-FIRST CONSEQUENCES v2.201

% BEGIN FRAMEWORK v2.220: BARYON SPIN--FLUX SPECTRAL READOUT BOUNDARY


\section{Standard extraction tools and physical comparison boundaries}
\label{p11:sec:v237-b7}
External scattering, lattice, chiral or factorization methods are applied only with their precise hypotheses; they do not retroactively select a Born--Hopf source member.

% Source Paper VII-B, original line 3419; inherited proof bodies unchanged.
\subsubsection{Imported standard machinery and physical-pole completion blocks}
\label{p11:sec:physical-pole-completion-blocks}

The preceding Blocks~B--F localize the remaining hard work to finite spectral execution.  This subsection records the standard QCD machinery that may be imported as theorem-bank material, and separates it from the Born--Hopf-native matrix elements that still have to be evaluated from \(\mathfrak C_L^{\rm QCD}\to\mathfrak M_L^{\rm QCD}\).  The purpose is to avoid two opposite errors: re-proving standard scattering, projection, anomaly, and factorization technology, and also importing lattice or phenomenological numbers as if they were native Born--Hopf derivations.

\begin{theorem}[Standard-method import boundary]
\label{p11:thm:standard-method-import-boundary}
\begin{sloppypar}
The following external results may be used in Paper~VII-B only as comparison/extraction machinery, not as selectors for Born--Hopf parameters:
\begin{enumerate}[label=(\roman*), leftmargin=2.2em]
\item Feshbach--Schur projection and energy-dependent effective Hamiltonians for same-channel pole dressing~\cite{Feshbach1958,Feshbach1962};
\item finite-volume scattering and resonance extraction from finite spectra in the sense of Luescher~\cite{Luscher1986Volume,Luscher1991TwoParticle,Luscher1991Unstable};
\item Witten--Veneziano/topological-susceptibility anomaly technology~\cite{Witten1979CurrentAlgebra,Veneziano1979};
\item chiral perturbation theory and resonance saturation of low-energy constants~\cite{GasserLeutwyler1984,EckerGasserPichDeRafael1989};
\item trace-anomaly scalar low-energy theorems and scalar glue-current normalization~\cite{NovikovShifmanVainshteinZakharov1981,ShifmanVainshteinZakharov1978};
\item factorization, light-cone evolution, GPD and angular-momentum
decompositions, and Euclidean-lattice quasi-distribution extraction as
downstream transport and comparison methods
\cite{EfremovRadyushkin1980,Ji1997,Radyushkin1997,Ji2013LaMET};
\item lattice glueball calculations and thermal QCD equation-of-state technology
as comparison maps~\cite{MorningstarPeardon1999,Borsanyi2014EoS}.
\end{enumerate}
These imports are therefore domain-stated extraction theorems: they may process native finite eigenvalues, residues, currents, or spectral measures, but they may not provide those quantities.
None of these imports supplies a value for a native Born--Hopf matrix element such as
\begin{equation}
M_G,
\qquad
\xi=\langle G|\mathcal M^2|S_1\rangle,
\qquad
\rho_G=\frac{\langle0|T^\mu{}_{\mu}|G\rangle}{M_G^2},
\qquad
\omega_A,
\qquad
M_{A,\perp},
\qquad
\Sigma_{\rm native}(s).
\end{equation}
Those quantities are fixed only by the compiled finite Born--Hopf QCD transfer and its operator residues.
\end{sloppypar}
\end{theorem}
\begin{proof}
The cited literature gives projection identities, finite-volume quantization
formulae, anomaly relations, chiral/resonance effective operators,
factorization/evolution equations, and lattice comparison procedures.  Each
statement is an equality or extraction theorem once its full domain hypotheses
and its input Hamiltonian, current, finite volume, or spectral density have
been specified.  Paper~VII-B specifies a different input object, namely the finite
Born--Hopf singlet strong Hamiltonian.  Therefore the imported theorem bank can
process native finite eigenvalues and residues into continued correlator poles,
PDF operators, thermal cumulants, and comparison observables.  Phase shifts,
decay rates, and cross sections additionally use the claim-indexed Paper-VII-C
scattering package.  None of these methods determines the native finite matrix
elements without executing the Born--Hopf compiler.
\end{proof}


\begin{proposition}[Hypothesis audit for standard-method imports]
\label{p11:prop:hypothesis-audit-standard-imports}
The external theorem bank in Theorem~\ref{p11:thm:standard-method-import-boundary} is admissible only under the following native-input checks.
\begin{enumerate}[label=(\alph*), leftmargin=2.2em]
\item Feshbach--Schur imports require a fixed block decomposition \(P\oplus Q\) of the compiled finite Hamiltonian and invertibility of \(z-H_{QQ}\) on the energy window.  The effective operator is then
\[
H_{\rm eff}(z)=H_{PP}+H_{PQ}(z-H_{QQ})^{-1}H_{QP},
\]
and pole positions and residues are Schur roots/readouts, not inserted masses.
\item Chiral perturbation theory imports require the native symmetry-breaking data
\[
SU(3)_L\times SU(3)_R\to SU(3)_V,
\qquad F_\pi,
\qquad B_0,
\qquad \Sigma_q,
\]
and the inherited renormalized physical quark-mass/source operator obtained from the derived charged preterminal source member through the certified G16 RG/threshold transport before any chiral expansion is invoked.
\item Witten--Veneziano and anomaly imports require native \(\chi_t\), \(F_\pi\), and \(N_f\).  The \(\eta'\) contribution is then an anomaly readout; it is not supplied by phenomenology.
\item Finite-volume scattering imports require native two-hadron levels
      \(E_{2h,L}(P)\), the finite geometry map, locality/finite-range control,
      a specified elastic window away from thresholds, and the channel
      unitarity/partial-wave hypotheses of the quantization theorem.  Calling
      the extracted quantity a physical phase shift or \(S\)-matrix pole also
      consumes the applicable clauses of
      Definition~\ref{p12:def:claim-indexed-smat-package}.
\item Factorization and light-cone imports require fixed native Wilson-line
      operators and finite hadron eigenstates.  DGLAP/ERBL/timelike kernels
      transport native initial measures; they do not fit those measures.  An
      operator evolution claim needs no \(S\)-matrix hypothesis, whereas an
      asymptotic process observable uses
      Theorem~\ref{p12:thm:scattering-hypothesis-firewall}.
\item Thermal and lattice-comparison imports require a fixed finite spectrum, charge operators, volume convention, and operator normalization before any comparison table or continuum-extrapolation technology is used.
\end{enumerate}
\end{proposition}
\begin{proof}
Each listed external method is a domain-stated extraction theorem.  It maps a specified Hamiltonian, current, finite-volume spectrum, topological susceptibility, or operator measure to a downstream observable.  If the required native input is absent, using the method would silently insert the missing quantity from outside.  Requiring the checks above preserves the readout-chronology criterion: standard QCD machinery may process Born--Hopf readouts, but it cannot select Born--Hopf matrix entries or residues.
\end{proof}

\begin{theorem}[Dimensionless scalar Hodge--Schur ledger]
\label{p11:thm:exact-scalar-hodge-schur-ledger}

The scalar Hodge spectral measure, its exact moments, and any finite Schur complement constructed solely from the dimensionless Hodge/incidence matrices are unconditional finite geometry.  A physical energy ledger is instead obtained by a generalized eigenproblem pairing the geometric metric with the supplied Yang--Mills quadratic form.  No multiplication by a universal \(\Lambda_{\rm BH}\) is permitted.
\phantomsection\label{p11:eq:scalar-hodge-resolvent}
\phantomsection\label{p11:eq:scalar-hodge-ledger}
\end{theorem}
\begin{proof}

The Hodge moments are traces/matrix elements of fixed finite incidence matrices.  Generalized eigenvalues of the pair \((Q_{\YM,L},G_H)\) are invariantly defined once the physical quadratic form is supplied.  Rescaling the latter changes energies while leaving every Hodge moment fixed, proving the separation.
\end{proof}

\begin{theorem}[Hodge Krylov/Ritz block and physical generalized eigenproblem]
\label{p11:thm:native-scalar-curvature-krylov-block}

Lanczos/Krylov coefficients generated from a normalized scalar Hodge source are exact dimensionless numbers.  Their Ritz values are variational approximants to the corresponding Hodge operator.  Physical scalar-curvature energies require applying the same subspace to the supplied Yang--Mills Hamiltonian/quadratic-form pair and solving the resulting generalized eigenvalue problem.
\phantomsection\label{p11:eq:scalar-hodge-central-moments-extended}
\phantomsection\label{p11:eq:scalar-krylov-first-vector}
\phantomsection\label{p11:eq:scalar-krylov-first-coupling}
\phantomsection\label{p11:eq:scalar-krylov-second-residual}
\phantomsection\label{p11:eq:scalar-curvature-three-state-krylov-block}
\phantomsection\label{p11:eq:scalar-curvature-three-state-ritz-values}
\phantomsection\label{p11:eq:scalar-curvature-three-state-ritz-energies}
\end{theorem}
\begin{proof}

The Lanczos recursion follows from repeated application and orthogonalization of the finite Hodge matrix.  Rayleigh--Ritz gives convergence/interlacing properties inside the generated subspace.  An independent physical quadratic form is necessary to attach energy units.
\end{proof}

\begin{theorem}[Compact scalar glue--flavor-singlet mixing from correlators]
\label{p11:thm:compact-scalar-glue-singlet-mixing}

In an exact flavor-symmetric normalization, a flavor-blind glue source couples to the normalized nonstrange/strange singlet directions with ratios fixed by their flavor multiplicities; departures are source-overlap effects.  The physical scalar mixing matrix is the Hermitian/Schur mass or inverse-propagator matrix extracted from the computed scalar correlators.  Its diagonal masses and glue couplings are not fixed by \(4\pi F_\pi\) or vector-sector splittings.
On the stated leading flavor-symmetric compact branch, choose the
two-channel phases so its real reduced glue coupling is $\xi$, and write
\begin{equation}
\label{p11:eq:compact-scalar-mixing-block}
 \mathcal M_P^2=
 \begin{pmatrix}
 M_G^2&\sqrt{2/3}\,\xi&\sqrt{1/3}\,\xi\\
 \sqrt{2/3}\,\xi&M_0^2&0\\
 \sqrt{1/3}\,\xi&0&M_0^2+2\Delta_s
 \end{pmatrix}.
\end{equation}
This is the compact ansatz consumed by the following cubic; a more general
computed correlator block retains its actual additional entries.
\phantomsection\label{p11:eq:xi-native-transfer}
\end{theorem}
\begin{proof}

Flavor symmetry determines only Clebsch/normalization factors.  The diagonal scalar energies and reduced glue coupling are independent correlator matrix elements.  Diagonalization of the resulting finite matrix gives the mixed poles, while arbitrary symmetry-compatible diagonal entries show that no unique numerical matrix follows from flavor counting.
\end{proof}

\begin{corollary}[Compact scalar cubic]
\label{p11:cor:compact-scalar-cubic}
With \(y=s=E^2\), the compact scalar poles before open-channel dressing obey
\begin{equation}
\label{p11:eq:compact-scalar-cubic}
(y-M_G^2)(y-M_0^2)(y-M_0^2-2\Delta_s)
-
\xi^2\left(y-M_0^2-\frac{4}{3}\Delta_s\right)=0.
\end{equation}
Equivalently, the glue-channel Schur equation is
\begin{equation}
\label{p11:eq:compact-scalar-schur-single}
y-M_G^2-
\xi^2\left(\frac{2/3}{y-M_0^2}+\frac{1/3}{y-M_0^2-2\Delta_s}\right)=0.
\end{equation}
When $M_G^2$ is distinct from both quarkonium entries, the analytic
glue-dominant branch has
\[
 y_G=M_G^2+\xi^2\left[
 {2/3\over M_G^2-M_0^2}+{1/3\over M_G^2-M_0^2-2\Delta_s}
 \right]+O(\xi^4).
\]
The shift is positive if the actual $M_G^2$ lies above both entries and
negative if it lies below both; between them its sign is the displayed
signed sum. No ordering is supplied by the finite carrier. A downward
shift need not come from an open channel unless the compact positive-shift
ordering has first been established.
\end{corollary}
\begin{proof}
Compute the determinant of \eqref{p11:eq:compact-scalar-mixing-block}
or eliminate its quarkonium block. At a nondegenerate $\xi=0$ glue root,
the implicit-function theorem gives an even analytic branch in $\xi$.
Inserting $y=M_G^2+O(\xi^2)$ into the Schur equation gives the coefficient
and its stated signs. Degenerate cases require the actual mixed block and
are not covered by this nondegenerate expansion.
\end{proof}

\begin{theorem}[Scalar two-Goldstone unitarity row with computed couplings]
\label{p11:thm:scalar-two-goldstone-trace-row}

Once scalar-to-two-Goldstone couplings and direct trace/glue residues are
computed, assume either a proved positive open-channel spectral discontinuity
for the correlator or the applicable channel-unitarity clauses of
Definition~\ref{p12:def:claim-indexed-smat-package}.  On that declared domain
the scalar self-energy has the corresponding cut and subtracted dispersion
representation.  Chiral couplings such as \(c_d,c_m\) remain independent
source-response coefficients; no theorem forces \(c_d=c_m=F_\pi/2\).
\phantomsection\label{p11:eq:scalar-quarkonium-to-pp}
\phantomsection\label{p11:eq:direct-glue-to-pp}
\phantomsection\label{p11:eq:scalar-imaginary-schur}
\end{theorem}
\begin{proof}

The assumed spectral discontinuity, or channel unitarity when an asymptotic
claim is made, fixes the imaginary part to the appropriate phase-space
quadratic form in on-shell couplings.  A subtracted dispersion relation then
reconstructs the real part up to subtraction data.  Symmetry classifies the
allowed couplings but leaves their reduced magnitudes to correlators.
\end{proof}

\begin{theorem}[Subtracted Schur real part, sheets, and residues]
\label{p11:thm:schur-real-part-sheets-residues}
Let
\begin{equation}
\Gamma_{ab}^{(\alpha)}(s)=-\operatorname{Im}\Sigma_{ab}^{(\alpha)}(s+i0).
\end{equation}
The nonlocal open-channel self-energy may be written in a no-double-counted subtracted form
\begin{equation}
\label{p11:eq:subtracted-schur-dispersion}
\widehat\Sigma_{ab}^{(\alpha)}(s)=
-\frac{(s-s_0)^N}{\pi}
\int_{s_{{\rm th},\alpha}}^\infty
\frac{\Gamma_{ab}^{(\alpha)}(s')}{(s'-s_0)^N(s'-s-i0)}\,ds'.
\end{equation}
The integral is asserted only on a certified spectral-tail domain. A
sufficient condition for its absolute tail is
$\int^\infty\|\Gamma(s')\|s'^{-N-1}\,ds'<\infty$.
If $\|\Gamma(s')\|\le C(s')^p$ for $s'\ge S$ and $N>p$, then for
$S\ge2\max(|s|,|s_0|)$ the omitted tail has the explicit bound
\[
 \left\|\widehat\Sigma(s)-\widehat\Sigma_{[s_{\rm th},S]}(s)\right\|
 \le {2^{N+1}C|s-s_0|^N\over\pi(N-p)}S^{p-N}.
\]
A finite actual spectral support or sufficiently decaying native form factor
is another sufficient domain. For an unrestricted derivative-coupling tail
$\Gamma(s')=O((s')^2)$, $N\ge3$ is sufficient; $N=2$ alone is
logarithmically insufficient. No removal of the protected native UV ceiling
is required: its actual tail/support, rather than a guessed unbounded
continuation, is retained. The choice $s_0=0$ is allowed only when it is
outside the cut on the analytic subtraction domain; a massless cut beginning
at zero instead uses an analytic off-cut subtraction point. Local
subtraction polynomials remain in the compact block and matched chiral
constants. The positive cut density and spectral-tail estimates alone imply \emph{no} analytic continuation across the real cut. For the following sheet formulas assume, in addition, that each relevant cut density admits a channel-local holomorphic continuation \(\Gamma^{(\alpha)}(s)\) on a specified simply connected crossing chart away from its threshold/endpoints, and that the boundary-jump convention \(\Sigma_{\rm crossed}-\Sigma_I=2i\Gamma^{(\alpha)}\) holds on that chart. Continuation around each cut is part of the supplied analytic data, not a finite-kernel conclusion. Only \emph{under these hypotheses}, on a sheet vector \(\boldsymbol\epsilon\),
\begin{equation}
\label{p11:eq:sheet-continuation}
\Sigma_{\boldsymbol\epsilon}(s)=\Sigma_I(s)+2i\sum_\alpha \epsilon_\alpha\Gamma^{(\alpha)}(s),
\qquad
\epsilon_\alpha\in\{0,1\}.
\end{equation}
A continued correlator/resolvent resonance pole is a zero of
\begin{equation}
\label{p11:eq:physical-pole-determinant-completed}
\det D_{\boldsymbol\epsilon}(s)=0,
\qquad
D_{\boldsymbol\epsilon}(s)
:=s\mathbf 1-M_P^2-\Sigma_{\boldsymbol\epsilon}(s),
\end{equation}
and if the zero \(s_p\) is simple, choose right/left null vectors
\(D_{\boldsymbol\epsilon}(s_p)v_p=0\) and
\(w_p^\dagger D_{\boldsymbol\epsilon}(s_p)=0\).  Then
\begin{equation}
\label{p11:eq:pole-residue-normalization}
\operatorname*{Res}_{s=s_p}D_{\boldsymbol\epsilon}(s)^{-1}
=\frac{v_pw_p^\dagger}
{w_p^\dagger D_{\boldsymbol\epsilon}'(s_p)v_p}
=\frac{v_pw_p^\dagger}
{w_p^\dagger[\mathbf 1-\Sigma_{\boldsymbol\epsilon}'(s_p)]v_p}.
\end{equation}
For source vectors \(g_A,g_B\), the observable residue is
\(g_A^\dagger(\operatorname*{Res}D^{-1})g_B\).  At a complex pole, component
``compositeness fractions'' formed from \(|(v_p)_a|^2\) are generally
basis- and scheme-dependent and are not probabilities.  A glue-dominant
resonance is therefore characterized only relative to a declared source basis
and by its matrix residues, not by the nearest real-axis mass.
\end{theorem}
\begin{proof}
At finite level the Schur term is a rational sum over \(Q\)-space eigenstates.
Under the stated open-channel spectral-density limit, the real part is the
Hilbert transform of the discontinuity up to local subtractions; the convention
records that local contact terms already lie in \(M_P^2\) and the chiral
constants.  Crossing a \emph{holomorphically continuable} channel cut with the declared jump orientation gives
\eqref{p11:eq:sheet-continuation}; positivity and tail decay by themselves do not justify evaluating a real-cut density at a complex pole.  For a simple zero, finite-dimensional
analytic matrix inversion gives the rank-one Laurent coefficient
\eqref{p11:eq:pole-residue-normalization}; the nonzero denominator is
equivalent to simplicity.
\end{proof}

\begin{theorem}[Isolated continued-pole and matrix-residue limit]
\label{p11:thm:isolated-resonance-pole-residue-limit}
Let \(U\Subset\CC\) be a disc on a specified continuation sheet and let
\(D_{L}^{\rm II},D^{\rm II}:U\to M_r(\CC)\) be holomorphic inverse-propagator
matrices.  Assume
\begin{equation}
\label{p11:eq:continued-inverse-propagator-convergence}
D_L^{\rm II}\longrightarrow D^{\rm II}
\quad\text{locally uniformly on a neighbourhood of }\overline U,
\end{equation}
that \(\det D^{\rm II}\) has no zero on \(\partial U\), and that it has exactly
one simple zero \(s_p\) in \(U\).  Then, for all sufficiently large \(L\),
\(\det D_L^{\rm II}\) has exactly one zero \(s_{p,L}\) in \(U\), counted with
multiplicity, and
\begin{equation}
\label{p11:eq:continued-pole-convergence}
s_{p,L}\longrightarrow s_p.
\end{equation}
Choose unit-norm right/left null vectors \(v_{p,L},w_{p,L}\), with phases
chosen continuously relative to limiting unit vectors \(v_p,w_p\).  Then
\begin{equation}
\label{p11:eq:continued-matrix-residue-convergence}
\frac{v_{p,L}w_{p,L}^\dagger}
{w_{p,L}^\dagger(D_L^{\rm II})'(s_{p,L})v_{p,L}}
\longrightarrow
\frac{v_pw_p^\dagger}
{w_p^\dagger(D^{\rm II})'(s_p)v_p}.
\end{equation}
If the continuation variable is energy and the limiting pole lies in the lower
half-plane, \(\Gamma_p=-2\operatorname{Im}s_p\) is the resonance width in that
energy convention.  The theorem proves a continued correlator/resolvent pole.
Calling it an \(S\)-matrix pole additionally requires the channel package of
Definition~\ref{p12:def:claim-indexed-smat-package}; asymptotic completeness is
not required for the analytic statement itself.
\end{theorem}
\begin{proof}
Uniform matrix convergence implies uniform convergence of the determinants on
\(\partial U\).  Since \(|\det D^{\rm II}|\) has a positive minimum there,
Rouch\'e's theorem gives the same number of zeros in \(U\), counted with
multiplicity.  Compactness then forces the unique zeros \(s_{p,L}\) to converge
to \(s_p\); otherwise a subsequence would have a different limiting zero of
\(\det D^{\rm II}\).  Simplicity makes the left and right kernels
one-dimensional and
\(w_p^\dagger(D^{\rm II})'(s_p)v_p\ne0\).  Finite-dimensional spectral
perturbation therefore permits continuous phases and normalizations with
\(v_{p,L}\to v_p\) and \(w_{p,L}\to w_p\).  Holomorphic local-uniform
convergence also gives derivative convergence by Cauchy's formula.  Substitution
in the simple-zero Laurent formula proves
\eqref{p11:eq:continued-matrix-residue-convergence}.
\end{proof}

\paragraph{Independent contour proof of the same conditional pole limit.}
Set $F_L(z)=\det D_L^{\rm II}(z)$ and
$F(z)=\det D^{\rm II}(z)$.  Local-uniform holomorphic convergence
also gives convergence of the derivatives on the closed contour
$\partial U$.  Because $F$ has no boundary zeros, the integers
\[
 N_L=\frac{1}{2\pi i}\oint_{\partial U}\frac{F_L'(z)}{F_L(z)}\,dz
 \quad\longrightarrow\quad
 \frac{1}{2\pi i}\oint_{\partial U}\frac{F'(z)}{F(z)}\,dz=1
\]
are eventually one.  Since the counted zero is unique and simple,
its position is the first logarithmic moment
\[
 s_{p,L}=\frac{1}{2\pi i}\oint_{\partial U}
 z\frac{F_L'(z)}{F_L(z)}\,dz\longrightarrow s_p.
\]
All $D_L^{\rm II}$ are invertible on $\partial U$ for large $L$,
so their inverses converge uniformly there.  Their matrix residues
therefore converge without selecting the null-vector phases:
\[
 R_L=\frac{1}{2\pi i}\oint_{\partial U}(D_L^{\rm II}(z))^{-1}\,dz
 \longrightarrow
 \frac{1}{2\pi i}\oint_{\partial U}(D^{\rm II}(z))^{-1}\,dz=R.
\]
At a simple determinant zero, the adjugate divided by the determinant
derivative identifies $R$ with the rank-one null-vector expression in
the theorem, and likewise for $R_L$.  This is an alternate analytic
proof of the matrix-residue limit, not a construction of the continued
sheet from the real cut, a physical channel, or an $S$-matrix pole.

\begin{theorem}[Tensor and pseudoscalar glue closures as computed correlator blocks]
\label{p11:thm:tensor-pseudoscalar-glue-closures}

Tensor polarization/projector algebra and phase-space factors are exact once the corresponding source is specified.  The pseudoscalar glue/anomaly block is determined by the computed topological susceptibility/cumulants and chiral singlet normalization.  Physical pole positions are eigenvalues/zeros of the coupled inverse-propagator blocks; no topology-count mass is inserted.
\phantomsection\label{p11:eq:tensor-pp-schur-cut}
\phantomsection\label{p11:eq:pseudoscalar-glue-rank-one-block}
\end{theorem}
\begin{proof}

The tensor statement is finite representation algebra.  The anomalous Ward identity relates pseudoscalar singlet sources to topological charge, but its magnitude is a correlator/cumulant output.  Standard Schur reduction then determines the coupled poles from those computed entries.
\end{proof}

\begin{corollary}[Transport endpoints without equal-share valence ansatz]
\label{p11:cor:transport-endpoint-closures}

The exact transport closures are the conserved flavor/charge/momentum sum rules, perturbative evolution equations, asymptotic singlet fixed points, and connected-cumulant identities under their stated domains.  Primitive color determinant support does not fix a boundary PDF to equal momentum fractions \(x_i=1/3\).

\end{corollary}
\begin{proof}

For any \(0<a<1/2\), the proton-like measures \(u_v=2\delta(x-a)\) and \(d_v=\delta(x-(1-2a))\) have the same valence counts and total momentum.  Therefore color support cannot determine the momentum split.  The remaining listed identities follow from conserved operators/evolution/cumulant algebra.
\end{proof}




\section{Execution, cross-volume handoff, and final readout scope}
\label{p11:sec:v237-b8}
This final ledger distinguishes completed finite algebraic readouts, state-preparation requirements, physical pole/measurement maps and the claim-indexed continuum-QFT conditions of Paper~VII-C.

% Source Paper VII-B, original line 3778; inherited proof bodies unchanged.
% Former logical part: Computation, Validation, and Dependency Control; integrated into the new thematic argument.
\label{p11:part:computation-validation}


% Source Paper VII-B, original line 3781; inherited proof bodies unchanged.
\subsection{Computation and validation path}
\label{p11:sec:minimal-upgrade-path}


% Source Paper VII-B, original line 3784; inherited proof bodies unchanged.
\subsubsection{Compute-first validation path}
\label{p11:subsec:compute-first-validation-path}


\begin{table}[htbp]
\centering
\footnotesize
\begin{tabular}{>{\raggedright\arraybackslash}p{0.34\textwidth} >{\raggedright\arraybackslash}p{0.56\textwidth}}
\toprule
Stage & Status \\
\midrule
Construct finite Hamiltonian/compiler & native finite-QCD construction \\
Diagonalize finite singlet blocks & native computation \\
Compute transfer kernels & native computation \\
Extract spectral, current, distribution, and topology readouts & native computation \\
Compare to lattice, experiment, PDF fits, and global data & validation only \\
Iterate numerical precision & computation refinement, not parameter fitting \\
\bottomrule
\end{tabular}
\caption[Computation and validation path]{Computation and validation path.}
\label{p11:tab:computation-validation-path}
\TableNotes{The validation path is ordered compute-first, compare-second.  The native finite Hamiltonian, transfer kernels, and readouts are constructed before any lattice, PDG, PDF, heavy-ion, or nuclear comparison can be used.}
\end{table}

After the finite matrix compiler is fixed, the remaining native task is the evaluation of a closed set of eigenvalues, residues, matrix elements, and thermal susceptibilities.  The computational chain is
\begin{equation}
\label{p11:eq:main-computational-chain}
\boxed{
\mathfrak C_L^{\rm QCD}
\xrightarrow{\operatorname{Compile}_{\rm QCD}}
\mathfrak M_L^{\rm QCD}
\xrightarrow{\operatorname{Spec/MatElem}}
\mathfrak r_L^{\rm QCD}
\xrightarrow{L\to\infty\;{\rm comparison}}
\mathfrak r_{\rm QCD}^{\rm BH}.
}
\end{equation}
The finite basis and admissible gauge domain are enumerated first; the singlet Hamiltonian blocks are then assembled, diagonalized, and evaluated before comparison.

\begin{proposition}[Finite-QCD benchmark package with proof-status typing]
\label{p11:prop:finite-qcd-benchmark-package}

The benchmark package consists of: exact finite geometry/algebra numbers; conditional source-member perturbative comparisons; computed finite spectral/correlator/thermal/topological outputs; and separately flagged continuum obligations.  Gap, string tension, \(F_\pi,B_0,\Sigma,T_c,\chi_t\), resonance poles and baryon observables are not placed in the unconditional benchmark column until their defining computations are executed.
\phantomsection\label{p11:eq:finite-benchmark-set}
\phantomsection\label{p11:tab:finite-qcd-benchmark-package}
\end{proposition}
\begin{proof}

This is a direct partition by logical type.  Each unconditional entry is a finite algebraic invariant; each physical observable listed in the second group depends on an explicit Hamiltonian/source/limit.  The acyclic handoff forbids comparison values from being substituted for missing computations.
\end{proof}



% Source Paper VII-B, original line 3836; inherited proof bodies unchanged.
\subsection{Readout compression and dependency graph}
\label{p11:sec:readout-compression-dependency-graph}


% Source Paper VII-B, original line 3839; inherited proof bodies unchanged.
\subsubsection{Dependency graph}
\label{p11:subsec:dependency-graph-forbidden-comparison-arrows}


The finite execution order is
\begin{equation}
\label{p11:eq:qcd-dependency-graph}
\boxed{\begin{gathered}
\Iqcd\longrightarrow\mathfrak C_L^{\rm QCD}
\xrightarrow{\operatorname{Compile}_{\rm QCD}}
\mathfrak M_L^{\rm QCD},\\
\mathfrak M_L^{\rm QCD}
\longrightarrow
\left(\mathcal R_{{\rm np},L},\Delta_{\YM,L},\sigma_{\YM,L},X_{t,L},
\Sigma_{q,L},B_{0,L},M_{h,L},\mathcal F_{h,L},\mathcal E_{{\rm QCD},L}\right)
\longrightarrow\mathfrak r_{\rm QCD}^{\rm BH}.
\end{gathered}}
\end{equation}


% Source Paper VII-B, original line 3858; inherited proof bodies unchanged.
\subsubsection{Comparison control and refinement stability}
\label{p11:subsec:comparison-control-refinement-stability}

\begin{theorem}[Uniform-gap finite-carrier comparison control]
\label{p11:thm:post-gap-finite-carrier-control}

Assume an admissible local/reflection-positive sequence of finite Hamiltonians with a level-independent gap \(g_*:=\inf_L(E_{1,L}-E_{0,L})>0\) and the convergence hypotheses of the chosen continuum construction.  Then connected correlators have uniform massive exponential control away from criticality, and isolated low-lying spectral data are stable under sufficiently small norm-resolvent perturbations.
\phantomsection\label{p11:eq:uniform-comparison-gap-condition}
\phantomsection\label{p11:eq:uniform-gap-comparison-corrections}
\end{theorem}
\begin{proof}

The spectral gap bounds the imaginary-time semigroup on the vacuum-orthogonal subspace by \(e^{-g_*t}\).  Locality/reflection positivity convert this spectral bound into the standard massive correlation estimates in the stated framework.  Perturbation stability of isolated spectral projections follows from the Riesz projection/norm-resolvent theorem.  The hypothesis \(g_*>0\) is not derived here.
\end{proof}

\begin{corollary}[Uniform pure-glue control remains a criterion, not a derived floor]
\label{p11:cor:native-pure-glue-comparison-control-derived}

Pure-glue comparison control follows if the admissible sequence satisfies a uniform positive gap/coercivity condition together with the locality/reflection-positivity hypotheses required by the chosen continuum theorem.  Current Paper~VII-A provides the exact criterion
\[
\inf_L\kappa_{{\rm sh},L}>0,
\]
but does not prove this inequality from finite counts.

\end{corollary}
\begin{proof}

The implication from a uniform gap to exponential clustering/isolated-spectrum control is standard under the stated hypotheses.  The sequence \(\kappa_{{\rm sh},L}=1/L\) shows that positive finite-level gaps do not imply a positive infimum, so the criterion cannot be silently discharged.
\end{proof}


\begin{theorem}[Isolated stable-pole and source-residue limit]
\label{p11:thm:admissible-refinement-stability-finite-readouts}
After isometric embedding into one Hilbert space, let \(H_L\) and \(H\) be
self-adjoint and suppose \(H_L\to H\) in norm resolvent sense.  Let
\(\Gamma\subset\rho(H)\) be a positively oriented rectifiable contour enclosing
one isolated finite-rank spectral island and no other spectrum.  Assume
\(\Gamma\subset\rho(H_L)\) for all sufficiently large \(L\), with no spectral
crossing through \(\Gamma\).  Define
\begin{equation}
\label{p11:eq:refinement-riesz-projectors}
P_L=\frac{1}{2\pi i}\oint_\Gamma(z-H_L)^{-1}\,dz,
\qquad
P=\frac{1}{2\pi i}\oint_\Gamma(z-H)^{-1}\,dz .
\end{equation}
Then
\begin{equation}
\label{p11:eq:refinement-projector-convergence}
\|P_L-P\|\longrightarrow0,
\qquad
\operatorname{rank}P_L=\operatorname{rank}P
\quad\text{eventually},
\end{equation}
and the enclosed eigenvalues converge as a multiset, with algebraic
multiplicity.

Let normalized vacua \(\Omega_L\to\Omega\) in norm and bounded sources
\(A_L\to A\), \(B_L\to B\) in operator norm.  Their island residues
\begin{equation}
\label{p11:eq:refinement-source-residue-limit}
\mathcal R_{AB,L}:=
\langle\Omega_L,A_LP_LB_L\Omega_L\rangle
\longrightarrow
\mathcal R_{AB}:=
\langle\Omega,APB\Omega\rangle .
\end{equation}
If the island is a simple eigenvalue \(E_L\to E\), then near the pole
\begin{equation}
\label{p11:eq:simple-resolvent-pole-decomposition}
G_{AB,L}(z):=
\langle\Omega_L,A_L(z-H_L)^{-1}B_L\Omega_L\rangle
=\frac{\mathcal R_{AB,L}}{z-E_L}+G^{\rm reg}_{AB,L}(z),
\end{equation}
so the operator-resolvent pole and the source residue both converge.  The
scalar function \(G_{AB,L}\) has a nonremovable limiting pole precisely when
\(\mathcal R_{AB}\ne0\); a vanishing source overlap makes that scalar pole
removable without removing the spectral eigenvalue.  For a degenerate island
the projector and matrix residue are canonical; individual vector residues
require an additional basis or splitting convention.  No locality,
scattering, LSZ, asymptotic completeness, or infrared hypothesis is used in
this theorem.
\phantomsection\label{p11:eq:refinement-transfer-convergence}
\phantomsection\label{p11:eq:refinement-operator-convergence}
\phantomsection\label{p11:eq:refinement-uniform-gap-floor}
\phantomsection\label{p11:eq:refinement-gap-liminf}
\end{theorem}
\begin{proof}
Norm-resolvent convergence is locally uniform on the compact contour
\(\Gamma\).  Integrating the norm difference of the resolvents gives
\begin{equation}
\|P_L-P\|
\le\frac{\operatorname{length}(\Gamma)}{2\pi}
\sup_{z\in\Gamma}
\|(z-H_L)^{-1}-(z-H)^{-1}\|\longrightarrow0.
\end{equation}
Two projections at norm distance smaller than one have isomorphic ranges, so
their finite ranks agree.  The usual Riesz reduction to these finite-dimensional
ranges then gives convergence of the enclosed eigenvalue multiset.  Expanding
\(\mathcal R_{AB,L}-\mathcal R_{AB}\) and inserting and subtracting
\(\Omega,A,B,P\) bounds it by a finite sum of the assumed vector/operator norm
differences and \(\|P_L-P\|\), proving
\eqref{p11:eq:refinement-source-residue-limit}.  For rank one, the spectral
theorem separates \((z-E_L)^{-1}P_L\) from the resolvent on the complementary
spectral subspace, which is holomorphic near \(E_L\), yielding
\eqref{p11:eq:simple-resolvent-pole-decomposition}.  This is a Riesz/spectral
argument only~\cite{Kato1950,ReedSimonII1975}.
\end{proof}

\begin{corollary}[Source-selected fixed-momentum hadron state]
\label{p11:cor:source-selected-fixed-momentum-state}
Assume \([H_L,P_{\rm phys,L}^+]=[H_L,G_L]=0\), and let \(P_{h,L}\) be the Riesz projector of an
isolated stable hadron island.  For a declared interpolating source \(B_L\),
put
\begin{equation}
\label{p11:eq:source-selected-fixed-momentum-vector}
 \Xi_{h,L}:=P_{h,L}E_{p_h,L}B_L\Omega_L .
\end{equation}
If \(\Xi_{h,L}\ne0\), then
\begin{equation}
\label{p11:eq:source-selected-fixed-momentum-state}
 \boxed{
 |h,L;p_h^+\rangle
 :=\frac{\Xi_{h,L}}{\|\Xi_{h,L}\|}}
\end{equation}
is a constructed normalized physical state in the required joint sector and
discharges the vector part of \(\LFState_{h,L}\).  The state/beam alignment
clause remains unless it follows from the same source and covariance data.

Under the norm-resolvent, source, vacuum, noncrossing, and nonzero-overlap
hypotheses of
Theorem~\ref{p11:thm:admissible-refinement-stability-finite-readouts}, together
with norm convergence of the compared fixed-momentum projectors, the
corresponding rank-one state projectors converge.  On a simple joint island
the ray is canonical relative to the declared source.  In a degenerate island
the displayed vector is source dependent; without a source or splitting rule,
only the projected density/matrix residue is canonical.  A continued resonance
residue does not define such a probability state.
\end{corollary}

\begin{proof}
Finite-dimensional strong commutation makes \(P_{h,L}\) commute with
\(E_{p_h,L}\).  Their product maps the source vector into the physical joint
island, and normalization proves the first claim.  The Riesz theorem gives
\(P_{h,L}\to P_h\) in norm; together with the assumed convergence of
\(E_{p_h,L}\), \(B_L\), and \(\Omega_L\), this gives convergence of
\(\Xi_{h,L}\).  A nonzero limiting norm makes normalization continuous, hence
the rank-one projectors converge.  The final qualifications distinguish
source selection from basis-independent spectral data and stable Hilbert
vectors from continued complex residues.
\end{proof}


\begin{table}[htbp]
\centering
\footnotesize
\begin{tabular}{>{\raggedright\arraybackslash}p{0.28\textwidth} >{\raggedright\arraybackslash}p{0.34\textwidth} >{\raggedright\arraybackslash}p{0.28\textwidth}}
\toprule
Open theorem/execution target & Native object to compute & Forbidden shortcut \\
\midrule
Vector octet harmonic trace & flag-averaged Schur trace of the Hurwitz octet block & fitting \(M_\rho\) or \(L_9\) \\
Scalar poles and \(L_5,L_8\) & \(PP\oplus q\bar q\oplus gg\) Schur poles and scalar-source residues & choosing scalar masses or chiral LECs from data \\
Heavy-source matching & \(m_Q^{\rm stat},m_Q^{\rm kin}\), contact density, heavy-light cloud energies & Cornell constants or quarkonium offsets \\
Full hadron spectroscopy & diagonalization and pole extraction of \(\mathfrak M_L^{\rm QCD}\) & PDG/lattice mass insertion \\
PDFs, FFs, and decay amplitudes & constructed light-front operator measures, Wilson-line limits, and claim-scoped out-state transition sums & global-fit boundary functions \\
\bottomrule
\end{tabular}
\caption[Open native QCD execution ledger]{Open native QCD execution ledger.}
\label{p11:tab:open-native-qcd-execution-ledger}
\TableNotes{An item in this table is not an observational input.  It is a sharply defined finite computation or theorem still to be executed inside the native compiler.}
\end{table}


% Source Paper VII-B, original line 4475; inherited proof bodies unchanged.
\subsection{Paper-VII-B completion boundary}
\label{p11:sec:completion-certificate-derivation-protocol}

\begin{theorem}[Paper-VII-B repaired finite readout completion certificate]
\label{p11:thm:paper-xi-finite-qcd-completion-certificate}

Paper~VII-B is complete as a finite readout-definition layer: given the repaired
Paper-VII-A finite member, the derived \(\LFRayGeom\), and a supplied physical
state certificate \(\LFState_{h,L}\), it constructs spectral masses, source
residues, current responses, light-ray/twist-two readouts, finite thermal traces,
topological cumulants, same-channel/Schur blocks, Hodge/Krylov diagnostics, and
multi-baryon cumulants without downstream fitting.  If the additional exact
\(\LFRep_L\) certificate is supplied, Appendix~\ref{p11:app:v264-exact-lfrep-positive-measures}
constructs the stronger finite fraction operators and positive number/momentum
operator measures.  It proves the isolated stable-pole/residue limit, the physical current/stress
PDF sum-rule bridge, and a separate sufficient continued-resonance limit.  The
optional appendix proves Hausdorff and bounded-Lipschitz positive-measure limits
conditional on exact \(\LFRep_L\).  It does not claim that every hadron sector
supplies \(\LFState_{h,L}\), that the consumed \(\LFLim\) clauses or out-state map
exist, or that the physical quark-mass operator, chiral phase
occurrence, uniform gap, string tension, or topology measure has been solved.

\end{theorem}
\begin{proof}

Every retained native object is either a finite matrix operation or a precisely typed conditional/limit interface.  The predecessor count-derived numerical chains have been removed, so no unresolved constructor is disguised as a completed datum.
\end{proof}

\begin{longtable}{>{\raggedright\arraybackslash}p{0.25\textwidth} >{\raggedright\arraybackslash}p{0.36\textwidth} >{\raggedright\arraybackslash}p{0.27\textwidth}}
\caption[Final Paper-VII-B closure boundary]{Final Paper~VII-B closure boundary.}\label{p11:tab:final-closure-boundary}\\
\toprule
Layer & Status at the Paper-VII-B boundary & Successor obligation \\
\midrule
\endfirsthead
\toprule
Layer & Status at the Paper-VII-B boundary & Successor obligation \\
\midrule
\endhead
Finite QCD core & repaired Paper-VII-A sequences \(H_L,\widehat K_L(t),\mathfrak A_L,\mathfrak A'_L,\mathcal O_L,W_L[J]\) plus exact Hodge geometry & execute the inherited source/action member and only the consumed state/limit certificates \\
Strong inputs & \(SU(3)_c\) representation and common-heat gauge matching exact; low-scale coupling and \(M_q^{\rm phys}(\mu)\) consume typed scale/RG transport of inherited owners; topological \(\theta\) is auxiliary, the native finite CP character and rooted determinant orientation are fixed by Paper~V-C, and finite action is regulator-representative & upstream owner/scale matching; claim-scoped action comparison only if representative-independence/canonicity is claimed; no downstream backflow \\
Confinement/gap & native \(E_{48}\) carrier gap derived; finite QCD gap/coercivity functionals, generated position area, and Wilson operators defined & use a certified phase cluster for the base export; prove QFT4$_C$ only for canonical same-phase refinement independence; on a specified pure-glue/unscreened Wilson sector the strict area exponent is conditional on actual same-owner APBI (APBC square specialization) or a direct multiface/flux estimate; it is not a universal full-QCD confinement gate; a continuous gap is optional and claim-specific \\
Chiral/topology & covered finite massive source responses and auxiliary \(Z_L(\theta)\)/cumulants; exact ordered condensate/Goldstone targets separately defined & execute finite current/winding responses; use the complete blocked/Duhamel source-jet tests of Theorems~\ref{p10:thm:v2210-lit-blocked-source-jets} and~\ref{p11:thm:v2210-lit-complete-duhamel-contacts}; prove ordered SSB and chiral-residue limits only for claims that consume \(\Sigma,F_\pi,B_0\) \\
Hadron blocks & finite real Schur poles plus stable and continued pole/residue limit theorems & reuse inherited quark/source packet; execute finite spectra with the source-frame/GEVP and return tests of Theorem~\ref{p11:thm:v2210-lit-gevp-frame-gap} and Proposition~\ref{p11:prop:v2210-lit-semigroup-leakage}; prove continuation hypotheses only for continued-pole claims \\
Heavy/source kernels & exact finite Hurwitz source form factors & reuse generated area/Wilson owner and certify the consumed physical heavy-propagator/potential and mass map \\
Distributions/out states & \(\LFRayGeom\), physical \(P_{\rm phys}^+\), and source-owned light-ray/twist-two operators give the physical readout chain; exact \(\LFRep_L\) gives an optional finite positive-measure realization & construct the physical state and claimed \(\LFLim\), hadronization, or out-state clauses; spacelike source jets may instead use Theorem~\ref{p11:thm:v2210-lit-spacelike-source-jets} with their actual matching; construct \(\LFRep_L\) only for the stronger finite POVM claim \\
Thermal/nuclear & finite analytic traces; OPE form; general contact/cumulant grammar & complete nonlinear source/metric contacts as in Theorem~\ref{p11:thm:v2210-lit-complete-duhamel-contacts}; thermodynamic limit and QCD matching of contact/multi-baryon operators \\
\bottomrule
\end{longtable}
\TableNotes{The successor column lists genuine constructors/limits, not observational parameters to be fitted into the finite kernel.}



% Source Paper VII-B, original line 4526; inherited proof bodies unchanged.
% Former logical part: Continuum-QFT Handoff and Final Export Ledger; integrated into the new thematic argument.
\label{p11:part:paper-xii-handoff}


% Source Paper VII-B, original line 4529; inherited proof bodies unchanged.
\subsection{Continuum-QFT handoff to Paper VII-C}
\label{p11:sec:clay-continuum-boundary}
\label{p11:sec:paper-xii-continuum-handoff}

\begin{theorem}[Two-layer Paper-VII-B export with explicit obligations]
\label{p11:thm:two-layer-paper-xi-export-paper-xii}
\label{p11:thm:paper-xi-xii-handoff}

The Paper-VII-B to Paper-VII-C handoff has two layers: (A) unconditional/conditionally supplied finite operator sequences \(\mathfrak D_{7B,\YM}^{SU(3)}\) and \(\mathfrak D_{7B}^{\rm QCD}\); (B) explicit obligations for any claimed continuum reconstruction, including convergence topology, reflection positivity where used, a uniform physical gap if required, APBC for promotion of the derived bounded-incidence heat area exponent to physical confinement, and the relevant source/operator limits.  Layer B is not part of the proved finite datum.

\end{theorem}
\begin{proof}

The finite tuple is defined by projection and matrix functional calculus.  Continuum statements require additional quantified limits and inequalities.  Separating the two prevents a conclusion of Paper~VII-C from being used retroactively as a premise of Paper~VII-B.
\end{proof}


% Source Paper VII-B, original line 4545; inherited proof bodies unchanged.
\subsubsection{Finite export datum}
\label{p11:subsec:finite-export-datum-paper-xii}

\begin{definition}[Paper-VII-B finite strong readout datum for Paper VII-C]
\label{p11:def:pure-g-wilson-continuum-datum}
The rigorous finite export is
\begin{equation}
\label{p11:eq:paper-xii-export-datum}
\mathfrak D_{7B}^{\rm QCD}
:=\left\{H_L,\widehat K_L(t),\mathfrak A_{L,{\rm str}},\mathfrak A'_{L,{\rm str}},\mathcal O_L,W_L[J],\omega_L\right\}_{L\in\mathcal R_{\rm adm}}
\cup\{\mathfrak T_{\rm chiral},\mathfrak T_{\rm top}\},
\end{equation}
with \(\omega_L(A)=\langle\Omega_L,A\Omega_L\rangle\).  Any light-front map, OS reconstruction, physical quark-mass member, uniform gap, or positive area-law estimate is exported only when separately constructed/proved; none is inserted into this datum by notation.  The reciprocal-position chart, metric area functional, and Wilson observable are already generated presentations rather than an extra pairing coefficient.
\end{definition}

\begin{definition}[Pure-glue finite subdatum for Paper VII-C]
\label{p11:def:pure-glue-wilson-curvature-subdatum-paper-xii}
The pure-glue projection is
\begin{equation}
\label{p11:eq:paper-xii-pure-glue-subdatum}
\mathfrak D_{7B,\YM}^{SU(3)}
:=\left\{\mathcal H_{L,\YM}^{\rm sing},\Omega_{L,\YM},H_{L,\YM},e^{-t(H_{L,\YM}-E_{0,L,\YM})},\mathfrak A_{W,L}^{SU(3)},\Delta_{2,H}\right\}_{L\in\mathcal R_{\rm adm}}.
\end{equation}
It is a \emph{typed submodel/restriction}, not a literal subset of the displayed
operator/state tuples: the Hilbert carriers, Hamiltonians, and vacua need not be identical. On an actual gauge-invariant pure-glue deletion/restriction map \(\mathcal R_{\YM}\) supplied by the finite compiler, with the pure-glue Hamiltonian and its own vacuum energy \(E_{0,L,\YM}=\min\operatorname{spec}H_{L,\YM}\), the handoff is
\begin{equation}
\label{p11:eq:pure-glue-subdatum-contained-in-full-qcd-export}
\mathfrak D_{7B,\YM}^{SU(3)}=\mathcal R_{\YM}(\mathfrak D_{7B}^{\rm QCD}),
\qquad \text{when }\mathcal R_{\YM}\text{ is defined}.
\end{equation}
No existence of a gauge-compatible reducing restriction or equality of full-QCD and pure-glue vacua is inferred by this notation.
Paper~VII-C must still prove any continuum reconstruction, OS positivity, uniform mass gap, position-space area law, or local-field limit that it uses.
\end{definition}


% Source Paper VII-B, original line 4576; inherited proof bodies unchanged.
\subsubsection{Handoff ledger}
\label{p11:subsec:paper-xi-xii-handoff-ledger}

\begin{longtable}{>{\raggedright\arraybackslash}p{0.28\textwidth} >{\raggedright\arraybackslash}p{0.38\textwidth} >{\raggedright\arraybackslash}p{0.22\textwidth}}
\caption[Paper-VII-B to Paper-VII-C handoff ledger]{Paper~VII-B to Paper~VII-C handoff ledger.}\label{p11:tab:paper12-handoff-ledger}\\
\toprule
Paper~VII-B finite export & Paper~VII-C target, not certified here & Handoff status \\
\midrule
\endfirsthead
\toprule
Paper~VII-B finite export & Paper~VII-C target, not certified here & Handoff status \\
\midrule
\endhead
Imported finite \(SU(3)_c\) source-transfer kernel \(\mathfrak C_L^{\rm QCD}\mapsto\mathfrak M_L^{\rm QCD}\) & admissible Born--Hopf spectral/momentum QFT datum & finite export here; admissibility of the continuum target remains a Paper-VII-C obligation \\
Finite \(m_x\)-side comparison charts \(\chi_L:\Gamma_L\to\mathbb R^4\) at fixed \(m_p=1\) & large-\(m_x\) quadrature and local flat chart recovery & closed in Paper~VII-C as presentation, not primitive lattice \\
Positive finite transfer family & candidate OS-positive continuum state & reflection positivity and reconstruction are Paper-VII-C obligations \\
Finite vacuum states \(\omega_L\) & candidate local continuum state & existence/uniqueness of the limit is a Paper-VII-C obligation \\
Finite plaquette/Hodge curvature blocks & contact-normalized local action and native spectral UCS & kernel block imported from Paper~VII-A; continuum local-field reconstruction is a Paper-VII-C target \\
Finite Haar/color identities and any separately certified finite quantum Ward identity & continuum Schwinger--Dyson and Ward identities & classical symmetry alone is insufficient; finite state/measure/regulator invariance and QFT4/QFT10 identity transport are Paper-VII-C obligations \\
Finite Euclidean transfer and state & OS/Wightman reconstruction with imported Lorentzian sign & Paper~VII-C target, not certified here \\
Finite gap sequence \(\gamma_L=E_{1,L}-E_{0,L}\) and uniform-gap criterion & optional continuum pure-glue mass-gap branch, not the native gap & if claimed, requires fixed scaling and an independent proof that \(\inf_L\gamma_L>0\) \\
Finite Wilson operators & continuum area law/string tension & fixed-loop convergence is QFT4; the heat-branch triality exponent is derived, while physical promotion requires APBC on the same action \\
Finite Wilson loops and local operator readouts & local \(F^2,\ F\widetilde F,\ T_{\mu\nu}\) fields from native residues & Paper~VII-C target, not certified here \\
\bottomrule
\end{longtable}
\TableNotes{Paper~VII-B supplies finite exports; Paper~VII-C is tasked with proving any continuum reconstruction it consumes.  The table records the acyclic handoff boundary and must not be read as advance certification of a Paper-VII-C continuum theorem.}

\begin{corollary}[Acyclic Papers VII-A--VII-C handoff]
\label{p11:cor:paper-x-xii-circularity-audit}

The dependency direction is \(\mathfrak F_{\le6B}\to\mathfrak D_{7A}^{\rm QCD,core}\to\mathfrak D_{7B}^{\rm QCD}\to\) Paper~VII-C continuum claims.  Paper~VII-C may consume the finite sequence and prove additional limits, but no Paper-VII-C continuum conclusion may be fed backward to select a Paper-VII-A or Paper-VII-B source coefficient.  Substantive continuum claims remain conditional until Paper~VII-C proves the obligations exported in Theorem~\ref{p11:thm:two-layer-paper-xi-export-paper-xii}.

\end{corollary}
\begin{proof}

Each arrow is functional dependence on previously defined data.  Backward use of a continuum result to choose an earlier coefficient would make the earlier object depend on its own downstream consequence, creating a logical cycle.  Keeping the continuum obligations as hypotheses preserves acyclicity.
\end{proof}

\begin{remark}[Status of numerical QCD readouts]\label{p11:rem:printed-qcd-numerics-source-branch}
Version~1.7 contains no unconditional hadronic/chiral/topological decimal derived from a finite count.  Perturbative comparison numbers are permitted only on their explicitly typed source/mass-threshold branch.  Nonperturbative masses, residues, \(F_\pi,B_0,\Sigma\), \(T_c\), \(\chi_t\), string tension, nuclear contacts and distribution moments are computation/limit outputs until their defining matrices are executed.
\end{remark}


% Source Paper VII-B, original line 4618; inherited proof bodies unchanged.
\subsection{flavor/readout handoff}
\label{pxi:sec:post-v125-flavor-handoff}

All quark-mass-dependent QCD readouts in this paper consume physical quark mass
eigenvalues only after the upstream charged preterminal source theorem plus G16 RG handoff has supplied the running quark-mass spurion.
The relevant finite flavor object is the simultaneous-conjugacy class of
$H_u=Y_uY_u^*$ and $H_d=Y_dY_d^*$, not the historical Paper-V-B terminal row or
comparison-exponent ledger.  Right-handed weak-basis frames and optional
terminal comparison atlases do not enter physical hadronic observables.



% Source Paper VII-B, original line 4629; inherited proof bodies unchanged.
\subsubsection{Finite compression and symmetry compatibility}
\label{p11:sec:v199-compression-handoff}
The exact finite LFRep target remains a stronger optional representation problem:
it requires the action-owned equal-time good-field seed, the actual-carrier
positive-polar certificate, exact R3 $P_{F,L}^+=P_{\rm phys,L}^+$, and the
Wilson-path/Gauss/species intertwiners.  The main physical PDF/GPD/DA chain no
longer waits on that certificate: it uses the source-owned light-ray/twist-two
operators and the physical total null generator.
The elapsed null generator itself is no longer a separate input: Paper~VII-A
constructs it from the same $H_L,P_{H,L}$ pair once the equal-time seed is
present.
Appendix~\ref{p11:app:v199-compression-defect} supplies the exact leakage
identity for a proposed finite compression.  It shows that a parent symmetry
is inherited exactly when the displayed leakage terms cancel; commutation of
the cutoff with the generator is a sufficient condition.  A small defect is
an approximation certificate, not exact finite LFRep or a new physical
generator construction.


% Source Paper VII-B, original line 4647; inherited proof bodies unchanged.
\subsection{Earlier theorem synthesis and scope}
Paper~VII-B is closed as a rigorous finite readout grammar rather than as a
catalogue of count-derived QCD numbers.  Exact color representation theory,
finite singlet/Hodge geometry, spectral/Schur algebra, thermal traces,
perturbative evolution endpoints, and boundary-rank mathematics are retained.
The physical distribution chain now runs through the generated light-ray and
twist-two operators.  Current charge fixes the first moments and the same-owner
stress tensor fixes the total second moment; exact finite species-fraction
operators and positive number/momentum measures are retained only as the
stronger Appendix~\ref{p11:app:v264-exact-lfrep-positive-measures} realization.
Paper~VII-A already closes conservation of the physical total null generator,
so physical fixed-total-plus sectors are invariant without R3.  The remaining
distributional obligations are the physical state occurrence and the declared
Wilson/light-ray renormalization and continuum intertwiners.  Finite
real poles with compressed spectral-projector residues, stable isolated Riesz
limits, and continued resonance
matrix residues are likewise distinct theorem layers.  Physical spectral
scales are computed from a supplied Hamiltonian/source family; current residues
are matrix elements; topology is a charge-covariance problem; chiral constants
are double-limit/correlator outputs; and nuclear forces are multi-baryon
operator/cumulant readouts.  Local and amputated kernels are computed from the
supplied Hamiltonian/source family; their interpretation as asymptotic
amplitudes or rates additionally uses the claim-indexed Paper-VII-C package.
Every quark-mass-dependent execution consumes the inherited
$M_q^{\rm phys}(\mu)$ supplied by the upstream charged preterminal source theorem plus G16 RG
handoff.  No new QCD mass constructor is required.  The remaining frontier is
explicit: transport the derived common-heat gauge boundary to any named
comparison scale, prove finite-action universality only for claimed representative-independence and QFT4$_C$ only for canonical same-phase export, execute covered finite source/spectral/topological readouts and preserve the
exact ordered chiral occurrence only for the condensate/Goldstone/GMOR claim;
prove the consumed physical light-ray/process and hadronization transport,
retain species/Fock correspondence only for the optional exact fraction
realization, and prove a positive area bound only for a selected pure-glue or
certified unscreened tension target.  A continuous
Yang--Mills gap is not on the active frontier.  No comparison datum is permitted
to close these interfaces by backflow.


% Source Paper VII-B, original line 4683; inherited proof bodies unchanged.
\subsection{revised QCD readout boundary}
\label{p11:sec:v164-revised-qcd-readout-boundary}

\begin{theorem}[Vacuum and gap readouts]
\label{p11:thm:v164-vacuum-gap-readouts}
The symmetric finite vacuum supplied to every readout is
\[
 \rho_{0,L}=P_{0,L}/\operatorname{rank}P_{0,L},
\]
the zero-temperature Gibbs limit of the compiled Hamiltonian.  The physical
finite gap is
\[
 \Delta_L^{\rm phys}=-a_L^{-1}\log q_L,
\]
not the raw one-step separation.  Hadron and glueball masses remain poles or
spectral energies of their specified same-action blocks.  If an enlarged
continuum branch claims a positive gap, it requires a uniform lower bound and
the Paper-VII-C limit theorems; the active framework makes no such claim.
\end{theorem}

\begin{theorem}[String-tension readout after factorization]
\label{p11:thm:v164-string-tension-readout}
The position metric area and Wilson observable are already generated on the
reciprocal-position chart; no monolithic momentum-to-position pairing
coefficient remains.  Fixed-loop convergence is part of QFT4, and the
local-flat chart exhaustion provides arbitrarily large rectangles.  After
taking the regulator limit at fixed \((R,T)\), positive transfer defines
\(V(R)\), and spatial reflection positivity gives the asymptotic slope
\[
 \sigma=\lim_{R\to\infty}V(R)/R.
\]
Strict positivity \(\sigma>0\) is a further claim on a specified pure-glue
or certified unscreened fundamental Wilson sector; it is not a required
property of every full-QCD screened ground-state or finite singlet readout.
Its actual complete-action area estimate is APBI or an alternative direct
surface/flux certificate.  Glueball-to-positive-string ratios retain this
sector identification and both computed readouts.
\end{theorem}

\begin{theorem}[Native carrier gap versus QCD readout gaps]
\label{p11:thm:v165-native-carrier-gap-qcd-boundary}
The mass gap predicted by the framework is the finite carrier-support gap
\[
 \mathbf1_{(E_{56},E_{40})}(H_{\rm car})=0,
 \qquad E_{48}=\sqrt{E_{40}E_{56}},
 \qquad P_{48}^{\rm act}=0.
\]
Finite QCD transfer gaps, glueball energies, and pole residues remain
well-defined readouts of their specified same-action blocks, but they do not
promote this structural theorem to a continuous Yang--Mills gap.  Accordingly
the positive continuum gap is removed from the active prediction frontier.
It remains only an optional comparison claim whose own uniform-scaling and
reconstruction hypotheses would have to be proved before use.
\end{theorem}


% Source Paper VII-B, original line 4738; inherited proof bodies unchanged.
\subsection{Effective rooted mass-source handoff}
\label{p11:sec:v2109-rooted-source-handoff}
The charged mass-source constructor is already closed on the occurring
rooted charged/action branch.  Import the complete mass/incidence sibling
packet and preterminal source of
Theorem~\ref{p5:thm:v397-charged-preterminal-srcval-realization} and
Corollary~\ref{p5:cor:v397-charged-preterminal-assembly-closure}:
\[
 R_s^{\rm ch}=\sum_r m_{s,r}^{\rm fin}\iota_{s,r},\qquad
 D_{F,s}^{\rm ch}
 =\Sigma_{\sigma_s}\otimes P_{Q_s}\otimes R_s^{\rm ch}+\mathrm{h.c.},
 \quad s\in\{u,d\}.
\]
The dependency is this inherited finite source followed by the applicable
\(\mathrm{G16}_{\rm scheme/RG}\) transport to \(M_q^{\rm phys}(\mu)\).
The direct finite coefficients and their family incidences remain siblings;
they are not redefined as singular values of $R_s^{\rm ch}(R_s^{\rm ch})^\dagger$.
The former noncentral SVD/terminal-binding gap is superseded on this branch,
not a live QCD constructor or a second source-member occurrence gate.
Actual scheme/threshold/running evaluation remains dependent execution.
Pole certification is consumed only by a pole claim; state/history, ordered
chiral limits and extensive bounds remain only where the readout uses them.


% BEGIN FRAMEWORK v2.228: CHARGED ORIGIN AND CHARACTER TESTS

% Source Paper VII-B, original line 4763; inherited proof bodies unchanged.
\subsection{Current charged-input origin in hadronic readouts}
\label{p7b:sec:v228-charged-origin}
Every hadronic readout consuming the noncentral charged source inherits
Definition~\ref{p5:def:v228-hierarchy-origin} and the generated-packet
scope of Proposition~\ref{p5:prop:v228-charged-origin-scope}.
The finite coefficient ledger is retained. A source-visible hadronic pole
requires the actual strong kernel and state; conditional charged assembly
alone does not discharge that calculation or identify a terminal coefficient
with a hadron rest mass.

% END FRAMEWORK v2.228: CHARGED ORIGIN AND CHARACTER TESTS


% Source Paper VII-B, original line 5680; inherited proof bodies unchanged.
\subsection{Native operator coefficients and hadronic readout scope}
\label{p11:sec:v226-hadron-coefficients}
\begin{corollary}[Complete finite momentum readout within an actual hadron frame]
\label{p11:cor:v226-hadron-coefficients}
After an actual hadron source/overlap frame and its finite action are supplied,
the native momentum part of its complete response can be reconstructed by
Theorem~\ref{p3f:thm:v226-response-reconstruction} and the tensor coefficient
inverse. Its admitted residues and source derivatives obey the same inverse.
A common positive kinetic normalization admits the error certificate of
Proposition~\ref{p2d:prop:v226-coefficient-spectral-error} on that native block.
\end{corollary}
\begin{proof}
Apply the complete response inverse before the actual state/source
contractions. The bounded finite sums commute with residue integration and
with derivatives on their admitted domain. The outer hadron Gram matrix,
source overlaps and normalization are unchanged inputs.
\end{proof}
The native 137 count is not the dimension of a full hadron, nuclear or
occupation space. The Hurwitz spin--flux, finite pole and thermal/nuclear
readouts retain their actual source frames, clocks and process-specific
matrices. A finite informationally complete momentum measurement does not
select those physical matrices or promote a finite pole to a continuum
scattering observable.

% END FRAMEWORK v2.226: KINEMATIC CLOSURES

% BEGIN FRAMEWORK v2.227: DOWNWARD CLOSURES

% Source Paper VII-B, original line 5707; inherited proof bodies unchanged.
\subsection{Hadronic pole and canonical-state handoffs}
\label{p11:sec:v227-hadron-response}
\begin{corollary}[Complete source response closes the finite pole comparison]
\label{p11:cor:v227-hadron-response}
For two actual finite Hermitian hadronic response words on one complete
retained/source target, the sample count in
Theorem~\ref{p1b:thm:v227-finite-response} suffices for equality of all
admitted same-source poles and residues. The dimension bound is the actual
complementary carrier size; 137 tests apply only to a genuinely native
momentum matrix factor. An unused reducing dark complement need not be
selected for this source-visible claim.
\end{corollary}
\begin{proof}
The sampled equality gives the whole meromorphic response. Its common
inverse and contour residues agree, and the already fixed source
contractions preserve that equality.
\end{proof}
The normalized full-ground/Gibbs state covariance is inherited from each
actual equivariant finite Hamiltonian. This propagates the finite state
antecedent to Wilson and hadron readouts without proving a pure-glue
continuum gap, scattering, or an area law. Odd baryonic/quark source
components retain the relative central offsets of
Proposition~\ref{p1b:prop:v227-central-boundary}; actual source incidence,
normalization and thermal/nuclear weights remain in the readout.

% END FRAMEWORK v2.227: DOWNWARD CLOSURES

% End of predecessor document.
\paragraph{Final theorem status and export boundary.}
The exact finite readouts are distinct from measured hadron masses, resonance widths, parton distributions or nuclear binding without their state, source and limiting certificates. The non-derivational Hurwitz numerical comparisons remain available in Paper~VII-D.

\paragraph{Technical-reference convention.} References in inherited theorem text to labels formerly described as appendices and now assigned to Paper~VII-D are references to that separate technical companion, not local appendices.



% BEGIN FRAMEWORK v2.338: MESON REST-ENERGY FALSIFIABILITY AND FINITE SOURCE RESIDUES
\subsection{Petersen spectral poles versus geometric-source residues}
\label{p11:sec:v338-source-poles-boundary}
The exact finite geometric family of Paper~VII-A provides a useful
model of source accessibility, not a predicted physical spectrum.

\begin{proposition}[Formal source resolvents and Hellmann--Feynman weights]
\label{p11:prop:v338-formal-source-residues}
For the moving roots $\lambda_a(t)$ of the three secular factors
$q_a(z,t)$ in Theorem~\ref{p10:thm:v338-full-occupancy-secular}, let
\[
 N_T=z^2-138z+3024,\quad N_S=z-18,\quad
 N_E=z^2-36z+216.
\]
The normalized formal source resolvent for each reduced channel is
$G_a(z,t)=N_a(z)/q_a(z,t)$.  At a simple pole $\lambda_a(t)$
its residue is
\[
 w_{a,\lambda}=\frac{N_a(\lambda)}{\partial_zq_a(\lambda,t)}
    =\frac{d\lambda_a(t)}{dt},
 \qquad 0\leq w_{a,\lambda}\leq1.
\]
The weights sum to one in each normalized source channel.
A spectral line dark to this finite source is not absent from the
full geometric operator; it is merely inaccessible to that probe.
\end{proposition}
\begin{proof}
Each determinant factor is linear in $t$ with
$\partial_tq_a=-N_a$.  Implicit differentiation at a simple root
gives the equality.  Since $\partial_tK=D_s$ is an orthogonal
projection, the finite Hellmann--Feynman formula bounds the slope
between zero and one and the spectral theorem normalizes the
residues.  No claim is made about physical correlators.
\end{proof}

\begin{proposition}[What remains to derive for ground-state mesons]
\label{p11:prop:v338-native-meson-no-mass-prediction}
The oriented-pair Hodge weights $(2/5,3/10,3/10)$, the rooted
Pfaffian phases, the three fifteen-dimensional two-source carriers,
and the finite incidence spectra have no unique action-selected
map to physical ground-state meson masses.  In particular:
(i) none fixes the actual pseudoscalar or vector meson source;
(ii) the possible geometric ratio $25/2$ compares different typed
norm readouts and cannot be installed as a pion/kaon squared-mass
identity; (iii) the rank-fifteen contact projector supplies only
one positive contact eigenlevel; and (iv) no charged--neutral
splitting follows from these operators alone.
The required observable test is the source-owned four-state spectrum
$\{\pi^0,\pi^{\pm},K^{\pm},K^0\}$, including the correct charge
splittings, derived before using their observed masses to select
operator coefficients.
\end{proposition}
\begin{proof}
The orthogonal Hodge idempotents permit independent geometric
source coefficients, as shown in Paper~III-F.  The contact
projector has spectrum $\{1^{15},0^{30}\}$.
Neither these geometric facts nor the representation-level rooted
comparison selects the missing color-singlet physical embedding
and actual action Hessian; different allowed coefficients have
different pole positions with identical incidence/representation
identities.  Consequently there is no unique mass formula under
the stated inputs.  The four observables supply the stringent
future discriminant, not present input data.
\end{proof}
\begin{remark}[Finite-native route, no assumption of QCD running]
A future direct source-owned finite transfer may provide a
different dynamical route from conventional RG descriptions.
Nothing in these exact group-theoretic derivations requires quark
running masses as \emph{axioms}; equally, refusing an RG input does not
by itself determine the finite physical Hamiltonian.  No spectrum
claim is made until the same-action source map and mass readout
have been evaluated.
\end{remark}
% END FRAMEWORK v2.338




% BEGIN SOURCE-INTEGRATED CROSS-PAPER MATHEMATICS v3_017, 15 November 2026.
%  Derived geometric comparison modules; physical source occurrence not promoted.
%  E3/E4 independent audit of these NEW results remains pending; inherited source statements retained.
\section{Additional exact mathematical consequences from the verified source geometry}

% -------- Research module Paper7B-NonetAlgebraBridge-v1_2 (latest supplied version) --------
% Proposed insert for future Paper7B-QCDReadouts-v1_159.tex.
% Research derivations, not in PapersMinimal v3_014.
% Imported exact mathematical inputs:
%   p11:thm:v1158-directed-two-core-nonet,
%   p10:prop:v1173-formal-two-core-meson-port,
%   p3c:thm:v179-clifford-projector-tangent,
%   p3c:thm:v237-generated-clifford-exceptional-resolution,
%   p1:thm:v285-ten-axis-quadrupole-exact-sequence.
% Physical source, mass, pole, width, coupling: NOT certified.

\subsection{A derived algebraic structure and a Clifford Gram on the two-core nonet}
\label{p11:subsec:v159-nonet-algebra-clifford}

\begin{theorem}[Jordan and compact Lie-algebra constructions on the nine-mode comparison]
\label{p11:thm:v159-nonet-jordan-lie-structure}
Let $\mathcal N_9=\mathbf1\oplus\mathbf3\oplus\mathbf5$
be the $A_5$-invariant orthogonal comparison module of
Theorem~\ref{p11:thm:v1158-directed-two-core-nonet}, and \emph{choose}
the equivariant real isometry
$\Upsilon:\mathcal N_9\to\operatorname{Herm}_3(\mathbb C)$
already admitted in
Proposition~\ref{p10:prop:v1173-formal-two-core-meson-port}.
Write $V=\operatorname{Im}\mathbb H\simeq\mathbb R^3$ and
$A(v)w=v\times w$ for the real skew-symmetric cross-product matrix.
In the induced decomposition, any image has the unique form
\[
\boxed{X=\Upsilon(\xi)=aI_3+Q+iA(v),\quad
 a\in\mathbb R,\ v\in V,\ Q\in\operatorname{Sym}_0^2(V).}
\]
The following operations are well-defined $A_5$-equivariant
\emph{mathematical transports} on $\mathcal N_9$:
\[
\xi\circ\eta=\Upsilon^{-1}\!\left(\frac{XY+YX}{2}\right),\qquad
\{\xi,\eta\}=\Upsilon^{-1}\!\bigl(-i[X,Y]\bigr)
\quad(\operatorname{tr}X=\operatorname{tr}Y=0).
\]
The first is the Euclidean Jordan algebra
$H_3(\mathbb C)$, whose connected Jordan automorphism group is
$PSU(3)$; the second is a compact simple real Lie algebra on the
traceless eight-dimensional subspace, isomorphic to $\mathfrak{su}(3)$.
Under the given $A_5\subset SO(3)$ action this Lie algebra splits as
$\mathbf3\oplus\mathbf5$. Its invariant scalar, quadratic and cubic
polynomials may be chosen as
\[
\boxed{\begin{aligned}
 \operatorname{tr}X&=3a,\\
 \operatorname{tr}X^2&=3a^2+\operatorname{tr}Q^2+2\|v\|^2,\\
 \det X&=a^3-\frac a2\bigl(\operatorname{tr}Q^2+2\|v\|^2\bigr)
             +\frac13\operatorname{tr}Q^3-v^{\mathsf T}Qv.
\end{aligned}}
\]
The isomorphism class of these transported structures is independent
of the chosen equivariant isometry, but their actual structure
constants on the labeled $20$-pair space depend on that choice.
\end{theorem}
\begin{proof}
The real decomposition
$\operatorname{Herm}_3(\mathbb C)=\mathbb RI
 \oplus i\mathfrak{so}(3)\oplus\operatorname{Sym}_0^2(\mathbb R^3)$
is exactly the $1+3+5$ character already proved in Paper VII-B.
For $R\in A_5\subset SO(3)$, one has
$A(Rv)=RA(v)R^{\mathsf T}$, so conjugation by $R$ commutes with
both the Jordan product and the commutator.
The matrix Jordan identity follows from associativity of matrix
multiplication. Traceless Hermitian matrices are closed under
$-i[X,Y]$, and multiplication by $-i$ identifies this bracket with
the ordinary bracket of $\mathfrak{su}(3)$.
The connected automorphisms $X\mapsto UXU^\dagger$ have
$U\in SU(3)$ and scalar-center kernel, giving $PSU(3)$.
Finally $A(v)^2=vv^{\mathsf T}-\|v\|^2I$ and
$\operatorname{tr}Q=0$ yield
$\operatorname{tr}(Q+iA(v))^2=\operatorname{tr}Q^2+2\|v\|^2$ and
$\operatorname{tr}(Q+iA(v))^3=\operatorname{tr}Q^3-3v^{\mathsf T}Qv$.
The determinant formula follows from the characteristic polynomial
of a traceless $3\times3$ matrix.
\end{proof}

\begin{proposition}[Explicit Jordan multiplication and octet fusion typing]
\label{p11:prop:v159-nonet-jordan-coordinate-product}
For $X=aI+Q+iA(v)$ and $Y=bI+R+iA(w)$, write
$X\circ Y=cI+T+iA(t)$. Then
\[
\boxed{\begin{aligned}
 c&=ab+\tfrac13\operatorname{tr}(QR)+\tfrac23v\cdot w,\\
 t&=aw+bv-\tfrac12(Qw+Rv),\\
 T&=aR+bQ+\tfrac12(QR+RQ)
          -\tfrac13\operatorname{tr}(QR)I\\
  &\hspace{2.5em}-\tfrac12(vw^{\mathsf T}+wv^{\mathsf T})
          +\tfrac13(v\cdot w)I.
\end{aligned}}
\]
Thus the comparison product has precisely typed channels
$\mathbf3\circ\mathbf3\subset\mathbf1\oplus\mathbf5$,
$\mathbf3\circ\mathbf5\subset\mathbf3$ and
$\mathbf5\circ\mathbf5\subset\mathbf1\oplus\mathbf5$.
Its cubic octet polynomial is
$\operatorname{tr}Q^3-3v^{\mathsf T}Qv$.
This is a \emph{flavor-tensor construction}, not an asserted cubic
pseudoscalar vertex (which must also satisfy parity, Lorentz and
source-action selection conditions).
\end{proposition}
\begin{proof}
Expand $(XY+YX)/2$ and use the matrix identities
$A(v)A(w)=wv^{\mathsf T}-(v\cdot w)I$ and
$QA(w)+A(w)Q=-A(Qw)$ for trace-free symmetric $Q$.
Sorting real scalar, real symmetric trace-free and imaginary skew
parts gives the displayed coefficients.
\end{proof}

\begin{proposition}[Balanced Clifford-to-nonet comparison Gram]
\label{p11:prop:v159-e36-nonet-triplet-gram}
Suppose a five-core $2I$ comparison packet has its ten shared unoriented
axes $u_a\in V$ and identify the nonet's matching $\mathbf3$ with
$V$ via the preceding $\Upsilon$. Let
$T:\mathcal N_9\to V$ be the \emph{isometric} triplet extraction,
zero on $\mathbf1\oplus\mathbf5$; explicitly, in the
Hilbert--Schmidt coordinates $X=aI+Q+iA(v)$ above one takes
$T(\xi)=\sqrt2\,v$ because
$\operatorname{tr}(iA(v))^2=2\|v\|^2$.
For each $a$, define the finite \emph{comparison} map
\[
 M_a(\xi)=E_{36}S_{u_a}\bigl((I-u_au_a^{\mathsf T})T\xi\bigr).
\]
Equip the target source block with the inherited whole-octet trace
metric $\tau_8$. Then
\[
 M_a^*M_a=\tfrac14E_{36}^2\,T^*
  (I-u_au_a^{\mathsf T})T,
 \qquad \operatorname{spec}(M_a^*M_a)
  =\{(E_{36}^2/4)^{(2)},0^{(7)}\}.
\]
Taking the orthogonal direct sum of the ten \emph{comparison} target
blocks, with normalization $1/\sqrt{10}$, gives
\[
\boxed{G_{36}^{(9)}:=\frac1{10}\sum_{a=1}^{10}M_a^*M_a
    =\frac{E_{36}^2}{6}P_{\mathbf3},\quad
 \operatorname{spec}G_{36}^{(9)}
    =\{(E_{36}^2/6)^{(3)},0^{(6)}\}.}
\]
Here $E_{36}^2=E_{32}E_{40}$ is the inherited transition
\emph{weight}, not a meson pole mass or a new eigenvalue of the
protected carrier Hamiltonian.
\end{proposition}
\begin{proof}
The polarised normalized source identity of
Theorem~\ref{p3c:thm:v179-clifford-projector-tangent}
gives the first Gram formula. Since $S_{u_a}$ is injective on
$u_a^\perp$ and $T$ is an isometry on the matching triplet, its
single-axis Gram has rank two and the indicated eigenvalues.
The ten-axis tight-frame identity from the generated Clifford
exceptional resolution gives
$\sum_{a=1}^{10}u_au_a^{\mathsf T}=\tfrac{10}{3}I_V$.
Substituting it yields
$\frac1{10}\sum_a(I-u_au_a^{\mathsf T})=\frac23 I_V$,
so the averaged Gram is $\frac{E_{36}^2}{6}P_{\mathbf3}$.
Direct-sum comparison avoids identifying the ten physical target
blocks or asserting a common QCD owner.
\end{proof}

\begin{proposition}[Fivefold axial resolution of the nonet quintet]
\label{p11:prop:v159-fivefold-axial-nonet-resolution}
Let $n\in V$ be a \emph{marked} unit fivefold axis of the supplied
$2I$ packet. On the induced $\mathbf5=\operatorname{Sym}_0^2(V)$,
define $\mathcal L_n(Q)=[A(n),Q]$ and the positive operator
$B_n=-\mathcal L_n^2$. Then
\[
\boxed{\operatorname{spec}B_n=\{0^{(1)},1^{(2)},4^{(2)}\},}
\]
with mutually orthogonal spectral projectors
\[
\boxed{\begin{aligned}
 P_0&=\frac{(B_n-I)(B_n-4I)}4,\\
 P_1&=\frac{B_n(4I-B_n)}3,\\
 P_2&=\frac{B_n(B_n-I)}{12}.
\end{aligned}}
\]
They give a purely geometric comparison decomposition
\[
\boxed{\mathcal N_9=\mathbf1\oplus\mathbf3
  \oplus\underbrace{(\mathbf2\oplus\mathbf2)}_{\mathbf4}
  \oplus\mathbf1,}
\]
where the four-dimensional quintet part is two \emph{inequivalent}
real $D_5$ doublets. This makes available a $3+4+1$ comparison
splitting of the octet $\mathbf3\oplus\mathbf5$ and a separate
nonet singlet, but neither the degeneracy of the four nor its physical
identification with kaons follows from the $D_5$ symmetry alone.
Indeed, $\mathbf3\downarrow D_5=\varepsilon\oplus\mathbf2_1$ and
$\mathbf5\downarrow D_5=\mathbf1\oplus\mathbf2_1\oplus\mathbf2_2$;
therefore $D_5$-equivariant physical correlators may mix the two
copies of $\mathbf2_1$ and may mix the two trivial lines from the
full nonet. A physically selected axis, source owner, kinetic
normalization, and the inherited $\mathbf3/\mathbf5$ block
protection are additional conditions for the displayed \emph{block}
reading.
\end{proposition}
\begin{proof}
Rotate coordinates so that $n=e_3$. The trace-free quadratic
forms $(x^2+y^2-2z^2)$, $(xz,yz)$ and
$(x^2-y^2,2xy)$ carry planar angular frequencies $0,1,2$
under rotations about $e_3$, respectively. The infinitesimal
rotation generator $\mathcal L_n$ therefore has imaginary
frequencies $0,\pm i,\pm2i$ on the real five-space.
Its negative square has eigenvalues $0,1,4$ with multiplicities
$1,2,2$, and Lagrange interpolation gives the three projectors.
The stabilizer of an unoriented icosahedral fivefold axis is $D_5$;
the stated restrictions and possible multiplicity-space mixings
follow from the real $D_5$ character decomposition.
\end{proof}

\begin{remark}[Physical occurrence and action ownership remain open]
\label{p11:rem:v159-nonet-algebra-physical-firewall}
All three constructions are defined using \emph{existing mathematical}
module identifications and the finite comparison response. An
$A_5$-isometric identification does not license the Jordan product
as the actual meson interaction Hamiltonian, does not establish a
continuous physical $SU(3)_f$ action, and does not identify the
quaternionic $E_{36}$ response with a QCD correlator residue.
The physical octet, singlet, $3+4+1$ chiral source channels,
quark-flavor spurions and positive kinetic norms must be connected by
one generated, source-visible map on the same action owner. In
particular none of the displayed invariants fixes a meson mass,
width or branching fraction.
\end{remark}

% BEGIN proposed v1_2 EXTENSION: explicit 20-pair polar isometry.
% Requires source adoption of p1:thm:cc12-directed-coherent-gram
% before this cross-reference is integrated into a canonical release.

\subsection{Source-ready directed-pair polar map for the nonet comparison}
\label{p11:sec:cc12-nonet-polar-map}
\begin{corollary}[The core-pair projector has a positive coherent Gram certificate]
\label{p11:cor:cc12-nonet-coherent-polar-map}
Let $\mathcal B_{20}$ be the ordered two-core pair module of
Theorem~\ref{p11:thm:v1158-directed-two-core-nonet}.
Using the same supplied exact $2I$ embedding, form the directed-axis
projector synthesis
\[
 S:\mathcal B_{20}\to\operatorname{Herm}_3(\mathbb C),
 \qquad Se_{ij}=\tfrac12\bigl(I-n_{ij}n_{ij}^T+iA(n_{ij})\bigr)
\]
of the proposed Paper-I-B theorem
\ref{p1:thm:cc12-directed-coherent-gram}. Its kernel and image
are the already specified representations:
\[
 \ker S=\mathbf3'\oplus2\mathbf4,
 \qquad \operatorname{im}S=\mathbf1\oplus\mathbf3\oplus\mathbf5.
\]
In particular the rank-nine projector of Paper VII-B is reconstructed
without a free nine-dimensional subspace choice:
\[
 \boxed{P_9=I_{20}
 -(I_{20}-\tfrac3{20}S^*S)
  (I_{20}-\tfrac3{10}S^*S)
  (I_{20}-\tfrac32S^*S).}
\]
The polar normalization
\[
 \boxed{\Upsilon_{\rm polar}
   =(SS^*)^{-1/2}S\big|_{\operatorname{im}P_9}}
\]
is an $A_5$-equivariant \emph{isometry} onto
$\operatorname{Herm}_3(\mathbb C)$, fixed by the quaternionic
coordinate/handedness data up to the overall complex-conjugation hand.
It supplies a specific mathematical candidate for the otherwise
chosen comparison isometry used above; it does not supply the
quark-flavor current, mass spurion, physical kinetic normalization,
interacting pole or width.
\end{corollary}
\begin{proof}
The coherent Gram has nonzero eigenvalues $20/3$ on the singlet,
$10/3$ on the selected $\mathbf3$ and $2/3$ on the quintet, and
vanishes on $\mathbf3'\oplus2\mathbf4$. The polynomial takes value
one on those three eigenspaces and zero on the kernel.
Since $SS^*$ is strictly positive on the whole nine-dimensional
Hermitian target, its positive inverse square root gives the polar
isometry. The $A_5$-covariance of the projectors implies the polar
operator commutes with the action; the global orientation reversal
acts by complex conjugation, preserving the real Hilbert--Schmidt
metric and Gram.
\end{proof}

\begin{remark}[New conditional source/observable contract]
If a single native VII-A meson-source owner is ever shown to
intertwine the selected two-core comparison and to implement the
auxiliary effects $E_{ij}=\tfrac3{20}Se_{ij}$, its measured
directed-label probabilities must obey all eleven homogeneous
comparison constraints of proposed
Theorem~\ref{p2c:thm:cc12-twenty-complex-povm} and
Proposition~\ref{p2c:prop:cc12-directed-probability-constraints}.
This would be a finite, directly falsifiable source test rather than
a guessed meson mass formula. Neither the intertwiner nor the
physical readout occurrence is asserted or derived by the
mathematical Gram itself.
\end{remark}
% END proposed v1_2 EXTENSION

% -------- Research module Paper7B-EnergyChargeZeroFreeGeometry-v1_1 (latest supplied version) --------
% Proposed research insertion for Paper VII-B. Source member fixed.
\subsection{Convex-spectral geometry of finite thermal zeros and charge response}
\label{p11:sec:cc13-zero-free-polytope}

\begin{theorem}[Exact spectral-polytope zero-free tube]
\label{p11:thm:cc13-zero-free-polytope}
Let $H$ be a fixed finite self-adjoint Hamiltonian and
$Q_1,\ldots,Q_r$ be commuting self-adjoint integer-spectrum charges
which commute with $H$.  This includes a charge-conserving fixed
member of the finite QCD source defined in
Definition~\ref{p11:def:native-thermal-partition-function}, on its
explicitly admitted vectorlike branch.  Choose a common eigenbasis
with energies $E_j\in\mathbb R$ and charge tuples
$q_j\in\mathbb Z^r$, counting multiplicities.  For complex source
coordinates $z=(\beta,\xi_1,\ldots,\xi_r)$ define
\[
 Z(z)=\operatorname{Tr}\exp(-\beta H+\sum_a\xi_aQ_a)
 =\sum_j e^{z\cdot s_j},\qquad s_j=(-E_j,q_j)\in\mathbb R^{r+1}.
\]
Let $\mathcal P=\operatorname{conv}\{s_j\}$ and define its directional
width $W_{\mathcal P}(y)=\max_{s\in\mathcal P}y\cdot s
-\min_{s\in\mathcal P}y\cdot s$.  For every $z=x+iy$ with
$W_{\mathcal P}(y)<\pi$, one has
\begin{equation}\label{p11:eq:cc13-zero-free-lower}
 \boxed{|Z(x+iy)|\ge
 \cos\!\left(\frac{W_{\mathcal P}(y)}2\right)Z(x)>0.}
\end{equation}
In particular $Z$ has no zeros on the open convex tube
$\{x+iy:W_{\mathcal P}(y)<\pi\}$, for \emph{every real source}
coordinate $x$.  The zero-free domain is certified by the actual
finite spectrum's convex hull, not by a fitted thermal scale or
an assumed physical mass gap.
\end{theorem}
\begin{proof}
Let $\theta_+=\max_j y\cdot s_j$,
$\theta_-=\min_j y\cdot s_j$, and
$\theta_0=(\theta_++\theta_-)/2$.  The rotated exponential sum
has real part
\[
 \operatorname{Re}(e^{-i\theta_0}Z(x+iy))
 =\sum_j e^{x\cdot s_j}\cos(y\cdot s_j-\theta_0)
 \ge\cos(W_{\mathcal P}(y)/2)\sum_j e^{x\cdot s_j}.
\]
All real weights are positive, and $W_{\mathcal P}(y)<\pi$
makes the last cosine strictly positive.  Taking absolute values
proves the bound.  The directional width is convex and positively
homogeneous as the support function of $\mathcal P-\mathcal P$,
so the displayed imaginary-source region is open and convex.
\end{proof}

\begin{corollary}[Energy strip, fugacity polynomial and integer-charge periodicity]
\label{p11:cor:cc13-fugacity-strips}
Write $\Delta E=E_{\max}-E_{\min}$.
For real $\beta_0$ and $|t|\Delta E<\pi$,
\[
 |Z(\beta_0+it,0)|\ge
 Z(\beta_0,0)\cos(|t|\Delta E/2)>0.
\]
For real $\beta>0$ and real chemical potentials $\mu$, the
commuting-integer-charge thermal trace is a finite positive-coefficient
Laurent polynomial in the fugacities $u_a=e^{\beta\mu_a}$:
\[
 \boxed{Z(\beta,\beta\mu)=
 \sum_{q\in\mathcal Q}Z_q(\beta)\,e^{\beta\mu\cdot q},\qquad
 Z_q(\beta)=\operatorname{Tr}_{\mathscr H_q}e^{-\beta H}>0.}
\]
It is exactly periodic under $\mu_a\mapsto\mu_a+2\pi i/\beta$.
The canonical fixed-charge partition functions satisfy
\[
 Z_q(\beta)=\frac1{(2\pi)^r}\int_{[-\pi,\pi]^r}
 Z(\beta,i\theta)e^{-i\theta\cdot q}\,d^r\theta.
\]
If $\theta_a=\beta\operatorname{Im}\mu_a$, the sufficient
zero-free condition is
$\sum_a|\theta_a|(q_{a,\max}-q_{a,\min})<\pi$;
using the actual joint support width
$\max_q\theta\cdot q-\min_q\theta\cdot q<\pi$ is sharper.
\end{corollary}
\begin{proof}
The first claim specializes the width theorem to $y=(t,0)$.
Commuting charge blocks provide the positive Laurent expansion.
Integer charges make each factor $e^{i2\pi q_a}$ equal one.
The Fourier inversion follows from character orthogonality on the
$r$-torus.  The width inequalities follow by taking maxima/minima of
the charge linear forms.
\end{proof}

\begin{corollary}[Finite charge-susceptibility bound]
\label{p11:cor:cc13-charge-variance}
At any real thermal source $\beta>0,\mu\in\mathbb R^r$,
$\partial_{\xi_a}\partial_{\xi_b}\log Z
=\operatorname{Cov}(Q_a,Q_b)$.  For every real vector $v$,
\[
 \boxed{0\le v^{\mathsf T}\nabla_\xi^2\log Z\,v
 =\operatorname{Var}(v\cdot Q)
 \le \tfrac14\left(\max_qv\cdot q-\min_qv\cdot q\right)^2.}
\]
For derivatives with respect to $\mu$ rather than $\xi=\beta\mu$,
multiply by $\beta^2$.  This is a finite source-support bound,
not a numerical susceptibility prediction without the Hamiltonian.
\end{corollary}
\begin{proof}
Differentiate the finite positive charge sum.  The variance upper
bound is the elementary bounded-real-variable inequality
$\operatorname{Var} X\le (\max X-\min X)^2/4$.
\end{proof}

\begin{remark}[Necessary scope protections]
The exact imaginary-chemical-potential periodicity requires
\emph{integer} charge normalization; fractional quark-number
coordinates have their correspondingly rescaled period.  Joint
charge eigenspaces require the stated commutation relations, not merely
individually self-adjoint sources.  Refinement-uniform nonzero limits
or real thermal phase transitions are not inferred: the spectral and
charge support widths can grow with the regulator.  The theorem
supplies a testable sufficient analytic chart on each actual finite
member only.
\end{remark}

% -------- Research module Paper7B-ChargeFanZeroTemperature-v1_1 (latest supplied version) --------
% Born--Hopf Geometry cross-DAG research insert, PROPOSED ONLY.
% Destination: Paper VII-B, after finite thermal trace and zero-free supplement.
\section{Zero-temperature charge fan and finite-gap bounds}
\label{p11:sec:cc14-charge-fan}
\begin{theorem}[Polyhedral ground-charge fan of a fixed finite member]
\label{p11:thm:cc14-charge-fan}
On a fixed finite-dimensional, charge-conserving sector of
Definition~\ref{p11:def:native-thermal-partition-function}, assume
$H$ and all real self-adjoint $Q_a$ mutually commute, and let
$(E_j,q_j)$, $j=1,\ldots,N$, be their simultaneous joint
energy--charge eigenvalues, repeated with multiplicity.  For real
chemical potentials $\mu\in\mathbb R^r$, set
\[
 k_j(\mu)=E_j-\mu\cdot q_j,\qquad
 m(\mu)=\min_j k_j(\mu),\qquad
 Z_\beta(\mu)=\sum_j e^{-\beta k_j(\mu)}.
\]
Then $m(\mu)$ is a concave, continuous, piecewise-affine function
whose linearity domains are normal-fan cells of the lower
energy--charge convex hull.  Ground-sector changes can occur only
on the explicitly computable hyperplanes
\[
 \boxed{\mu\cdot(q_i-q_j)=E_i-E_j.}
\]
On every open cell with a unique ground \emph{charge vector}
$q_*$, the zero-temperature charge readout is $q_*$; it is a
piecewise constant plateau and changes only upon reaching a
cell boundary.  All finite $\beta>0$ partition functions remain
strictly positive and real analytic in $\mu$.
\end{theorem}
\begin{proof}
A pointwise minimum of finitely many affine functions is concave,
continuous and piecewise affine.  Dominance changes only at equality
of two affine functions.  The active affine functions correspond to
exposed lower faces of the joint convex hull.  Within a fixed
strict cell, all ground states have the same charge vector by
hypothesis; the finite Gibbs weights concentrate on them as
$\beta\to\infty$.  At finite positive $\beta$ the positive finite
exponential sum and its logarithm are real analytic.
\end{proof}

\begin{theorem}[Exact finite-temperature approach to the charge fan]
\label{p11:thm:cc14-charge-fan-error}
Fix $\mu$ and let $m=m(\mu)$ have ground multiplicity $g\ge1$.
Suppose the next distinct value has gap $\delta>0$ (if all states
are ground states, interpret the remainder as zero).  Then
\[
 \boxed{Z_\beta=e^{-\beta m}(g+R_\beta),\quad
 0\le R_\beta\le(N-g)e^{-\beta\delta}.}
\]
For the free energy $\mathcal F_\beta=-\beta^{-1}\log Z_\beta$,
\[
 \boxed{m-\frac{\log N}{\beta}\le\mathcal F_\beta
 \le m-\frac{\log g}{\beta},\qquad
 0\le m-\frac{\log g}{\beta}-\mathcal F_\beta
 \le\frac{N-g}{g\beta}e^{-\beta\delta}.}
\]
If $\rho_\beta=e^{-\beta K}/Z_\beta$ and
$\rho_0=P_{\min}/g$, the normalized ground projector, then
\[
 \boxed{\|\rho_\beta-\rho_0\|_1
 =\frac{2R_\beta}{g+R_\beta}
 \le\frac{2(N-g)}g e^{-\beta\delta}.}
\]
If every ground state has the same charge $q_*$, then with
$D_Q=\max_{i,j}\|q_i-q_j\|_2$, one further has
\[
 \|\langle Q\rangle_\beta-q_*\|_2
 \le D_Q\frac{N-g}{g}e^{-\beta\delta},
 \qquad
 \operatorname{Var}_\beta(v\cdot Q)
 \le \|v\|_2^2D_Q^2\frac{N-g}{g}e^{-\beta\delta}.
\]
\end{theorem}
\begin{proof}
Factor $e^{-\beta m}$ out of the finite thermal sum.
Every excited factor is at most $e^{-\beta\delta}$, giving the
remainder bound.  The log bounds follow from
$g\le g+R_\beta\le N$ and $\log(1+x)\le x$.
The two density matrices are diagonal in a common eigenbasis;
subtracting their normalized ground and excited probabilities gives
exact trace distance $2R_\beta/(g+R_\beta)$.
For ground charge $q_*$, the excited Gibbs probability is
$R_\beta/(g+R_\beta)$; all charge vectors differ from $q_*$
by at most $D_Q$.  Thus the expectation deviation is at most
$D_Q R_\beta/(g+R_\beta)$.  Similarly
$\operatorname{Var}(v\cdot Q)\le
\mathbb E[(v\cdot(Q-q_*))^2]
\le\|v\|^2D_Q^2 R_\beta/(g+R_\beta)$.
\end{proof}
\begin{remark}[No finite-size phase-transition claim]
The charge fan is the zero-temperature limit of a \emph{fixed finite}
Hamiltonian/charge packet; its edges are not thermodynamic-limit
phase transitions or measured QCD chemical-potential thresholds.
The equal-charge hypothesis is required for a single charge plateau,
and joint commutation is required for the simultaneous-eigenvalue
convex-hull formula.  The full QCD/refinement physical realization
requires the source's existing action, current and transfer ports.
\end{remark}

% -------- Research module Paper7B-EnergyChargeInformationGeometry-v1_1 (latest supplied version) --------
% Born--Hopf Geometry research insert v1.5. PROPOSED ONLY.
% Destination: Paper VII-B after native finite thermal trace.
\section{Energy--charge moment polytope and finite information geometry}
\label{p11:sec:cc15-energy-charge-information}
\begin{theorem}[Exact finite thermal moment map]
\label{p11:thm:cc15-energy-charge-moment-map}
On a fixed finite, mutually commuting energy--charge source chart of
Definition~\ref{p11:def:native-thermal-partition-function}, let
$v_j=(-E_j,q_j)\in\mathbb R^d$ enumerate the \emph{distinct} joint
energy--charge support points, with multiplicity $g_j>0$, and let
$\mathcal P=\operatorname{conv}\{v_j\}$ and
$L=\operatorname{span}\{v_j-v_k\}$. Choose one support point $v_*$ and put $z_j=v_j-v_*\in L$.
The centered trace below differs from the native trace by the nonzero
factor $\exp(-\theta\cdot v_*)$; the normalized probabilities and
zero set are unchanged. For any real natural parameter $\theta\in L$ define
\[
 Z(\theta)=\sum_jg_je^{\theta\cdot z_j},\quad
 A(\theta)=\log Z(\theta),\quad
 p_j(\theta)=\frac{g_je^{\theta\cdot z_j}}{Z(\theta)}.
\]
Then
\[
 \boxed{v_*+\nabla_L A=\sum_jp_jv_j=:\bar v,
 \qquad \nabla_L^2 A=\operatorname{Cov}_{p}(v)|_L>0.}
\]
The affine moment map $\theta\mapsto v_*+\nabla_L A(\theta)$ is
a smooth bijection, with smooth inverse, from $L$ onto the relative
interior $\operatorname{relint}\mathcal P$.
The convex dual $A^*(\bar v-v_*)$ is the finite relative-entropy problem
\[
 \boxed{A^*(y-v_*)=\min_{\substack{p_j\ge0,\ \sum_jp_j=1\\
              \sum_jp_jv_j=y}}
       \sum_jp_j\log\frac{p_j}{g_j},\qquad
       y\in\operatorname{relint}\mathcal P.}
\]
The minimizing distribution is the unique exponential-family member
$p_j(\theta(y))$; the inverse moment-map Jacobian is the inverse
covariance on $L$.
\end{theorem}
\begin{proof}
Differentiating the positive finite exponential sum gives the mean
and covariance formulas. For $u\in L\setminus\{0\}$, the scalar
$u\cdot v_j$ cannot be constant on all support points by the
definition of $L$. Every $p_j$ is positive, so its variance is
strictly positive. Therefore $A$ is strictly convex on $L$, and
its gradient is injective. The affine moment map $v_*+\nabla A$
takes values in the relative interior of the support polytope.
Conversely fix $y$ in that relative interior. The function
$A(\theta)-\theta\cdot(y-v_*)$ is coercive on $L$: its recession growth in
unit direction $u$ is $h_{\mathcal P}(u)-u\cdot y>0$ and the compact
unit sphere has a positive minimum. Thus it admits a unique critical
point, where $v_*+\nabla A=y$. The inverse-function theorem gives a smooth
inverse because the covariance is nonsingular on $L$.
Finally the elementary finite Gibbs variational identity
$\log\sum_jg_je^{\theta\cdot z_j}
=\max_{p\in\Delta}[\theta\cdot\sum_jp_jz_j
 -\sum_jp_j\log(p_j/g_j)]$
and convex duality give the constrained entropy formula and uniqueness.
\end{proof}

\begin{corollary}[One polytope controls fluctuations, complex zeros and cold charge fans]
\label{p11:cor:cc15-polytope-triple-control}
With the same fixed finite source support, let its width in any
$u\in L$ be
$W_{\mathcal P}(u)=\max_j(u\cdot v_j)-\min_j(u\cdot v_j)$.
For all real $\theta$,
\[
 \boxed{0<u^{\mathsf T}\nabla^2 A(\theta)u
 \le\tfrac14 W_{\mathcal P}(u)^2\quad(u\ne0).}
\]
Moreover, whenever $W_{\mathcal P}(\eta)<\pi$, the complexified finite
partition function satisfies
\[
 \boxed{|Z(\theta+i\eta)|\ge
 Z(\theta)\cos\bigl(W_{\mathcal P}(\eta)/2\bigr)>0.}
\]
The same support polytope supplies the lower supporting faces of the
zero-temperature charge fan of
Theorem~\ref{p11:thm:cc14-charge-fan}. These are three different
finite mathematical consequences of one spectral support object:
real covariance, complex zero exclusion and extreme-source plateaus.
\end{corollary}
\begin{proof}
The Hessian quadratic form is the variance of $u\cdot v$. For a
random variable in an interval of width $W$, the elementary
Popoviciu bound gives variance at most $W^2/4$; strict positivity
is established in the theorem. For complex $\theta+i\eta$ rotate
the sum of phasors $p_j e^{i\eta\cdot z_j}$ by the midpoint phase.
All rotated phases lie in $[-W/2,W/2]$, so the real part is at
least $\cos(W/2)$; multiply by $Z(\theta)$. Finally the
$\beta\to\infty$ limit minimizes $E_j-\mu\cdot q_j$, which is
exactly the lower-supporting-face problem already identified in
the cited charge-fan theorem.
\end{proof}
\begin{remark}[The full mathematical natural-parameter space is larger than physical temperature]
The bijection is on the auxiliary full real natural-parameter space
$L$. In a physical thermal application one restricts the energy
coordinate to $\theta_E=\beta>0$ (with $v_E=-E$) and the charge
coordinates to $\theta_Q=\beta\mu$, followed by projection onto $L$; this restricted chart need not realize all
of $\operatorname{relint}\mathcal P$. The theorem does not establish
a physical state selection, thermodynamic-limit phase transition,
mechanical pressure, or measured QCD susceptibilities.
\end{remark}

% -------- Research module Paper7B-EnergyChargeInverseStability-v1_1 (latest supplied version) --------
% CROSS-DAG RESEARCH v1.6. Proposed Paper VII-B insertion after finite
% thermal moment chart. Assumes a FIXED finite H and jointly commuting
% conserved charges; distinct from volume/metric response or physical T>0 coverage.
\subsection{Quantitative stability of the finite energy--charge inverse chart}
\label{p11:sec:cc16-thermal-inverse-stability}
\begin{theorem}[Certified inverse Lipschitz constant on bounded thermal source charts]
\label{p11:thm:cc16-thermal-inverse-stability}
Let a fixed finite Hamiltonian and the charges of the finite thermal
trace of Definition~\ref{p11:def:native-thermal-partition-function}
commute strongly and have joint distinct support points
$v_1,\ldots,v_s\in\mathbb R^d$ with positive multiplicities
$g_1,\ldots,g_s$. Put $L=\operatorname{span}\{v_j-v_1\}$,
$G=\sum_jg_j$, $\bar v=s^{-1}\sum_jv_j$,
\[
 D=\max_{j,k}\|v_j-v_k\|_2,\quad
 C_L=\left.\sum_{j=1}^s(v_j-\bar v)(v_j-\bar v)^{\mathsf T}
 \right|_L,
\]
where $C_L$ is positive definite on the nonzero affine span $L$.
For $\theta\in L$ define
\[
 p_j(\theta)=\frac{g_j e^{\theta\cdot v_j}}
 {\sum_k g_ke^{\theta\cdot v_k}},\qquad
 \mu(\theta)=\sum_j p_j(\theta)v_j.
\]
For every $B\ge0$ set
\[
 \boxed{\kappa_B=\frac{\min_j g_j}{G}
 e^{-BD}\,\lambda_{\min}(C_L)>0.}
\]
For all $\theta,\theta'\in L$ with $\|\theta\|,\|\theta'\|\le B$,
\begin{equation}\label{p11:eq:cc16-thermal-inverse-stability}
\boxed{(\mu(\theta)-\mu(\theta'))\cdot(\theta-\theta')
 \ge\kappa_B\|\theta-\theta'\|^2,
\quad\|\theta-\theta'\|
 \le\kappa_B^{-1}\|\mu(\theta)-\mu(\theta')\|.}
\end{equation}
Thus errors in an interior mean energy--charge measurement propagate
to uniquely identifiable auxiliary source parameters with a
computable spectrum-dependent bound on every bounded chart.
No bound is asserted uniformly as $B\to\infty$, across rank changes,
or for physical parameters outside the stipulated commuting fixed-owner
thermal chart.
\end{theorem}
\begin{proof}
For $\|\theta\|\le B$, every pair satisfies
$\theta\cdot(v_k-v_j)\le BD$, so
$p_j(\theta)\ge(g_j/G)e^{-BD}\ge p_*= (\min_jg_j/G)e^{-BD}$.
For $u\in L$, the covariance quadratic form is
\[
 u^{\mathsf T}\nabla_L\mu(\theta)u
 =\min_{c\in\mathbb R}\sum_jp_j(\theta)(u\cdot v_j-c)^2
 \ge p_*\min_{c\in\mathbb R}\sum_j(u\cdot v_j-c)^2
 =p_*u^{\mathsf T}C_Lu\ge\kappa_B\|u\|^2.
\]
The convex segment from $\theta$ to $\theta'$ remains in the
radius-$B$ ball. Integrating the Hessian lower bound along it proves
the strong-monotonicity inequality. Cauchy--Schwarz then gives the
inverse Lipschitz estimate. If $L=\{0\}$, there are no independent
source parameters and the conclusion is vacuous; no positive
$\kappa_B$ is assigned in that case.
\end{proof}
% END SOURCE-INTEGRATED CROSS-PAPER MATHEMATICS v3_017.

\printbibliography

% BEGIN FRAMEWORK v2.344: SAME-ACTION BARYON CURVATURE OBSERVABLE
\section{Matched baryon curvature as a single-action source readout}
\label{p11:sec:v344-curvature-correlator}
Let $B_{qqQ,J}$, $B_{qsQ,J}$ and $B_{ssQ,J}$ be three genuine
color-singlet interpolating operators with matched spin $J$ and
flavor/charge projections, all evaluated using \emph{one} specified
physical Hamiltonian $H_L$ and vacuum $\Omega_L$.  Define the
positive Euclidean correlators
\[
 C_a(t)=\langle\Omega_L,B_a e^{-t(H_L-E_{0,L})}
                 B_a^\dagger\Omega_L\rangle,
 \qquad a\in\{qqQ,qsQ,ssQ\}.
\]

\begin{proposition}[Spectral extraction without additive quark fits]
\label{p11:prop:v344-same-action-curvature-extraction}
Suppose for each channel $a$ the operator $B_a^\dagger\Omega_L$
has nonzero overlap with its unique lowest isolated selected state
of energy $M_a>0$ relative to the vacuum and the finite/gapped
spectral approximation is valid at the stated $t\to\infty$ limit.
Then
\[
 \boxed{K_{Q,J}:=M_{qqQ,J}-2M_{qsQ,J}+M_{ssQ,J}
   =-\lim_{t\to\infty}\frac{d}{dt}
      \log\frac{C_{qqQ}(t)C_{ssQ}(t)}{C_{qsQ}(t)^2}.}
\]
Every constant interpolation-source normalization cancels.
A strictly additive consistent constituent law
$M_{n_s}=A+Bn_s$, $n_s=0,1,2$, would give $K_{Q,J}=0$.
\end{proposition}
\begin{proof}
The spectral decomposition gives
$C_a(t)=Z_ae^{-M_at}(1+o(1))$ with $Z_a>0$ and exponentially
suppressed subleading terms under the stated isolated-gap
assumption.  Taking the logarithmic derivative proves the
identity, and the second finite difference annihilates any
function affine in $n_s$.
\end{proof}

\begin{corollary}[Petersen eligibility does not determine curvature]
\label{p11:cor:v344-character-no-hadron-mass}
The nonzero $C_3$-projected three-face mode of
Proposition~\ref{p10:prop:v344-threeface-character-witness} and
the exact geometric transfer norm $1/25$ of
Theorem~\ref{p1:thm:v344-twisted-25-gram} do not determine
$K_{Q,J}$ or the sign of any physical response.  A physical
flavor-to-character map, normalized source Gram, full auxiliary
resolvent, and pole identification must come from the same
constituted $H_L$.  No equality $K_{Q,J}=-E_{36}/25$ follows.
\end{corollary}
\begin{proof}
The color determinant is flavor blind, and a geometric rank-five
character representation cannot by itself select which of
$qq$, $qs$, or $ss$ couples to it.  Inserting a common correction
into all three masses leaves the curvature unchanged.  Other
flavor-dependent couplings compatible with the same finite
character identities yield different values and signs of
$K_{Q,J}$; hence the numerical coefficient is underdetermined.
\end{proof}

\begin{remark}[Resonances require pole extraction]
The large-$t$ identity isolates the lowest stable/gapped state of
each selected channel.  For broad or excited resonances the
physical complex pole and its coupled-channel residue, rather than
a nominal Breit--Wigner mass or the ground-state large-$t$
limit, are the required observable.
\end{remark}
% END FRAMEWORK v2.344


% BEGIN FRAMEWORK v2.345: F40 OBSERVABLE DISTINCTION
\clearpage
\section[Source-visible F40 mirrors: a conditional hadronic test]{Source-visible F40 mirrors:\\a conditional hadronic test}
\label{p11:sec:v345-f40-readout-hadrons}
\begin{proposition}[Three geometric probes versus one physical mass owner]
\label{p11:prop:v345-f40-no-mass-inference}
The right and left fivefold quotient maps and the ordered
five-core arrow readout jointly resolve the entire
$21$-dimensional F40 cycle carrier, with squared frame lower bound
$4-\sqrt{10}$.  Yet none of their dimensionless eigenvalues
supplies an energy pole of the physical Paper~VII-A Hamiltonian
without a color-singlet physical embedding and a source residue.
A source factored through $T_R+T_L$ is insensitive to both
Hodge mirror components.  A source with a nonzero
$T_R-T_L$ component is \emph{eligible}, but not guaranteed, to
see them.  The physically relevant comparison is a common-action
source correlation, for example
\[
 \langle\Omega_L,O_L^{\dagger}(z-H_L)^{-1}
 O_L\Omega_L\rangle,
\]
with the actual action domain and source normalization fixed;
not a selected root of the geometric frame polynomial.
\end{proposition}
\begin{proof}
These are the direct rank/kernel consequences of
Proposition~\ref{p1:prop:v345-f40-balanced-mirror}, together
with the same-action pole and leakage conditions in
Proposition~\ref{p10:prop:v345-f40-qcd-source-obligation}.
\end{proof}
% END FRAMEWORK v2.345

% BEGIN FRAMEWORK v2.349: NATIVE F40 RECONSTRUCTION VERSUS MASS POLES
\clearpage
\section{Verified F40 comparison data versus hadronic pole information}
\label{p11:sec:v349-f40-pole-firewall}
\begin{proposition}[A complete twenty-one-mode packet is not an energy spectrum]
\label{p11:prop:v349-f40-unique-packet-no-pole}
A native 44-coordinate F40 comparison triple can be
reconstructed from a unique twenty-one-dimensional cycle
precisely when the established 23 linear compatibility
conditions hold.  Its right/left Hodge mirror response is
related to the skew core-arrow response by the explicit
operator $K$ of Theorem~\ref{p1:thm:v349-f40-mirror-gluing},
with condition number $(1+\sqrt5)/2$.
This reconstruction and its comparison Laplacian eigenvalues
$8/R^2$ and $24/R^2$ provide no mass poles by themselves.
A hadronic claim requires an action-selected embedding and
source $O_L$, and must extract energy poles and residues from
\[
 G_{O_L}(z)=\langle\Omega_L,
 O_L^\dagger(z-H_L)^{-1}O_L\Omega_L\rangle
\]
for the \emph{full} Hamiltonian, retaining the leakage block
of Proposition~\ref{p10:prop:v345-f40-qcd-source-obligation}.
\end{proposition}
\begin{proof}
The finite reconstruction follows from the rank of the
native frame, its rational polynomial inverse and its exact
compatibility projector.  The geometric Laplacians act on
comparison coefficient spaces, whereas the resolvent acts on
the physical Hamiltonian domain.  Without a specified source
embedding and operator identity, no equality of their poles
or spectra is implied.
\end{proof}
% END FRAMEWORK v2.349

% BEGIN FRAMEWORK v2.350: PHYSICAL READOUT AND CUMULATIVE NUCLEAR COUNT REPAIR
\clearpage
\section{Degree-wise nuclear ranks, physical level counting and light-ray closure tests}
\label{p11:sec:v350-physical-readout-repairs}
The preceding Gram theorem and Paper~I-A boundary quotients yield
precise \emph{degree-wise} finite ranks. To identify physically cumulative
magic numbers, a further same-Hamiltonian nesting and ordering test is
required; this is not a new primitive nuclear projector or a new
compact momentum-truncation level.

\begin{proposition}[Degree-wise versus cumulative Hurwitz shell capacities]
\label{p11:prop:v350-nuclear-cumulative-obstruction}
The seven spaces $\mathcal N_r=\mathcal B_r/\mathcal L_r^\sigma$
are degree-$r$ quotients, with
\[
 (d_0,\ldots,d_6)=(1,4,10,14,25,41,63).
\]
Twofold occupation of each \emph{individual} image produces
$(2,8,20,28,50,82,126)$ capacity values. If the seven degree
spaces are instead orthogonal independent consecutive shells,
their cumulative twofold counts are
\[
 \boxed{(2,10,30,58,108,190,316).}
\]
Thus the rank and CAR arguments alone do not prove that the first
sequence is a list of cumulative physical shell closures.
\end{proposition}
\begin{proof}
Each $\mathcal N_r$ is homogeneous of degree $r$ by its
construction, not the filtered sum of all prior quotients.
CAR gives $2d_r$ individual spin occupations. For mutually
independent successive shells the filled count after degree $r$
is $2\sum_{s=0}^rd_s$, giving the displayed second sequence.
No physical level ordering follows from rank-nullity.
\end{proof}

\begin{theorem}[Exact conditional cumulative nuclear spectral certificate]
\label{p11:thm:v350-nuclear-physical-spectral-count}
Let $h^{(1)}_{\rm nuc}$ be the \emph{actual} gauge-reduced finite
one-body nuclear-response Hamiltonian on one common one-body
Hilbert space. Suppose there are increasing thresholds
$\lambda_0<\cdots<\lambda_6$ in genuine spectral gaps with
\begin{equation}
\label{p11:eq:v350-nuclear-spectral-count}
 \operatorname{rank}\mathbf1_{(-\infty,\lambda_r)}(h^{(1)}_{\rm nuc})
 =d_r,\qquad
 \operatorname{dist}(\lambda_r,\operatorname{spec}h^{(1)}_{\rm nuc})>0.
\end{equation}
If the native source maps intertwine these subspaces with the
selected coordinate quotients, and each orbital level has the
licensed twofold spin occupation without additional splitting,
the cumulative closed-shell particle numbers are exactly
$2d_r=(2,8,20,28,50,82,126)$.
The required incremental orbital degeneracies are
$(1,3,6,4,11,16,22)$. The theorem is a \emph{test} on the supplied
physical Hamiltonian, not an assertion that it has already passed it.
\end{theorem}
\begin{proof}
The spectral projections of one self-adjoint operator onto
nested intervals $(-\infty,\lambda_r)$ are nested. Their
ranks are $d_r$ by hypothesis; differences of consecutive
ranks give the listed orbital degeneracies. The level gaps
make the occupied subspaces spectrally identifiable, and
twofold CAR occupancy makes the cumulative capacities $2d_r$.
Compatibility with the geometric quotients is an independent
intertwining condition; dimensions alone never supply it.
\end{proof}

\begin{proposition}[Physical null sectors and light-ray Ward obligations]
\label{p11:prop:v350-light-ray-owner-qualification}
If the inherited longitudinal action satisfies
$[H_L,P_{H,L}]=0$, then for
$P^+_{\rm phys,L}=(H_L/c_{\rm op}+P_{H,L})/\sqrt2$,
\[
 [H_L,P^+_{\rm phys,L}]=0.
\]
This assertion does not consume the optional finite
$P^+_{\rm Fock,L}=P^+_{\rm phys,L}$ certificate.
The valence, charge, baryon and momentum PDF sum rules follow from
same-owner renormalized light-ray operators only when (i) the
Wilson-line/zero-mode distributional limits and Mellin exchanges
exist, (ii) the first local Mellin moments agree with the actual
flavor Noether currents, and (iii) the second moments, including
gluon mixing and contact terms, equal the same-action full
Belinfante $T^{++}$.
\end{proposition}
\begin{proof}
The commutator identity follows by linearity. Under the listed
light-ray certificates, the first two Mellin moments are the
matrix elements of the conserved charges and of the total null
momentum generator. Their normalized forward matrix elements
give the indicated sum rules. Conservation of $P^+$ alone
supplies none of the extra support or renormalization identities.
\end{proof}

\begin{theorem}[Correct leading-order second Mellin fixed point]
\label{p11:thm:v350-dglap-fixed-nf-correction}
On the imported leading-order $SU(3)$ singlet evolution with
\emph{fixed} active flavor number $n_f$ and
$Q+G=1$, the second-moment system is
\[
 \frac{d}{d\ln\mu^2}\binom QG
 =\frac{\alpha_s}{2\pi}
 \begin{pmatrix}-16/9&n_f/3\\16/9&-n_f/3\end{pmatrix}
 \binom QG.
\]
The zero eigenvector gives
\[
 \boxed{Q_*=\frac{3n_f}{16+3n_f},\qquad
 G_*=\frac{16}{16+3n_f}.}
\]
For $n_f=3$ this is $(Q_*,G_*)=(9/25,16/25)$;
the corresponding leading-order total angular momentum fractions
are $(J_q,J_g)=(9/50,8/25)$.
The ratio $9/17$ is the \emph{six}-flavor formal fixed point,
not an alternative three-flavor scheme. An infinite-scale limit
is a formal evolution fixed point when the physical UV
momentum ceiling remains finite.
\end{theorem}
\begin{proof}
The displayed matrix has zero column sums and eigenvalues
$0$ and $-(16+3n_f)/9$. The stationary equation is
$16Q=3n_fG$; together with $Q+G=1$ it gives the fractions.
For $n_f=3$ the nonzero eigenvalue is $-25/9$. At leading
order the total Ji angular momentum moments obey the same
mixing equation and sum to $1/2$, so their stationary values
are half the momentum fractions. Changes of active flavors
are separate threshold matching operations, not competing
three-flavor formulas.
\end{proof}
% END FRAMEWORK v2.350

% BEGIN FRAMEWORK v2.356: finite correlator and closed-cell Wilson readouts
\section{Quantitative spectral readout and closed gauge-surface diagnostics}
\label{p11:sec:v356-closed-gauge-readouts}

\begin{proposition}[Finite-time hadron source-overlap bound]
\label{p11:prop:v356-hadron-effective-mass-error}
For a positive finite spectral correlator with lowest \emph{source-visible}
energy $M_0$ and nonzero residue $w_0$, write
\[
 C(t)=w_0e^{-M_0t}+\sum_{n>0}w_ne^{-(M_0+\delta_n)t},
 \qquad w_n\ge0,\quad\delta_n>0.
\]
Put $\delta=\min_{w_n>0,n>0}\delta_n$,
$\delta_{\max}=\max_{w_n>0,n>0}\delta_n$, and
$W_+=\sum_{n>0}w_n$ (the empty-sum case is exact).
Then $m_{\rm eff}(t)=-\partial_t\log C(t)$ is nonincreasing,
$m_{\rm eff}(t)\ge M_0$, and
\[
 \boxed{0\le m_{\rm eff}(t)-M_0
       \le\delta_{\max}\frac{W_+}{w_0}e^{-\delta t}.}
\]
\end{proposition}
\begin{proof}
After cancelling the common $e^{-M_0t}$, the excess effective
mass is a positive weighted mean of the $\delta_n$.
Its derivative is minus the weighted variance of the full energy
list, and is therefore nonpositive.  Bound the numerator by
$\delta_{\max} W_+e^{-\delta t}$ and the denominator below by
$w_0$.  Without the actual source residue $w_0>0$, a lowest
Hamiltonian eigenvalue need not be the lowest correlator exponent.
\end{proof}

\begin{theorem}[Exact isolated-octahedron heat-model Wilson functions]
\label{p11:thm:v356-octa-heat-wilson}
Take one eight-triangle octahedral boundary of the fixed Hurwitz
complex, give each face the same central $SU(3)$ heat weight $K_t$,
and integrate its twelve links with normalized Haar measure.
This defines an \emph{isolated conditional gauge submodel}, for which
\[
 \boxed{Z_{\rm oct}(t)=K_{8t}(e)
      =\sum_R d_R^2e^{-4tC_R}.}
\]
For a normalized fundamental Wilson loop surrounding $k$ of its
faces, $1\le k\le7$, its expectation is
\[
 \boxed{W_k(t)=\frac{1}{3Z_{\rm oct}(t)}
 \sum_{R,S}d_Rd_S N_{\mathbf3,R}^{S}
 e^{-\frac t2(kC_R+(8-k)C_S)}.}
\]
Here $N_{\mathbf3,R}^{S}$ is the exact fusion multiplicity in
$\mathbf3\otimes R$.  The loop is gauge-invariant; neither its
area ratio nor the heat time is an experimentally calibrated string
tension or a hadron mass.
\end{theorem}
\begin{proof}
Heat-kernel character expansion and repeated Haar link integration
reduce the sphere partition to one face with area time $8t$.
Inserting a separating fundamental character splits the sphere
into two discs with times $kt$ and $(8-k)t$.
The common boundary integration projects the tensor product to
its trivial component; equivalently it applies the fusion
coefficient $N_{\mathbf3,R}^{S}$.  The dimensions and the
normalization by $3Z_{\rm oct}$ give the formula.
\end{proof}

\begin{proposition}[First full-shell character coefficient]
\label{p11:prop:v356-full-hurwitz-character-leading}
Impose the additional uniform heat-face specialization on all
$96$ triangles of the native Hurwitz boundary, and integrate all
$96$ edge variables.  For this \emph{conditional} finite gauge model,
\[
 Z_{96}(t)=\int_{SU(3)^{96}}\prod_{f=1}^{96}
   K_t(U_{\partial f})\prod_{e=1}^{96}dU_e,
\qquad
 \boxed{Z_{96}(t)=1+432e^{-16t/3}+O(e^{-6t})}\quad(t\to\infty).
\]
\end{proposition}
\begin{proof}
The trivial character sector gives $1$.  Any nontrivial term
surviving Haar edge projection has no edge incident to precisely
one represented face.  By the minimum-support theorem of
Paper~I-A, Theorem~\ref{p1:thm:v356-hurwitz-minimal-faces},
its first possible support has eight faces, and there are exactly
$24$ such octahedral boundary supports.  The least nonzero Casimir
is $C_{\mathbf3}=4/3$, with both fundamental and antifundamental
choices.  Haar projection on an eight-face sphere forces a common
representation.  Each minimal term therefore carries
$d_{\mathbf3}^2e^{-8tC_{\mathbf3}/2}=9e^{-16t/3}$,
so the coefficient is $24\cdot2\cdot9=432$.
The next possible nontrivial nine-face minimal-Casimir cost is at
least $9(4/3)t/2=6t$; higher irreducible Casimirs have larger
cost.  At positive $t$ the finite-complex character sum converges,
so the remaining large-$t$ terms give the stated asymptotic order.
\end{proof}

\begin{remark}[Numerical comparison and action-selection firewall]
\label{p11:rem:v356-closed-shell-status}
On the same \emph{declared} isotropic heat time
$t_{\rm iso}=5.753295846633\ldots$, the exact isolated-octahedron
series gives $W_1=0.021589880610665\ldots$ and
$W_2=0.000466123046099\ldots$.  A separate full-$96$-face
self-normalized importance-sampling diagnostic reported
$W_1\simeq0.02146286\pm0.00010491$ and
$W_2\simeq0.00044313\pm0.00004677$ from $200{,}000$
independent sampled configurations (approximate standard errors,
not rigorous confidence intervals).  The two calculations concern
different finite actions and their agreement within sampling
error is not an equality theorem.  Neither is a physical QCD
prediction: the action-preserving presentation bridge, interacting
same-action quark member, transfer/vacuum source and contact
responses remain separate conditions.
\end{remark}
% END FRAMEWORK v2.356


% BEGIN FRAMEWORK v2.357: connected baryon-pair shell-gap readout
\clearpage
\section{Two-neutron shell gaps as same-action baryon pair-transfer observables}
\label{p11:sec:v357-pair-shell-gap}
\begin{theorem}[Exact finite pair-transfer curvature readout]
\label{p11:thm:v357-pair-transfer-shell-gap}
Let $H_{Z,N}$ be the actual self-adjoint finite QCD-derived
nuclear-sector Hamiltonian specified by the same-action transfer,
and let $E_{Z,N}$ denote its lowest eigenvalue. For neighboring
allowed even-neutron sectors choose gauge-invariant pair sources
$Q_N^\dagger:\mathcal H_{Z,N}\to\mathcal H_{Z,N+2}$.
Write $|\Psi_{Z,N}\rangle$ for a normalized ground state and set
\[
 C_{Z,N}(t)=\langle\Psi_{Z,N}|
 Q_Ne^{-t(H_{Z,N+2}-E_{Z,N}I)}Q_N^\dagger
 |\Psi_{Z,N}\rangle,\qquad t>0.
\]
If $Q_N^\dagger|\Psi_{Z,N}\rangle$ has nonzero overlap with the
lowest destination eigenspace for both sources indexed by $N$ and
$N-2$, then
\[
 \boxed{\delta_{2n}(Z,N):=
 E_{Z,N+2}-2E_{Z,N}+E_{Z,N-2}
 =-\lim_{t\to\infty}\partial_t
 \log\frac{C_{Z,N}(t)}{C_{Z,N-2}(t)}.}
\]
For each source with a positive visible excitation gap $\Delta_M$
and positive ground residue $w_{0,M}$, the effective estimator
$\delta_{2n}^{\rm eff}(t):=-\partial_t\log
(C_{Z,N}(t)/C_{Z,N-2}(t))$ obeys
\[
 |\delta_{2n}^{\rm eff}(t)-\delta_{2n}|
 \le A_Ne^{-\Delta_Nt}+A_{N-2}e^{-\Delta_{N-2}t},\qquad
 A_M=\Delta_{\max,M}W_{+,M}/w_{0,M},
\]
where the constants are the actual spectral quantities of
Proposition~\ref{p11:prop:v356-hadron-effective-mass-error}.
\end{theorem}
\begin{proof}
Insertion of a complete orthonormal destination eigenbasis
gives $C_{Z,N}(t)=\sum_jw_{j,N}
 e^{-t(E^{(j)}_{Z,N+2}-E_{Z,N})}$ with $w_{j,N}\ge0$.
The nonzero overlap hypothesis makes its lowest exponent exactly
$E_{Z,N+2}-E_{Z,N}$. Subtracting the two effective exponents
produces the displayed second difference. Apply the finite-time
spectral bound of Proposition~\ref{p11:prop:v356-hadron-effective-mass-error}
to each correlator and use the triangle inequality. The energy
reference shift does not change excitation gaps or residue ratios.
\end{proof}
\begin{proposition}[Connected-energy curvature and common quark-source response]
\label{p11:prop:v357-shell-gap-curvature-source}
For any common spin-paired one-body reference energies $e_a$ with
$E^{(0)}_{Z,2K}=E_Z^{\rm fixed}+2\sum_{a=1}^Ke_a$, put
$U_Z(N)=E_{Z,N}-E^{(0)}_{Z,N}$. Then
\[
 \boxed{\delta_{2n}(Z,2K)=2(e_{K+1}-e_K)
 +\bigl[U_Z(2K+2)-2U_Z(2K)+U_Z(2K-2)\bigr].}
\]
Here $U_Z$ is an exact \emph{residual energy}; identifying its terms
with particular connected two- and higher-baryon cumulants consumes
the actual cluster/M\"obius decomposition of the native transfer.
For a differentiable common quark-mass source $m_q$ and the
sector sigma responses $\sigma_{q,Z,N}=m_q\partial_{m_q}E_{Z,N}$,
\[
 \boxed{m_q\,\partial_{m_q}\delta_{2n}(Z,N)
 =\sigma_{q,Z,N+2}-2\sigma_{q,Z,N}+\sigma_{q,Z,N-2}.}
\]
Any separately additive proton and neutron rest-mass terms cancel
from both energy curvatures. Neither identity asserts a numerical
QCD prediction or a realized magic-number peak.
\end{proposition}
\begin{proof}
Take the second difference in steps of two of the stated one-body
reference and of the definition of $U_Z$. In a differentiable
finite-eigenvalue branch, the Feynman--Hellmann identity gives the
sector sigma response; differentiation commutes with the finite
second difference. Affine particle-number contributions cancel.
\end{proof}
\begin{remark}[Existing physical certificates remain open]
\label{p11:rem:v357-nuclear-pair-scope}
The pair sources, their nonvanishing destination ground overlaps,
and the actual QCD-derived sector spectra require execution of the
existing hadronic/nuclear transfer. This proposition does not
generate the degree-three normal boundary kernel. In particular,
Theorem~\ref{p11:thm:v2211-nuclear-gram-kernel-certificate} and
Theorem~\ref{p11:thm:v350-nuclear-physical-spectral-count}
retain their conditional physical interpretation; neither a
comparison Hamiltonian nor geometric matching of ranks discharges
their antecedents.
\end{remark}
% END FRAMEWORK v2.357


% BEGIN FRAMEWORK v2.365: HADRON POLE AND COMPLETE DECAY READOUT CONTRACT
\section{Observable-complete hadron pole and weak-decay evaluation contract}
\label{p11:sec:v365-hurwitz-decay-test}
The geometric Hurwitz edge fan is an eligible finite source only
after the action-preserving presentation and physical fermion-source
requirements of Definition~\ref{p10:def:v365-apbi-hadron-presentation}
are met.  Neither an adjacency eigenvalue nor the common terminal
mass grammar alone identifies an observable hadron pole.

\begin{definition}[Source-visible hadronic pole and channel hyperedge]
\label{p11:def:v365-hadron-pole-channel}
For the actual same-action finite color-singlet $H_L$, normalized
preparation $|\Omega\rangle$ and a source with definite baryon
number, charge, spin and flavor, define
\[
G_O(z)=\langle\Omega|O^\dagger(z-H_L)^{-1}O|\Omega\rangle.
\]
A physical-candidate pole is a spectral pole of $G_O$ with nonzero
residue and the requisite detector/volume/continuum interpretation;
a source-absent eigenvalue is not a visible pole.  For parent
$h$ and daughter multiset $D=(d_1,\ldots,d_n)$, a channel is the
directed hyperedge $h\to D$, evaluated with the complete
same-action transition and outgoing projectors of
Definition~\ref{p5c:def:v365-decay-source-contract}.  Distinct
antiparticles, neutrino flavors, pions, photons and radiative
states are not silently omitted when claiming completeness.
\end{definition}

\begin{proposition}[APBI leakage is the first necessary pole test]
\label{p11:prop:v365-pole-leakage-test}
For a proposed normalized fan source map $\iota$, a root of
$\det(z-\iota^\dagger H_L\iota)$ need not be a pole of the
actual full resolvent.  The source-visible test must retain
$\Sigma(z)=B^\dagger(z-QH_LQ)^{-1}B$ with
$B=(I-\iota\iota^\dagger)H_L\iota$; if $B=0$ the compressed
subspace is reducing for Hermitian $H_L$ and the pole candidates
are exact in that decoupled sector, subject to nonzero source residues.
\end{proposition}
\begin{proof}
The $P\oplus Q$ block inverse identity gives the Schur return.
A nonzero off-diagonal $B$ modifies the compression's inverse
through the energy-dependent self-energy.  If $B=0$, Hermiticity
also removes the reverse block and the subspace is reducing.
The residue of $G_O$ at a full-space eigenvalue is the squared
projection of $O|\Omega\rangle$ onto the corresponding eigenspace,
which may still vanish.
\end{proof}


% BEGIN FRAMEWORK v3.025: ACTION-SENSITIVE SOURCE-GRAM SCOPE
\begin{remark}[Nonlinear face response is not a hadron energy]
\label{p11:rem:v3025-action-sensitive-gram-firewall}
In the aligned-spin, ordered-orbital three-face \emph{reference}
family of
Proposition~\ref{p10:prop:v3023-marked-baryon-source-overlap},
the two-root averaged baryon-source Gram is
\[
 G_{B,\lambda}=\begin{pmatrix}
 I_{56}&B_{t,\lambda}I_{56}\\
 B_{t,\lambda}I_{56}&I_{56}
 \end{pmatrix},\quad
 B_{t,\lambda}=a_{t,\lambda}^3.
\]
Its eigenvalues are $1+B_{t,\lambda}$ and $1-B_{t,\lambda}$,
each $56$-fold.  The nonzero derivative proved in
Proposition~\ref{p10:prop:v3025-threeface-action-sensitivity}
therefore changes reference source overlaps, \emph{not} hadron masses.
A physical pole requires the selected locally Gauss-invariant
$H_L(\lambda)$, the admitted source $J_L(\lambda)$ and its complete
Schur return of
Proposition~\ref{p11:prop:v365-pole-leakage-test}.
On differentiable isolated eigenbranches of a common finite
Hamiltonian and vacuum, Feynman--Hellmann gives
\[
 \frac{d}{d\lambda}(E_h-E_\Omega)
 =\langle h|\partial_\lambda H_L|h\rangle
 -\langle\Omega|\partial_\lambda H_L|\Omega\rangle.
\]
The two right-hand-side matrix elements are not determined by the
reference Wilson derivative.  Decay widths additionally require a
licensed open-channel/continuation readout.
\end{remark}
% END FRAMEWORK v3.025


\begin{remark}[Actual observational frontier]
\label{p11:rem:v365-observable-decay-frontier}
The necessary next calculations are (i) an action- and
source-preserving chamber/fan-to-$\mathcal H_L$ map, (ii) true
$O_M,O_B,O_G^{2++}$ source residues and C-even stress/Wilson
mixed correlators, and (iii) native $C_{ud}$ current contractions
between actual parent and daughter pole states with charge,
angular momentum, center and energy support.  The 35 valence
triples, 90 schematic flavor substitutions and a selected
43-channel comparator are not exhaustive predicted hadrons or
decays.  A geometric $2/5:3/5$ spin weighting is not a branching
fraction, and a finite Hermitian carrier without a detector or
controlled asymptotic regime does not yield irreversible widths.
\end{remark}
% END FRAMEWORK v2.365


% BEGIN FRAMEWORK v2.366: OBSERVABLE-READY BUT CONDITIONAL HURWITZ READOUTS
\section{Conditional Hurwitz-source angular and decay-amplitude diagnostics}
\label{p11:sec:v366-hurwitz-decay-diagnostics}
This section supplies \emph{tests of source hypotheses} for the
existing color-singlet Hamiltonian and its outgoing channel
projectors.  None of the formulae below represents a measured
branching fraction or a new hadron pole.

\begin{proposition}[Prepared Hurwitz spin-source angular intensities]
\label{p11:prop:v366-prepared-hurwitz-angular-law}
\emph{Assume} the source has an action-owned physical rotation
intertwiner into the spin-$2$ carrier, an aligned Hurwitz frame,
the pure $Q_{\chi}$ tensor and an unpolarized spin-$1/2$ factor.
Project the resulting density onto $J=3/2$ or $J=5/2$ as in
Proposition~\ref{p10:prop:v366-hurwitz-spin2-character-source}.
Assume a positive-parity strong $N\pi$ final state and the single
lowest allowed relative partial wave, $L=1$ for $3/2^+$ and $L=3$
for $5/2^+$.  After summing over the daughter nucleon spin,
its normalized angular intensities in the aligned source axes are
\[
 \boxed{W_{3/2}(\mathbf n)=\frac1{4\pi},\qquad
 W_{5/2}(\mathbf n)=
 \frac{1+5(n_x^4+n_y^4+n_z^4)}{16\pi}},\quad
 \frac{W_{5/2}(\mathrm{axis})}{W_{5/2}(\mathrm{diagonal})}
 =\frac94.
\]
The corresponding threshold powers are $p^{2L+1}=p^3,p^7$.
Orientation averaging of a randomly rotated production frame
makes both densities isotropic.  There is no claim that the actual
$N(1720)$, $N(1680)$ or any other resonance realizes this
special source preparation.
\end{proposition}
\begin{proof}
Form the normalized density
$\rho_J=P_J(|Q_{\chi}\rangle\langle Q_{\chi}|
 \otimes I_2/2)P_J/\operatorname{tr}(P_J\rho)$.
The $SU(2)$ Clebsch--Gordan formula and normalized spherical
harmonics give the stated expressions after summing the final
spin.  The $J=3/2$ density is isotropic; the $J=5/2$ density is
tetrahedrally invariant with quartic term as displayed.
The explicit symbolic spin-matrix and independent spherical-harmonic
checks are in \texttt{Proofs/Paper00-HadronChamberSplitCertificate-v1\_12.py}.
Rotational Haar averaging of a normalized density on the sphere
is the uniform density $1/(4\pi)$.
\end{proof}

\begin{proposition}[Actual pole/residue test, not geometric-count prediction]
\label{p11:prop:v366-hurwitz-decay-residue-gate}
Let the proposed Hurwitz-tagged source be admitted under
Definition~\ref{p10:def:v366-hurwitz-stress-baryon-interface} and
let $G_O(z)$ be the \emph{actual} color-singlet source resolvent of
Definition~\ref{p11:def:v365-hadron-pole-channel}.  A physical
Hurwitz influence on a specified hadron decay requires (i) nonzero
same-action source overlap, (ii) the full Feshbach/Schur return
when leakage is nonzero, (iii) correctly conserved quantum numbers
and physical $SU(2)$ amplitudes, and (iv) nonzero pole-to-channel
residues in the actual outgoing-state operators.  The phase-space,
centrifugal barrier and isospin factors must be removed before
comparing proposed reduced geometric amplitudes.
The rank-three cross-core pentagonal filter kills center-odd and
nontrivial triality coefficients; a physical transition in those
sectors therefore requires a different action-owned phase- or
orientation-sensitive operator.  Neither $600$, $120$, $25$ nor
$2/5:3/5$ by itself determines a number of decay channels or
branching fractions.
\end{proposition}
\begin{proof}
The full-space resolvent identity of
Proposition~\ref{p11:prop:v365-pole-leakage-test}
requires the Schur self-energy.  Conservation laws and the
completeness of the outgoing channel hyperedge enter the already
established physical decay contract.  The geometric kernel
selection rule follows from
Proposition~\ref{p1:prop:v366-cross-core-pentagonal-kernel},
but only an action-owned intertwiner can transport it to the
physical transition amplitude.
\end{proof}

\begin{remark}[No additional Standard Model or hadron identity]
\label{p11:rem:v366-decay-physical-frontier}
The meson $q\bar q$ and baryon $qqq$ \emph{unprojected}
spin-carrier decompositions in
Proposition~\ref{p1:prop:v366-twisted-quintet-spinorial-branch}
are not physical spectroscopy until color, flavor, Pauli and
source-ownership selection are applied.  The $N(1680)$/$N(1720)$
and $\Delta(1905)$/$\Delta(1920)$ pairs may motivate experiments,
but are not identified by the finite geometry.  No existing
Hamiltonian, mass relation, weak-decay count or physical partial
width has been modified in this release.
\end{remark}
% END FRAMEWORK v2.366


% BEGIN PAPER VII-B v1.158: 2I TWO/THREE CORE STRUCTURAL GROUND STATES
\section{Polyhedral two-core and three-core ground-source comparison}
\label{p11:sec:v1158-two-three-core-ground-structure}
This section integrates finite group/tensor results without changing
the physical spectral-pole, width, source and continuum definitions above.
A low-lying finite source representation is not automatically a
physical ground state, and an excited mode is not automatically a resonance.

\begin{theorem}[Directed two-core nonet candidate and its complement]
\label{p11:thm:v1158-directed-two-core-nonet}
For a supplied $2I$ packet with five tetrahedral cores and its
$20$ directed core pairs, the real pair module decomposes as
\[
 \boxed{\mathcal B_{20}=(\mathbf1\oplus\mathbf4\oplus\mathbf5)
              \oplus(\mathbf4\oplus\mathbf3\oplus\mathbf3').}
\]
The unoriented $A_5$-equivariant Petersen projector selects a
singlet and quintet, and its signed Hodge operator selects the
three-dimensional hand matching $\operatorname{Im}\mathbb H$.
Hence there is an exactly defined rank-nine comparison projector
$P_9$ onto $\mathbf1\oplus\mathbf3\oplus\mathbf5$ and an
orthogonal rank-eleven complement $\mathbf3'\oplus2\mathbf4$.
The projected nine is $A_5$-isometric to
$\operatorname{Herm}_3(\mathbb C)$ with its conjugation action
under the embedded real icosahedral triplet.
\end{theorem}
\begin{proof}
The directed-pair permutation character on the five core labels
has values $(20,0,2,0,0)$ on the five $A_5$ classes, yielding the
stated irreducible decomposition.  The unsigned Petersen
adjacency has eigenspaces $\mathbf1\oplus\mathbf5\oplus\mathbf4$;
the signed Petersen Hodge theorem supplies the two oriented
$\mathbf3,\mathbf3'$ eigenspaces.  The embedded momentum triplet
selects the matching Hodge character.  Finally
$\operatorname{Herm}_3(\mathbb C)=\mathbb RI\oplus
 i\mathfrak{so}(3)\oplus\operatorname{Sym}_0^2(\mathbb R^3)$,
with dimensions $1+3+5$ and the same $A_5$ action.
\end{proof}
\begin{remark}[Meson source and nonet firewall]
\label{p11:rem:v1158-meson-ground-pole-gate}
The exact $P_9$ is a mathematical comparison projector, not a
physical meson-ground-state or force projector.  Paper VII-A's
color-singlet meson source can be labelled by it formally, but
the native flavor action, its charge and mass spurion,
kinetic normalization and interacting correlator must select the
physical port.  The native conditional chiral $8$ directions and
its $3+4+1$ pion/kaon/neutral-octet source pattern remain distinct
from the ninth anomalous singlet, and neither a vector-meson mass
nor a radial resonance follows from the finite $A_5$ character.
\end{remark}

\begin{theorem}[Three-core cyclic selection and S-wave baryon source]
\label{p11:thm:v1158-three-core-ground-source}
Let $F=\mathbb C^3$ be the supplied light three-flavor source and
let $C$ cyclically permute the three positions in $F^{\otimes3}$.
Then
\[
 \boxed{Q_{\rm rel}=\tfrac13(2I-C-C^{-1})
       =P_{[21]}^{\rm flavor}},\qquad
 \boxed{F^{\otimes3}=\mathbf{10}\oplus
                   (\mathbf8\otimes\mathbf2)\oplus\mathbf1.}
\]
Thus $\operatorname{rank}Q_{\rm rel}=16$: it selects the two
mixed-positional flavor octets and annihilates the symmetric decuplet
and alternating singlet.  With the existing three-color antisymmetric
source and the minimal symmetric S-wave spin--flavor construction,
\[
 \boxed{\operatorname{Sym}^3(\mathbb C^2_{\rm spin}\otimes
          \mathbb C^3_{\rm flavor})
       =(\mathbf4,\mathbf{10})\oplus(\mathbf2,\mathbf8),
       \quad 56=40+16.}
\]
On this restricted source the cyclic relative-flavor projector
coincides with the spin-$\tfrac12$ octet projector, and the
symmetric spin-$\tfrac32$ decuplet is dark to this particular
relative junction.  These are representation/source identities,
not eigenvalue ordering or lifetime theorems.
\end{theorem}
\begin{proof}
Schur--Weyl decomposition under $SU(3)\times S_3$ yields the
$10\oplus(8\otimes2)\oplus1$ expression.  A three-cycle has
character $1$ on the symmetric and alternating $S_3$ lines,
and eigenvalues $\omega,\bar\omega$ on $[21]$, so
$(2I-C-C^{-1})/3$ is exactly its standard-sector projector.
The symmetric cube of the $2\times3$ spin--flavor product has
Schur partitions $[3]$ and $[21]$, giving $(4,10)$ and $(2,8)$;
the color epsilon tensor requires this symmetry in the minimal
spatial S-wave.  Both source blocks have the stated spin and
flavor eigenvalues.
\end{proof}

% BEGIN FRAMEWORK v3.069 -- verified five-core baryon comparison.
\begin{theorem}[Five-core pair and cubic chart for the 56-state baryon source]
\label{p11:thm:v369-five-core-baryon-chart}
Retain the \emph{existing} symmetric three-core source of
Theorem~\ref{p11:thm:v1158-three-core-ground-source} and restrict its
three-dimensional flavor representation to the selected real $A_5$
triplet in Theorem~\ref{p1b:thm:v369-five-core-cubic-intertwiner}.
On the $20$ ordered distinct core pairs, let $\mathcal R|a,b\rangle
=|b,a\rangle$ and let $P_{\mathbf3}^{-},P_{\mathbf3'}^{-}$ and
$P_{\mathbf5}^{+}$ denote the central character projectors for their
respective irreducible components.  Put
\[
 F_8=\operatorname{im}(P_{\mathbf3}^{-}+P_{\mathbf5}^{+})
 \cong\mathbf3\oplus\mathbf5,\quad
 F_{10}=\operatorname{im}(P_{\mathbf3}^{-}+P_{\mathbf3'}^{-})
 \oplus\operatorname{im}P_{\mathbf4}^{(\operatorname{Sym}^3V)}
 \cong\mathbf3\oplus\mathbf3'\oplus\mathbf4.
\]
The same odd $\mathbf3$ can label the two distinct \emph{spin-typed}
summands; they are not disjoint subspaces of the bare pair Hilbert
space.  There is an explicit $A_5\times\operatorname{Spin}(3)$
representation isometry
\[
 \boxed{W_{56}:\operatorname{Sym}^3(\mathbb C^2_{\rm spin}
 \otimes V_{\mathbb C})\ \xrightarrow{\ \simeq\ }
 (\mathbb C^2_{\rm spin}\otimes F_8)
 \oplus(\mathbb C^4_{\rm spin}\otimes F_{10}),\quad
 56=2\cdot8+4\cdot10,\quad W_{56}^\dagger W_{56}=I_{56}.}
\]
The map intertwines the independently specified target spin generators
and the $A_5$ flavor action.  Its highest-spin cubic-$\mathbf4$
component can be chosen to agree with $W_3$ up to an overall phase.
\end{theorem}
\begin{proof}
The symmetric-power/Cauchy decomposition gives source spin types
$(\mathbf2,\mathbf8)\oplus(\mathbf4,\mathbf{10})$.
Under the selected $A_5$ flavor restriction the octet is
$\mathbf3\oplus\mathbf5$, and the decuplet is
$\mathbf3\oplus\mathbf3'\oplus\mathbf4$.  The directed-pair
characters and Theorem~\ref{p1b:thm:v369-five-core-cubic-intertwiner}
provide exactly these target representations.  Match the pair and
cubic isotypic components by the $A_5$-averaged polar maps, and extend
over independent spin weights by standard spin lowering.  This
constructs a unitary $A_5\times\operatorname{Spin}(3)$ intertwiner.
Direct $56\times56$ construction checks both unitary normalization and
the $60$ effective $A_5$ matrices below $3\cdot10^{-14}$; independent
target spin generators, rather than only transported generators, obey
the full commutator test at floating-point precision.
\end{proof}

\begin{corollary}[Scoped fiber-capacity bound for a rank-five core port]
\label{p11:cor:v369-five-core-factorized-capacity}
If a five-core outgoing port of $A_5$ type $\mathbf1\oplus\mathbf4$
is required to host the entire $56$-state source through a
\emph{factorized} module $(\mathbf1\oplus\mathbf4)\otimes F$, then
$\dim F\ge18$.  The lower bound is attained at the representation
capacity level by $F=6\mathbf3$, since
\[
 (\mathbf1\oplus\mathbf4)\otimes6\mathbf3
 =6(\mathbf3\oplus\mathbf3'\oplus\mathbf4\oplus\mathbf5)
 \supset6\mathbf3\oplus4\mathbf3'\oplus4\mathbf4\oplus2\mathbf5.
\]
This is not a universal minimal dimension for a nonfactorized physical
Born/Clifford embedding.
\end{corollary}
\begin{proof}
The flavor-restricted source has six copies of $\mathbf3$.  For any
irreducible $A_5$ summand $R\subseteq F$, the tensor product
$(\mathbf1\oplus\mathbf4)\otimes R$ contains at most one
$\mathbf3$.  Thus at least six nontrivial irreducible summands are
required and their dimensions are at least three.  The explicit
$6\mathbf3$ tensor product shows sharpness within the stated class.
\end{proof}

\begin{proposition}[Finite orientation intermediary and direct-vertex obstruction]
\label{p11:prop:v369-oriented-core-auxiliary}
Let $E=\{(a,b):1\le a,b\le5,\ a\ne b\}$ and define isometries
$U_{\rm out},U_{\rm in}:\mathbb C^5\to\mathbb C^E$ by
\[
 U_{\rm out}|a\rangle={1\over2}\sum_{b\ne a}|a,b\rangle,
 \qquad U_{\rm in}|b\rangle={1\over2}\sum_{a\ne b}|a,b\rangle.
\]
Set $P=U_{\rm out}U_{\rm out}^\dagger$,
$Q=U_{\rm in}U_{\rm in}^\dagger$,
$K=[P,Q]$, $H_{\rm ori}=16iK/\sqrt{15}$, and
$P_{\rm ori}=-256K^2/15$.  Then
\[
 \boxed{K^3=-{15\over256}K,\qquad
 H_{\rm ori}^2=P_{\rm ori},\quad
 \operatorname{spec}H_{\rm ori}=\{-1^{(4)},0^{(12)},1^{(4)}\}.}
\]
Define $W_e=\sqrt{2/3}(U_{\rm out}+U_{\rm in})P_\perp$,
$W_o=\sqrt{2/5}(U_{\rm out}-U_{\rm in})P_\perp$, and
$J_e=W_eW_3$, $J_o=W_oW_3$, with $W_3$ as above.  They satisfy
\[
 \boxed{J_e^\dagger J_e=J_o^\dagger J_o=P_{\mathbf4}^{(\operatorname{Sym}^3V)},\quad
 J_e^\dagger J_o=0,\quad
 H_{\rm ori}J_e=-iJ_o,\quad H_{\rm ori}J_o=iJ_e.}
\]
For $z\not\in\{-1,0,1\}$ the full finite resolvent obeys
\[
 \boxed{J_e^\dagger(z-H_{\rm ori})^{-1}J_e
 ={z\over z^2-1}P_{\mathbf4}^{(\operatorname{Sym}^3V)},\quad
 J_o^\dagger(z-H_{\rm ori})^{-1}J_e
 =-{i\over z^2-1}P_{\mathbf4}^{(\operatorname{Sym}^3V)}.}
\]
But $H_{\rm ori}$ annihilates the $\mathbf3$, $\mathbf3'$ and
$\mathbf5$ pair components.  In particular, for the pair-octet and
pair-decuplet projectors
$P_8^{\rm pair}=P_{\mathbf3}^{-}+P_{\mathbf5}^{+}$ and
$P_6^{\rm pair}=P_{\mathbf3}^{-}+P_{\mathbf3'}^{-}$,
\[
 \boxed{P_8^{\rm pair}H_{\rm ori}P_8^{\rm pair}=0,\quad
 P_6^{\rm pair}H_{\rm ori}P_6^{\rm pair}=0,\quad
 P_8^{\rm pair}H_{\rm ori}P_6^{\rm pair}=0.}
\]
Thus this bare $20$-pair Hamiltonian cannot \emph{by itself} be the
physical octet--decuplet axial decay vertex.  It may be an auxiliary
orientation sector of a larger action-derived current.
\end{proposition}
\begin{proof}
The cross Gram is $U_{\rm out}^\dagger U_{\rm in}
=J_5/4-I_5/4=P_0-P_\perp/4$ with $P_0=J_5/5$.
Its principal-angle cosine on the augmentation is $1/4$.
The two-projector calculus gives $K^3=-15K/256$ and rank-eight
orientation support, and yields the stated even/odd quarter-turn.
The spectral projectors on $\operatorname{span}\{J_e,J_o\}$
give the rational resolvent contractions.  The same pair
permutation character decomposition shows that the orientation
support consists only of its two $\mathbf4$ summands, orthogonal to
the $\mathbf3$, $\mathbf3'$ and $\mathbf5$ components; hence the
three displayed zero blocks follow directly.  All formulas have an
independent $20$-pair finite-matrix numerical replay.
\end{proof}

\begin{corollary}[Axial spin--flavor strength benchmark in the new chart]
\label{p11:cor:v369-axial-transition-benchmark}
Let $G^{ia}=\sum_{r=1}^3(\sigma_r^i/2)(\lambda_r^a/2)$ be the
previously constructed $24$-component primitive one-body
spin--flavor operator on the symmetric $56$ source, and let
$P_8,P_{10}$ denote its physical spin-half octet and spin-three-half
decuplet \emph{representation} projectors.  Then the established
identity and its off-diagonal strength are
\[
 \boxed{\sum_{i=1}^3\sum_{a=1}^8G^{ia}G^{ia}
 ={19\over4}P_8+{7\over2}P_{10},\qquad
 \sum_{ia}\|P_8G^{ia}P_{10}\|_{\rm HS}^2=40.}
\]
Unitary transport through $W_{56}$ preserves these equalities.
This is a \emph{consistency benchmark}, not an independent derivation
of the axial current from the five-core Hamiltonian and not a decay
width or lifetime formula.
\end{corollary}
\begin{proof}
The symmetric three-slot Pauli and flavor Fierz identities give the
first equality.  Its octet-to-decuplet part has contraction
$\sum_{ia}P_8G^{ia}P_{10}G^{ia}P_8=(5/2)P_8$;
tracing gives $(5/2)\operatorname{rank}P_8=40$.
The identities survive unitary conjugation by $W_{56}$.
\end{proof}

% BEGIN v3.071 channel-resolved 56-state axial algebra (no physical-current promotion).
\begin{theorem}[Unique covariant decuplet--octet axial tensor and twelve exact strengths]
\label{p11:thm:v371-unique-decuplet-current-strengths}
Let $\mathcal H_{56}=\operatorname{Sym}^3(\mathbb C^2_{\rm spin}
\otimes\mathbb C^3_{\rm flavor})$ be the existing symmetric
three-quark source.  Its spin and exact flavor comparison group is
$SU(2)\times SU(3)$, with
\[
 \mathcal H_{56}\cong(\mathbf4,\mathbf{10})\oplus
 (\mathbf2,\mathbf8).
\]
For an axial-current \emph{representation slot} transforming as
$(\mathbf3,\mathbf8)$, the decuplet-to-octet intertwiner space has
multiplicity one:
\[
 \boxed{\dim\operatorname{Hom}_{SU(2)\times SU(3)}
 ((\mathbf4,\mathbf{10}),
 (\mathbf2,\mathbf8)\otimes(\mathbf3,\mathbf8))=1.}
\]
Thus every such exactly covariant current, including higher-body
candidates, has the same normalized decuplet-to-octet spin--flavor
matrix packet as the already existing $G^{ia}$, up to \emph{one}
common factor $g$ on this sector.  This statement presupposes the
full flavor-$SU(3)$ covariance of the \emph{current}, and is not a
consequence of its $A_5$ restriction alone.

For the normalized primitive packet
\[
 G^{ia}=\sum_{r=1}^3(\sigma^i_r/2)(\lambda^a_r/2),
 \qquad i=1,2,3,\quad a=1,\ldots,8,
\]
let $P_D$ and $P_B$ project onto the indicated entire decuplet and
octet spin--isospin multiplets within $\mathcal H_{56}$.  Define
\[
 S_{D\to B;\mathcal A}:=
 \frac1{\operatorname{rank}P_D}
 \sum_{i=1}^3\sum_{a\in\mathcal A}
 \|P_BG^{ia}P_D\|_{\rm HS}^2,
 \quad
 \mathcal A_\pi=\{1,2,3\},\ 
 \mathcal A_K=\{4,5,6,7\},\ 
 \mathcal A_\eta=\{8\}.
\]
Every nonzero strength is shown below; all omitted classes vanish:
\[
\begin{array}{c|c|c|r|c}
 D&B&\mathcal A&\sum_{ia}\|P_BG^{ia}P_D\|_{\rm HS}^2
             &S_{D\to B;\mathcal A}\\\hline
 \Delta&N&\pi&8&1/2\\
 \Delta&\Sigma&K&8&1/2\\
 \Sigma^*&N&K&2&1/6\\
 \Sigma^*&\Lambda&\pi&3&1/4\\
 \Sigma^*&\Sigma&\pi&2&1/6\\
 \Sigma^*&\Sigma&\eta_8&3&1/4\\
 \Sigma^*&\Xi&K&2&1/6\\
 \Xi^*&\Lambda&K&2&1/4\\
 \Xi^*&\Sigma&K&2&1/4\\
 \Xi^*&\Xi&\pi&2&1/4\\
 \Xi^*&\Xi&\eta_8&2&1/4\\
 \Omega&\Xi&K&4&1
\end{array}
\]
In particular, the four pion-class strength ratios are
\[
 \boxed{S_{\Delta N\pi}:S_{\Sigma^*\Lambda\pi}:
        S_{\Sigma^*\Sigma\pi}:S_{\Xi^*\Xi\pi}=6:3:2:3.}
\]
For an exactly covariant current with coefficient $g$, every entry
is multiplied by $|g|^2$ with no independent channel weight.
This is a \emph{complete representation-level current-strength
classification}, not an action-derived pion current or partial width.
\end{theorem}
\begin{proof}
In the exact product group,
$\mathbf2\otimes\mathbf3=\mathbf2\oplus\mathbf4$ and
$\mathbf8\otimes\mathbf8=\mathbf1\oplus\mathbf8_s
\oplus\mathbf8_a\oplus\mathbf{10}\oplus\overline{\mathbf{10}}
\oplus\mathbf{27}$. Each desired factor has multiplicity one,
so the Hom space has dimension one.  The primitive packet is a
nonzero intertwiner, as certified by its source total norm $40$;
it therefore spans this Hom space.

For exact evaluation use the six-mode bosonic polynomial realization
of the \emph{symmetric} cube, with occupation vectors
$n\in\mathbb N^6$, $\sum n_\alpha=3$, and generator action
$a^\dagger_\alpha a_\beta z^n=n_\beta
z^{n-e_\beta+e_\alpha}$.  There are exactly $56$ such
monomials, with Gram weights $\prod_\alpha n_\alpha!$.
Writing total spin Casimir $\mathbf S^2$, the exact orthogonal
projectors are
\[
 P_{10}=(\mathbf S^2-3I/4)/3,\qquad P_8=I-P_{10}.
\]
Let $Q_s$ be the commuting projector onto precisely $s$ strange
quarks.  Then $P_\Delta=Q_0P_{10}$,
$P_{\Sigma^*}=Q_1P_{10}$, $P_{\Xi^*}=Q_2P_{10}$,
$P_\Omega=Q_3P_{10}$, and $P_N=Q_0P_8$,
$P_\Xi=Q_2P_8$.  On $Q_1P_8$, the exact isospin
Casimir $\mathbf I^2$ has eigenvalues zero and two; hence
\[
 P_\Lambda=Q_1P_8(I-\mathbf I^2/2),\qquad
 P_\Sigma=Q_1P_8\mathbf I^2/2.
\]
Substituting these projectors and the eight Gell--Mann matrices
into the stated finite trace evaluates each displayed entry over
$\mathbb Q$; the rational-matrix evaluation is independently
reproducible by the accompanying exact verifier.  The projector
ranks are $(16,12,8,4)$ and $(4,2,6,4)$, respectively.  The
sum of the twelve total strengths is exactly $40$, matching the
already proved global Casimir identity and excluding unreported
nonzero classes.
\end{proof}

\begin{corollary}[Two-contraction exact test and strength-error bound for the occurring axial packet]
\label{p11:cor:v371-axial-current-gram-test}
Work on the same $56$-source with the same fixed normalization,
and let $J^{ia}:\mathcal H_{10}\to\mathcal H_8$ denote \emph{any}
independently evaluated actual action-current packet, when such a
packet exists.  Put $B^{ia}:=P_8G^{ia}P_{10}$ and define the
same-owner Hilbert--Schmidt contractions
\[
 A:=\sum_{ia}\operatorname{tr}[(B^{ia})^\dagger J^{ia}],
 \qquad N:=\sum_{ia}\operatorname{tr}[(J^{ia})^\dagger J^{ia}],
 \qquad g:=A/40.
\]
Then the \emph{complete} current-packet matching criterion is
\[
 \boxed{\sum_{ia}\|J^{ia}-gB^{ia}\|_{\rm HS}^2
 =N-\frac{|A|^2}{40}\ge0,\qquad
 J^{ia}=gB^{ia}\ (\forall i,a)
 \ \Longleftrightarrow\ N=\frac{|A|^2}{40}.}
\]
No twenty-four independent channel fits are required.  Moreover,
writing $\varepsilon:=N-|A|^2/40$ and $S_{D\to B;\mathcal A}(J)$
for the initial-rank-normalized squared norm of the corresponding
physical-current subpacket, one has the exact stable estimate
\[
 \boxed{\left|\sqrt{S_{D\to B;\mathcal A}(J)}-
 |g|\sqrt{S_{D\to B;\mathcal A}(G)}\right|
 \le\sqrt{\frac{\varepsilon}{\operatorname{rank}P_D}}.}
\]
The certificate is a \emph{test on an actual current supplied elsewhere},
not an evaluation of $J$, a proof of $SU(3)$ flavor symmetry, or a
width formula.
\end{corollary}
\begin{proof}
Treat the twenty-four blocks as one element of their direct-sum
Hilbert--Schmidt space.  The source identity
$\sum_{ia}\|B^{ia}\|_{\rm HS}^2=40$ makes $gB$ the orthogonal
projection of $J$ onto the one-dimensional complex span of $B$.
Pythagoras yields the displayed exact residual and equality test.
Every individual channel extraction is a contraction on that direct
sum (two orthogonal state projectors and a subset of the eight
flavor labels), hence its residual norm cannot exceed
$\sqrt\varepsilon$.  The reverse triangle inequality, divided by
$\sqrt{\operatorname{rank}P_D}$, gives the last bound.
\end{proof}

\begin{corollary}[Exact pseudoscalar two-body support and unique orbital partial wave]
\label{p11:cor:v371-strong-decuplet-threshold-wave}
For any same-owner spectral masses $M_D,M_B,m_\mathcal A>0$ supplied
for the above representation classes, a positive-energy on-shell
$D\to B+\phi_\mathcal A$ two-body final state is kinematically
accessible precisely when
\[
 M_D> M_B+m_\mathcal A,\qquad
 p_{DB\mathcal A}=\frac{\sqrt{\lambda(M_D^2,M_B^2,m_\mathcal A^2)}}{2M_D},
 \quad \lambda(x,y,z)=(x-y-z)^2-4yz.
\]
Threshold equality has vanishing two-body phase volume. The preceding
representation theorem supplies the exhaustive \emph{candidate} octet
pseudoscalar classes, and multiplication by the threshold indicator
$\mathbf1_{\{M_D>M_B+m_\mathcal A\}}$ supplies a unique
\emph{kinematically open candidate} list without fitting a channel weight.
For the familiar experimentally observed light-hadron mass ordering,
the open classes among the preceding twelve are exactly
\[
 \boxed{\Delta\to N\pi,\quad
 \Sigma^*\to\Lambda\pi,\quad
 \Sigma^*\to\Sigma\pi,\quad
 \Xi^*\to\Xi\pi.}
\]
In particular, $\Omega\to\Xi K$ is an algebraically supported
representation transition but closed at that mass ordering. These
mass inequalities are \emph{input checks}, not values derived in
this corollary.

If the occurring interaction conserves spatial rotations and parity,
the spin-parity labels $D(3/2^+)$, $B(1/2^+)$ and
$\phi(0^-)$ force the two-body final relative orbital momentum to
be \emph{exactly} $\ell=1$: angular momentum allows $\ell\in\{1,2\}$,
while parity permits odd $\ell$ only. For an analytic nonzero
partial-wave amplitude near threshold, the phase-volume and
centrifugal factors consequently give $\Gamma_{DB\mathcal A}
=O(p_{DB\mathcal A}^{3})$, with leading power $p^3$ whenever the
on-shell $P$-wave coefficient is nonzero. This fixes the threshold
power, \emph{not} the current coefficient, the momentum-dependent
form factors, the physical partial widths, or any lifetime.
\end{corollary}
\begin{proof}
In the parent rest frame conservation of four-momentum and positive
daughter energies give $M_D=\sqrt{M_B^2+p^2}+\sqrt{m_\mathcal A^2+p^2}$,
whose positive real solution exists only at or above the mass sum and
has the stated K\"all\'en expression. Crossing to a negative-energy
branch is not a physical two-body decay. Applying this inequality to
the twelve nonzero classes gives the stated list at the declared mass
ordering. The final intrinsic spin is $s=1/2$; coupling it to orbital
$\ell$ can produce $j=3/2$ only for $\ell=1,2$. The final parity is
$(+1)(-1)(-1)^\ell=(-1)^{\ell+1}$; matching the positive parent
parity excludes $\ell=2$. The centrifugal amplitude scales as $p^\ell$
under the stated analyticity assumption and two-body phase space as
$p$, proving the $p^{2\ell+1}=p^3$ leading law when its coefficient
is nonzero.
\end{proof}

\begin{corollary}[Fixed primitive octet axial $F/D$ ratio]
\label{p11:cor:v371-primitive-octet-FD}
On the \emph{same primitive one-body packet} $G^{ia}$, the
usual $SU(3)$ octet reduced-coupling convention of
Corollary~\ref{p11:cor:native-octet-axial-current-block}
gives
\[
 \boxed{F=\frac23,\qquad D=1,\qquad \frac FD=\frac23,
 \qquad F+D=\frac53.}
\]
These are parameter-free $56$-source representation benchmarks.
The preceding multiplicity-one argument does \emph{not} extend to the
octet-to-octet block: there are two independent $\mathbf8$ copies in
$\mathbf8\otimes\mathbf8$.  Thus a general physical current need
not have the primitive $F/D$ value.
\end{corollary}
\begin{proof}
The exact rank-one spin-up highest-weight projectors onto the
$N$ proton and $\Sigma^+$ source vectors, evaluated against
$4G^{3,3}=\sum_r\sigma_r^3\lambda_r^3$, have diagonal
matrix elements $5/3$ and $4/3$, respectively.  In the standard
octet convention these equal $F+D$ and $2F$.  Hence
$F=2/3$ and $D=1$.  Both matrix elements follow exactly from
the occupation/projector calculus in the preceding theorem, with
no phenomenological input.
\end{proof}

\begin{remark}[Exact symmetry versus occurring decay current]
\label{p11:rem:v371-axial-strength-scope}
The twelve source strengths do not assert that the actual $A_5$
carrier or QCD action enjoys unbroken flavor $SU(3)$, that its
pseudoscalar physical meson occupies the declared octet current
slot, or that its \emph{action-derived} axial current equals this
operator.  Real flavor breaking, phase space, pole residues,
finite-volume/continuum limits and absorptive open channels remain
separate source-owned obligations.  No physical branching ratio,
absolute width or lifetime is derived by the purely finite tensor.
\end{remark}
% END v3.071 channel-resolved 56-state axial algebra.
% BEGIN v3.072 aligned E36 mass/current and transition-status firewall.
\begin{proposition}[Zeroth-moment obstruction to identifying dressed axial and mass ports]
\label{p11:prop:v372-internal-external-axial-gram-no-go}
On the \emph{specified} symmetric $56$ spin--flavor source, the
primitive packet of Theorem~\ref{p11:thm:v371-unique-decuplet-current-strengths}
has the exact quadratic Gram
\[
 \boxed{\sum_{ia}G^{ia}G^{ia}=\frac{19}{4}P_8+\frac72P_{10}.}
\]
Define, for any real $t$, the octet-diagonally dressed packet
$J_t^{ia}=G^{ia}+(t-1)P_8G^{ia}P_8$. Its inter-multiplet blocks are
unchanged but its zeroth moment is
\[
 \boxed{\sum_{ia}J_t^{ia}J_t^{ia}
 =\left(\frac52+\frac94t^2\right)P_8+\frac72P_{10}.}
\]
In the \emph{conditional} common-$E_{36}$ scalar response
$H[C]=\tfrac{22}{5}E_{36}I-\tfrac{16}{15}E_{36}C$,
substituting the bare packet gives a decuplet--octet eigenvalue
gap $4E_{36}/3$, but substituting $C=\sum J_t^2$ gives instead
\[
 \boxed{\Delta_{10-8}(J_t)
 =\frac{16}{15}E_{36}\left(\frac94t^2-1\right),\quad
 \Delta_{10-8}(J_{3/4})=\frac{17}{60}E_{36}.}
\]
The latter is only $17/80$ of the bare gap. Therefore a dressed
physical axial-current \emph{candidate} cannot simply replace the
internal mass port at a common unmodified Schur denominator. A single
physical parent may generate both as different field variations,
but no occurring physical mass or current is established by this
finite comparison or by the choice of $E_{36}$ coefficient.
\end{proposition}
\begin{proof}
The primitive operator has octet-to-octet contribution
$\sum_{ia}(P_8G^{ia}P_8)^2=(9/4)P_8$ and cross-octet return
$\sum_{ia}P_8G^{ia}P_{10}G^{ia}P_8=(5/2)P_8$.
The decuplet return is $(7/2)P_{10}$. Dressing only the octet
diagonal multiplies its squared contribution by $t^2$, leaving
both inter-multiplet blocks intact. Substitution into the stated
scalar mass functional yields the displayed gaps by subtraction.
\end{proof}

\begin{corollary}[Dressed octet axial benchmarks and invariant inter-multiplet budget]
\label{p11:cor:v372-octet-dressing-moment-firewall}
The explicit comparison packet with \emph{only} its octet block
modified to
\[
 P_8J_{\rm loc}^{ia}P_8
 =\frac45P_8G^{ia}P_8-\frac1{15}P_8S^iT^aP_8,
 \qquad P_8J_{\rm loc}^{ia}P_{10}=P_8G^{ia}P_{10}
\]
(and Hermitian conjugate, with unchanged decuplet block)
satisfies the exact representation identities
\[
 \boxed{F_{\rm loc}=\frac7{15},\qquad D_{\rm loc}=\frac45,
 \qquad \frac{F_{\rm loc}}{D_{\rm loc}}=\frac7{12},}
\]
\[
 \boxed{\sum_{ia}J_{\rm loc}^{ia}J_{\rm loc}^{ia}
 =\frac{379}{100}P_8+\frac72P_{10},\qquad
 \Delta_{10-8}(J_{\rm loc})=\frac{116}{375}E_{36}.}
\]
Every inter-multiplet strength in
Theorem~\ref{p11:thm:v371-unique-decuplet-current-strengths}
is unchanged by this octet-only modification. Moreover, if
$D^{ia}=P_8D^{ia}P_8$ is \emph{any} exactly
$SU(2)\times SU(3)$-covariant spin-vector/flavor-octet packet,
$J^{ia}=G^{ia}+D^{ia}$, and $H=h_8P_8+h_{10}P_{10}$, then
\[
 \boxed{\frac12\sum_{ia}[J^{ia},[H,J^{ia}]]
 =\frac12\sum_{ia}[G^{ia},[H,G^{ia}]].}
\]
Thus the cross-multiplet energy-weighted first moment can agree
while the diagonal zeroth moment and octet weak-current ratios differ.
The displayed $F,D$ and $E_{36}$ values are benchmarks of the
\emph{specified} finite operators, not a derivation of the
physically occurring weak current or mass coefficient.
\end{corollary}
\begin{proof}
The symmetric source identities
$\sum(P_8GP_8)^2=(9/4)P_8$,
$\sum(P_8GP_8)(P_8STP_8)=(3/2)P_8$ and
$\sum(P_8STP_8)^2=(9/4)P_8$ give the octet-diagonal
square $129P_8/100$. The unchanged cross return contributes
$5P_8/2$, yielding $379P_8/100$; the decuplet square remains
$7P_{10}/2$. The octet $SU(3)$ reduced-coupling rules applied
to the displayed tensor yield the stated $F,D$ values.
For the first-moment identity, $[H,D^{ia}]=0$, so the
Jacobi identity reduces the difference of the double-commutator
packets to $\tfrac12[H,\sum_{ia}[D^{ia},G^{ia}]]$.
The contracted sum transforms as a scalar of $SU(2)\times SU(3)$,
so by the multiplicity-one decomposition
$(\mathbf2,\mathbf8)\oplus(\mathbf4,\mathbf{10})$ it is a
linear combination of $P_8,P_{10}$ and commutes with $H$.
\end{proof}

\begin{remark}[Comparison of exact symmetry and on-shell spectral trials]
\label{p11:rem:v372-mass-decay-trial-scope}
The multiplicity-one decuplet-to-octet tensor is unique under the
\emph{full} comparison $SU(2)\times SU(3)$, not under the
five-core $A_5$ alone. The finite five-core pair-plus-cubic chart
has unfixed irrep phases until a source-selected physical embedding
exists. At a \emph{common} spectral argument, full-symmetry
covariance selects one reduced tensor; observed-mass on-shell
channels can nevertheless probe different arguments of an
\emph{independently derived} common form factor. The exploratory
map $z_{B\to B'\pi}=2-E_\pi/(4E_{36})$ and its directed
Born-resolvent ratio $R_{\Sigma^*}\simeq7.493915$, versus the
recoil-corrected constant-current trial $6.9040865$, are
\emph{post-hoc conditional comparisons}, not same-action predictions.
The common-width trials also use external phase-space masses,
$F_\pi$ and an unproved physical endpoint normalization. A
finite Hermitian orientation propagator has no standalone
absorptive lifetime; the actual source must supply a retarded
open-channel continuation, its pole residues and chiral Ward
completion before any channel-dependent width is claimed.
\end{remark}
% END v3.072 aligned E36 mass/current and transition-status firewall.


% BEGIN v3.074 -- sharp positive-connector contact bound on the axial 56.
\begin{theorem}[Unique least affine positive connector and its sharp contact bound]
\label{p11:thm:v374-axial-minimal-positive-contact}
Let $B_0:\mathbb C^n\to\mathbb C^m$ be injective,
$K=B_0^\dagger B_0>0$, and let $J_a=J_a^\dagger$ be specified
field derivatives of a positive response
$\Sigma(\Phi)=B(\Phi)^\dagger B(\Phi)$ with an affine connector
$B(\Phi)=B_0+\sum_a\Phi_aX_a$.  Among \emph{all} (including
rectangular and non-Hermitian) $X_a$ satisfying
\[
 B_0^\dagger X_a+X_a^\dagger B_0=J_a,
\]
the unique Hilbert--Schmidt least connector is
\[
 \boxed{X_a^{\min}=B_0\Lambda_a,\qquad
 K\Lambda_a+\Lambda_a K=J_a,\quad\Lambda_a=\Lambda_a^\dagger.}
\]
Every other admissible $X_a=X_a^{\min}+N_a$ obeys the exact
orthogonal decomposition
\[
 \boxed{\|X_a\|_{\rm HS}^2=\|X_a^{\min}\|_{\rm HS}^2
       +\|N_a\|_{\rm HS}^2.}
\]
Writing $k_r>0$ for the eigenvalues of $K$, the quadratic diagonal
response $C_{aa}=\partial_{\Phi_a}^2\Sigma(0)=2X_a^\dagger X_a$
has the sharp, auxiliary-dimension-independent trace inequality
\[
 \boxed{\operatorname{tr}C_{aa}\ge
 \sum_{r,s}\frac{|(J_a)_{rs}|^2}{k_r+k_s},}
\]
with equality exactly at $X_a^{\min}$. The complete two-field
contact of the \emph{least affine construction} is fixed by
$C_{ab}^{\min}=\Lambda_aK\Lambda_b+\Lambda_bK\Lambda_a$.
This is a mathematical minimization within the declared positive
affine class, \emph{not} an asserted new GP3 axiom or selection of
the occurring physical QCD port.
\end{theorem}
\begin{proof}
Since $K>0$, the Hermitian Sylvester equation has the unique
solution $(\Lambda_a)_{rs}=(J_a)_{rs}/(k_r+k_s)$.
The proposed $X_a^{\min}$ satisfies the first-derivative constraint.
For any $N_a$ preserving that constraint,
$B_0^\dagger N_a+N_a^\dagger B_0=0$, and hence
$\operatorname{Re}\operatorname{tr}[(X_a^{\min})^\dagger N_a]
=\operatorname{Re}\operatorname{tr}(\Lambda_a B_0^\dagger N_a)=0$.
Pythagoras proves the unique minimum. Further,
$2\|X_a^{\min}\|_{\rm HS}^2
=\operatorname{tr}(\Lambda_aJ_a)
=\sum_{r,s}|(J_a)_{rs}|^2/(k_r+k_s)$.
The affine product rule gives the stated diagonal and mixed contacts.
\end{proof}

\begin{corollary}[Exact 56-carrier current budgets, contact bound and affine no-go]
\label{p11:cor:v374-native-56-contact-budget-and-nonuniqueness}
For the \emph{existing} Hermitian primitive packet $G^{ia}$ of
Corollary~\ref{p11:cor:v369-axial-transition-benchmark}, the four
full spin/flavor Hilbert--Schmidt budgets are exactly
\[
 \boxed{\sum_{ia}\|P_8G^{ia}P_8\|_{\rm HS}^2=36,\quad
 \sum_{ia}\|P_8G^{ia}P_{10}\|_{\rm HS}^2=40,}
\]
\[
 \boxed{\sum_{ia}\|P_{10}G^{ia}P_8\|_{\rm HS}^2=40,\quad
 \sum_{ia}\|P_{10}G^{ia}P_{10}\|_{\rm HS}^2=100.}
\]
For an $SU(2)\times SU(3)$-scalar positive comparison response
$K=k_8P_8+k_{10}P_{10}$ with $k_8,k_{10}>0$ and an assigned
primitive first current $J^{ia}=\alpha G^{ia}$ ($\alpha\in\mathbb R$),
the sharp summed affine-contact trace bound is
\[
 \boxed{\sum_{ia}\operatorname{tr}C_{ia,ia}\ge
 \alpha^2\left(\frac{18}{k_8}+\frac{50}{k_{10}}
 +\frac{80}{k_8+k_{10}}\right).}
\]
The earlier \emph{Hermitian} square-root Lyapunov connector
$X_{rs}^{\rm H}=(J)_{rs}/(\sqrt{k_r}+\sqrt{k_s})$ obeys the
same first derivative but has summed contact
\[
 \boxed{\alpha^2\left(\frac{18}{k_8}+\frac{50}{k_{10}}
 +\frac{160}{(\sqrt{k_8}+\sqrt{k_{10}})^2}\right),}
\]
which exceeds the least value by exactly
\[
 \boxed{\frac{80\alpha^2(\sqrt{k_8}-\sqrt{k_{10}})^2}
 {(k_8+k_{10})(\sqrt{k_8}+\sqrt{k_{10}})^2}.}
\]
Even \emph{positivity, affinity, full spin/flavor covariance, the
same $K$, and the same entire $J$ do not by themselves select the
contact}: for $B_0=K^{1/2}$, all real $\beta$ give the admissible
covariant family
\[
 \boxed{X_a(\beta)=X_a^{\min}+i\beta K^{-1/2}J_a,}
\]
with identical first derivative and an additional summed contact trace
$2\beta^2\alpha^2(76/k_8+140/k_{10})$. The minimal solution is
uniquely selected \emph{if} least Hilbert--Schmidt connector norm is
adopted; the actual action must decide whether that criterion applies.
This conditional comparison does not determine physical two-pion
contacts, chiral seagulls, $F_\pi$, partial widths or lifetimes.
\end{corollary}
\begin{proof}
The source Casimir identity gives
$\operatorname{tr}_{P_8}\sum G^2=16(19/4)=76$ and
$\operatorname{tr}_{P_{10}}\sum G^2=40(7/2)=140$.
The known cross budget is $40$ in each direction. Subtraction gives
the diagonal budgets $76-40=36$ and $140-40=100$.
Insert the four blocks into the preceding sharp trace formula;
the cross directions contribute $80/(k_8+k_{10})$.
For the Hermitian Lyapunov lift, the diagonal contributions agree,
while the cross contributions total
$160/(\sqrt{k_8}+\sqrt{k_{10}})^2$. Their difference is the
nonnegative displayed expression. Finally
$B_0^\dagger(i\beta K^{-1/2}J_a)=i\beta J_a$ and its adjoint
is $-i\beta J_a$, so the current is unchanged. The null direction
is Hilbert--Schmidt orthogonal to the minimum; its summed extra
squared norm uses $\sum\operatorname{tr}(J_aK^{-1}J_a)
=\alpha^2(76/k_8+140/k_{10})$. Covariance is preserved because
$K$ is a scalar under the given comparison group.
\end{proof}
% END v3.074 sharp positive-connector contact bound on the axial 56.


% BEGIN FULL-CARRIER MESON--BARYON SOURCE CROSSWALK v1.1.
\begin{proposition}[Geometric flavor source map and symmetric-56 current crosswalk]
\label{p11:prop:fullcarrier-geometric-flavor-source-crosswalk}
On the finite representation source, let $F_A$ be the eight
standard-normalized traceless Hermitian matrices generated by the
single dodecahedral one-hop lift in
Theorem~\ref{p1b:thm:fullcarrier-protected-local-flavor}.
Their primitive symmetric-three-constituent spin--flavor current is
$G^{iA}=\sum_{r=1}^3(\sigma_i/2)_r(F_A)_r$.
The established octet/decuplet identity of
Corollary~\ref{p11:cor:v369-axial-transition-benchmark} decomposes
under the geometric $A_5$ triplet/quintet as
\[
 \boxed{\sum_{i,A\in\mathbf3}\|P_8G^{iA}P_{10}\|_{\rm HS}^2=15,
 \qquad\sum_{i,A\in\mathbf5}\|P_8G^{iA}P_{10}\|_{\rm HS}^2=25,
 \qquad 15+25=40.}
\]
The constant, linear and traceless-quadratic dodecahedral vertex
harmonics give a normalized real isometry
$\mathcal Y:\operatorname{Herm}(3)\to
\mathcal N_9=\mathbf1\oplus\mathbf3\oplus\mathbf5$.
For $H=cI+T(v)+Q_0$ (traceless real symmetric $Q_0$), it is
\[
 \mathcal Y(H)_j=\sqrt{3/20}\,c
 +\sqrt{2/5}\,u_j\cdot v
 +\sqrt{2/3}\,u_j^{\mathsf T}Q_0u_j,
 \qquad\|\mathcal Y(H)\|^2=\operatorname{tr}(H^2),
\]
with $T(v)$ the normalized Hermitian rotation matrix of the
vector $v$.  For orthonormal traceless sources $\xi_A$ and their
corresponding primitive current maps $\mathcal M_i$,
\[
 \boxed{\sum_i\langle P_{10}\mathcal M_i(\xi_A)P_8,
 P_{10}\mathcal M_i(\xi_B)P_8\rangle_{\rm HS}=5\delta_{AB}.}
\]
These are \emph{reference representation} matrix strengths; the
geometric $A_5$ triplet is not automatically the physical isospin
triplet.  A physically chosen light-flavor plane determines its
orientation invariant $\kappa=|n^{\mathsf T}n|^2$, not independent
partial-width weights.  No physical hadron current, pion vertex,
branching fraction, pole or lifetime is inferred without the
same-action gate of Definition~\ref{p10:def:fullcarrier-current-action-comparison}.
\end{proposition}
\begin{proof}
The dodecahedral second and fourth moments give the stated isometry.
The symmetric three-slot spin/flavor Fierz contractions yield the
triplet and quintet strengths $15,25$ and the eight-source Gram
$5I_8$; the integer and quadratic-number projector identities
are recorded in the native finite geometry certificates.  The
source-to-physical implication is explicitly conditional.
\end{proof}
% END FULL-CARRIER MESON--BARYON SOURCE CROSSWALK v1.1.

\begin{remark}[Status of decay support versus physical action]
\label{p11:rem:v369-five-core-decay-scope}
The new $W_3$, $J_{\rm geom}$, $W_{56}$ and $H_{\rm ori}$ establish
finite geometric and representation-level support, not the unique
field-dependent physical baryon current.  The relative reduced
couplings are not fixed by $A_5$ alone: in particular,
$\dim\operatorname{Hom}_{A_5}(\mathbf{10},\mathbf8\otimes\mathbf8)=11$
and $\dim\operatorname{Hom}_{A_5}(\mathbf{10},\mathbf8\otimes\mathbf3)=4$.
A geometric $\mathbf3$ is not automatically the marked-axis physical
pion isotriplet.  Five independent irrep phases remain before further
same-action matching.  The bare Hermitian finite resolvent cannot by
itself produce an irreversible open-channel lifetime, a CP-violating
observable, or a physical hadron pole.
\end{remark}
% END FRAMEWORK v3.069 -- verified five-core baryon comparison.

\begin{proposition}[Conditional antisymmetric light-pair stabilization]
\label{p11:prop:v1158-light-pair-conditional-ground-order}
Let $L_i$ select the already typed $u,d$ $E_{40}$ source on
constituent slot $i$ and let $\Pi_{ij}^{(f)}$ exchange flavor slots.
The positive comparison operator
\[
 A_{\ell\ell}=\sum_{i<j}L_iL_j(I-\Pi_{ij}^{(f)})
\]
has, on the above 56-dimensional symmetric spin--flavor source,
\[
 \boxed{\operatorname{spec}A_{\ell\ell}
        =\{3^{(4)},\;2^{(2)},\;0^{(50)}\}.}
\]
Its values on the octet families are $N:3$, $\Lambda:2$,
$\Sigma:0$, $\Xi:0$; the decuplet is dark.  If and only if a
same-action, kinetically normalized connector return produces
$-gA_{\ell\ell}$ with $g>0$ and a positive first-order
strange-source increment $b>0$, the conditional octet masses obey
\[
 M_N=m_0-3g,\quad M_\Lambda=m_0+b-2g,\quad
 M_\Sigma=m_0+b,\quad M_\Xi=m_0+2b,
\]
so $M_N<M_\Lambda<M_\Sigma<M_\Xi$ when $b>g>0$, and
$2(M_N+M_\Xi)=3M_\Lambda+M_\Sigma$ in that comparison model.
The exchange-even pair channel is also symmetry-allowed; the actual
action must select its relative coefficient.  No positive $g$ or
observed mass is inferred merely from $A_{\ell\ell}$.
\end{proposition}
\begin{proof}
On a three-position flavor space, each factor
$I-\Pi_{ij}^{(f)}$ is twice the light antisymmetric-pair
projector.  The mixed-symmetry nucleon source gives the eigenvalue
three, the $ud$ isosinglet in $\Lambda$ gives two,
and other minimal S-wave families give zero.  The spin degeneracies
are four and two, respectively, and the other fifty states are
in the kernel.  Substitution in the displayed conditional mass
operator proves both the ordering when $b>g>0$ and the relation.
\end{proof}

\begin{remark}[Full-action odd/even connector requirement]
\label{p11:rem:v1158-connector-action-frontier}
The color-adjoint relative three-face Gram of
the color-current Gram proposition in Paper~VII-A
and the native heat Hessian share the projector
$I_3-J_3/3$.  The full action still has to produce the
$E_{40}$-gated, color/spin/flavor-compatible mixed source.  One
finite \emph{comparison} odd-channel return is
$A_{\ell\ell}+A_{\ell\ell}^{2}/6$, with eigenvalues
$9/2$ on the nucleon, $8/3$ on $\Lambda$ and zero on the others;
its physical vertex and coefficient are not established.
The octet may be stabilized virtually by a lower eigenbranch,
whereas other allowed states can be excited or resonant with
threshold-dependent widths.  A light-pair count or large geometric
source norm alone is not a binding energy.
\end{remark}
% END PAPER VII-B v1.158


% BEGIN FRAMEWORK v3.023: MOBILE 2I-SOURCE READOUT AND POLE GUARD
\section{Reference $2I$ source embeddings and conditional hadron pole readout}
\label{p11:sec:v3023-mobile-2i-hadron-readout}
Paper VII-A's marked-Hurwitz-fan CAR sources use the \emph{existing}
color singlets and inherited two-core nine-component, respectively
three-core $56$-component, source labels. They are reference-sector
operator embeddings, not $9$ or $56$ new elementary matter fields.
In particular the $56$-dimensional source still decomposes as
\[
\operatorname{Sym}^3(\mathbb C^2\otimes\mathbb C^3)
 \cong(\mathbf4,\mathbf{10})\oplus(\mathbf2,\mathbf8)
\]
on the existing minimal symmetric spin--flavor domain.

\begin{proposition}[Conditional same-action mobile-source pole interface]
\label{p11:prop:v3023-mobile-2i-pole-interface}
Suppose the marked, link-dressed baryon/meson sources of Paper VII-A
are admitted into the \emph{same} finite gauge-invariant physical
Hilbert owner and measure as the existing QCD Hamiltonian $H_L$,
with the required spin, orbital, charge and source pullbacks. Let
$|J_{v,e}(\chi)\rangle=O_{v,e}^\dagger(\chi)|\Omega\rangle$.
Then their physical two-root response is the full resolvent
\[
\boxed{G_e^{uv}(z;\eta,\chi)=
\langle J_{u,e}(\eta)|\,(z-H_L)^{-1}\,
|J_{v,e}(\chi)\rangle.}
\]
For any source projection $P$ and $Q=I-P$, its exact compressed
resolvent includes the leakage self-energy
\[
\boxed{P(z-H_L)^{-1}P=
 [z-PH_LP-PH_LQ(z-QH_LQ)^{-1}QH_LP]^{-1}}
\]
whenever the indicated inverses exist. A physical hadron-pole claim
requires a nonzero residue in the actual source channel. The fixed-link
reference Gram $W_eI_{56}$, or its spin-transport correction, does
not establish such a pole and does not determine a binding energy,
physical mass or decay width.
\end{proposition}
\begin{proof}
The first formula is the definition of the matched source-visible
resolvent on the assumed common owner. The second follows by the
block Schur-complement formula for $z-H_L$ and retains all
source-to-complement couplings. The spectral theorem expresses the
source residue at an isolated eigenvalue as the norm of the
corresponding spectral projection on the source; that norm may
vanish independently of the free-CAR reference Gram.
\end{proof}

\begin{remark}[RG and hadronic observables]
\label{p11:rem:v3023-beta-hadron-source-firewall}
The one-loop $SU(3)_c$ curvature beta coefficient and the
quark-singlet color factors use the same generated color
representation, but their equality of underlying Casimir data
is not a hadron-mass prediction. Exact same-action RG blocking of a
retained source, with induced multiface operators and source contacts,
may preserve its full correlator while the effective gauge coefficients
run. No identification between compact-group heat time, $\ln\mu$,
physical hierarchy energy and a hadron pole is made here.
\end{remark}
% END FRAMEWORK v3.023


% BEGIN BERGER READOUT GATE v1.17
\section{Source-induced shape effects in hadronic and transition readouts}
\label{bh:berger:v117:p7b:observable-gate}
\begin{remark}[Physical source first, spectral observation second]
\label{bh:berger:v117:p7b:spectral-gate}
A matter-induced homogeneous Milnor/axisymmetric Berger deformation
can enter the already established finite QCD pole and current readouts
only through a single actual metric-coupled and locally Gauss-reduced
Hamiltonian $H_L(B)$, its generated occupied state, and its full
source/measure/contact variations. The five-component source Hessian
and the current spectral measure must refer to that \emph{same} member.
A geometric spin-five transfer-strength ratio, such as the mathematical
$84{:}63{:}8$ angular comparison, is not a physical branching ratio,
width, lifetime, pole mass or selection of an on-shell hadronic channel.
Likewise an icosahedral sextic angular invariant does not set the
amplitude of a laboratory or cosmological observable without
source-to-spatial transport and detector normalization.
The finite-matrix geometry supplies a well-typed diagnostic, not a
new adjustable matter coupling or a replacement for the existing
color-singlet/physical Gauss and channel conditions.
\end{remark}
% END BERGER READOUT GATE v1.17



% BEGIN PAPER00 v1.5 -- exact rooted and dressed action currents
% New exact consequence for Paper VII-B (proposed v1.171).
\section{Dressed octet currents and exact same-action transition readouts}
\label{p11:sec:v375-dressed-baryon-transition}
\begin{corollary}[All-orders symmetry protection of the decuplet--octet tensor]
\label{p11:cor:v375-dressed-axial-reduced-tensor}
Suppose the actual Gauss-reduced finite strong Hamiltonian and
source projection of the degenerate dressed-current theorem of Paper~VII-A
are equivariant under \emph{exact} $SU(3)_{\rm flavour}\times SU(2)_{\rm spin}$,
and the physical field-derived current transforms as an octet spin
vector $(\mathbf8,\mathbf3)$.
Between specified isolated pole multiplets (each radial multiplicity resolved) of types
$(\mathbf{10},\mathbf4)$ and $(\mathbf8,\mathbf2)$ the complete
Schur-dressed current, including all complementary-state and residue
corrections, is proportional to one fixed reduced tensor:
\[
\boxed{\dim\operatorname{Hom}_{SU(3)\times SU(2)}
 ((\mathbf{10},\mathbf4),
  (\mathbf8,\mathbf2)\otimes(\mathbf8,\mathbf3))=1.}
\]
Hence the exact primitive Clebsch--Gordan strength ratios
$6:3:2:3$ (initial-state averaged) remain valid across one such
fully symmetric source/pole family, up to a \emph{common} reduced
matrix element. The corresponding unaveraged block sums are
$8:3:2:2$, as in the existing finite source ledger.
This statement is conditional on \emph{full} flavour covariance and
matched radial/kinematic branches: the observed unequal quark masses,
channel-dependent thresholds and spin-orbital/current corrections
are not symmetry-protected, and physical partial-width ratios cannot
be read from these numbers alone.
\end{corollary}
\begin{proof}
The Schur complementary projector, $C$, $B$, and each resolvent are
intertwiners under the stated exact group action. The $J$ insertion
transforms as $(\mathbf8,\mathbf3)$, so every term of
$\mathcal V_J(z_f,z_i)$ has the same covariance. Its pole normalization
$Z(E)$ is an intertwiner and cannot change the tensor multiplicity.
The flavour tensor product $\mathbf8\otimes\mathbf8$ contains
$\mathbf{10}$ exactly once, and the spin product
$\mathbf2\otimes\mathbf3$ contains $\mathbf4$ exactly once.
Schur's lemma then leaves one reduced coefficient for the entire
transition tensor. The existing CAR/primitive axial computation
fixes its relative component normalization on the specified
source convention, not its physical overall strength.
\end{proof}
\begin{remark}[Exact physical-current matching interface after the Schur closure]
\label{p11:rem:v375-action-current-remaining-edge}
The rooted one-body source of the exact CAR insertion theorem of Paper~VII-A
fixes its \emph{bare} colour-singlet spin--flavour tensor and Wilson
transport without per-channel weights. The dressed-current theorem
then fixes every finite transition from one \emph{actual} action and
its derivative, even if physical pole states contain sea, gluon or
other orbital components. The remaining occurrence obligations are:
(i) a selected Gauss-compatible Hamiltonian and axial-source member
with its full link/spin/source derivatives; (ii) actual parent and
daughter spectral projectors; and (iii) the controlled open-channel
continuum/meromorphic continuation and phase-space normalization for
irreversible widths. The outgoing finite five-core resolvent is an
eligible geometric comparison, not an independently imposed decay
clock or physical absorptive operator.
\end{remark}

% END PAPER00 v1.5

% BEGIN PAPER00 v1.6 -- transfer inversion and relative spectral port
% Proposed v1.6 VII-B addition: finite exact symmetry + relative-port obstruction.
\section{All-energy axial tensor reduction and the relative spectral port}
\label{p11:sec:v376-symmetry-moment-port}
\begin{theorem}[Entire finite mixed-moment axial packet reduces to four scalar families]
\label{p11:thm:v376-all-energy-axial-moment-factorization}
Let $G=SU(3)_{\rm flavor}\times SU(2)_{\rm spin}$ act unitarily
on a finite Gauss-compatible full carrier. Assume $H$ commutes
with this action and the source isometry
$\iota:\mathcal W\hookrightarrow\mathcal H$ intertwines it,
where
$\mathcal W=(\mathbf8,\mathbf2)\oplus(\mathbf{10},\mathbf4)$
is the inherited $56$-state carrier with one copy of each summand.
Let $J^{ia}$ be a covariant current transforming as
$(\mathbf8,\mathbf3)$. Then for every $k,m,n\ge0$
\[
 \boxed{M_k=\iota^\dagger H^k\iota
 =m_k^{(8)}P_8+m_k^{(10)}P_{10},}
 \qquad
 N_{mn}^{ia}=\iota^\dagger H^mJ^{ia}H^n\iota.
\]
Every decuplet-to-octet block has \emph{one and the same}
spin--flavor tensor $B^{ia}=P_8G^{ia}P_{10}$,
\[
 \boxed{P_8N_{mn}^{ia}P_{10}=\nu_{mn}B^{ia},\quad
 P_8\iota^\dagger(z_f-H)^{-1}J^{ia}(z_i-H)^{-1}
 \iota P_{10}=\nu(z_f,z_i)B^{ia}.}
\]
The octet-to-octet block has two (the usual $F,D$) reduced
scalar kernels, and the decuplet-to-decuplet block has one;
the reverse octet-to-decuplet block is Hermitian-conjugate after
interchanging moment and energy arguments. Thus the entire finite
24-component one-current transition packet reduces to four
independent scalar-kernel \emph{types} on this exact-symmetry branch.
The primitive one-body CAR source fixes the tensor normalization
$B^{ia}$, not the dynamical scalar kernels $\nu_{mn}$.
\end{theorem}
\begin{proof}
The powers and resolvents of $H$ commute with $G$. Since
$\iota$ intertwines $G$, Schur's lemma on the multiplicity-free
source gives the two scalar mass moments. The current insertion
transforms as $(\mathbf8,\mathbf3)$ for every pair of powers
or resolvents. In the tensor decompositions
$\mathbf8\otimes\mathbf8$, the output $\mathbf{10}$ occurs
once, the output $\mathbf8$ occurs twice, and in
$\mathbf{10}\otimes\mathbf8$, the output $\mathbf{10}$ occurs
once. The corresponding $SU(2)$ spin-vector products contain
the required target spins once each. The reduced-tensor counts
follow for every $m,n$ and off-spectrum pair of energies, without
restricting radial multiplicities in the full Hamiltonian. The
adjoint relation of a Hermitian current identifies reversed blocks.
\end{proof}

\begin{proposition}[Mass moments sharply bound but do not select the axial port]
\label{p11:prop:v376-sharp-relative-port-moment-bound}
For a common auxiliary radial Hilbert space and source unit vector
$e$, let $h_8,h_{10}$ be self-adjoint radial Hamiltonians and
let the off-diagonal current act as the radial identity times
$B^{ia}$. Write
$\mu_r=\langle e,h_re\rangle$,
$v_r=\langle e,h_r^2e\rangle-\mu_r^2$, and define
$\nu_{11}=\langle e,h_8h_{10}e\rangle$. Then
\[
\boxed{\nu_{11}=\mu_8\mu_{10}
+\langle(h_8-\mu_8)e,(h_{10}-\mu_{10})e\rangle,
\qquad |\nu_{11}-\mu_8\mu_{10}|\le\sqrt{v_8v_{10}}.}
\]
The bound is sharp for aligned or oppositely aligned auxiliary
leakage vectors. If the physical current has additional
source--auxiliary blocks, the complete mixed-moment formula of
Proposition~\ref{p10:prop:v376-first-mixed-moment-leakage}
replaces this restricted bound.
\end{proposition}
\begin{proof}
Subtract the two source means and apply Cauchy--Schwarz.
Equality holds if the two residual radial vectors are linearly
dependent by a nonzero scalar of the appropriate sign; for real
aligned and antialigned vectors both endpoints are attained.
\end{proof}

\begin{theorem}[In-domain finite symmetry counterexample: same masses and bare axial, different transitions]
\label{p11:thm:v376-identical-mass-distinct-current-moments}
On $\mathbb C^2\otimes\mathcal W$ put
\[
 h_8=\begin{pmatrix}1&1/3\\1/3&2\end{pmatrix},\quad
 h_{10}^{\pm}=\begin{pmatrix}2&\pm1/2\\\pm1/2&3\end{pmatrix},
 \quad H_\pm=h_8\otimes P_8+h_{10}^\pm\otimes P_{10},
 \quad J^{ia}=I_2\otimes G^{ia},
 \quad \iota x=e_1\otimes x.
\]
Both Hamiltonians are self-adjoint and fully
$SU(3)\times SU(2)$ invariant; both currents are the same
covariant Hermitian 24-component packet. Their complete spectra
and \emph{all} source-visible mass moments
$\iota^\dagger H_\pm^k\iota$ coincide, as does the full
bare source current $\iota^\dagger J^{ia}\iota=G^{ia}$.
Nevertheless
\[
 \boxed{P_8N_{11}^{ia,\pm}P_{10}
 =\left(2\pm\frac16\right)B^{ia},\qquad
 P_8(N_{11}^{ia,+}-N_{11}^{ia,-})P_{10}
 =\frac13 B^{ia}.}
\]
The sharp moment bound of the preceding proposition is saturated:
$v_8=1/9$, $v_{10}=1/4$. Moreover, the squared radial overlap
of the respective lowest octet and decuplet eigenvectors is
precisely $\frac12(1+5/\sqrt{26})\simeq0.99029034$ in the
plus case and $\frac12(1+1/\sqrt{26})\simeq0.59805807$ in the
minus case, although their pole energies are identical.
These are \emph{finite mathematical source models}, not two
certified physical Born--Hopf QCD action members. They disprove
any inference from mass spectral moments, exact symmetry and a
fixed primitive source-current normalization \emph{alone} to
the full dressed current: the missing datum is the action-owned
relative auxiliary-port orientation.
\end{theorem}
\begin{proof}
For $D=\operatorname{diag}(1,-1)$,
$h_{10}^-=Dh_{10}^+D$ and $De_1=e_1$; hence
$\langle e_1,(h_{10}^-)^ke_1\rangle=
\langle e_1,(h_{10}^+)^ke_1\rangle$ for every $k$.
The octet radial block is unchanged, so the full source moments
and spectra coincide. The current is identical by construction.
In the $8\leftarrow10$ sector, the mixed moment is
$(h_8h_{10}^{\pm})_{11}B^{ia}$; its coefficient equals
$2\pm1/6$. The source residual norms are $1/3$ and $1/2$,
so Cauchy--Schwarz is saturated. Direct diagonalization of the
two real $2\times2$ matrices gives the stated squared overlaps.
The pair is \emph{not} an equivalence of complete $(H,J)$ pairs:
the radial reflection transporting $H_+$ to $H_-$ does not
preserve the auxiliary octet--decuplet current. This is precisely
why the mixed moments distinguish the pairs.
\end{proof}

\begin{proposition}[Canonical symmetry-breaking defect and all-energy comparison bound]
\label{p11:prop:v376-haar-symmetry-response-bound}
On any fixed finite carrier with a compact spin--flavor action,
let $H_0$ be the Haar average of the actual $H$ and let $J_0$
be the orthogonal Haar projection of its current tuple onto
the covariant $(\mathbf8,\mathbf3)$ tensor subspace. These
are the unique Hilbert--Schmidt nearest symmetry-compatible
components of the given data; they are \emph{diagnostics}, not
replacement physical actions. Put $\Delta H=H-H_0$ and
$\Delta J=J-J_0$, and for a fixed source isometry $\iota$ set
\[
 C=\iota^\dagger R_H(z_f)J R_H(z_i)\iota,
 \quad C_0=\iota^\dagger R_{H_0}(z_f)J_0 R_{H_0}(z_i)\iota.
\]
For $\delta_a=\operatorname{dist}(z_a,\operatorname{spec}H)$
and $\delta_a^0=\operatorname{dist}(z_a,\operatorname{spec}H_0)$,
if all these distances are nonzero, then for each current component
\[
\boxed{\begin{aligned}
\|C-C_0\|_{\rm op}\le {}&\|\Delta H\|_{\rm op}
 \left(\frac{\|J\|_{\rm op}}{\delta_f\delta_f^0\delta_i}
 +\frac{\|J_0\|_{\rm op}}{\delta_f^0\delta_i\delta_i^0}\right)
 +\frac{\|\Delta J\|_{\rm op}}{\delta_f^0\delta_i}.
\end{aligned}}
\]
If the source intertwines the same group action, $C_0$ obeys
the exact reduced tensor classification of
Theorem~\ref{p11:thm:v376-all-energy-axial-moment-factorization}.
This gives a no-fitted-weight \emph{off-shell} bound on its
symmetry-breaking violation. The bound by itself does not control
an on-shell resonance at a vanishing spectral gap or establish
physical widths.
\end{proposition}
\begin{proof}
Compact Haar averaging is the orthogonal projection in the
finite Hilbert--Schmidt operator/tensor space onto the fixed
subspace of the group representation. Use the resolvent identity
$R_H-R_{H_0}=R_H\Delta H R_{H_0}$ and telescope the two-energy
product into
$(R_f-R_f^0)JR_i+R_f^0\Delta J R_i+
 R_f^0J_0(R_i-R_i^0)$.
The self-adjoint resolvent norms are inverse spectral distances;
submultiplicativity and $\|\iota\|=1$ give the displayed
bound. Group covariance of $H_0,J_0$ proves the last assertion.
\end{proof}

% END PAPER00 v1.6
% BEGIN PAPER00 v1.7 -- native spin-flavour Gram response boundary
\section{Primitive axial Gram and a positive auxiliary mass-order certificate}
\label{p11:sec:v377-positive-gram-boundary}
\begin{proposition}[Primitive $E_{36}$ Gram has no hidden auxiliary-current correction]
\label{p11:prop:v377-primitive-gram-zero-centered}
On the \emph{complete declared} symmetric $56$-state primitive
spin--flavor source $\mathcal W$, let
$A:=\sum_{ia}G^{ia}G^{ia}=(19/4)P_8+(7/2)P_{10}$, as already
established in Corollary~\ref{p11:cor:v369-axial-transition-benchmark}.
For the restricted algebraic trial $H_0=\kappa A$ and
$J_0^{ia}=G^{ia}$, with $\mathcal W$ itself the whole source space,
\[
\boxed{M_n=\left(\tfrac{19\kappa}{4}\right)^nP_8
 +\left(\tfrac{7\kappa}{2}\right)^nP_{10},\qquad
 P_8N_{mn}^{ia}P_{10}
 =\left(\tfrac{19\kappa}{4}\right)^m
  \left(\tfrac{7\kappa}{2}\right)^n P_8G^{ia}P_{10}.}
\]
All centered auxiliary-current moments of
the all-order centered auxiliary-moment theorem of proposed Paper~VII-A v1.187
vanish exactly, $\Xi_{mn}^{ia}=0$. In particular the
$\chi_{8,10}$ double-leakage scalar is exactly zero on this
\emph{isolated primitive model}. A nonzero physical auxiliary
correction must arise from actual same-action degrees outside this
primitive complete $56$-source, from moving-source derivatives, or
from a different physical current. This result does \emph{not}
assert the true QCD source has no leakage.
\end{proposition}
\begin{proof}
$H_0$ is scalar on each spin--flavor representation projector,
so its powers and the $8\leftarrow10$ current block give the
claimed formulas. With source projector $P=I_{\mathcal W}$
there is no complementary space: $Q=0$, $E_n=QH_0^n=0$, and
all the exact centered moments vanish.
\end{proof}

\begin{theorem}[Positive common-pole Schur parent forces the octet mass-order sign]
\label{p11:thm:v377-conditional-positive-gram-ordering}
Consider the \emph{conditional} common auxiliary member on
$\mathcal W\oplus\mathcal W$
\[
H_{g,\Delta,k}=\begin{pmatrix}
 kI&\sqrt g\,A^{1/2}\\
 \sqrt g\,A^{1/2}&\Delta I
\end{pmatrix},\qquad
J^{ia}=\operatorname{diag}(G^{ia},G^{ia}),
\]
where $g>0$, $k,\Delta\in\mathbb R$ are shared comparison
coordinates, \emph{not selected physical constants}. Then the
full positive connector return is
$\Sigma(z)=gA/(z-\Delta)$ and the exact source-visible pole
energies in the octet/decuplet sectors $r=8,10$ are
\[
 \boxed{E_r^{\pm}=\frac{k+\Delta\pm
 \sqrt{(\Delta-k)^2+4g\lambda_r}}2,
 \quad \lambda_8=19/4,\quad\lambda_{10}=7/2.}
\]
For all $g>0$ the lower-pole ordering is \emph{strict}:
\[
 \boxed{E_8^-<E_{10}^-,\qquad
 E_{10}^--E_8^-
 =\tfrac12\!\left(\sqrt{(\Delta-k)^2+19g}
 -\sqrt{(\Delta-k)^2+14g}\right)>0.}
\]
The full normalized $10\to8$ lower-pole current has the
primitive tensor multiplied by one common radial factor; its
\emph{squared} multiplier is
\[
 \boxed{\mathcal R_{8,10}
 =\frac12\left[1+
 \frac{(\Delta-k)^2+4g\sqrt{\lambda_8\lambda_{10}}}
 {\sqrt{(\Delta-k)^2+4g\lambda_8}
 \sqrt{(\Delta-k)^2+4g\lambda_{10}}}\right]\le1.}
\]
Thus the initial-state-averaged primitive
$6:3:2:3$ decuplet-to-octet axial-strength pattern is preserved
on this exactly covariant common-parent branch. For $\Delta>k$,
these are the primarily source-connected low poles.
The theorem derives the \emph{sign} and the common radial response
\emph{within the stated mathematical parent}; it does not show
that the actual Wilson/Clifford QCD action selects this member,
its values of $g,\Delta,k$, or a physical absorptive decay width.
\end{theorem}
\begin{proof}
The primitive Gram identity makes $A$ positive and scalar on
both representation summands. Hence $H_{g,\Delta,k}$ splits
into two-by-two radial Hamiltonians with off-diagonal
$\sqrt{g\lambda_r}$. Their characteristic polynomials give the
pole formula. Since $\lambda_8>\lambda_{10}$, the lower root
is strictly decreasing in $\lambda$; this fixes the ordering
without an independently chosen negative hyperfine coefficient.
For each radial lower eigenvector write its source and auxiliary
components in the same positive-norm convention. Their squared
mutual overlap is the displayed $\mathcal R_{8,10}$, obtained
from the Schur pole-residue formula in
the normalized finite Schur pole-state formula.
The current is the same primitive $G^{ia}$ in both radial copies,
so its $10\to8$ block is the single common radial overlap times
$P_8G^{ia}P_{10}$. Cauchy--Schwarz bounds the squared overlap by
one; the previously established primitive channel ratios follow.
\end{proof}
% END PAPER00 v1.7

% Proposed Paper VII-B v1.174: conditional reference-source consequence.
\begin{corollary}[Full-Hurwitz heat root factor cannot reweight the primitive strong channels]
\label{p11:cor:v378-heat-root-common-factor-no-go}
On the finite, aligned-spin, locally Gauss-invariant heat-reference
baryon source built from the existing rooted CAR source and
homogeneous $96$-face heat measure, the cross-root primitive axial
packet is $b_tG^{ia}$ with a single scalar $b_t$, independent of
the spin--flavor channel. Consequently the twelve nonzero channel
strengths in
Theorem~\ref{p11:thm:v371-unique-decuplet-current-strengths}
are all multiplied by $|b_t|^2$ and retain their exact relative
ratios. In particular,
\[
 \boxed{S_{\Delta N\pi}:S_{\Sigma^*\Lambda\pi}:
 S_{\Sigma^*\Sigma\pi}:S_{\Xi^*\Xi\pi}=6:3:2:3.}
\]
Moreover the same scalar Wilson factor multiplies the two-root
source Gram and the bare-current matrix. Canonical Gram
orthonormalization removes it completely, leaving
$I_2\otimes G^{ia}$ on the two supported root sectors (or the
corresponding support projector in a singular case). Thus the
nonzero reference overlap is not itself a new physical vertex.
The full finite heat-action reference constructs a nonzero shared
baryon overlap at leading order, but \emph{cannot by itself}
generate a daughter-dependent modification of the
$\Sigma^*\to\Lambda\pi$ versus $\Sigma^*\to\Sigma\pi$
primitive strength ratio. Such modifications must be obtained from
the actual action's spin/flavor/source jets, dressed pole residues,
momentum-dependent form factors or other declared effects. The
heat-reference scalar is not a physical decay-vertex normalization.
\end{corollary}
\begin{proof}
The rooted CAR identity gives a single scalar multiple of the
entire $24$-component primitive current. Every squared transition
matrix contraction is homogeneous of degree two in that current,
so each carries the same $|b_t|^2$ factor. The primitive ratios
are those of the existing exact $56$-state theorem; the common
factor cancels. This argument is independent of the finite heat
time and does not impose any physical axial current.
\end{proof}
% END proposed v1.174 heat source readout.

\end{document}
